\input{preamble}

% OK, start here.
%
\begin{document}

\title{Géométrie algébrique et formelle}


\maketitle

\phantomsection
\label{section-phantom}

\tableofcontents

\section{Introduction}
\label{section-introduction}

\noindent
Ce chapitre poursuit l'étude de la géométrie algébrique formelle,
et notamment de la question de savoir si un objet formel est
le complété d'un objet algébrique. Une référence fondamentale est \cite{SGA2}.
Voici une liste des résultats déjà abordés
dans le projet Champs :
\begin{enumerate}
\item Le théorème des fonctions formelles, voir
Cohomologie des schémas, section \ref{coherent-section-theorem-formal-functions}.
\item Les modules formels cohérents, voir
Cohomologie des schémas, section \ref{coherent-section-coherent-formal}.
\item Le théorème d'existence de Grothendieck, voir
Cohomologie des schémas, sections \ref{coherent-section-existence},
\ref{coherent-section-existence-proper} et
\ref{coherent-section-existence-proper-support}.
\item Le théorème d'algébrisation de Grothendieck, voir
Cohomologie des schémas, section \ref{coherent-section-algebraization}.
\item Le théorème d'existence de Grothendieck dans un cadre plus général, voir
Compléments sur la platitude, sections \ref{flat-section-existence} et
\ref{flat-section-existence-derived}.
\end{enumerate}
Donnons un aperçu du contenu de ce chapitre.

\medskip\noindent
Soient $X$ un schéma et $\mathcal{I} \subset \mathcal{O}_X$
un faisceau quasi-cohérent d'idéaux de type fini. De nombreuses questions
de ce chapitre portent sur des systèmes projectifs $(\mathcal{F}_n)$
de $\mathcal{O}_X$-modules quasi-cohérents tels que
$\mathcal{F}_n = \mathcal{F}_{n + 1}/\mathcal{I}^n\mathcal{F}_{n + 1}$.
Un cas particulier important est celui où $X$ est un schéma sur un anneau
noethérien $A$ et où $\mathcal{I} = I \mathcal{O}_X$ pour un idéal $I \subset A$.
Dans Cohomologie, sections \ref{cohomology-section-inverse-systems},
\ref{cohomology-section-inverse-systems-bis} et
\ref{cohomology-section-inverse-systems-tri}, nous avons établi des résultats
généraux. Dans ce chapitre, les sections \ref{section-ML-degree-zero} et
\ref{section-formal-functions-principal} contiennent des résultats propres
aux schémas et aux modules quasi-cohérents. Dans la
section \ref{section-formal-sections-cd-one}, nous démontrons
que la topologie de limite sur $\lim H^p(X, \mathcal{F}_n)$ est $I$-adique
lorsque $\text{cd}(A, I) = 1$. L'un des thèmes
de ce chapitre sera de montrer que les résultats démontrés lorsque
l'idéal est principal, $I = (f)$,
restent valables sous la seule hypothèse $\text{cd}(A, I) = 1$.

\medskip\noindent
Dans la section \ref{section-derived-completion}, nous étudions la complétion dérivée
des modules sur un site annelé $(\mathcal{C}, \mathcal{O})$
relativement à un faisceau d'idéaux de type fini $\mathcal{I}$.
Cette section prolonge naturellement la théorie de la complétion dérivée
en algèbre commutative, telle qu'elle est exposée dans
Compléments d'algèbre, section \ref{more-algebra-section-derived-completion}.
Le premier résultat principal affirme l'existence de la complétion dérivée.
Le second affirme que, pour un morphisme $f$ de sites annelés,
la complétion dérivée commute à l'image directe dérivée :
$$
(Rf_*K)^\wedge = Rf_*(K^\wedge)
$$
si le faisceau d'idéaux à la source est localement engendré par des sections provenant
de l'idéal sur la base, voir le
lemme \ref{lemma-pushforward-commutes-with-derived-completion}.
Soulignons que ces deux résultats principaux sont très élémentaires lorsque
les idéaux en question sont globalement de type fini, ce qui sera
le cas dans toutes les applications de cette théorie faites dans ce chapitre.
L'égalité ci-dessus est la version « correcte » du
théorème des fonctions formelles ; voir la discussion de la
section \ref{section-formal-functions}.

\medskip\noindent
Soient $A$ un anneau noethérien et $I, J$ deux idéaux de $A$.
Soit $M$ un $A$-module de type fini.
Le sujet suivant de ce chapitre est l'application
$$
R\Gamma_J(M) \longrightarrow R\Gamma_J(M)^\wedge
$$
de la cohomologie locale de $M$ dans son complété $I$-adique dérivé.
Il se trouve que, sous des conditions de profondeur convenables,
cette application induit un isomorphisme en cohomologie dans un intervalle de degrés.
Dans la section \ref{section-algebraization-sections-general},
nous travaillons essentiellement dans le cadre général qui vient d'être indiqué.
Dans la section \ref{section-algebraization-punctured},
nous supposons que $A$ est un anneau local et que $J = \mathfrak m$ est un idéal maximal.
Nous conseillons au lecteur de lire cette section avant les deux autres
de cette partie du chapitre.
Enfin, dans la section \ref{section-bootstrap}, nous exploitons
le cas local pour obtenir des résultats plus forts dans le cas général.

\medskip\noindent
Dans la partie suivante de ce chapitre, nous utilisons les résultats sur
la complétion de la cohomologie locale pour obtenir une liste non exhaustive de résultats sur
la cohomologie du complété des modules cohérents.
Plus précisément, soient $A$ un anneau noethérien, $I \subset A$
un idéal et $U \subset \Spec(A)$ un sous-schéma ouvert.
Si $\mathcal{F}$ est un $\mathcal{O}_U$-module cohérent, alors
nous pouvons considérer les applications
$$
H^i(U, \mathcal{F}) \longrightarrow \lim H^i(U, \mathcal{F}/I^n\mathcal{F})
$$
et nous demander si elles sont des isomorphismes dans un certain intervalle de degrés.
Dans la section \ref{section-algebraization-sections},
nous étudions quelques exemples où $U$ est le spectre épointé
d'un anneau local. Dans la section \ref{section-algebraization-sections-coherent},
nous traitons le cas général.
Dans la section \ref{section-connected}, nous appliquons certains des résultats
obtenus à des questions de connexité en géométrie algébrique.

\medskip\noindent
Les dernières sections de ce chapitre sont consacrées à l'étude
de l'algébrisation des modules formels cohérents. Autrement dit, étant donné
un système projectif de modules cohérents $(\mathcal{F}_n)$ sur $U$
comme ci-dessus, avec
$\mathcal{F}_n = \mathcal{F}_{n + 1}/I^n\mathcal{F}_{n + 1}$
nous demandons s'il existe un $\mathcal{O}_U$-module cohérent
$\mathcal{F}$ tel que
$\mathcal{F}_n = \mathcal{F}/I^n\mathcal{F}$
pour tout $n$. Nous conseillons au lecteur de consulter la
section \ref{section-algebraization-modules}
pour un énoncé précis de la question, un résultat général utile
(lemme \ref{lemma-when-done}) et une application non triviale
(lemme \ref{lemma-algebraization-principal-variant}).
La démonstration d'un résultat dépassant substantiellement ce cas
exige de développer une théorie sensiblement plus étendue.
On trouvera dans la section \ref{section-algebraization-modules-conclusion}
les résultats de ce type les plus forts obtenus dans ce chapitre.



\section{sections formelles, I}
\label{section-ML-degree-zero}

\noindent
Nous conseillons de consulter d'abord
Cohomologie, section \ref{cohomology-section-inverse-systems}.

\begin{lemma}
\label{lemma-properties-system}
Soient $X$ un schéma et $\mathcal{I} \subset \mathcal{O}_X$
un faisceau quasi-cohérent d'idéaux. Soit
$$
\ldots \to \mathcal{F}_3 \to \mathcal{F}_2 \to \mathcal{F}_1
$$
un système projectif de $\mathcal{O}_X$-modules quasi-cohérents
tel que
$\mathcal{F}_n = \mathcal{F}_{n + 1}/\mathcal{I}^n\mathcal{F}_{n + 1}$.
Posons $\mathcal{F} = \lim \mathcal{F}_n$. Alors
\begin{enumerate}
\item $\mathcal{F} = R\lim \mathcal{F}_n$,
\item pour tout ouvert affine $U \subset X$, on a
$H^p(U, \mathcal{F}) = 0$ pour $p > 0$, et
\item pour tout $p$, il existe une suite exacte courte
$0 \to R^1\lim H^{p - 1}(X, \mathcal{F}_n) \to
H^p(X, \mathcal{F}) \to \lim H^p(X, \mathcal{F}_n) \to 0$.
\end{enumerate}
Si, de plus, $\mathcal{I}$ est de type fini, alors
\begin{enumerate}
\item[(4)]
$\mathcal{F}_n = \mathcal{F}/\mathcal{I}^n\mathcal{F}$, et
\item[(5)]
$\mathcal{I}^n \mathcal{F} = \lim_{m \geq n} \mathcal{I}^n\mathcal{F}_m$.
\end{enumerate}
\end{lemma}

\begin{proof}
Les assertions (1), (2) et (3) sont des faits généraux sur les systèmes projectifs de
modules quasi-cohérents dont les applications de transition sont surjectives ; voir
Catégories dérivées des schémas, lemme \ref{perfect-lemma-Rlim-quasi-coherent}
et Cohomologie, lemme \ref{cohomology-lemma-RGamma-commutes-with-Rlim}.
Supposons maintenant que $\mathcal{I}$ est de type fini.
Soit $U \subset X$ un ouvert affine. Écrivons $U = \Spec(A)$ et supposons que $\mathcal{I}|_U$
correspond à $I \subset A$. Remarquons que $I$ est un idéal de type fini.
L'équivalence entre la catégorie des $\mathcal{O}_U$-modules quasi-cohérents
et celle des $A$-modules (Schémas, lemme \ref{schemes-lemma-equivalence-quasi-coherent})
montre que $M_n = \mathcal{F}_n(U)$ est un système projectif
de $A$-modules tel que $M_n = M_{n + 1}/I^nM_{n + 1}$. Ainsi
$$
M = \mathcal{F}(U) = \lim \mathcal{F}_n(U) = \lim M_n
$$
est un module $I$-adiquement complet, avec $M/I^nM = M_n$, d'après
Algèbre, lemme \ref{algebra-lemma-limit-complete}. Cela prouve (4).
L'assertion (5) se traduit par l'égalité
$\lim_{m \geq n} I^nM/I^mM = I^nM$.
Puisque $I^mM = I^{m - n} \cdot I^nM$, cela signifie simplement que
$I^mM$ est $I$-adiquement complet. Cela résulte
d'Algèbre, lemme \ref{algebra-lemma-hathat-finitely-generated},
et du fait que $M$ est complet.
\end{proof}





\section{sections formelles, II}
\label{section-formal-functions-principal}

\noindent
Nous conseillons de consulter d'abord
Cohomologie, sections
\ref{cohomology-section-inverse-systems-bis} et
\ref{cohomology-section-inverse-systems-tri}.

\begin{lemma}
\label{lemma-equivalent-f-good}
Soient $X$ un schéma et $f \in \Gamma(X, \mathcal{O}_X)$. Soit
$$
\ldots \to \mathcal{F}_3 \to \mathcal{F}_2 \to \mathcal{F}_1
$$
un système projectif de $\mathcal{O}_X$-modules quasi-cohérents.
Les conditions suivantes sont équivalentes :
\begin{enumerate}
\item pour tout $n \geq 1$, l'application
$f : \mathcal{F}_{n + 1} \to \mathcal{F}_{n + 1}$ se factorise
par $\mathcal{F}_{n + 1} \to \mathcal{F}_n$ et définit une
suite exacte courte
$0 \to \mathcal{F}_n \to \mathcal{F}_{n + 1} \to \mathcal{F}_1 \to 0$,
\item pour tout $n \geq 1$, l'application
$f^n : \mathcal{F}_{n + 1} \to \mathcal{F}_{n + 1}$
se factorise par $\mathcal{F}_{n + 1} \to \mathcal{F}_1$
et définit une suite exacte courte
$0 \to \mathcal{F}_1 \to \mathcal{F}_{n + 1} \to \mathcal{F}_n \to 0$
\item il existe un $\mathcal{O}_X$-module $\mathcal{G}$
qui est $f$-divisible et tel que $\mathcal{F}_n = \mathcal{G}[f^n]$.
\item il existe un $\mathcal{O}_X$-module $\mathcal{F}$ qui est
sans $f$-torsion et tel que $\mathcal{F}_n = \mathcal{F}/f^n\mathcal{F}$.
\end{enumerate}
\end{lemma}

\begin{proof}
L'équivalence de (1), (2), (3) et l'implication (4) $\Rightarrow$ (1)
sont démontrées dans Cohomologie, lemme \ref{cohomology-lemma-equivalent-f-good}.
Supposons (1) vérifiée. Posons
$\mathcal{F} = \lim \mathcal{F}_n$. D'après le lemme \ref{lemma-properties-system},
assertion (4), on a $\mathcal{F}_n = \mathcal{F}/f^n\mathcal{F}$.
Soient $U \subset X$ un ouvert et
$s = (s_n) \in \mathcal{F}(U) = \lim \mathcal{F}_n(U)$.
Choisissons $n \geq 1$. Si $fs = 0$, alors $s_{n + 1}$ appartient au noyau de
$\mathcal{F}_{n + 1} \to \mathcal{F}_n$ par la condition (1).
Ainsi $s_n = 0$. Comme $n$ est arbitraire, on obtient $s = 0$.
Donc $\mathcal{F}$ est sans $f$-torsion.
\end{proof}

\begin{lemma}
\label{lemma-formal-functions-principal}
\begin{reference}
Version légèrement améliorée de \cite[Lemma 1.6]{Bhatt-local}
\end{reference}
Soient $A$ un anneau et $f \in A$. Soit $X$ un schéma sur $A$.
Soit $\mathcal{F}$ un $\mathcal{O}_X$-module quasi-cohérent.
Supposons que la suite $\mathcal{F}[f^n] = \Ker(f^n : \mathcal{F} \to \mathcal{F})$
stationne. Alors
$$
R\Gamma(X, \lim \mathcal{F}/f^n\mathcal{F}) =
R\Gamma(X, \mathcal{F})^\wedge
$$
où le membre de droite désigne le complété dérivé
relativement à l'idéal $(f) \subset A$. Par conséquent, pour
$p \in \mathbf{Z}$, on obtient un diagramme commutatif
$$
\xymatrix{
& 0 & 0 \\
0 \ar[r] &
\widehat{H^p(X, \mathcal{F})} \ar[r] \ar[u] &
\lim H^p(X, \mathcal{F}/f^n\mathcal{F}) \ar[r] \ar[u] &
T_f(H^{p + 1}(X, \mathcal{F})) \ar[r] &
0 \\
0 \ar[r] &
H^0(H^p(X, \mathcal{F})^\wedge) \ar[r] \ar[u] &
H^p(X, \lim \mathcal{F}/f^n\mathcal{F}) \ar[r] \ar[u] &
T_f(H^{p + 1}(X, \mathcal{F})) \ar[r] \ar@{=}[u] &
0 \\
&
R^1\lim H^p(X, \mathcal{F})[f^n] \ar[u] \ar[r]^\cong &
R^1\lim H^{p - 1}(X, \mathcal{F}/f^n\mathcal{F}) \ar[u] \\
& 0 \ar[u] & 0 \ar[u]
}
$$
à lignes et colonnes exactes, où
$\widehat{H^p(X, \mathcal{F})} =
\lim H^p(X, \mathcal{F})/f^n H^p(X, \mathcal{F})$
est le complété $f$-adique usuel
et où $T_f(-)$ désigne le module de Tate $f$-adique comme dans
Compléments d'algèbre, exemple
\ref{more-algebra-example-spectral-sequence-principal}.
\end{lemma}

\begin{proof}
D'après le lemme \ref{lemma-properties-system}, on a
$\lim \mathcal{F}/f^n\mathcal{F} = R\lim \mathcal{F}/f^n \mathcal{F}$.
Le reste découle de Cohomologie, exemple
\ref{cohomology-example-formal-functions-principal}.
\end{proof}









\section{sections formelles, III}
\label{section-formal-sections-cd-one}

\noindent
Dans cette section, nous démontrons le lemme \ref{lemma-topology-I-adic},
qui, dans le cadre des schémas noethériens et des modules cohérents,
est l'analogue de Cohomologie, lemme \ref{cohomology-lemma-topology-I-adic-f},
lorsque l'idéal $I$ n'est pas supposé principal, mais
vérifie $\text{cd}(A, I) = 1$.

\begin{lemma}
\label{lemma-cd-one}
Soit $I = (f_1, \ldots, f_r)$ un idéal d'un anneau noethérien $A$.
Si $\text{cd}(A, I) = 1$, il existe $c \geq 1$ et des applications
$\varphi_j : I^c \to A$ telles que $\sum f_j \varphi_j : I^c \to I$
soit l'inclusion.
\end{lemma}

\begin{proof}
Puisque $\text{cd}(A, I) = 1$, le complémentaire $U = \Spec(A) \setminus V(I)$
est affine (Cohomologie locale, lemme \ref{local-cohomology-lemma-cd-is-one}).
Écrivons $U = \Spec(B)$. Alors $IB = B$
et l'on peut écrire $1 = \sum_{j = 1, \ldots, r} f_j b_j$
pour des $b_j \in B$. D'après
Cohomologie des schémas, lemme \ref{coherent-lemma-homs-over-open},
on peut représenter les $b_j$ par des applications $\varphi_j : I^c \to A$
pour un certain $c \geq 0$. Alors $\sum f_j \varphi_j : I^c \to I \subset A$
est l'inclusion canonique, quitte à remplacer $c$ par un entier
plus grand, d'après le même lemme.
\end{proof}

\begin{lemma}
\label{lemma-cd-one-extend}
Soit $I = (f_1, \ldots, f_r)$ un idéal d'un anneau noethérien $A$
tel que $\text{cd}(A, I) = 1$. Soient $c \geq 1$ et $\varphi_j : I^c \to A$,
$j = 1, \ldots, r$, comme dans le lemme \ref{lemma-cd-one}.
Il existe alors un unique homomorphisme de $A$-algèbres graduées
$$
\Phi : \bigoplus\nolimits_{n \geq 0} I^{nc} \to A[T_1, \ldots, T_r]
$$
tel que $\Phi(g) = \sum \varphi_j(g) T_j$ pour $g \in I^c$.
De plus, la composée de $\Phi$ avec l'application
$A[T_1, \ldots, T_r] \to \bigoplus_{n \geq 0} I^n$,
$T_j \mapsto f_j$, est l'inclusion
$\bigoplus_{n \geq 0} I^{nc} \to \bigoplus_{n \geq 0} I^n$.
\end{lemma}

\begin{proof}
Pour tout $j$ et tout $m \geq c$, la restriction de $\varphi_j$ à
$I^m$ est une application $\varphi_j : I^m \to I^{m - c}$.
Étant donnés $j_1, \ldots, j_n \in \{1, \ldots, r\}$, montrons que la
composée
$$
\varphi_{j_1} \ldots \varphi_{j_n} :
I^{nc} \to I^{(n - 1)c} \to \ldots \to I^c \to A
$$
est indépendante de l'ordre des indices $j_1, \ldots, j_n$.
En effet, si $g = g_1 \ldots g_n$ avec $g_i \in I^c$, on
voit que
$$
(\varphi_{j_1} \ldots \varphi_{j_n})(g) =
\varphi_{j_1}(g_1) \ldots \varphi_{j_n}(g_n)
$$
est indépendant de l'ordre, car la multiplication dans $A$ est commutative.
On peut donc définir $\Phi$ en envoyant $g \in I^{nc}$ sur
$$
\Phi(g) = \sum\nolimits_{e_1 + \ldots + e_r = n}
(\varphi_1^{e_1} \circ \ldots \circ \varphi_r^{e_r})(g)
T_1^{e_1} \ldots T_r^{e_r}
$$
On vérifie directement qu'il s'agit d'un homomorphisme de $A$-algèbres
graduées possédant la propriété voulue. L'unicité est immédiate,
ainsi que la dernière propriété. Le lemme est démontré.
\end{proof}

\begin{lemma}
\label{lemma-cd-one-extend-to-module}
Soit $I = (f_1, \ldots, f_r)$ un idéal d'un anneau noethérien $A$
tel que $\text{cd}(A, I) = 1$. Soient $c \geq 1$ et $\varphi_j : I^c \to A$,
$j = 1, \ldots, r$, comme dans le lemme \ref{lemma-cd-one}.
Soit $A \to B$ un homomorphisme d'anneaux, avec $B$ noethérien, et soit $N$
un $B$-module de type fini. Alors, quitte à augmenter $c$
et à modifier les $\varphi_j$ en conséquence, il existe un unique
homomorphisme de $B$-modules gradués
$$
\Phi_N : \bigoplus\nolimits_{n \geq 0} I^{nc}N \to N[T_1, \ldots, T_r]
$$
tel que $\Phi_N(g x) = \Phi(g) x$ pour $g \in I^{nc}$ et $x \in N$,
où $\Phi$ est comme dans le lemme \ref{lemma-cd-one-extend}.
La composée de $\Phi_N$ avec l'application
$N[T_1, \ldots, T_r] \to \bigoplus_{n \geq 0} I^nN$,
$T_j \mapsto f_j$, est l'inclusion
$\bigoplus_{n \geq 0} I^{nc}N \to \bigoplus_{n \geq 0} I^nN$.
\end{lemma}

\begin{proof}
L'unicité résulte clairement de la formule et de l'unicité de $\Phi$ dans le
lemme \ref{lemma-cd-one-extend}. Considérons la $A$-algèbre noethérienne
$B' = B \oplus N$, où $N$ est un idéal de carré nul. Pour établir
l'existence de $\Phi_N$, il suffit
(d'après le lemme \ref{lemma-cd-one}) de montrer que $\varphi_j$ se prolonge en
une application $\varphi'_j : I^cB' \to B'$, quitte à augmenter $c$
en un certain $c'$ (et à remplacer $\varphi_j$ par la composée de l'inclusion
$I^{c'} \to I^c$ avec $\varphi_j$). Rappelons que $\varphi_j$ correspond à une
section
$$
h_j \in \Gamma(\Spec(A) \setminus V(I), \mathcal{O}_{\Spec(A)})
$$
voir Cohomologie des schémas, lemme \ref{coherent-lemma-homs-over-open}.
(C'est d'ailleurs ainsi que nous avons choisi les $\varphi_j$ dans la démonstration du
lemme \ref{lemma-cd-one}.) Utilisons ce même lemme pour représenter l'image inverse
$$
h'_j \in \Gamma(\Spec(B') \setminus V(IB'), \mathcal{O}_{\Spec(B')})
$$
de $h_j$ par une application $B'$-linéaire
$\varphi'_j : I^{c'}B' \to B'$ pour un certain $c' \geq c$.
La compatibilité avec $\varphi_j$ sera vérifiée pour $c'$
assez grand, par une nouvelle application du lemme :
il suffit en effet de la vérifier sur une famille finie de générateurs de $I^{c'}$.
Nous omettons ce détail.
\end{proof}

\begin{lemma}
\label{lemma-cd-is-one-for-system}
Soit $I = (f_1, \ldots, f_r)$ un idéal d'un anneau noethérien $A$ tel que
$\text{cd}(A, I) = 1$. Soient $c \geq 1$ et $\varphi_j : I^c \to A$,
$j = 1, \ldots, r$, comme dans le lemme \ref{lemma-cd-one}.
Soit $X$ un schéma noethérien sur $\Spec(A)$. Soit
$$
\ldots \to \mathcal{F}_3 \to \mathcal{F}_2 \to \mathcal{F}_1
$$
un système projectif de $\mathcal{O}_X$-modules cohérents
tel que $\mathcal{F}_n = \mathcal{F}_{n + 1}/I^n\mathcal{F}_{n + 1}$.
Posons $\mathcal{F} = \lim \mathcal{F}_n$.
Alors, quitte à augmenter $c$ et à modifier les $\varphi_j$ en conséquence,
il existe un unique homomorphisme de $\mathcal{O}_X$-modules gradués
$$
\Phi_\mathcal{F} :
\bigoplus\nolimits_{n \geq 0} I^{nc}\mathcal{F}
\longrightarrow
\mathcal{F}[T_1, \ldots, T_r]
$$
tel que $\Phi_\mathcal{F}(g s) = \Phi(g) s$ pour $g \in I^{nc}$ et
$s$ une section locale de $\mathcal{F}$, où $\Phi$ est comme dans le
lemme \ref{lemma-cd-one-extend}. La composée de $\Phi_\mathcal{F}$
avec l'application
$\mathcal{F}[T_1, \ldots, T_r] \to \bigoplus_{n \geq 0} I^n\mathcal{F}$,
$T_j \mapsto f_j$
est l'inclusion canonique
$\bigoplus_{n \geq 0} I^{nc}\mathcal{F} \to
\bigoplus_{n \geq 0} I^n\mathcal{F}$.
\end{lemma}

\begin{proof}
L'unicité résulte immédiatement de la $\mathcal{O}_X$-linéarité
et de la condition $\Phi_\mathcal{F}(g s) = \Phi(g) s$ pour
$g \in I^{nc}$ et $s$ une section locale de $\mathcal{F}$.
On peut donc supposer $X = \Spec(B)$ affine.
Remarquons que $(\mathcal{F}_n)$ est un objet de la catégorie
$\textit{Coh}(X, I\mathcal{O}_X)$ introduite
dans Cohomologie des schémas, section \ref{coherent-section-coherent-formal}.
Soit $B' = B^\wedge$ le complété $I$-adique de $B$.
D'après Cohomologie des schémas, lemme \ref{coherent-lemma-inverse-systems-affine},
l'objet $(\mathcal{F}_n)$ correspond à un $B'$-module de type fini $N$,
au sens où $\mathcal{F}_n$ est le module cohérent
associé au $B$-module de type fini $N/I^n N$.
En appliquant le lemme \ref{lemma-cd-one-extend-to-module}
à $I \subset A \to B'$ et à $N$,
on voit que, quitte à augmenter $c$ et à modifier les
$\varphi_j$ en conséquence, on obtient des applications uniques
$$
\Phi_N : \bigoplus\nolimits_{n \geq 0} I^{nc}N \to N[T_1, \ldots, T_r]
$$
ayant les propriétés correspondantes. Notons qu'en degré $n$, on obtient
un système projectif d'applications $N/I^mN \to \bigoplus_{e_1 + \ldots + e_r = n}
N/I^{m - nc}N \cdot T_1^{e_1} \ldots T_r^{e_r}$ pour $m \geq nc$.
En revenant au langage des faisceaux
cohérents, on voit que $\Phi_N$ correspond à un système d'applications
$$
\Phi^n_m :
I^{nc}\mathcal{F}_m
\longrightarrow
\bigoplus\nolimits_{e_1 + \ldots + e_r = n}
\mathcal{F}_{m - nc} \cdot T_1^{e_1} \ldots T_r^{e_r}
$$
pour $m \geq nc$ et $n \geq 1$ variables. En prenant la limite projective
de ces applications suivant $m$, on obtient $\Phi_\mathcal{F} = \bigoplus_n \lim_m \Phi^n_m$.
Notons que $\lim_m I^t\mathcal{F}_m = I^t \mathcal{F}$, comme on peut le voir
en évaluant sur les ouverts affines, par exemple ; mais ce fait est inutile ici, car
il existe clairement une application $I^t\mathcal{F} \to \lim_m I^t\mathcal{F}_m$.
\end{proof}

\begin{lemma}
\label{lemma-topology-I-adic}
Soient $I$ un idéal d'un anneau noethérien $A$ et $X$ un schéma noethérien
sur $\Spec(A)$. Soit
$$
\ldots \to \mathcal{F}_3 \to \mathcal{F}_2 \to \mathcal{F}_1
$$
un système projectif de $\mathcal{O}_X$-modules cohérents
tel que $\mathcal{F}_n = \mathcal{F}_{n + 1}/I^n\mathcal{F}_{n + 1}$.
Si $\text{cd}(A, I) = 1$, alors, pour tout $p \in \mathbf{Z}$, la topologie de limite sur
$\lim H^p(X, \mathcal{F}_n)$ est $I$-adique.
\end{lemma}

\begin{proof}
Il est d'abord clair que l'image de $I^t \lim H^p(X, \mathcal{F}_n)$
dans $H^p(X, \mathcal{F}_t)$ est nulle. La topologie $I$-adique
est donc plus fine que la topologie de limite. Pour la réciproque, posons
$\mathcal{F} = \lim \mathcal{F}_n$, choisissons des générateurs $f_1, \ldots, f_r$
de $I$, choisissons $c \geq 1$, puis choisissons
$\Phi_\mathcal{F}$ comme dans le lemme \ref{lemma-cd-is-one-for-system}.
Nous utiliserons les résultats du lemme \ref{lemma-properties-system}
sans autre mention. En particulier, on a une suite exacte
courte
$$
0 \to R^1\lim H^{p - 1}(X, \mathcal{F}_n) \to H^p(X, \mathcal{F})
\to \lim H^p(X, \mathcal{F}_n) \to 0
$$
On peut donc relever tout élément $\xi$ de $\lim H^p(X, \mathcal{F}_n)$
en un élément $\xi' \in H^p(X, \mathcal{F})$. Supposons que $\xi$ ait une image nulle
dans $H^p(X, \mathcal{F}_{nc})$ pour un certain $n$ ; autrement
dit, supposons que $\xi$ soit « petit » pour la topologie de limite. On a une
suite exacte courte
$$
0 \to I^{nc}\mathcal{F} \to \mathcal{F} \to \mathcal{F}_{nc} \to 0
$$
et l'hypothèse signifie donc que l'on peut relever $\xi'$ en un élément
$\xi'' \in H^p(X, I^{nc}\mathcal{F})$. En appliquant $\Phi_\mathcal{F}$,
on obtient
$$
\Phi_\mathcal{F}(\xi'') = \sum\nolimits_{e_1 + \ldots + e_r = n}
\xi'_{e_1, \ldots, e_r} \cdot T_1^{e_1} \ldots T_r^{e_r}
$$
pour certains $\xi'_{e_1, \ldots, e_r} \in H^p(X, \mathcal{F})$.
En notant $\xi_{e_1, \ldots, e_r} \in \lim H^p(X, \mathcal{F}_n)$
leurs images et en utilisant la dernière assertion du
lemme \ref{lemma-cd-is-one-for-system},
on conclut que
$$
\xi = \sum f_1^{e_1} \ldots f_r^{e_r} \xi_{e_1, \ldots, e_r}
$$
appartient à $I^n \lim H^p(X, \mathcal{F}_n)$, comme voulu.
\end{proof}

\begin{example}
\label{example-not-I-adic}
Soit $k$ un corps. Posons $A = k[x, y][[s, t]]/(xs - yt)$.
Posons $I = (s, t)$ et $\mathfrak a = (x, y, s, t)$.
Posons $X = \Spec(A) - V(\mathfrak a)$ et
$\mathcal{F}_n = \mathcal{O}_X/I^n\mathcal{O}_X$.
Remarquons que la fonction rationnelle
$$
g = \frac{t}{x} = \frac{s}{y}
$$
est régulière dans un voisinage ouvert $V \subset X$ de
$V(I\mathcal{O}_X)$. Toute puissance $g^e$ définit donc une section
$g^e \in M = \lim H^0(X, \mathcal{F}_n)$. Remarquons que
$g^e \to 0$ lorsque $e \to \infty$ pour la topologie de limite sur $M$,
puisque l'image de $g^e$ dans $\mathcal{F}_e$ est nulle.
En revanche, $g^e \not \in IM$ pour tout $e$,
comme le lecteur pourra le voir en calculant $H^0(U, \mathcal{F}_n)$ ;
nous omettons ce calcul. Remarquons que $\text{cd}(A, I) = 2$.
Le résultat du lemme \ref{lemma-topology-I-adic} est donc optimal.
\end{example}








\section{Conditions de Mittag-Leffler}
\label{section-ML}

\noindent
Lorsque l'on prend la cohomologie locale relativement à l'idéal maximal
d'un anneau local noethérien, la condition de Mittag-Leffler est souvent
automatiquement satisfaite. Il en va donc de même des groupes de cohomologie
supérieure d'un système projectif de faisceaux cohérents à applications de transition
surjectives sur le spectre épointé.

\begin{lemma}
\label{lemma-descending-chain}
Soit $(A, \mathfrak m)$ un anneau local noethérien.
\begin{enumerate}
\item Soit $M$ un $A$-module de type fini. Alors le $A$-module
$H^i_\mathfrak m(M)$ satisfait à la condition des chaînes descendantes
pour tout $i$.
\item Soit $U = \Spec(A) \setminus \{\mathfrak m\}$ le
spectre épointé de $A$.
Soit $\mathcal{F}$ un $\mathcal{O}_U$-module cohérent.
Alors le $A$-module $H^i(U, \mathcal{F})$
satisfait à la condition des chaînes descendantes pour $i > 0$.
\end{enumerate}
\end{lemma}

\begin{proof}
Démontrons (1) par récurrence sur la dimension du support de $M$.
L'énoncé vaut si $M = 0$ ; on peut donc supposer, et l'on suppose, $M$ non nul.

\medskip\noindent
Initialisation de la récurrence.
Si $\dim(\text{Supp}(M)) = 0$, le support
de $M$ est $\{\mathfrak m\}$, et l'on voit que $H^0_\mathfrak m(M) = M$
et $H^i_\mathfrak m(M) = 0$ pour $i > 0$, comme il résulte de la
construction de la cohomologie locale ; voir
Complexes dualisants, section \ref{dualizing-section-local-cohomology}.
Puisque $M$ est de longueur finie (Algèbre, lemme \ref{algebra-lemma-length-finite}),
il satisfait à la condition des chaînes descendantes.

\medskip\noindent
Étape de récurrence. Supposons $\dim(\text{Supp}(M)) > 0$.
D'après le cas initial, le module de type fini $H^0_\mathfrak m(M) \subset M$
satisfait à la condition des chaînes descendantes.
D'après Complexes dualisants, lemme \ref{dualizing-lemma-divide-by-torsion},
on peut remplacer $M$ par $M/H^0_\mathfrak m(M)$.
Alors $H^0_\mathfrak m(M) = 0$, c'est-à-dire que $M$ est de profondeur $\geq 1$ ; voir
Complexes dualisants, lemme \ref{dualizing-lemma-depth}.
Choisissons $x \in \mathfrak m$ tel que $x : M \to M$ soit injective.
D'après Algèbre, lemme \ref{algebra-lemma-one-equation-module}, on a
$\dim(\text{Supp}(M/xM)) = \dim(\text{Supp}(M)) - 1$ et l'hypothèse
de récurrence s'applique. Fixons un indice $i$ et considérons la
suite exacte
$$
H^{i - 1}_\mathfrak m(M/xM) \to H^i_\mathfrak m(M) \xrightarrow{x}
H^i_\mathfrak m(M)
$$
issue de la suite exacte courte $0 \to M \xrightarrow{x} M \to M/xM \to 0$.
Il s'ensuit que le sous-module de $x$-torsion $H^i_\mathfrak m(M)[x]$
est quotient d'un module satisfaisant à la condition des chaînes descendantes,
et satisfait donc lui-même à cette condition. Ainsi, le sous-module
de $\mathfrak m$-torsion $H^i_\mathfrak m(M)[\mathfrak m]$ satisfait
à la condition des chaînes descendantes (et est donc de dimension finie
sur $A/\mathfrak m$). On en conclut que le module de torsion par les puissances de $\mathfrak m$
$H^i_\mathfrak m(M)$ satisfait à la condition des chaînes descendantes
d'après Complexes dualisants, lemme
\ref{dualizing-lemma-describe-categories}.

\medskip\noindent
L'assertion (2) découle de (1) par Cohomologie locale,
lemme \ref{local-cohomology-lemma-finiteness-pushforwards-and-H1-local}.
\end{proof}

\begin{lemma}
\label{lemma-ML-local}
Soit $(A, \mathfrak m)$ un anneau local noethérien.
\begin{enumerate}
\item Soit $(M_n)$ un système projectif de $A$-modules de type fini. Alors le
système projectif $H^i_\mathfrak m(M_n)$ satisfait à la condition de Mittag-Leffler
pour tout $i$.
\item Soit $U = \Spec(A) \setminus \{\mathfrak m\}$ le
spectre épointé de $A$.
Soit $\mathcal{F}_n$ un système projectif de
$\mathcal{O}_U$-modules cohérents.
Alors le système projectif $H^i(U, \mathcal{F}_n)$
satisfait à la condition de Mittag-Leffler pour $i > 0$.
\end{enumerate}
\end{lemma}

\begin{proof}
Cela résulte immédiatement du lemme \ref{lemma-descending-chain}.
\end{proof}

\begin{lemma}
\label{lemma-terrific}
Soit $(A, \mathfrak m)$ un anneau local noethérien.
Soit $(M_n)$ un système projectif de $A$-modules de type fini.
Soit $M \to \lim M_n$ une application, où $M$ est un $A$-module de type fini,
telle que, pour un certain $i$, l'application
$H^i_\mathfrak m(M) \to \lim H^i_\mathfrak m(M_n)$
soit un isomorphisme.
Alors le système projectif $H^i_\mathfrak m(M_n)$
est essentiellement constant, de valeur $H^i_\mathfrak m(M)$.
\end{lemma}

\begin{proof}
D'après le lemme \ref{lemma-ML-local}, le système projectif $H^i_\mathfrak m(M_n)$
satisfait à la condition de Mittag-Leffler. Soit $E_n \subset H^i_\mathfrak m(M_n)$
l'image de $H^i_\mathfrak m(M_{n'})$ pour $n' \gg n$.
Alors $(E_n)$ est un système projectif à applications de transition surjectives
et $H^i_\mathfrak m(M) = \lim E_n$. Puisque $H^i_\mathfrak m(M)$
satisfait à la condition des chaînes descendantes d'après le
lemme \ref{lemma-descending-chain},
il ne peut y avoir qu'un nombre fini de noyaux non triviaux
parmi ceux des surjections $H^i_\mathfrak m(M) \to E_n$.
Ainsi $E_n \to E_{n - 1}$ est un isomorphisme pour tout $n \gg 0$,
comme voulu.
\end{proof}

\begin{lemma}
\label{lemma-local-cohomology-derived-completion}
Soit $(A, \mathfrak m)$ un anneau local noethérien.
Soient $I \subset A$ un idéal et $M$ un $A$-module de type fini.
Alors
$$
H^i(R\Gamma_\mathfrak m(M)^\wedge) = \lim H^i_\mathfrak m(M/I^nM)
$$
pour tout $i$, où $R\Gamma_\mathfrak m(M)^\wedge$ désigne
le complété $I$-adique dérivé.
\end{lemma}

\begin{proof}
Appliquons Complexes dualisants, lemme \ref{dualizing-lemma-completion-local},
et le lemme \ref{lemma-ML-local} pour obtenir l'annulation des termes $R^1\lim$.
\end{proof}








\section{Complétion dérivée sur un site annelé}
\label{section-derived-completion}

\noindent
Nous recommandons au lecteur de sauter cette section à la première lecture.

\medskip\noindent
La version algébrique de ce qui suit figure dans
Compléments d'algèbre, section \ref{more-algebra-section-derived-completion}.
Soit $\mathcal{O}$ un faisceau d'anneaux sur un site $\mathcal{C}$.
Soit $f$ une section globale de $\mathcal{O}$. Notons
$\mathcal{O}_f$ le faisceau associé au préfaisceau de localisés
$U \mapsto \mathcal{O}(U)_f$.

\begin{lemma}
\label{lemma-map-twice-localize}
Soit $(\mathcal{C}, \mathcal{O})$ un site annelé. Soit $f$ une section
globale de $\mathcal{O}$.
\begin{enumerate}
\item Pour $L, N \in D(\mathcal{O}_f)$, on a
$R\SheafHom_\mathcal{O}(L, N) = R\SheafHom_{\mathcal{O}_f}(L, N)$.
En particulier, les deux structures de $\mathcal{O}_f$-module sur
$R\SheafHom_\mathcal{O}(L, N)$ coïncident.
\item Pour $K \in D(\mathcal{O})$ et
$L \in D(\mathcal{O}_f)$, on a
$$
R\SheafHom_\mathcal{O}(L, K) =
R\SheafHom_{\mathcal{O}_f}(L, R\SheafHom_\mathcal{O}(\mathcal{O}_f, K))
$$
En particulier
$R\SheafHom_\mathcal{O}(\mathcal{O}_f,
R\SheafHom_\mathcal{O}(\mathcal{O}_f, K)) =
R\SheafHom_\mathcal{O}(\mathcal{O}_f, K)$.
\item Si $g$ est une seconde section
globale de $\mathcal{O}$, alors
$$
R\SheafHom_\mathcal{O}(\mathcal{O}_f, R\SheafHom_\mathcal{O}(\mathcal{O}_g, K))
= R\SheafHom_\mathcal{O}(\mathcal{O}_{gf}, K).
$$
\end{enumerate}
\end{lemma}

\begin{proof}
Démonstration de (1). Soit $\mathcal{J}^\bullet$ un complexe K-injectif de
$\mathcal{O}_f$-modules représentant $N$. D'après Cohomologie sur les sites, lemme
\ref{sites-cohomology-lemma-K-injective-flat}, il s'ensuit que
$\mathcal{J}^\bullet$ est aussi un complexe K-injectif de
$\mathcal{O}$-modules. Soit $\mathcal{F}^\bullet$ un complexe de
$\mathcal{O}_f$-modules représentant $L$. Alors
$$
R\SheafHom_\mathcal{O}(L, N) =
R\SheafHom_\mathcal{O}(\mathcal{F}^\bullet, \mathcal{J}^\bullet) =
R\SheafHom_{\mathcal{O}_f}(\mathcal{F}^\bullet, \mathcal{J}^\bullet)
$$
d'après
Modules sur les sites, lemme \ref{sites-modules-lemma-epimorphism-modules},
puisque $\mathcal{J}^\bullet$ est un complexe K-injectif de $\mathcal{O}$-modules
et de $\mathcal{O}_f$-modules.

\medskip\noindent
Démonstration de (2). Soit $\mathcal{I}^\bullet$ un complexe K-injectif de
$\mathcal{O}$-modules représentant $K$.
Alors $R\SheafHom_\mathcal{O}(\mathcal{O}_f, K)$ est représenté par
$\SheafHom_\mathcal{O}(\mathcal{O}_f, \mathcal{I}^\bullet)$, qui est
un complexe K-injectif de $\mathcal{O}_f$-modules et de
$\mathcal{O}$-modules d'après
Cohomologie sur les sites, lemmes \ref{sites-cohomology-lemma-hom-K-injective} et
\ref{sites-cohomology-lemma-K-injective-flat}.
Soit $\mathcal{F}^\bullet$ un complexe de $\mathcal{O}_f$-modules
représentant $L$. Alors
$$
R\SheafHom_\mathcal{O}(L, K) =
R\SheafHom_\mathcal{O}(\mathcal{F}^\bullet, \mathcal{I}^\bullet) =
R\SheafHom_{\mathcal{O}_f}(\mathcal{F}^\bullet,
\SheafHom_\mathcal{O}(\mathcal{O}_f, \mathcal{I}^\bullet))
$$
d'après Modules sur les sites, lemme \ref{sites-modules-lemma-adjoint-hom-restrict},
et puisque $\SheafHom_\mathcal{O}(\mathcal{O}_f, \mathcal{I}^\bullet)$ est un
complexe K-injectif de $\mathcal{O}_f$-modules.

\medskip\noindent
Démonstration de (3). Cela résulte du fait que
$R\SheafHom_\mathcal{O}(\mathcal{O}_g, \mathcal{I}^\bullet)$
est K-injectif comme complexe de $\mathcal{O}$-modules et du fait que
$\SheafHom_\mathcal{O}(\mathcal{O}_f,
\SheafHom_\mathcal{O}(\mathcal{O}_g, \mathcal{H})) = 
\SheafHom_\mathcal{O}(\mathcal{O}_{gf}, \mathcal{H})$
pour tout faisceau de $\mathcal{O}$-modules $\mathcal{H}$.
\end{proof}

\noindent
Soit $K \in D(\mathcal{O})$. Notons
$T(K, f)$ une limite dérivée (Catégories dérivées, définition
\ref{derived-definition-derived-limit}) du système projectif
$$
\ldots \to K \xrightarrow{f} K \xrightarrow{f} K
$$
dans $D(\mathcal{O})$.

\begin{lemma}
\label{lemma-hom-from-Af}
Soit $(\mathcal{C}, \mathcal{O})$ un site annelé. Soit $f$ une section
globale de $\mathcal{O}$. Soit $K \in D(\mathcal{O})$.
Les conditions suivantes sont équivalentes :
\begin{enumerate}
\item $R\SheafHom_\mathcal{O}(\mathcal{O}_f, K) = 0$,
\item $R\SheafHom_\mathcal{O}(L, K) = 0$ pour tout $L$ dans $D(\mathcal{O}_f)$,
\item $T(K, f) = 0$.
\end{enumerate}
\end{lemma}

\begin{proof}
Il est clair que (2) implique (1). L'implication (1) $\Rightarrow$ (2)
résulte du lemme \ref{lemma-map-twice-localize}.
Une résolution libre du $\mathcal{O}$-module $\mathcal{O}_f$ est donnée par
$$
0 \to \bigoplus\nolimits_{n \in \mathbf{N}} \mathcal{O} \to
\bigoplus\nolimits_{n \in \mathbf{N}} \mathcal{O}
\to \mathcal{O}_f \to 0
$$
où la première application envoie une section locale $(x_0, x_1, \ldots)$ sur
$(x_0, x_1 - fx_0, x_2 - fx_1, \ldots)$, et la seconde envoie
$(x_0, x_1, \ldots)$ sur $x_0 + x_1/f + x_2/f^2 + \ldots$.
En appliquant $\SheafHom_\mathcal{O}(-, \mathcal{I}^\bullet)$,
où $\mathcal{I}^\bullet$ est un complexe K-injectif de $\mathcal{O}$-modules
représentant $K$, on obtient une suite exacte courte de complexes
$$
0 \to \SheafHom_\mathcal{O}(\mathcal{O}_f, \mathcal{I}^\bullet) \to
\prod \mathcal{I}^\bullet \to \prod \mathcal{I}^\bullet \to 0
$$
puisque $\mathcal{I}^n$ est un $\mathcal{O}$-module injectif.
Les produits sont des produits dans $D(\mathcal{O})$ ; voir
Injectifs, lemme \ref{injectives-lemma-derived-products}.
Cela signifie que l'objet $T(K, f)$ est un représentant de
$R\SheafHom_\mathcal{O}(\mathcal{O}_f, K)$ dans $D(\mathcal{O})$.
D'où l'équivalence de (1) et (3).
\end{proof}

\begin{lemma}
\label{lemma-ideal-of-elements-complete-wrt}
Soit $(\mathcal{C}, \mathcal{O})$ un site annelé. Soit $K \in D(\mathcal{O})$.
La règle qui associe à $U$ l'ensemble $\mathcal{I}(U)$
des sections $f \in \mathcal{O}(U)$ telles que $T(K|_U, f) = 0$
est un faisceau d'idéaux de $\mathcal{O}$.
\end{lemma}

\begin{proof}
Nous utiliserons les résultats du lemme \ref{lemma-hom-from-Af} sans autre
mention. Si $f \in \mathcal{I}(U)$ et $g \in \mathcal{O}(U)$, alors
$\mathcal{O}_{U, gf}$ est un $\mathcal{O}_{U, f}$-module,
donc $R\SheafHom_\mathcal{O}(\mathcal{O}_{U, gf}, K|_U) = 0$, d'où
$gf \in \mathcal{I}(U)$. Supposons $f, g \in \mathcal{O}(U)$.
On a alors une suite exacte courte
$$
0 \to \mathcal{O}_{U, f + g} \to
\mathcal{O}_{U, f(f + g)} \oplus \mathcal{O}_{U, g(f + g)} \to
\mathcal{O}_{U, gf(f + g)} \to 0
$$
puisque $f, g$ engendrent l'idéal unité de $\mathcal{O}(U)_{f + g}$.
Cela résulte
d'Algèbre, lemme \ref{algebra-lemma-standard-covering},
et du fait immédiat que la dernière flèche est surjective.
Puisque $R\SheafHom_\mathcal{O}( - , K|_U)$ est un foncteur exact
entre catégories triangulées, l'annulation de
$R\SheafHom_{\mathcal{O}_U}(\mathcal{O}_{U, f(f + g)}, K|_U)$,
$R\SheafHom_{\mathcal{O}_U}(\mathcal{O}_{U, g(f + g)}, K|_U)$ et
$R\SheafHom_{\mathcal{O}_U}(\mathcal{O}_{U, gf(f + g)}, K|_U)$,
entraîne l'annulation de
$R\SheafHom_{\mathcal{O}_U}(\mathcal{O}_{U, f + g}, K|_U)$.
Nous omettons la vérification de la condition de faisceau.
\end{proof}

\noindent
On peut donner la définition suivante pour tout site annelé.

\begin{definition}
\label{definition-derived-complete}
Soit $(\mathcal{C}, \mathcal{O})$ un site annelé.
Soit $\mathcal{I} \subset \mathcal{O}$ un faisceau d'idéaux.
Soit $K \in D(\mathcal{O})$. On dit que $K$ est
{\it complet au sens dérivé relativement à $\mathcal{I}$}
si, pour tout objet $U$ de $\mathcal{C}$ et tout $f \in \mathcal{I}(U)$,
l'objet $T(K|_U, f)$ de $D(\mathcal{O}_U)$ est nul.
\end{definition}

\noindent
Il est clair que la sous-catégorie pleine
$D_{comp}(\mathcal{O}) = D_{comp}(\mathcal{O}, \mathcal{I}) \subset
D(\mathcal{O})$ formée des objets complets au sens dérivé
est une sous-catégorie triangulée saturée ; voir
Catégories dérivées, définitions
\ref{derived-definition-triangulated-subcategory} et
\ref{derived-definition-saturated}. Cette sous-catégorie est stable
par produits et limites homotopiques dans $D(\mathcal{O})$.
En revanche, elle n'est en général pas stable par sommes directes dénombrables.

\begin{lemma}
\label{lemma-derived-complete-internal-hom}
Soit $(\mathcal{C}, \mathcal{O})$ un site annelé.
Soit $\mathcal{I} \subset \mathcal{O}$ un faisceau d'idéaux.
Si $K \in D(\mathcal{O})$ et $L \in D_{comp}(\mathcal{O})$, alors
$R\SheafHom_\mathcal{O}(K, L) \in D_{comp}(\mathcal{O})$.
\end{lemma}

\begin{proof}
Soient $U$ un objet de $\mathcal{C}$ et $f \in \mathcal{I}(U)$.
Rappelons que
$$
\Hom_{D(\mathcal{O}_U)}(\mathcal{O}_{U, f}, R\SheafHom_\mathcal{O}(K, L)|_U)
=
\Hom_{D(\mathcal{O}_U)}(
K|_U \otimes_{\mathcal{O}_U}^\mathbf{L} \mathcal{O}_{U, f}, L|_U)
$$
d'après Cohomologie sur les sites, lemme \ref{sites-cohomology-lemma-internal-hom}.
Le membre de droite est nul par le lemme \ref{lemma-hom-from-Af}
et la relation entre Hom interne et Hom ; voir
Cohomologie sur les sites, lemme \ref{sites-cohomology-lemma-section-RHom-over-U}.
La même annulation vaut pour tout $U'/U$. Ainsi, l'objet
$R\SheafHom_{\mathcal{O}_U}(\mathcal{O}_{U, f},
R\SheafHom_\mathcal{O}(K, L)|_U)$ de $D(\mathcal{O}_U)$ a un faisceau de cohomologie
de degré $0$ nul (loc. cit.). Il en va de même des autres
faisceaux de cohomologie, autrement dit $R\SheafHom_{\mathcal{O}_U}(\mathcal{O}_{U, f},
R\SheafHom_\mathcal{O}(K, L)|_U)$ est nul dans $D(\mathcal{O}_U)$.
On conclut par le lemme \ref{lemma-hom-from-Af}.
\end{proof}

\begin{lemma}
\label{lemma-restriction-derived-complete}
Soit $\mathcal{C}$ un site. Soit $\mathcal{O} \to \mathcal{O}'$
un homomorphisme de faisceaux d'anneaux. Soit $\mathcal{I} \subset \mathcal{O}$
un faisceau d'idéaux. L'image réciproque de $D_{comp}(\mathcal{O}, \mathcal{I})$
par le foncteur de restriction $D(\mathcal{O}') \to D(\mathcal{O})$ est
$D_{comp}(\mathcal{O}', \mathcal{I}\mathcal{O}')$.
\end{lemma}

\begin{proof}
En utilisant le lemme \ref{lemma-ideal-of-elements-complete-wrt},
on voit que $K' \in D(\mathcal{O}')$ appartient à
$D_{comp}(\mathcal{O}', \mathcal{I}\mathcal{O}')$
si et seulement si $T(K'|_U, f)$ est nul pour toute section locale
$f \in \mathcal{I}(U)$. Remarquons que les faisceaux de cohomologie de
$T(K'|_U, f)$ se calculent dans la catégorie des faisceaux abéliens ;
il importe donc peu que l'on considère $f$ comme une section de
$\mathcal{O}$ ou que l'on prenne l'image de $f$ comme section de $\mathcal{O}'$.
Le lemme résulte immédiatement de cette remarque et de la
définition des objets complets au sens dérivé.
\end{proof}

\begin{lemma}
\label{lemma-pushforward-derived-complete}
Soit $f : (\Sh(\mathcal{D}), \mathcal{O}') \to (\Sh(\mathcal{C}), \mathcal{O})$
un morphisme de topos annelés. Soient $\mathcal{I} \subset \mathcal{O}$
et $\mathcal{I}' \subset \mathcal{O}'$ des faisceaux d'idéaux tels
que $f^\sharp$ envoie $f^{-1}\mathcal{I}$ dans $\mathcal{I}'$.
Alors $Rf_*$ envoie $D_{comp}(\mathcal{O}', \mathcal{I}')$
dans $D_{comp}(\mathcal{O}, \mathcal{I})$.
\end{lemma}

\begin{proof}
On peut supposer que $f$ est donné par un morphisme de sites annelés correspondant
à un foncteur continu $\mathcal{C} \to \mathcal{D}$
(Modules sur les sites, lemme
\ref{sites-modules-lemma-morphism-ringed-topoi-comes-from-morphism-ringed-sites}
).
Soit $U$ un objet de $\mathcal{C}$ et soit $g$ une section de
$\mathcal{I}$ sur $U$. Il faut montrer que
$\Hom_{D(\mathcal{O}_U)}(\mathcal{O}_{U, g}, Rf_*K|_U) = 0$
lorsque $K$ est complet au sens dérivé relativement à $\mathcal{I}'$.
En effet, d'après Cohomologie sur les sites, lemme
\ref{sites-cohomology-lemma-section-RHom-over-U}
cette annulation, appliquée à tous les objets au-dessus de $U$ et à tous les décalages de $K$,
entraînera que $R\SheafHom_{\mathcal{O}_U}(\mathcal{O}_{U, g}, Rf_*K|_U)$
est nul, donc que $T(Rf_*K|_U, g)$ est nul
(lemme \ref{lemma-hom-from-Af}), ce qui est l'assertion à démontrer
(définition \ref{definition-derived-complete}).
Soit $V$ dans $\mathcal{D}$ l'image de $U$. Alors
$$
\Hom_{D(\mathcal{O}_U)}(\mathcal{O}_{U, g}, Rf_*K|_U) =
\Hom_{D(\mathcal{O}'_V)}(\mathcal{O}'_{V, g'}, K|_V) = 0
$$
où $g' = f^\sharp(g) \in \mathcal{I}'(V)$. La seconde égalité
résulte de ce que $K$ est complet au sens dérivé, et la première de ce que
l'image inverse dérivée de $\mathcal{O}_{U, g}$ est $\mathcal{O}'_{V, g'}$,
ainsi que de
Cohomologie sur les sites, lemme \ref{sites-cohomology-lemma-adjoint}.
\end{proof}

\noindent
Le lemme suivant donne le cas le plus simple où l'on dispose de la complétion dérivée.

\begin{lemma}
\label{lemma-derived-completion}
Soit $(\mathcal{C}, \mathcal{O})$ un site annelé. Soient $f_1, \ldots, f_r$
des sections globales de $\mathcal{O}$. Soit $\mathcal{I} \subset \mathcal{O}$
le faisceau d'idéaux engendré par $f_1, \ldots, f_r$.
Alors le foncteur d'inclusion $D_{comp}(\mathcal{O}) \to D(\mathcal{O})$
admet un adjoint à gauche ; autrement dit, pour tout objet $K$ de $D(\mathcal{O})$,
il existe une application $K \to K^\wedge$, avec $K^\wedge$ dans $D_{comp}(\mathcal{O})$,
telle que l'application
$$
\Hom_{D(\mathcal{O})}(K^\wedge, E) \longrightarrow \Hom_{D(\mathcal{O})}(K, E)
$$
soit bijective dès que $E$ appartient à $D_{comp}(\mathcal{O})$. En fait,
on a
$$
K^\wedge =
R\SheafHom_\mathcal{O}
(\mathcal{O} \to \prod\nolimits_{i_0} \mathcal{O}_{f_{i_0}} \to
\prod\nolimits_{i_0 < i_1} \mathcal{O}_{f_{i_0}f_{i_1}} \to
\ldots \to \mathcal{O}_{f_1\ldots f_r}, K)
$$
fonctoriellement en $K$.
\end{lemma}

\begin{proof}
Définissons $K^\wedge$ par la dernière formule affichée du lemme.
On a un morphisme de complexes
$$
(\mathcal{O} \to \prod\nolimits_{i_0} \mathcal{O}_{f_{i_0}} \to
\prod\nolimits_{i_0 < i_1} \mathcal{O}_{f_{i_0}f_{i_1}} \to
\ldots \to \mathcal{O}_{f_1\ldots f_r}) \longrightarrow \mathcal{O}
$$
qui induit une application $K \to K^\wedge$. Il suffit de montrer que
$K^\wedge$ est complet au sens dérivé et que $K \to K^\wedge$ est un
isomorphisme si $K$ est complet au sens dérivé.

\medskip\noindent
Soit $f$ une section globale de $\mathcal{O}$.
D'après le lemme \ref{lemma-map-twice-localize}, l'objet
$R\SheafHom_\mathcal{O}(\mathcal{O}_f, K^\wedge)$
est égal à
$$
R\SheafHom_\mathcal{O}(
(\mathcal{O}_f \to \prod\nolimits_{i_0} \mathcal{O}_{ff_{i_0}} \to
\prod\nolimits_{i_0 < i_1} \mathcal{O}_{ff_{i_0}f_{i_1}} \to
\ldots \to \mathcal{O}_{ff_1\ldots f_r}), K)
$$
Si $f = f_i$ pour un certain $i$, alors $f_1, \ldots, f_r$ engendrent
l'idéal unité de $\mathcal{O}_f$ ; le complexe de
{\v C}ech alterné augmenté
$$
\mathcal{O}_f \to \prod\nolimits_{i_0} \mathcal{O}_{ff_{i_0}} \to
\prod\nolimits_{i_0 < i_1} \mathcal{O}_{ff_{i_0}f_{i_1}} \to
\ldots \to \mathcal{O}_{ff_1\ldots f_r}
$$
est donc nul (et même homotope à zéro). On voit ainsi que $K^\wedge$
est complet au sens dérivé.

\medskip\noindent
Si $K$ est complet au sens dérivé, alors $R\SheafHom_\mathcal{O}(\mathcal{O}_f, K)$
est nul pour tout $f = f_{i_0} \ldots f_{i_p}$, $p \geq 0$. Ainsi
$K \to K^\wedge$ est un isomorphisme dans $D(\mathcal{O})$.
\end{proof}

\noindent
Expliquons maintenant en quoi la complétion dérivée est une complétion.

\begin{lemma}
\label{lemma-derived-completion-koszul}
Soit $(\mathcal{C}, \mathcal{O})$ un site annelé. Soient $f_1, \ldots, f_r$
des sections globales de $\mathcal{O}$. Soit $\mathcal{I} \subset \mathcal{O}$
le faisceau d'idéaux engendré par $f_1, \ldots, f_r$. Soit $K \in D(\mathcal{O})$.
Le complété dérivé $K^\wedge$ du lemme \ref{lemma-derived-completion}
est donné par la formule
$$
K^\wedge = R\lim K \otimes^\mathbf{L}_\mathcal{O} K_n
$$
où $K_n = K(\mathcal{O}, f_1^n, \ldots, f_r^n)$
est le complexe de Koszul de $f_1^n, \ldots, f_r^n$ sur $\mathcal{O}$.
\end{lemma}

\begin{proof}
Dans Compléments d'algèbre, lemme
\ref{more-algebra-lemma-extended-alternating-Cech-is-colimit-koszul}
nous avons vu que le complexe de {\v C}ech alterné augmenté
$$
\mathcal{O} \to \prod\nolimits_{i_0} \mathcal{O}_{f_{i_0}} \to
\prod\nolimits_{i_0 < i_1} \mathcal{O}_{f_{i_0}f_{i_1}} \to
\ldots \to \mathcal{O}_{f_1\ldots f_r}
$$
est une limite inductive des complexes de Koszul
$K^n = K(\mathcal{O}, f_1^n, \ldots, f_r^n)$ placés en
degrés $0, \ldots, r$. Notons que $K^n$ est un complexe de chaînes fini
de $\mathcal{O}$-modules libres de type fini, dont le dual est
$\SheafHom_\mathcal{O}(K^n, \mathcal{O}) = K_n$, où $K_n$ est le complexe de Koszul
de cochaînes placé en degrés $-r, \ldots, 0$ (comme d'habitude). D'après le
lemme \ref{lemma-derived-completion},
le foncteur $E \mapsto E^\wedge$ s'obtient en prenant
$R\SheafHom$ du complexe de {\v C}ech alterné augmenté dans $E$ :
$$
E^\wedge = R\SheafHom(\colim K^n, E)
$$
Cette expression est égale à $R\lim (E \otimes_\mathcal{O}^\mathbf{L} K_n)$
d'après
Cohomologie sur les sites, lemme \ref{sites-cohomology-lemma-colim-and-lim-of-duals}.
\end{proof}

\begin{lemma}
\label{lemma-all-rings}
Il existe une construction des données suivantes :
\begin{enumerate}
\item pour tout couple $(A, I)$ formé d'un anneau $A$ et d'un idéal
de type fini $I \subset A$, un complexe $K(A, I)$ de $A$-modules,
\item un morphisme $K(A, I) \to A$ de complexes de $A$-modules,
\item pour tout homomorphisme d'anneaux $A \to B$ et tout idéal de type fini $I \subset A$,
un morphisme de complexes $K(A, I) \to K(B, IB)$,
\end{enumerate}
telle que
\begin{enumerate}
\item[(a)] pour $A \to B$ et $I \subset A$ de type fini, le diagramme
$$
\xymatrix{
K(A, I) \ar[r] \ar[d] & A \ar[d] \\
K(B, IB) \ar[r] & B
}
$$
soit commutatif,
\item[(b)] pour $A \to B \to C$ et $I \subset A$ de type fini,
la composée des applications
$K(A, I) \to K(B, IB) \to K(C, IC)$ soit l'application $K(A, I) \to K(C, IC)$.
\item[(c)] pour $A \to B$ et un idéal de type fini $I \subset A$,
l'application induite $K(A, I) \otimes_A^\mathbf{L} B \to K(B, IB)$
soit un isomorphisme dans $D(B)$, et
\item[(d)] si $I = (f_1, \ldots, f_r) \subset A$, il existe un diagramme
commutatif
$$
\xymatrix{
(A \to \prod\nolimits_{i_0} A_{f_{i_0}} \to
\prod\nolimits_{i_0 < i_1} A_{f_{i_0}f_{i_1}} \to
\ldots \to A_{f_1\ldots f_r}) \ar[r] \ar[d] &  K(A, I) \ar[d] \\
A \ar[r]^1 & A
}
$$
dans $D(A)$ dont les flèches horizontales sont des isomorphismes.
\end{enumerate}
\end{lemma}

\begin{proof}
Soit $S$ l'ensemble des anneaux $A_0$ de la forme
$A_0 = \mathbf{Z}[x_1, \ldots, x_n]/J$.
Toute $\mathbf{Z}$-algèbre de type fini est isomorphe à
un élément de $S$. Soit $\mathcal{A}_0$ la catégorie dont les objets
sont les couples $(A_0, I_0)$, où $A_0 \in S$ et $I_0 \subset A_0$
est un idéal, et dont les morphismes $(A_0, I_0) \to (B_0, J_0)$ sont
les homomorphismes d'anneaux $\varphi : A_0 \to B_0$ tels que $J_0 = \varphi(I_0)B_0$.

\medskip\noindent
Supposons que l'on puisse construire $K(A_0, I_0) \to A_0$ fonctoriellement pour
les objets de $\mathcal{A}_0$, avec les propriétés (a), (b), (c) et (d).
On pose alors
$$
K(A, I) = \colim_{\varphi : (A_0, I_0) \to (A, I)} K(A_0, I_0)
$$
où la limite inductive porte sur les homomorphismes d'anneaux $\varphi : A_0 \to A$ tels
que $\varphi(I_0)A = I$, avec $(A_0, I_0)$ dans $\mathcal{A}_0$.
Un morphisme de $(A_0, I_0) \to (A, I)$ vers $(A_0', I_0') \to (A, I)$
est donné par un morphisme $(A_0, I_0) \to (A_0', I_0')$ dans $\mathcal{A}_0$
commutant aux applications vers $A$.
La catégorie de ces $(A_0, I_0) \to (A, I)$ est filtrante
(nous omettons les détails). De plus, $\colim_{\varphi : (A_0, I_0) \to (A, I)} A_0 = A$,
de sorte que $K(A, I)$ est un complexe de $A$-modules.
Enfin, étant donnés $\varphi : A \to B$ et $I \subset A$,
pour chaque $(A_0, I_0) \to (A, I)$ intervenant dans la limite inductive, la composée
$(A_0, I_0) \to (B, IB)$ intervient dans celle relative à $(B, IB)$.
On obtient ainsi une application entre les limites inductives. Les propriétés (a), (b), (c) et (d)
en découlent aussitôt, compte tenu des propriétés correspondantes
des complexes $K(A_0, I_0)$.

\medskip\noindent
Munissons $\mathcal{C}_0 = \mathcal{A}_0^{opp}$ de la topologie chaotique.
Munissons $\mathcal{C}_0$ du faisceau d'anneaux
$\mathcal{O} : (A, I) \mapsto A$. Les idéaux $I$ définissent ensemble un
faisceau d'idéaux $\mathcal{I} \subset \mathcal{O}$.
Choisissons une résolution injective $\mathcal{O} \to \mathcal{J}^\bullet$.
Considérons l'objet
$$
\mathcal{F}^\bullet = \bigcup\nolimits_n \mathcal{J}^\bullet[\mathcal{I}^n]
$$
Soit $U = (A, I) \in \Ob(\mathcal{C}_0)$.
Puisque la topologie de $\mathcal{C}_0$ est chaotique, la valeur
$\mathcal{J}^\bullet(U)$ est une résolution de $A$ par des
$A$-modules injectifs. La valeur $\mathcal{F}^\bullet(U)$ est donc un
objet de $D(A)$ représentant l'image de $R\Gamma_I(A)$ dans $D(A)$ ; voir
Complexes dualisants, section \ref{dualizing-section-local-cohomology}.
Choisissons un complexe de $\mathcal{O}$-modules $\mathcal{K}^\bullet$
et un diagramme commutatif
$$
\xymatrix{
\mathcal{O} \ar[r] & \mathcal{J}^\bullet \\
\mathcal{K}^\bullet \ar[r] \ar[u] & \mathcal{F}^\bullet \ar[u]
}
$$
dont les flèches horizontales sont des quasi-isomorphismes. Cela est possible
par la construction de la catégorie dérivée $D(\mathcal{O})$.
Posons $K(A, I) = \mathcal{K}^\bullet(U)$, où $U = (A, I)$.
Les propriétés (a) et (b) sont claires, et les propriétés (c) et (d)
résultent de Complexes dualisants, lemmes
\ref{dualizing-lemma-compute-local-cohomology-noetherian} et
\ref{dualizing-lemma-local-cohomology-change-rings}.
\end{proof}

\begin{lemma}
\label{lemma-global-extended-cech-complex}
Soit $(\mathcal{C}, \mathcal{O})$ un site annelé. Soit
$\mathcal{I} \subset \mathcal{O}$ un faisceau d'idéaux de type fini.
Il existe une application $K \to \mathcal{O}$ dans $D(\mathcal{O})$
telle que, pour tout $U \in \Ob(\mathcal{C})$ tel que
$\mathcal{I}|_U$ soit engendré par $f_1, \ldots, f_r \in \mathcal{I}(U)$,
il existe un isomorphisme
$$
(\mathcal{O}_U \to \prod\nolimits_{i_0} \mathcal{O}_{U, f_{i_0}} \to
\prod\nolimits_{i_0 < i_1} \mathcal{O}_{U, f_{i_0}f_{i_1}} \to
\ldots \to \mathcal{O}_{U, f_1\ldots f_r}) \longrightarrow K|_U
$$
compatible avec les applications vers $\mathcal{O}_U$.
\end{lemma}

\begin{proof}
Soit $\mathcal{C}' \subset \mathcal{C}$ la sous-catégorie pleine
des objets $U$ tels que $\mathcal{I}|_U$ soit engendré par
un nombre fini de sections. Alors $\mathcal{C}' \to \mathcal{C}$
est un foncteur cocontinu spécial
(Sites, définition \ref{sites-definition-special-cocontinuous-functor}).
Il suffit donc de travailler avec $\mathcal{C}'$ ; voir
Sites, lemme \ref{sites-lemma-equivalence}.
Autrement dit, on peut supposer que, pour tout
objet $U$ de $\mathcal{C}$, il existe un idéal de type fini
$I \subset \mathcal{I}(U)$ tel que
$\mathcal{I}|_U = \Im(I \otimes \mathcal{O}_U \to \mathcal{O}_U)$.
Nous dirons que $I$ engendre $\mathcal{I}|_U$.
Attention : on ne sait pas si $\mathcal{I}(U)$ est un idéal de type fini
de $\mathcal{O}(U)$.

\medskip\noindent
Soient $U$ un objet et $I \subset \mathcal{O}(U)$ un idéal de type
fini qui engendre $\mathcal{I}|_U$.
Sur la catégorie $\mathcal{C}/U$, considérons le complexe de préfaisceaux
$$
K_{U, I}^\bullet : U'/U \longmapsto K(\mathcal{O}(U'), I\mathcal{O}(U'))
$$
où $K(-, -)$ est comme dans le lemme \ref{lemma-all-rings}.
Montrons que son faisceau associé est indépendant
du choix de $I$. En effet, si $I' \subset \mathcal{O}(U)$
est un idéal de type fini engendrant lui aussi $\mathcal{I}|_U$,
il existe un recouvrement $\{U_j \to U\}$ tel que
$I\mathcal{O}(U_j) = I'\mathcal{O}(U_j)$. (Indication : cela tient au fait que
$I$ et $I'$ sont tous deux de type fini et engendrent $\mathcal{I}|_U$.)
Ainsi $K_{U, I}^\bullet$ et $K_{U, I'}^\bullet$ sont {\it identiques}
sur tout objet au-dessus de l'un des $U_j$. L'assertion
sur les faisceaux associés s'ensuit. Notons $K_U^\bullet$ leur valeur commune.

\medskip\noindent
L'indépendance par rapport au choix de $I$ montre également que
$K_U^\bullet|_{\mathcal{C}/U'} = K_{U'}^\bullet$
dès que l'on se donne un morphisme
$U' \to U$, et donc un morphisme de localisation
$\mathcal{C}/U' \to \mathcal{C}/U$. Ainsi, les complexes
$K_U^\bullet$ se recollent en un unique complexe bien défini $K^\bullet$
de $\mathcal{O}$-modules. L'existence de l'application $K^\bullet \to \mathcal{O}$
et du quasi-isomorphisme du lemme résulte immédiatement
des propriétés correspondantes des complexes $K(-, -)$ dans le
lemme \ref{lemma-all-rings}.
\end{proof}

\begin{proposition}
\label{proposition-derived-completion}
Soit $(\mathcal{C}, \mathcal{O})$ un site annelé.
Soit $\mathcal{I} \subset \mathcal{O}$ un faisceau d'idéaux
de type fini. Le foncteur d'inclusion
$D_{comp}(\mathcal{O}) \to D(\mathcal{O})$ admet un adjoint à gauche.
\end{proposition}

\begin{proof}
Soit $K \to \mathcal{O}$ dans $D(\mathcal{O})$ l'application construite dans le
lemme \ref{lemma-global-extended-cech-complex}. Soit $E \in D(\mathcal{O})$.
Alors $E^\wedge = R\SheafHom(K, E)$, muni de l'application $E \to E^\wedge$,
convient. En effet, localement sur le site $\mathcal{C}$, on
retrouve l'adjoint du lemme \ref{lemma-derived-completion}.
Cela montre que $E^\wedge$ est toujours complet au sens dérivé et que
$E \to E^\wedge$ est un isomorphisme si $E$ est complet au sens dérivé.
\end{proof}

\begin{remark}[Comparaison avec la complétion]
\label{remark-compare-with-completion}
Soit $(\mathcal{C}, \mathcal{O})$ un site annelé.
Soit $\mathcal{I} \subset \mathcal{O}$ un faisceau d'idéaux
de type fini. Soit $K \mapsto K^\wedge$ le foncteur de complétion dérivée
de la proposition \ref{proposition-derived-completion}.
Pour tout $n \geq 1$, l'objet
$K \otimes_\mathcal{O}^\mathbf{L} \mathcal{O}/\mathcal{I}^n$
est complet au sens dérivé, car il est annulé par des puissances des
sections locales de $\mathcal{I}$. On a donc une factorisation canonique
$$
K \to K^\wedge \to K \otimes_\mathcal{O}^\mathbf{L} \mathcal{O}/\mathcal{I}^n
$$
de l'application canonique
$K \to K \otimes_\mathcal{O}^\mathbf{L} \mathcal{O}/\mathcal{I}^n$.
Ces applications sont compatibles lorsque $n$ varie, et l'on obtient une application de comparaison
$$
K^\wedge
\longrightarrow
R\lim \left(K \otimes_\mathcal{O}^\mathbf{L} \mathcal{O}/\mathcal{I}^n\right)
$$
Le membre de droite s'interprète plus manifestement comme une sorte de complété.
En général, cette application de comparaison n'est pas un isomorphisme.
\end{remark}

\begin{remark}[Localisation et complétion dérivée]
\label{remark-localization-and-completion}
Soit $(\mathcal{C}, \mathcal{O})$ un site annelé.
Soit $\mathcal{I} \subset \mathcal{O}$ un faisceau d'idéaux
de type fini. Soit $K \mapsto K^\wedge$ le foncteur de complétion dérivée
de la proposition \ref{proposition-derived-completion}. La construction
donnée dans la démonstration de la proposition montre que $K^\wedge|_U$
est le complété dérivé de $K|_U$ pour tout $U \in \Ob(\mathcal{C})$.
On peut aussi le démontrer comme suit. La définition
des objets complets au sens dérivé montre que $K^\wedge|_U$ est complet au sens dérivé.
On obtient ainsi une application canonique $a : (K|_U)^\wedge \to K^\wedge|_U$.
D'autre part, si $E$ est un objet complet au sens dérivé de
$D(\mathcal{O}_U)$, alors $Rj_*E$ est un objet complet au sens dérivé de
$D(\mathcal{O})$ d'après le lemme \ref{lemma-pushforward-derived-complete}.
Ici $j$ est le morphisme de localisation
(Modules sur les sites, section \ref{sites-modules-section-localize}).
On obtient donc aussi une application
canonique $b : K^\wedge \to Rj_*((K|_U)^\wedge)$. Nous omettons la vérification formelle
du fait que l'adjoint de $b$ est l'inverse de $a$.
\end{remark}

\begin{remark}[Produit tensoriel complété]
\label{remark-completed-tensor-product}
Soit $(\mathcal{C}, \mathcal{O})$ un site annelé. Soit
$\mathcal{I} \subset \mathcal{O}$ un faisceau d'idéaux de type fini.
Notons $K \mapsto K^\wedge$ l'adjoint de la
proposition \ref{proposition-derived-completion}.
Posons alors
$$
K \otimes^\wedge_\mathcal{O} L = (K \otimes_\mathcal{O}^\mathbf{L} L)^\wedge
$$
Ce {\it produit tensoriel complété} définit un foncteur
$D_{comp}(\mathcal{O}) \times D_{comp}(\mathcal{O}) \to D_{comp}(\mathcal{O})$
tel que l'on ait
$$
\Hom_{D_{comp}(\mathcal{O})}(K, R\SheafHom_\mathcal{O}(L, M))
=
\Hom_{D_{comp}(\mathcal{O})}(K \otimes_\mathcal{O}^\wedge L, M)
$$
pour $K, L, M \in D_{comp}(\mathcal{O})$. Notons que
$R\SheafHom_\mathcal{O}(L, M) \in D_{comp}(\mathcal{O})$ d'après le
lemme \ref{lemma-derived-complete-internal-hom}.
\end{remark}

\begin{lemma}
\label{lemma-map-identifies-koszul-and-cech-complexes}
Soit $\mathcal{C}$ un site.
Supposons que $\varphi : \mathcal{O} \to \mathcal{O}'$ soit un homomorphisme plat
de faisceaux d'anneaux. Soient $f_1, \ldots, f_r$ des sections globales
de $\mathcal{O}$ telles que $\mathcal{O}/(f_1, \ldots, f_r) \cong
\mathcal{O}'/(f_1, \ldots, f_r)\mathcal{O}'$.
Alors le morphisme de complexes de {\v C}ech alternés augmentés
$$
\xymatrix{
\mathcal{O} \to
\prod_{i_0} \mathcal{O}_{f_{i_0}} \to
\prod_{i_0 < i_1} \mathcal{O}_{f_{i_0}f_{i_1}} \to \ldots \to
\mathcal{O}_{f_1\ldots f_r} \ar[d] \\
\mathcal{O}' \to
\prod_{i_0} \mathcal{O}'_{f_{i_0}} \to
\prod_{i_0 < i_1} \mathcal{O}'_{f_{i_0}f_{i_1}} \to \ldots \to
\mathcal{O}'_{f_1\ldots f_r}
}
$$
est un quasi-isomorphisme.
\end{lemma}

\begin{proof}
Remarquons que le second complexe est le produit tensoriel du premier
avec $\mathcal{O}'$. On peut écrire le premier complexe
de {\v C}ech alterné augmenté comme limite inductive des complexes de Koszul
$K_n = K(\mathcal{O}, f_1^n, \ldots, f_r^n)$ ; voir
Compléments d'algèbre, lemme
\ref{more-algebra-lemma-extended-alternating-Cech-is-colimit-koszul}.
Il suffit donc de montrer que $K_n \to K_n \otimes_\mathcal{O} \mathcal{O}'$
est un quasi-isomorphisme. Puisque $\mathcal{O} \to \mathcal{O}'$ est plat,
il suffit de montrer que $H^i \to H^i \otimes_\mathcal{O} \mathcal{O}'$
est un isomorphisme, où $H^i$ est le faisceau de cohomologie de degré $i$,
$H^i = H^i(K_n)$. Ces faisceaux sont annulés par $f_1^n, \ldots, f_r^n$ ; voir
Compléments d'algèbre, lemme \ref{more-algebra-lemma-homotopy-koszul}.
Ils sont donc annulés par $(f_1, \ldots, f_r)^m$
pour un certain $m \gg 0$. Ainsi $H^i \to H^i \otimes_\mathcal{O} \mathcal{O}'$
est un isomorphisme d'après Modules sur les sites, lemme
\ref{sites-modules-lemma-neighbourhood-isomorphism}.
\end{proof}

\begin{lemma}
\label{lemma-restriction-derived-complete-equivalence}
Soit $\mathcal{C}$ un site. Soit $\mathcal{O} \to \mathcal{O}'$ un
homomorphisme de faisceaux d'anneaux. Soit $\mathcal{I} \subset \mathcal{O}$
un faisceau d'idéaux de type fini.
Si $\mathcal{O} \to \mathcal{O}'$ est plat et
$\mathcal{O}/\mathcal{I} \cong \mathcal{O}'/\mathcal{I}\mathcal{O}'$,
alors le foncteur de restriction $D(\mathcal{O}') \to D(\mathcal{O})$
induit une équivalence
$D_{comp}(\mathcal{O}', \mathcal{I}\mathcal{O}') \to
D_{comp}(\mathcal{O}, \mathcal{I})$.
\end{lemma}

\begin{proof}
Le lemme \ref{lemma-pushforward-derived-complete} entraîne que le foncteur de
restriction $r : D(\mathcal{O}') \to D(\mathcal{O})$
envoie $D_{comp}(\mathcal{O}', \mathcal{I}\mathcal{O}')$
dans $D_{comp}(\mathcal{O}, \mathcal{I})$. Nous allons construire un
quasi-inverse $E \mapsto E'$.

\medskip\noindent
Soit $K \to \mathcal{O}$ le morphisme de $D(\mathcal{O})$
construit dans le lemme \ref{lemma-global-extended-cech-complex}.
Posons $K' = K \otimes_\mathcal{O}^\mathbf{L} \mathcal{O}'$ dans $D(\mathcal{O}')$.
Alors $K' \to \mathcal{O}'$ est une application dans $D(\mathcal{O}')$ qui
satisfait aux conclusions du lemme \ref{lemma-global-extended-cech-complex}
relativement à $\mathcal{I}' = \mathcal{I}\mathcal{O}'$.
L'application $K \to r(K')$ est un quasi-isomorphisme d'après le
lemme \ref{lemma-map-identifies-koszul-and-cech-complexes}.
Pour $E \in D_{comp}(\mathcal{O}, \mathcal{I})$, posons maintenant
$$
E' = R\SheafHom_\mathcal{O}(r(K'), E)
$$
considéré comme un objet de $D(\mathcal{O}')$ à l'aide de la structure de $\mathcal{O}'$-module
de $K'$. Puisque $E$ est complet au sens dérivé,
on a $E = R\SheafHom_\mathcal{O}(K, E)$ ; voir la
démonstration de la proposition \ref{proposition-derived-completion}.
D'autre part, puisque $K \to r(K')$ est un isomorphisme,
on voit qu'il existe un isomorphisme
$E \to r(E')$ dans $D(\mathcal{O})$. Pour achever la démonstration, il
faut montrer que, si $E = r(M')$ pour un objet $M'$ de
$D_{comp}(\mathcal{O}', \mathcal{I}')$, alors
$E' \cong M'$. Pour obtenir une application, on utilise
$$
M' = R\SheafHom_{\mathcal{O}'}(\mathcal{O}', M') \to
R\SheafHom_\mathcal{O}(r(\mathcal{O}'), r(M')) \to
R\SheafHom_\mathcal{O}(r(K'), r(M')) = E'
$$
où la seconde flèche utilise l'application $K' \to \mathcal{O}'$.
Pour voir que c'est un isomorphisme, on montre que $r$ appliqué
à cette flèche coïncide avec l'isomorphisme $E \to r(E')$ ci-dessus.
Nous omettons les détails.
\end{proof}

\begin{lemma}
\label{lemma-pushforward-derived-complete-adjoint}
Soit $f : (\Sh(\mathcal{D}), \mathcal{O}') \to (\Sh(\mathcal{C}), \mathcal{O})$
un morphisme de topos annelés. Soient $\mathcal{I} \subset \mathcal{O}$
et $\mathcal{I}' \subset \mathcal{O}'$
des faisceaux d'idéaux de type fini tels que $f^\sharp$ envoie
$f^{-1}\mathcal{I}$ dans $\mathcal{I}'$.
Alors $Rf_*$ envoie $D_{comp}(\mathcal{O}', \mathcal{I}')$
dans $D_{comp}(\mathcal{O}, \mathcal{I})$ et admet un adjoint à gauche
$Lf_{comp}^*$, qui est $Lf^*$ suivi de la complétion dérivée.
\end{lemma}

\begin{proof}
La première assertion a été établie dans le
lemme \ref{lemma-pushforward-derived-complete}.
Notons que la seconde a un sens, car on dispose d'un foncteur de complétion
dérivée $D(\mathcal{O}') \to D_{comp}(\mathcal{O}', \mathcal{I}')$
d'après la proposition \ref{proposition-derived-completion}.
Soient donc $K \in D_{comp}(\mathcal{O}, \mathcal{I})$
et $M \in D_{comp}(\mathcal{O}', \mathcal{I}')$. On a alors
$$
\Hom(K, Rf_*M) = \Hom(Lf^*K, M) = \Hom(Lf_{comp}^*K, M)
$$
par la propriété universelle de la complétion dérivée.
\end{proof}

\begin{lemma}
\label{lemma-pushforward-commutes-with-derived-completion}
\begin{reference}
Généralisation de \cite[Lemma 6.5.9 (2)]{BS}. Comparer à
\cite[Theorem 6.5]{HL-P} dans le cadre des modules quasi-cohérents
et des morphismes de champs algébriques (dérivés).
\end{reference}
Soit $f : (\Sh(\mathcal{D}), \mathcal{O}') \to (\Sh(\mathcal{C}), \mathcal{O})$
un morphisme de topos annelés. Soit $\mathcal{I} \subset \mathcal{O}$
un faisceau d'idéaux de type fini. Soit $\mathcal{I}' \subset \mathcal{O}'$
l'idéal engendré par $f^\sharp(f^{-1}\mathcal{I})$.
Alors $Rf_*$ commute à la complétion dérivée, c'est-à-dire
$Rf_*(K^\wedge) = (Rf_*K)^\wedge$.
\end{lemma}

\begin{proof}
Les foncteurs de complétion dérivée existent d'après la proposition \ref{proposition-derived-completion}.
D'après le lemme \ref{lemma-pushforward-derived-complete}, l'objet
$Rf_*(K^\wedge)$ est complet au sens dérivé ; on obtient donc une application canonique
$(Rf_*K)^\wedge \to Rf_*(K^\wedge)$ par la propriété universelle de la complétion
dérivée. On peut vérifier que c'est un isomorphisme localement sur $\mathcal{C}$.
Ainsi, puisque la complétion dérivée commute à la localisation
(remarque \ref{remark-localization-and-completion}), on peut supposer
que $\mathcal{I}$ est engendré par des sections globales $f_1, \ldots, f_r$.
Alors $\mathcal{I}'$ est engendré par $g_i = f^\sharp(f_i)$. D'après le
lemme \ref{lemma-derived-completion-koszul},
il faut montrer que
$$
R\lim \left(
Rf_*K \otimes^\mathbf{L}_\mathcal{O} K(\mathcal{O}, f_1^n, \ldots, f_r^n)
\right)
=
Rf_*\left(
R\lim
K \otimes^\mathbf{L}_{\mathcal{O}'} K(\mathcal{O}', g_1^n, \ldots, g_r^n)
\right)
$$
Puisque $Rf_*$ commute à $R\lim$
(Cohomologie sur les sites, lemme
\ref{sites-cohomology-lemma-Rf-commutes-with-Rlim})
il suffit de montrer que
$$
Rf_*K \otimes^\mathbf{L}_\mathcal{O} K(\mathcal{O}, f_1^n, \ldots, f_r^n) =
Rf_*\left(
K \otimes^\mathbf{L}_{\mathcal{O}'} K(\mathcal{O}', g_1^n, \ldots, g_r^n)
\right)
$$
Cela résulte de la formule de projection (Cohomologie sur les sites, lemme
\ref{sites-cohomology-lemma-projection-formula}) et du fait que
$Lf^*K(\mathcal{O}, f_1^n, \ldots, f_r^n) =
K(\mathcal{O}', g_1^n, \ldots, g_r^n)$.
\end{proof}

\begin{lemma}
\label{lemma-formal-functions-general}
Soient $A$ un anneau et $I \subset A$ un idéal de type fini.
Soit $\mathcal{C}$ un site et soit $\mathcal{O}$ un faisceau
de $A$-algèbres. Soit $\mathcal{F}$ un faisceau de $\mathcal{O}$-modules.
On a alors
$$
R\Gamma(\mathcal{C}, \mathcal{F})^\wedge =
R\Gamma(\mathcal{C}, \mathcal{F}^\wedge)
$$
dans $D(A)$, où $\mathcal{F}^\wedge$ est le complété
dérivé de $\mathcal{F}$ relativement à $I\mathcal{O}$, tandis que le membre de
gauche est le complété dérivé relativement à $I$.
On obtient ainsi deux suites spectrales
$$
E_2^{i, j} = H^i(H^j(\mathcal{C}, \mathcal{F})^\wedge)
\quad\text{et}\quad
E_2^{p, q} = H^p(\mathcal{C}, H^q(\mathcal{F}^\wedge))
$$
convergeant toutes deux vers
$H^*(R\Gamma(\mathcal{C}, \mathcal{F})^\wedge) =
H^*(\mathcal{C}, \mathcal{F}^\wedge)$
\end{lemma}

\begin{proof}
Appliquons le lemme \ref{lemma-pushforward-commutes-with-derived-completion}
au morphisme de topos annelés $(\mathcal{C}, \mathcal{O}) \to (pt, A)$
et prenons la cohomologie pour obtenir la première assertion. La seconde suite spectrale
est la seconde suite spectrale de
Catégories dérivées, lemme \ref{derived-lemma-two-ss-complex-functor}.
La première suite spectrale est celle de
Compléments d'algèbre, exemple
\ref{more-algebra-example-derived-completion-spectral-sequence}
appliqué à $R\Gamma(\mathcal{C}, \mathcal{F})^\wedge$.
\end{proof}

\begin{remark}
\label{remark-local-calculation-derived-completion}
Soit $(\mathcal{C}, \mathcal{O})$ un site annelé.
Soit $\mathcal{I} \subset \mathcal{O}$ un faisceau d'idéaux
de type fini. Soit $K \mapsto K^\wedge$ la complétion dérivée de la
proposition \ref{proposition-derived-completion}.
Soit $U \in \Ob(\mathcal{C})$ un objet tel que $\mathcal{I}$
soit engendré comme faisceau d'idéaux par $f_1, \ldots, f_r \in \mathcal{I}(U)$.
Posons $A = \mathcal{O}(U)$ et $I = (f_1, \ldots, f_r) \subset A$.
Attention : il se peut que l'égalité $I = \mathcal{I}(U)$ ne soit pas vérifiée.
On a alors
$$
R\Gamma(U, K^\wedge) = R\Gamma(U, K)^\wedge
$$
où le membre de droite est le complété dérivé de
l'objet $R\Gamma(U, K)$ de $D(A)$ relativement à $I$.
Cela résulte de la commutation de la complétion dérivée à la localisation
(remarque \ref{remark-localization-and-completion}) et du
lemme \ref{lemma-formal-functions-general}.
\end{remark}








\section{Le théorème des fonctions formelles}
\label{section-formal-functions}

\noindent
Interrompons le fil de l'exposé pour dire quelques mots de
la complétion dérivée dans le cadre des modules quasi-cohérents sur les schémas
et l'utiliser pour donner une démonstration quelque peu différente du théorème
des fonctions formelles. Quelques indications bibliographiques figurent dans la
remarque \ref{remark-references}.

\medskip\noindent
Le lemme \ref{lemma-pushforward-commutes-with-derived-completion} est une
version dérivée (très formelle) du théorème des fonctions formelles
(Cohomologie des schémas, théorème \ref{coherent-theorem-formal-functions}).
Pour l'expliciter, supposons que $f : X \to S$ soit un morphisme de schémas,
que $\mathcal{I} \subset \mathcal{O}_S$ soit un faisceau quasi-cohérent d'idéaux
de type fini
et que $\mathcal{F}$ soit un faisceau quasi-cohérent sur $X$. Le lemme affirme alors que
\begin{equation}
\label{equation-formal-functions}
Rf_*(\mathcal{F}^\wedge) = (Rf_*\mathcal{F})^\wedge
\end{equation}
où $\mathcal{F}^\wedge$ est le complété dérivé de $\mathcal{F}$
relativement à $f^{-1}\mathcal{I} \cdot \mathcal{O}_X$, et le membre
de droite est le complété dérivé de $Rf_*\mathcal{F}$
relativement à $\mathcal{I}$. Pour retrouver ainsi le théorème
des fonctions formelles, il reste quelques vérifications à faire.

\begin{lemma}
\label{lemma-sections-derived-completion-pseudo-coherent}
Soit $X$ un schéma localement noethérien. Soit $\mathcal{I} \subset \mathcal{O}_X$
un faisceau quasi-cohérent d'idéaux. Soit $K$ un
objet pseudo-cohérent de $D(\mathcal{O}_X)$, de complété dérivé
$K^\wedge$. Alors
$$
H^p(U, K^\wedge) = \lim H^p(U, K)/I^nH^p(U, K) =
H^p(U, K)^\wedge
$$
pour tout ouvert affine $U \subset X$,
où $I = \mathcal{I}(U)$ et où le membre de droite est le complété
dérivé relativement à $I$.
\end{lemma}

\begin{proof}
Écrivons $U = \Spec(A)$. L'anneau $A$ est noethérien,
donc $I \subset A$ est de type fini. On a alors
$$
R\Gamma(U, K^\wedge) = R\Gamma(U, K)^\wedge
$$
d'après la remarque \ref{remark-local-calculation-derived-completion}.
Or $R\Gamma(U, K)$ est un complexe pseudo-cohérent de $A$-modules
(Catégories dérivées des schémas, lemme
\ref{perfect-lemma-pseudo-coherent-affine}).
D'après Compléments d'algèbre, lemme
\ref{more-algebra-lemma-derived-completion-pseudo-coherent}
on en conclut que le module de cohomologie de degré $p$ de $R\Gamma(U, K^\wedge)$
est égal au complété $I$-adique de $H^p(U, K)$.
Cela prouve la première égalité. La seconde, moins importante,
résulte immédiatement d'une nouvelle application du lemme que l'on vient d'utiliser.
\end{proof}

\begin{lemma}
\label{lemma-derived-completion-pseudo-coherent}
Soit $X$ un schéma localement noethérien. Soit $\mathcal{I} \subset \mathcal{O}_X$
un faisceau quasi-cohérent d'idéaux.
Soit $K$ un objet de $D(\mathcal{O}_X)$. Alors
\begin{enumerate}
\item le complété dérivé $K^\wedge$ est égal à
$R\lim (K \otimes_{\mathcal{O}_X}^\mathbf{L} \mathcal{O}_X/\mathcal{I}^n)$.
\end{enumerate}
Soit $K$ un objet pseudo-cohérent de $D(\mathcal{O}_X)$. Alors
\begin{enumerate}
\item[(2)] le faisceau de cohomologie $H^q(K^\wedge)$ est égal à
$\lim H^q(K)/\mathcal{I}^nH^q(K)$.
\end{enumerate}
Soit $\mathcal{F}$ un $\mathcal{O}_X$-module cohérent\footnote{Par exemple
$H^q(K)$ pour $K$ pseudo-cohérent sur notre schéma localement noethérien $X$.}. Alors
\begin{enumerate}
\item[(3)] le complété dérivé $\mathcal{F}^\wedge$ est égal à
$\lim \mathcal{F}/\mathcal{I}^n\mathcal{F}$,
\item[(4)]
$\lim \mathcal{F}/\mathcal{I}^n \mathcal{F} =
R\lim \mathcal{F}/\mathcal{I}^n \mathcal{F}$,
\item[(5)] $H^p(U, \mathcal{F}^\wedge) = 0$ pour $p \not = 0$ et pour tout
ouvert affine $U \subset X$.
\end{enumerate}
\end{lemma}

\begin{proof}
Démonstration de (1). On a une application canonique
$$
K \longrightarrow
R\lim (K \otimes_{\mathcal{O}_X}^\mathbf{L} \mathcal{O}_X/\mathcal{I}^n),
$$
voir la remarque \ref{remark-compare-with-completion}.
La complétion dérivée commute au passage à un sous-schéma ouvert
(remarque \ref{remark-localization-and-completion}).
La formation de $R\lim$ commute au passage à un sous-schéma ouvert.
Pour vérifier que notre application est un isomorphisme, on peut donc travailler localement.
On peut ainsi supposer $X = U = \Spec(A)$.
Écrivons $I = (f_1, \ldots, f_r)$. Soit
$K_n = K(A, f_1^n, \ldots, f_r^n)$ le complexe de Koszul.
D'après Compléments d'algèbre, lemme \ref{more-algebra-lemma-sequence-Koszul-complexes},
nous avons vu que les pro-systèmes $\{K_n\}$ et
$\{A/I^n\}$ de $D(A)$ sont isomorphes.
En utilisant l'équivalence $D(A) = D_{\QCoh}(\mathcal{O}_X)$
de Catégories dérivées des schémas, lemme
\ref{perfect-lemma-affine-compare-bounded}
on voit que les pro-systèmes $\{K(\mathcal{O}_X, f_1^n, \ldots, f_r^n)\}$
et $\{\mathcal{O}_X/\mathcal{I}^n\}$ sont isomorphes dans
$D(\mathcal{O}_X)$. Cela prouve la seconde égalité dans
$$
K^\wedge = R\lim \left(
K \otimes_{\mathcal{O}_X}^\mathbf{L} K(\mathcal{O}_X, f_1^n, \ldots, f_r^n)
\right) =
R\lim (K \otimes_{\mathcal{O}_X}^\mathbf{L} \mathcal{O}_X/\mathcal{I}^n)
$$
La première égalité est donnée par le
lemme \ref{lemma-derived-completion-koszul}.

\medskip\noindent
Supposons $K$ pseudo-cohérent. Pour un ouvert affine $U \subset X$,
on a $H^q(U, K^\wedge) = \lim H^q(U, K)/\mathcal{I}^n(U)H^q(U, K)$
d'après le lemme \ref{lemma-sections-derived-completion-pseudo-coherent}.
Comme cela vaut pour tout $U$, on voit que
$H^q(K^\wedge) = \lim H^q(K)/\mathcal{I}^nH^q(K)$ comme faisceaux.
Cela prouve (2).

\medskip\noindent
L'assertion (3) est un cas particulier de (2).
Les assertions (4) et (5) résultent de
Catégories dérivées des schémas, lemme
\ref{perfect-lemma-Rlim-quasi-coherent}.
\end{proof}

\begin{lemma}
\label{lemma-formal-functions}
Soient $A$ un anneau noethérien et $I \subset A$ un idéal. Soit $X$ un
schéma noethérien sur $A$. Soit $\mathcal{F}$ un
$\mathcal{O}_X$-module cohérent. Supposons que $H^p(X, \mathcal{F})$ soit
un $A$-module de type fini pour tout $p$. Il existe alors des suites exactes courtes
$$
0 \to R^1\lim H^{p - 1}(X, \mathcal{F}/I^n\mathcal{F}) \to
H^p(X, \mathcal{F})^\wedge \to \lim H^p(X, \mathcal{F}/I^n\mathcal{F}) \to 0
$$
de $A$-modules, où $H^p(X, \mathcal{F})^\wedge$ est le complété $I$-adique
usuel. Si $f$ est propre, le terme $R^1\lim$ est nul.
\end{lemma}

\begin{proof}
Considérons les deux suites spectrales du
lemme \ref{lemma-formal-functions-general}.
La première dégénère d'après Compléments d'algèbre, lemme
\ref{more-algebra-lemma-derived-completion-pseudo-coherent}.
On obtient $H^p(X, \mathcal{F})^\wedge$ en degré $p$.
C'est ici qu'intervient l'hypothèse que $H^p(X, \mathcal{F})$ est
un $A$-module de type fini. La seconde dégénère puisque
$$
\mathcal{F}^\wedge = \lim \mathcal{F}/I^n\mathcal{F} =
R\lim \mathcal{F}/I^n\mathcal{F}
$$
est un faisceau d'après le lemme \ref{lemma-derived-completion-pseudo-coherent}.
On obtient $H^p(X, \lim \mathcal{F}/I^n\mathcal{F})$ en degré $p$.
Puisque $R\Gamma(X, -)$ commute aux limites dérivées
(Injectifs, lemme \ref{injectives-lemma-RF-commutes-with-Rlim}),
on obtient aussi
$$
R\Gamma(X, \lim \mathcal{F}/I^n\mathcal{F}) =
R\Gamma(X, R\lim \mathcal{F}/I^n\mathcal{F}) =
R\lim R\Gamma(X, \mathcal{F}/I^n\mathcal{F})
$$
D'après Compléments d'algèbre, remarque
\ref{more-algebra-remark-how-unique}
on obtient des suites exactes
$$
0 \to
R^1\lim H^{p - 1}(X, \mathcal{F}/I^n\mathcal{F}) \to
H^p(X, \lim \mathcal{F}/I^n\mathcal{F}) \to
\lim H^p(X, \mathcal{F}/I^n\mathcal{F}) \to 0
$$
de $A$-modules. En réunissant ces résultats, on obtient la première assertion du lemme.
L'annulation du terme $R^1\lim$ résulte de
Cohomologie des schémas, lemme \ref{coherent-lemma-ML-cohomology-powers-ideal}.
\end{proof}

\begin{remark}
\label{remark-references}
Voici quelques références bibliographiques sur des questions voisines.
Il semble qu'un « théorème dérivé des fonctions formelles » pour les morphismes propres
remonte à \cite[théorème 6.3.1]{lurie-thesis}.
On trouve aussi un exposé dans \cite{dag12}, en particulier
au chapitre 4, qui traite du cas affine ; voir
Compléments d'algèbre, section \ref{more-algebra-section-derived-completion}.
Dans \cite[section 2.9]{G-R}, le changement de base propre et la
complétion dérivée sont étudiés à l'aide de modules (ind-)cohérents.
Un analogue de (\ref{equation-formal-functions})
pour les complexes de modules quasi-cohérents figure dans
\cite[théorème 6.5]{HL-P}
\end{remark}







\section{Algébrisation de la cohomologie locale, I}
\label{section-algebraization-sections-general}

\noindent
Soit $A$ un anneau noethérien et soient $I$ et $J$ deux idéaux de $A$.
Soit $M$ un $A$-module de type fini. Dans cette section, nous étudions les
groupes de cohomologie de l'objet
$$
R\Gamma_J(M)^\wedge
\quad\text{de}\quad
D(A)
$$
où ${}^\wedge$ désigne la complétion $I$-adique dérivée. Rappelons que, dans
Complexes dualisants, lemme \ref{dualizing-lemma-completion-local-H0},
nous avons montré que, si $A$ est complet relativement à $I$,
il existe un isomorphisme
$$
\colim H^0_Z(M) \longrightarrow H^0(R\Gamma_J(M)^\wedge)
$$
où la limite inductive filtrante porte sur les parties fermées $Z = V(J')$
avec $J' \subset J$ et $V(J') \cap V(I) = V(J) \cap V(I)$.
La réunion de ces parties fermées est
\begin{equation}
\label{equation-associated-subset}
T = \{\mathfrak p \in \Spec(A) :
V(\mathfrak p) \cap V(I) \subset V(J) \cap V(I)\}
\end{equation}
C'est une partie de $\Spec(A)$ stable par spécialisation.
Le résultat précédent signifie alors que
$$
H^0_T(M) \longrightarrow H^0(R\Gamma_J(M)^\wedge)
$$
est un isomorphisme si $A$ est complet relativement à $I$ ; voir
Cohomologie locale, lemme \ref{local-cohomology-lemma-adjoint-ext} et
remarque \ref{local-cohomology-remark-upshot}.
Notre méthode pour étendre cet isomorphisme aux groupes de cohomologie supérieure
repose sur le lemme suivant.

\begin{lemma}
\label{lemma-kill-completion-general}
Soient $I, J$ des idéaux d'un anneau noethérien $A$.
Soit $M$ un $A$-module de type fini. Soit $\mathfrak p \subset A$ un idéal premier.
Soient $s$ et $d$ des entiers. Supposons que
\begin{enumerate}
\item $A$ possède un complexe dualisant,
\item $\mathfrak p \not \in V(J) \cap V(I)$,
\item $\text{cd}(A, I) \leq d$, et
\item pour tout idéal premier $\mathfrak p' \subset \mathfrak p$,
on ait
$\text{depth}_{A_{\mathfrak p'}}(M_{\mathfrak p'}) +
\dim((A/\mathfrak p')_\mathfrak q) > d + s$
pour tout $\mathfrak q \in V(\mathfrak p') \cap V(J) \cap V(I)$.
\end{enumerate}
Alors il existe un $f \in A$, $f \not \in \mathfrak p$, qui annule
$H^i(R\Gamma_J(M)^\wedge)$ pour $i \leq s$, où ${}^\wedge$
désigne la complétion $I$-adique.
\end{lemma}

\begin{proof}
Nous utiliserons l'égalité $R\Gamma_J = R\Gamma_{V(J)}$ et de même pour
$I + J$ ; voir
Complexes dualisants, lemme \ref{dualizing-lemma-local-cohomology-noetherian}.
Remarquons que
$R\Gamma_J(M)^\wedge = R\Gamma_I(R\Gamma_J(M))^\wedge =
R\Gamma_{I + J}(M)^\wedge$ ; voir
Complexes dualisants, lemmes
\ref{dualizing-lemma-complete-and-local} et
\ref{dualizing-lemma-local-cohomology-ss}.
On peut donc remplacer $J$ par $I + J$ et supposer $I \subset J$
et $\mathfrak p \not \in V(J)$.
Rappelons que
$$
R\Gamma_J(M)^\wedge = R\Hom_A(R\Gamma_I(A), R\Gamma_J(M))
$$
d'après la description de la complétion dérivée dans
Compléments d'algèbre, lemme \ref{more-algebra-lemma-derived-completion},
combinée à la description de la cohomologie locale dans
Complexes dualisants, lemme
\ref{dualizing-lemma-compute-local-cohomology-noetherian}.
L'hypothèse (3) signifie que la cohomologie de $R\Gamma_I(A)$ n'est non nulle
qu'en degrés $\leq d$. En utilisant les troncatures canoniques de
$R\Gamma_I(A)$, on voit qu'il suffit de montrer que
$$
\text{Ext}^i(N, R\Gamma_J(M))
$$
est annulé par un $f \in A$, $f \not \in \mathfrak p$, pour
$i \leq s + d$ et pour tout $A$-module $N$.
En utilisant à présent les troncatures canoniques de $R\Gamma_J(M)$,
on voit qu'il suffit de montrer que
$H^i_J(M)$ est annulé par un $f \in A$, $f \not \in \mathfrak p$,
pour $i \leq s + d$.
Cela résulte de Cohomologie locale, lemme
\ref{local-cohomology-lemma-kill-local-cohomology-at-prime}.
\end{proof}

\begin{lemma}
\label{lemma-kill-colimit-weak-general}
Soient $I, J$ des idéaux d'un anneau noethérien. Soit $M$ un $A$-module de type fini.
Soient $s$ et $d$ des entiers. Avec $T$ comme dans
(\ref{equation-associated-subset}), supposons que
\begin{enumerate}
\item $A$ possède un complexe dualisant,
\item si $\mathfrak p \in V(I)$, aucune condition n'est imposée,
\item si $\mathfrak p \not \in V(I)$, $\mathfrak p \in T$, alors
$\dim((A/\mathfrak p)_\mathfrak q) \leq d$ pour un certain
$\mathfrak q \in V(\mathfrak p) \cap V(J) \cap V(I)$,
\item si $\mathfrak p \not \in V(I)$, $\mathfrak p \not \in T$, alors
$$
\text{depth}_{A_\mathfrak p}(M_\mathfrak p) \geq s
\quad\text{ou}\quad
\text{depth}_{A_\mathfrak p}(M_\mathfrak p) +
\dim((A/\mathfrak p)_\mathfrak q) > d + s
$$
pour tout $\mathfrak q \in V(\mathfrak p) \cap V(J) \cap V(I)$.
\end{enumerate}
Alors il existe un idéal $J_0 \subset J$ tel que
$V(J_0) \cap V(I) = V(J) \cap V(I)$ et tel que, pour tout $J' \subset J_0$ avec
$V(J') \cap V(I) = V(J) \cap V(I)$, l'application
$$
R\Gamma_{J'}(M) \longrightarrow R\Gamma_{J_0}(M)
$$
induise un isomorphisme en cohomologie en degrés $\leq s$ ;
de plus, ces modules sont annulés par une puissance de $J_0I$.
\end{lemma}

\begin{proof}
Considérons l'ensemble
$$
B = \{\mathfrak p \not \in V(I),\ \mathfrak p \in T,\text{ et }
\text{depth}(M_\mathfrak p) \leq s\}
$$
Choisissons $J_0 \subset J$ tel que $V(J_0)$ soit l'adhérence de $B \cup V(J)$.

\medskip\noindent
Assertion I : $V(J_0) \cap V(I) = V(J) \cap V(I)$.

\medskip\noindent
Démonstration de l'assertion I. L'inclusion $\supset$ résulte de la construction.
Soit $\mathfrak p$ un idéal premier minimal de $V(J_0)$.
Si $\mathfrak p \in B \cup V(J)$, alors ou bien $\mathfrak p \in T$,
ou bien $\mathfrak p \in V(J)$, et, dans les deux cas,
$V(\mathfrak p) \cap V(I) \subset V(J) \cap V(I)$, comme voulu.
Si $\mathfrak p \not \in B \cup V(J)$, alors
$V(\mathfrak p) \cap B$ est dense, donc infini, et l'on en conclut que
$\text{depth}(M_\mathfrak p) < s$ d'après
Cohomologie locale, lemme \ref{local-cohomology-lemma-depth-function}.
En fait, posons
$V(\mathfrak p) \cap B = \{\mathfrak p_\lambda\}_{\lambda \in \Lambda}$.
Choisissons $\mathfrak q_\lambda \in V(\mathfrak p_\lambda) \cap V(J) \cap V(I)$
comme dans (3).
Soit $\delta : \Spec(A) \to \mathbf{Z}$ la fonction de dimension
associée à un complexe dualisant $\omega_A^\bullet$ de $A$.
Puisque $\Lambda$ est infini et que $\delta$ est bornée,
il existe une partie infinie $\Lambda' \subset \Lambda$ sur laquelle
$\delta(\mathfrak q_\lambda)$ est constante. Pour
$\lambda \in \Lambda'$, on a
$$
\text{depth}(M_{\mathfrak p_\lambda}) +
\delta(\mathfrak p_\lambda) - \delta(\mathfrak q_\lambda) =
\text{depth}(M_{\mathfrak p_\lambda}) +
\dim((A/\mathfrak p_\lambda)_{\mathfrak q_\lambda})
\leq d + s
$$
d'après (3) et la définition de $B$. Par la semi-continuité de
la fonction $\text{depth} + \delta$ établie dans
Dualité pour les schémas, lemme \ref{duality-lemma-sitting-in-degrees},
on en conclut que
$$
\text{depth}(M_\mathfrak p) +
\dim((A/\mathfrak p)_{\mathfrak q_\lambda}) =
\text{depth}(M_\mathfrak p) + \delta(\mathfrak p) - \delta(\mathfrak q_\lambda)
\leq d + s
$$
Puisque l'on a aussi $\mathfrak p \not \in V(I)$, on déduit de (4) que
$\mathfrak p \in T$, c'est-à-dire
$V(\mathfrak p) \cap V(I) \subset V(J) \cap V(I)$. Cela achève la
démonstration de l'assertion I.

\medskip\noindent
Assertion II : $H^i_{J_0}(M) \to H^i_J(M)$ est un isomorphisme pour $i \leq s$
et $J' \subset J_0$ tel que $V(J') \cap V(I) = V(J) \cap V(I)$.

\medskip\noindent
Démonstration de l'assertion II. Choisissons $\mathfrak p \in V(J')$ n'appartenant pas à $V(J_0)$.
Il suffit de montrer que $H^i_{\mathfrak pA_\mathfrak p}(M_\mathfrak p) = 0$
pour $i \leq s$ ; voir
Cohomologie locale, lemme \ref{local-cohomology-lemma-isomorphism}.
Remarquons que $\mathfrak p \in T$. Puisque $\mathfrak p$ n'appartient pas à $B$,
on a donc $\text{depth}(M_\mathfrak p) > s$, et les groupes s'annulent d'après
Complexes dualisants, lemme \ref{dualizing-lemma-depth}.

\medskip\noindent
Assertion III. La dernière assertion du lemme est vraie.

\medskip\noindent
D'après l'assertion II, pour $i \leq s$, on a
$$
H^i_T(M) = H^i_{J_0}(M) = H^i_{J'}(M)
$$
pour tout idéal $J' \subset J_0$ tel que $V(J')  \cap V(I) = V(J) \cap V(I)$.
Voir Cohomologie locale, lemme \ref{local-cohomology-lemma-adjoint-ext}.
Vérifions les hypothèses de Cohomologie locale,
proposition \ref{local-cohomology-proposition-annihilator},
pour les parties $T \subset T \cup V(I)$, le module $M$ et l'entier $s$.
Il faut montrer qu'étant donnés $\mathfrak p \subset \mathfrak q$
tels que $\mathfrak p \not \in T \cup V(I)$ et $\mathfrak q \in T$,
on a
$$
\text{depth}_{A_\mathfrak p}(M_\mathfrak p) +
\dim((A/\mathfrak p)_\mathfrak q) > s
$$
Si $\text{depth}(M_\mathfrak p) \geq s$, c'est vrai puisque
la dimension de $(A/\mathfrak p)_\mathfrak q$ est au moins $1$.
On peut donc supposer $\text{depth}(M_\mathfrak p) < s$.
Si $\mathfrak q \in V(I)$, alors $\mathfrak q \in V(J) \cap V(I)$
et l'inégalité résulte de (4). Si $\mathfrak q \not \in V(I)$,
on peut utiliser (3) pour choisir
$\mathfrak q' \in V(\mathfrak q) \cap V(J) \cap V(I)$ tel que
$\dim((A/\mathfrak q)_{\mathfrak q'}) \leq d$.
L'hypothèse (4) donne alors
$$
\text{depth}_{A_\mathfrak p}(M_\mathfrak p) +
\dim((A/\mathfrak p)_{\mathfrak q'}) > s + d
$$
Puisque $A$ est caténaire, on en déduit l'inégalité voulue.
En appliquant Cohomologie locale,
proposition \ref{local-cohomology-proposition-annihilator}, on
trouve $J'' \subset A$ tel que $V(J'') \subset T \cup V(I)$
et tel que $J''$ annule $H^i_T(M)$ pour $i \leq s$.
On peut alors écrire $V(J'') \cup V(J_0) \cup V(I) = V(J'I)$
pour un certain $J' \subset J_0$ avec $V(J') \cap V(I) = V(J) \cap V(I)$.
En remplaçant $J_0$ par $J'$, on achève la démonstration.
\end{proof}

\begin{lemma}
\label{lemma-kill-colimit-general}
Dans le lemme \ref{lemma-kill-colimit-weak-general}, si, au lieu de la condition
vide (2), on suppose que
\begin{enumerate}
\item[(2')] si $\mathfrak p \in V(I)$, $\mathfrak p \not \in V(J) \cap V(I)$,
alors
$\text{depth}_{A_\mathfrak p}(M_\mathfrak p) +
\dim((A/\mathfrak p)_\mathfrak q) > s$
pour tout $\mathfrak q \in V(\mathfrak p) \cap V(J) \cap V(I)$,
\end{enumerate}
ces conditions entraînent aussi que $H^i_{J_0}(M)$ est un
$A$-module de type fini pour $i \leq s$.
\end{lemma}

\begin{proof}
Rappelons que $H^i_{J_0}(M) = H^i_T(M)$ ; voir la démonstration du
lemme \ref{lemma-kill-colimit-weak-general}. Il suffit donc de
vérifier que, pour $\mathfrak p \not \in T$ et $\mathfrak q \in T$
avec $\mathfrak p \subset \mathfrak q$, on a
$\text{depth}_{A_\mathfrak p}(M_\mathfrak p) +
\dim((A/\mathfrak p)_\mathfrak q) > s$ ; voir Cohomologie locale,
proposition \ref{local-cohomology-proposition-finiteness}.
La condition (2') l'assure pour $\mathfrak p \in V(I)$.
Puisque l'on sait que $H^i_T(M)$ est annulé par une puissance de $IJ_0$,
on sait que la condition vaut si $\mathfrak p \not \in V(IJ_0)$
d'après Cohomologie locale, proposition \ref{local-cohomology-proposition-annihilator}.
Tous les cas sont ainsi couverts, et la démonstration est achevée.
\end{proof}

\begin{lemma}
\label{lemma-kill-colimit-support-general}
Si, dans le lemme \ref{lemma-kill-colimit-weak-general}, on suppose en outre que
\begin{enumerate}
\item[(6)] si $\mathfrak p \not \in V(I)$, $\mathfrak p \in T$, alors
$\text{depth}_{A_\mathfrak p}(M_\mathfrak p) > s$,
\end{enumerate}
alors $H^i_{J_0}(M) = H^i_J(M) = H^i_{J + I}(M)$ pour $i \leq s$, et ces
modules sont annulés par une puissance de $I$.
\end{lemma}

\begin{proof}
Choisissons $\mathfrak p \in V(J)$ ou $\mathfrak p \in V(J_0)$, mais
$\mathfrak p \not \in V(J + I) = V(J_0 + I)$.
Il suffit de montrer que $H^i_{\mathfrak pA_\mathfrak p}(M_\mathfrak p) = 0$
pour $i \leq s$ ; voir
Cohomologie locale, lemme \ref{local-cohomology-lemma-isomorphism}.
Ces groupes s'annulent d'après la condition (6) et
Complexes dualisants, lemme \ref{dualizing-lemma-depth}.
La dernière assertion résulte de
Cohomologie locale, proposition \ref{local-cohomology-proposition-annihilator}.
\end{proof}

\begin{lemma}
\label{lemma-algebraize-local-cohomology-general}
Soient $I, J$ des idéaux d'un anneau noethérien $A$.
Soit $M$ un $A$-module de type fini.
Soient $s$ et $d$ des entiers. Avec $T$ comme dans
(\ref{equation-associated-subset}), supposons que
\begin{enumerate}
\item $A$ est complet pour la topologie $I$-adique et possède un complexe dualisant,
\item si $\mathfrak p \in V(I)$, aucune condition n'est imposée,
\item $\text{cd}(A, I) \leq d$,
\item si $\mathfrak p \not \in V(I)$, $\mathfrak p \not \in T$, alors
$$
\text{depth}_{A_\mathfrak p}(M_\mathfrak p) \geq s
\quad\text{ou}\quad
\text{depth}_{A_\mathfrak p}(M_\mathfrak p) +
\dim((A/\mathfrak p)_\mathfrak q) > d + s
$$
pour tout $\mathfrak q \in V(\mathfrak p) \cap V(J) \cap V(I)$,
\item si $\mathfrak p \not \in V(I)$, $\mathfrak p \not \in T$,
$V(\mathfrak p) \cap V(J) \cap V(I) \not = \emptyset$ et
$\text{depth}(M_\mathfrak p) < s$, alors l'une
des conditions suivantes est satisfaite\footnote{Notre méthode
impose cette condition supplémentaire. Nous y reviendrons
(insérer une référence ultérieure).} :
\begin{enumerate}
\item $\dim(\text{Supp}(M_\mathfrak p)) < s + 2$\footnote{Par exemple,
si $M$ satisfait la condition $(S_s)$ de Serre
sur le complémentaire de $V(I) \cup T$.}, ou
\item  $\delta(\mathfrak p) > d + \delta_{max} - 1$
où $\delta$ est une fonction de dimension et $\delta_{max}$
est le maximum de $\delta$ sur $V(J) \cap V(I)$, ou
\item $\text{depth}_{A_\mathfrak p}(M_\mathfrak p) +
\dim((A/\mathfrak p)_\mathfrak q) > d + s + \delta_{max} - \delta_{min} - 2$
pour tout $\mathfrak q \in V(\mathfrak p) \cap V(J) \cap V(I)$.
\end{enumerate}
\end{enumerate}
Alors il existe un idéal $J_0 \subset J$ tel que
$V(J_0) \cap V(I) = V(J) \cap V(I)$
et tel que, pour tout $J' \subset J_0$ avec
$V(J') \cap V(I) = V(J) \cap V(I)$, l'application
$$
R\Gamma_{J'}(M) \longrightarrow R\Gamma_J(M)^\wedge
$$
induise un isomorphisme en cohomologie en degrés $\leq s$.
Ici ${}^\wedge$ désigne la complétion $I$-adique dérivée.
\end{lemma}

\noindent
Nous invitons le lecteur à lire d'abord la démonstration dans le cas local
(lemme \ref{lemma-algebraize-local-cohomology}), car elle explique la structure
de la démonstration sans avoir à traiter toutes les inégalités.

\begin{proof}
Pour un idéal $\mathfrak a \subset A$, on a
$R\Gamma_\mathfrak a = R\Gamma_{V(\mathfrak a)}$ ; voir
Complexes dualisants, lemme \ref{dualizing-lemma-local-cohomology-noetherian}.
Remarquons ensuite que
$$
R\Gamma_J(M)^\wedge =
R\Gamma_I(R\Gamma_J(M))^\wedge =
R\Gamma_{I + J}(M)^\wedge =
R\Gamma_{I + J'}(M)^\wedge =
R\Gamma_I(R\Gamma_{J'}(M))^\wedge =
R\Gamma_{J'}(M)^\wedge
$$
d'après Complexes dualisants, lemmes \ref{dualizing-lemma-local-cohomology-ss} et
\ref{dualizing-lemma-complete-and-local}.
Cela explique comment on définit la flèche de l'énoncé du lemme.

\medskip\noindent
Nous affirmons que les hypothèses du lemme \ref{lemma-kill-colimit-weak-general}
résultent de nos hypothèses actuelles sur $M$.
La seule chose à vérifier est l'hypothèse (3).
Soit donc $\mathfrak p \not \in V(I)$, $\mathfrak p \in T$.
Alors $V(\mathfrak p) \cap V(I)$ est non vide, puisque $I$ est
contenu dans le radical de Jacobson de $A$
(Algèbre, lemme \ref{algebra-lemma-radical-completion}).
Puisque $\mathfrak p \in T$, on a
$V(\mathfrak p) \cap V(I) = V(\mathfrak p) \cap V(J) \cap V(I)$.
Soit $\mathfrak q \in V(\mathfrak p) \cap V(I)$ le
point générique d'une composante irréductible.
On a $\text{cd}(A_\mathfrak q, I_\mathfrak q) \leq d$
d'après Cohomologie locale, lemme \ref{local-cohomology-lemma-cd-local}.
On a $V(\mathfrak pA_\mathfrak q) \cap V(I_\mathfrak q) =
\{\mathfrak qA_\mathfrak q\}$ par notre choix de $\mathfrak q$,
et l'on en conclut que $\dim((A/\mathfrak p)_\mathfrak q) \leq d$
d'après Cohomologie locale, lemme \ref{local-cohomology-lemma-cd-bound-dim-local}.

\medskip\noindent
Remarquons que le lemme vaut pour $s < 0$. Ce cas n'est pas trivial, car
il n'est pas clair a priori que $H^i(R\Gamma_J(M)^\wedge)$
soit nul pour $i < 0$. Cette annulation a toutefois été établie dans
Complexes dualisants, lemme \ref{dualizing-lemma-completion-local}.
Nous démontrerons le lemme par récurrence pour $s \geq 0$.

\medskip\noindent
Le lemme pour $s = 0$ résulte immédiatement de
la conclusion du lemme \ref{lemma-kill-colimit-weak-general}
et de Complexes dualisants, lemme \ref{dualizing-lemma-completion-local-H0}.

\medskip\noindent
Supposons $s > 0$ et supposons le lemme démontré pour les valeurs plus petites de $s$.
Soit $M' \subset M$ le sous-module maximal dont le support est contenu
dans $V(I) \cup T$. Alors $M'$ est un $A$-module de type fini dont le support
est contenu dans $V(J') \cup V(I)$ pour un certain idéal $J' \subset J$
tel que $V(J') \cap V(I) = V(J) \cap V(I)$.
Nous affirmons que
$$
R\Gamma_{J'}(M') \to R\Gamma_J(M')^\wedge
$$
est un isomorphisme quel que soit le choix de $J'$.
En effet, on peut choisir une suite exacte courte
$0 \to M_1 \oplus M_2 \to M' \to N \to 0$, où
$M_1$ est annulé par une puissance de $J'$, où $M_2$ est annulé
par une puissance de $I$, et où $N$ est annulé par une puissance de $I + J'$.
Il suffit donc de montrer que l'assertion vaut pour $M_1$, $M_2$ et $N$.
Dans le cas de $M_1$, on voit que $R\Gamma_{J'}(M_1) = M_1$ et,
puisque $M_1$ est un $A$-module de type fini et complet pour la topologie $I$-adique,
on a $M_1^\wedge = M_1$. Cela prouve l'assertion pour $M_1$
d'après les remarques initiales de la démonstration. Dans le cas de $M_2$, on voit que
$H^i_J(M_2) = H^i_{I + J}(M) = H^i_{I + J'}(M) = H^i_{J'}(M_2)$
est annulé par une puissance de $I$, et est donc complet au sens dérivé.
L'assertion vaut donc aussi dans ce cas. Pour $N$, on peut utiliser l'un ou l'autre
des arguments que l'on vient de donner. En considérant la suite exacte courte
$0 \to M' \to M \to M/M' \to 0$
on voit qu'il suffit de démontrer le lemme pour $M/M'$.
On peut donc supposer $\text{Ass}(M) \cap (V(I) \cup T) = \emptyset$.

\medskip\noindent
Soit $\mathfrak p \in \text{Ass}(M)$ tel que
$V(\mathfrak p) \cap V(J) \cap V(I) = \emptyset$.
Puisque $I$ est contenu dans le radical de Jacobson de $A$, cela entraîne
que $V(\mathfrak p) \cap V(J') = \emptyset$ pour tout
$J' \subset J$ tel que $V(J') \cap V(I) = V(J) \cap V(I)$.
En posant donc $N = H^0_\mathfrak p(M)$, on voit que
$R\Gamma_J(N) = R\Gamma_{J'}(N) = 0$ pour tout
$J' \subset J$ tel que $V(J') \cap V(I) = V(J) \cap V(I)$.
En particulier, $R\Gamma_J(N)^\wedge = 0$.
On peut donc remplacer $M$ par $M/N$, puisque cela ne modifie la
structure de $M$ qu'aux idéaux premiers qui ne jouent aucun
rôle dans les conditions (4) ou (5). En répétant l'opération, on peut supposer que
$V(\mathfrak p) \cap V(J) \cap V(I) \not = \emptyset$
pour tout $\mathfrak p \in \text{Ass}(M)$.

\medskip\noindent
Supposons $\text{Ass}(M) \cap (V(I) \cup T) = \emptyset$ et que
$V(\mathfrak p) \cap V(J) \cap V(I) \not = \emptyset$
pour tout $\mathfrak p \in \text{Ass}(M)$.
Soit $\mathfrak p \in \text{Ass}(M)$. Nous voulons montrer que l'on peut appliquer le
lemme \ref{lemma-kill-completion-general}.
C'est dans cette vérification que nous utiliserons la condition supplémentaire
(5). Choisissons $\mathfrak p' \subset \mathfrak p$
et $\mathfrak q' \subset V(\mathfrak p) \cap V(J) \cap V(I)$.
\begin{enumerate}
\item Si $M_{\mathfrak p'} = 0$, alors
$\text{depth}(M_{\mathfrak p'}) = \infty$ et
$\text{depth}(M_{\mathfrak p'}) +
\dim((A/\mathfrak p')_{\mathfrak q'}) > d + s$.
\item Si $\text{depth}(M_{\mathfrak p'}) < s$, alors
$\text{depth}(M_{\mathfrak p'}) +
\dim((A/\mathfrak p')_{\mathfrak q'}) > d + s$ d'après (4).
\end{enumerate}
Dans les cas restants, on a $M_{\mathfrak p'} \not = 0$ et
$\text{depth}(M_{\mathfrak p'}) \geq s$. En particulier, on voit que
$\mathfrak p'$ appartient au support de $M$, et l'on peut choisir
$\mathfrak p'' \subset \mathfrak p'$ avec $\mathfrak p'' \in \text{Ass}(M)$.
\begin{enumerate}
\item[(a)] Remarquons que
$\dim((A/\mathfrak p'')_{\mathfrak p'}) \geq \text{depth}(M_{\mathfrak p'})$
d'après Algèbre, lemme \ref{algebra-lemma-depth-dim-associated-primes}.
Si l'égalité a lieu, alors on a
$$
\text{depth}(M_{\mathfrak p'}) + \dim((A/\mathfrak p')_{\mathfrak q'}) =
\text{depth}(M_{\mathfrak p''}) + \dim((A/\mathfrak p'')_{\mathfrak q'})
> s + d
$$
d'après (4) appliquée à $\mathfrak p''$, et l'on a terminé. Cela signifie que la
seule difficulté se présente lorsque
$\dim((A/\mathfrak p'')_{\mathfrak p'}) > \text{depth}(M_{\mathfrak p'})$.
Cela entraîne que $\dim(M_\mathfrak p) \geq s + 2$.
Ainsi, si (5)(a) est satisfaite, ce cas ne se présente pas.
\item[(b)] Si (5)(b) est satisfaite, on obtient
$$
\text{depth}(M_{\mathfrak p'}) + \dim((A/\mathfrak p')_{\mathfrak q'})
\geq s + \delta(\mathfrak p') - \delta(\mathfrak q')
\geq s + 1 + \delta(\mathfrak p) - \delta_{max}
> s + d
$$
comme voulu.
\item[(c)] Si (5)(c) est satisfaite, on obtient
\begin{align*}
\text{depth}(M_{\mathfrak p'}) + \dim((A/\mathfrak p')_{\mathfrak q'})
& \geq
s + \delta(\mathfrak p') - \delta(\mathfrak q') \\
& \geq
s + 1 + \delta(\mathfrak p) - \delta(\mathfrak q') \\
& =
s + 1 + \delta(\mathfrak p) - \delta(\mathfrak q) +
\delta(\mathfrak q) - \delta(\mathfrak q') \\
& >
s + 1 + (s + d + \delta_{max} - \delta_{min} - 2) +
\delta(\mathfrak q) - \delta(\mathfrak q') \\
& \geq 
2s + d - 1 \geq s + d
\end{align*}
comme voulu. Remarquons que cet argument fonctionne parce que
nous savons qu'un idéal premier $\mathfrak q \in V(\mathfrak p) \cap V(J) \cap V(I)$
existe.
\end{enumerate}
Nous sommes maintenant prêts à effectuer l'étape de récurrence.

\medskip\noindent
Choisissons un idéal $J_0$ comme dans le lemme \ref{lemma-kill-colimit-weak-general}
et un entier $t > 0$ tel que $(J_0I)^t$ annule $H^s_J(M)$.
Les hypothèses du lemme \ref{lemma-kill-completion-general}
sont satisfaites pour tout $\mathfrak p \in \text{Ass}(M)$
(voir le paragraphe précédent).
Ainsi, l'annulateur $\mathfrak a \subset A$ de
$H^s(R\Gamma_J(M)^\wedge)$
n'est pas contenu dans $\mathfrak p$ pour $\mathfrak p \in \text{Ass}(M)$.
On peut donc trouver un $f \in \mathfrak a(J_0I)^t$
qui n'appartient à aucun idéal premier associé de $M$ et qui annule
à la fois $H^s(R\Gamma_J(M)^\wedge)$ et $H^s_J(M)$.
Alors $f$ est un non-diviseur de zéro dans $M$, et l'on peut considérer la
suite exacte courte
$$
0 \to M \xrightarrow{f} M \to M/fM \to 0
$$
Notre choix de $f$ montre que l'on obtient
$$
\xymatrix{
H^{s - 1}_{J'}(M) \ar[d] \ar[r] &
H^{s - 1}_{J'}(M/fM) \ar[d] \ar[r] &
H^s_{J'}(M) \ar[d] \ar[r] & 0 \\
H^{s - 1}(R\Gamma_J(M)^\wedge) \ar[r] &
H^{s - 1}(R\Gamma_J(M/fM)^\wedge) \ar[r] &
H^s(R\Gamma_J(M)^\wedge) \ar[r] & 0
}
$$
pour tout $J' \subset J_0$ tel que $V(J') \cap V(I) = V(J) \cap V(I)$.
Ainsi, si l'on choisit $J'$ de sorte qu'il convienne pour
$M$, $M/fM$ et $s - 1$ (ce qui est possible par l'hypothèse de récurrence ---
voir le paragraphe suivant), on en conclut que le lemme est vrai.

\medskip\noindent
Pour achever la démonstration, il faut montrer que le module
$M/fM$ satisfait les hypothèses (4) et (5) pour $s - 1$.
Soit donc $\mathfrak p$ un idéal premier du support
de $M/fM$, avec $\text{depth}((M/fM)_\mathfrak p) < s - 1$
et avec $V(\mathfrak p) \cap V(J) \cap V(I)$ non vide.
Alors $\dim(M_\mathfrak p) = \dim((M/fM)_\mathfrak p) + 1$
et $\text{depth}(M_\mathfrak p) = \text{depth}((M/fM)_\mathfrak p) + 1$.
En particulier, on sait que (4) et (5) sont satisfaites pour $\mathfrak p$ et $M$
avec la valeur initiale de $s$.
Les inégalités voulues s'en déduisent alors immédiatement.
\end{proof}

\begin{example}
\label{example-no-ML}
Dans le lemme \ref{lemma-algebraize-local-cohomology-general},
nous ne savons pas si les systèmes projectifs $H^i_J(M/I^nM)$ satisfont la
condition de Mittag-Leffler.
Par exemple, supposons $A = \mathbf{Z}_p[[x, y]]$, $I = (p)$,
$J = (p, x)$ et $M = A/(xy - p)$. Alors l'image de
$H^0_J(M/p^nM) \to H^0_J(M/pM)$
est l'idéal engendré par $y^n$ dans $M/pM = A/(p, xy)$.
\end{example}





\section{Algébrisation de la cohomologie locale, II}
\label{section-algebraization-punctured}

\noindent
Dans cette section, nous reprenons les arguments de la
section \ref{section-algebraization-sections-general}
lorsque $(A, \mathfrak m)$ est un anneau local et que l'on considère la cohomologie locale
$R\Gamma_\mathfrak m$ relativement à $\mathfrak m$. Comme précédemment, notre
outil principal est le lemme suivant.

\begin{lemma}
\label{lemma-kill-completion}
Soit $(A, \mathfrak m)$ un anneau local noethérien.
Soit $I \subset A$ un idéal. Soit $M$ un $A$-module de type fini et
soit $\mathfrak p \subset A$ un idéal premier. Soient $s$ et $d$ des entiers. Supposons que
\begin{enumerate}
\item $A$ possède un complexe dualisant,
\item $\text{cd}(A, I) \leq d$, et
\item
$\text{depth}_{A_\mathfrak p}(M_\mathfrak p) + \dim(A/\mathfrak p) > d + s$.
\end{enumerate}
Alors il existe un $f \in A \setminus \mathfrak p$ qui annule
$H^i(R\Gamma_\mathfrak m(M)^\wedge)$ pour $i \leq s$, où ${}^\wedge$
désigne la complétion $I$-adique.
\end{lemma}

\begin{proof}
D'après Cohomologie locale, lemme
\ref{local-cohomology-lemma-sitting-in-degrees}
la fonction
$$
\mathfrak p' \longmapsto
\text{depth}_{A_{\mathfrak p'}}(M_{\mathfrak p'}) + \dim(A/\mathfrak p')
$$
est semi-continue inférieurement sur $\Spec(A)$. La valeur
de cette fonction en tout $\mathfrak p' \subset \mathfrak p$
est donc $> s + d$. Notre lemme est ainsi un cas particulier du
lemme \ref{lemma-kill-completion-general},
pourvu que $\mathfrak p \not = \mathfrak m$.
Si $\mathfrak p = \mathfrak m$,
alors $H^i_\mathfrak m(M) = 0$ pour $i \leq s + d$ d'après
la relation entre la profondeur et la cohomologie locale
(Complexes dualisants, lemme \ref{dualizing-lemma-depth}).
Ainsi, l'argument donné dans la démonstration du
lemme \ref{lemma-kill-completion-general}
montre que $H^i(R\Gamma_\mathfrak m(M)^\wedge) = 0$
pour $i \leq s$ dans ce cas (dégénéré).
\end{proof}

\begin{lemma}
\label{lemma-kill-colimit-weak}
Soit $(A, \mathfrak m)$ un anneau local noethérien.
Soit $I \subset A$ un idéal. Soit $M$ un $A$-module de type fini.
Soient $s$ et $d$ des entiers. Supposons que
\begin{enumerate}
\item $A$ possède un complexe dualisant,
\item si $\mathfrak p \in V(I)$, aucune condition n'est imposée,
\item si $\mathfrak p \not \in V(I)$ et
$V(\mathfrak p) \cap V(I) = \{\mathfrak m\}$, alors
$\dim(A/\mathfrak p) \leq d$,
\item si $\mathfrak p \not \in V(I)$ et
$V(\mathfrak p) \cap V(I) \not = \{\mathfrak m\}$, alors
$$
\text{depth}_{A_\mathfrak p}(M_\mathfrak p) \geq s
\quad\text{ou}\quad
\text{depth}_{A_\mathfrak p}(M_\mathfrak p) + \dim(A/\mathfrak p) > d + s
$$
\end{enumerate}
Alors il existe un idéal $J_0 \subset A$ tel que
$V(J_0) \cap V(I) = \{\mathfrak m\}$ et tel que, pour tout $J \subset J_0$ avec
$V(J) \cap V(I) = \{\mathfrak m\}$, l'application
$$
R\Gamma_J(M) \longrightarrow R\Gamma_{J_0}(M)
$$
induise un isomorphisme en cohomologie en degrés $\leq s$ ;
de plus, ces modules sont annulés par une puissance de $J_0I$.
\end{lemma}

\begin{proof}
C'est un cas particulier du lemme \ref{lemma-kill-colimit-weak-general}.
\end{proof}

\begin{lemma}
\label{lemma-kill-colimit}
Dans le lemme \ref{lemma-kill-colimit-weak}, si, au lieu de la condition
vide (2), on suppose que
\begin{enumerate}
\item[(2')] si $\mathfrak p \in V(I)$ et $\mathfrak p \not = \mathfrak m$,
alors $\text{depth}_{A_\mathfrak p}(M_\mathfrak p) + \dim(A/\mathfrak p) > s$,
\end{enumerate}
alors ces conditions entraînent aussi que $H^i_{J_0}(M)$ est un
$A$-module de type fini pour $i \leq s$.
\end{lemma}

\begin{proof}
C'est un cas particulier du lemme \ref{lemma-kill-colimit-general}.
\end{proof}

\begin{lemma}
\label{lemma-kill-colimit-support}
Si, dans le lemme \ref{lemma-kill-colimit-weak}, on suppose en outre que
\begin{enumerate}
\item[(6)] si $\mathfrak p \not \in V(I)$ et
$V(\mathfrak p) \cap V(I) = \{\mathfrak m\}$, alors
$\text{depth}_{A_\mathfrak p}(M_\mathfrak p) > s$,
\end{enumerate}
alors $H^i_{J_0}(M) = H^i_J(M) = H^i_\mathfrak m(M)$ pour $i \leq s$,
et ces modules sont annulés par une puissance de $I$.
\end{lemma}

\begin{proof}
C'est un cas particulier du lemme \ref{lemma-kill-colimit-support-general}.
\end{proof}

\begin{lemma}
\label{lemma-algebraize-local-cohomology}
Soit $(A, \mathfrak m)$ un anneau local noethérien.
Soit $I \subset A$ un idéal. Soit $M$ un $A$-module de type fini.
Soient $s$ et $d$ des entiers. Supposons que
\begin{enumerate}
\item $A$ est complet pour la topologie $I$-adique et possède un complexe dualisant,
\item si $\mathfrak p \in V(I)$, aucune condition n'est imposée,
\item $\text{cd}(A, I) \leq d$,
\item si $\mathfrak p \not \in V(I)$ et
$V(\mathfrak p) \cap V(I) \not = \{\mathfrak m\}$, alors
$$
\text{depth}_{A_\mathfrak p}(M_\mathfrak p) \geq s
\quad\text{ou}\quad
\text{depth}_{A_\mathfrak p}(M_\mathfrak p) + \dim(A/\mathfrak p) > d + s
$$
\end{enumerate}
Alors il existe un idéal $J_0 \subset A$ tel que
$V(J_0) \cap V(I) = \{\mathfrak m\}$ et tel que, pour tout $J \subset J_0$ avec
$V(J) \cap V(I) = \{\mathfrak m\}$, l'application
$$
R\Gamma_J(M) \longrightarrow
R\Gamma_J(M)^\wedge = R\Gamma_\mathfrak m(M)^\wedge
$$
induise un isomorphisme en cohomologie en degrés $\leq s$.
Ici ${}^\wedge$ désigne la complétion $I$-adique dérivée.
\end{lemma}

\begin{proof}
Ce lemme est un cas particulier du
lemme \ref{lemma-algebraize-local-cohomology-general},
puisque la condition (5)(c) résulte de la condition (4),
car $\delta_{max} = \delta_{min} = \delta(\mathfrak m)$.
Nous allons donner la démonstration de ce cas particulier important,
car elle est un peu plus simple (il y a moins de choses à vérifier).

\medskip\noindent
Il n'y a pas de différence entre $R\Gamma_\mathfrak a$ et
$R\Gamma_{V(\mathfrak a)}$ dans la situation présente ; voir
Complexes dualisants, lemme \ref{dualizing-lemma-local-cohomology-noetherian}.
Remarquons ensuite que
$$
R\Gamma_\mathfrak m(M)^\wedge =
R\Gamma_I(R\Gamma_J(M))^\wedge =
R\Gamma_J(M)^\wedge
$$
d'après Complexes dualisants, lemmes \ref{dualizing-lemma-local-cohomology-ss} et
\ref{dualizing-lemma-complete-and-local},
ce qui explique le signe d'égalité dans l'énoncé du lemme.

\medskip\noindent
Remarquons que le lemme vaut pour $s < 0$. Ce cas n'est pas trivial, car
il n'est pas clair a priori que $H^s(R\Gamma_\mathfrak m(M)^\wedge)$
soit nul pour $s$ négatif. Cette annulation a toutefois été établie
dans le lemme \ref{lemma-local-cohomology-derived-completion}.
Nous démontrerons le lemme par récurrence pour $s \geq 0$.

\medskip\noindent
Les hypothèses du lemme \ref{lemma-kill-colimit-weak}
sont satisfaites d'après Cohomologie locale, lemme
\ref{local-cohomology-lemma-cd-bound-dim-local}.
Le lemme pour $s = 0$ résulte du lemme \ref{lemma-kill-colimit-weak} et de
Complexes dualisants, lemme \ref{dualizing-lemma-completion-local-H0}.

\medskip\noindent
Supposons $s > 0$ et le lemme vrai pour les valeurs plus petites de $s$.
Soit $M' \subset M$ le sous-module des éléments dont le
support est contenu dans $V(I) \cup V(J)$ pour un certain
idéal $J$ tel que $V(J) \cap V(I) = \{\mathfrak m\}$.
Alors $M'$ est un $A$-module de type fini.
Nous affirmons que
$$
R\Gamma_J(M') \to R\Gamma_\mathfrak m(M')^\wedge
$$
est un isomorphisme quel que soit le choix de $J$.
En effet, pour tout module de ce type, il existe une suite exacte courte
$0 \to M_1 \oplus M_2 \to M' \to N \to 0$, où
$M_1$ est annulé par une puissance de $J$, où $M_2$ est annulé
par une puissance de $I$ et où $N$ est annulé par une puissance de $\mathfrak m$.
Dans le cas de $M_1$, on voit que $R\Gamma_J(M_1) = M_1$ et,
puisque $M_1$ est un $A$-module de type fini et complet pour la topologie $I$-adique,
on a $M_1^\wedge = M_1$. L'assertion vaut donc pour $M_1$.
Dans le cas de $M_2$, on voit que $H^i_J(M_2)$ est annulé
par une puissance de $I$, et est donc complet au sens dérivé. Ainsi, l'assertion vaut
pour $M_2$. Les mêmes arguments montrent que l'assertion vaut pour $N$,
et l'on en conclut qu'elle est vraie. En considérant la
suite exacte courte $0 \to M' \to M \to M/M' \to 0$,
on voit qu'il suffit de démontrer le lemme pour $M/M'$.
On peut donc supposer que $\mathfrak p \in \text{Ass}(M)$
entraîne $V(\mathfrak p) \cap V(I) \not = \{\mathfrak m\}$, c'est-à-dire que
$\mathfrak p$ est un idéal premier comme dans (4).

\medskip\noindent
Choisissons un idéal $J_0$ comme dans le lemme \ref{lemma-kill-colimit-weak}
et un entier $t > 0$ tel que $(J_0I)^t$ annule $H^s_J(M)$.
Ici, $J$ désigne un idéal arbitraire $J \subset J_0$ tel que
$V(J) \cap V(I) = \{\mathfrak m\}$.
Les hypothèses du lemme \ref{lemma-kill-completion}
sont satisfaites pour tout $\mathfrak p \in \text{Ass}(M)$
(voir le paragraphe précédent). Ainsi, l'annulateur $\mathfrak a \subset A$ de
$H^s(R\Gamma_\mathfrak m(M)^\wedge)$
n'est pas contenu dans $\mathfrak p$ pour $\mathfrak p \in \text{Ass}(M)$.
On peut donc trouver un $f \in \mathfrak a(J_0I)^t$
qui n'appartient à aucun idéal premier associé de $M$ et qui annule
à la fois $H^s(R\Gamma_\mathfrak m(M)^\wedge)$ et $H^s_J(M)$.
Alors $f$ est un non-diviseur de zéro dans $M$, et l'on peut considérer la
suite exacte courte
$$
0 \to M \xrightarrow{f} M \to M/fM \to 0
$$
Notre choix de $f$ montre que l'on obtient
$$
\xymatrix{
H^{s - 1}_J(M) \ar[d] \ar[r] &
H^{s - 1}_J(M/fM) \ar[d] \ar[r] &
H^s_J(M) \ar[d] \ar[r] & 0 \\
H^{s - 1}(R\Gamma_\mathfrak m(M)^\wedge) \ar[r] &
H^{s - 1}(R\Gamma_\mathfrak m(M/fM)^\wedge) \ar[r] &
H^s(R\Gamma_\mathfrak m(M)^\wedge) \ar[r] & 0
}
$$
pour tout $J \subset J_0$ tel que $V(J) \cap V(I) = \{\mathfrak m\}$.
Ainsi, si l'on choisit $J$ de sorte qu'il convienne pour
$M$, $M/fM$ et $s - 1$ (ce qui est possible par l'hypothèse de récurrence),
on en conclut que le lemme est vrai.
\end{proof}








\section{Algébrisation de la cohomologie locale, III}
\label{section-bootstrap}

\noindent
Dans cette section, nous nous appuyons sur les résultats des
sections \ref{section-algebraization-sections-general} et
\ref{section-algebraization-punctured}
pour obtenir un résultat plus fort dans la situation suivante.

\begin{situation}
\label{situation-bootstrap}
Ici, $A$ est un anneau noethérien. Nous avons un idéal $I \subset A$,
un $A$-module de type fini $M$ et un sous-ensemble $T \subset V(I)$ stable par
spécialisation. Nous avons des entiers $s$ et $d$. Nous supposons que
\begin{enumerate}
\item[(1)] $A$ possède un complexe dualisant,
\item[(3)] $\text{cd}(A, I) \leq d$,
\item[(4)] pour des idéaux premiers $\mathfrak p \subset \mathfrak r \subset \mathfrak q$
tels que $\mathfrak p \not \in V(I)$,
$\mathfrak r \in V(I) \setminus T$,
$\mathfrak q \in T$, on a
$$
\text{depth}_{A_\mathfrak p}(M_\mathfrak p) \geq s
\quad\text{ou}\quad
\text{depth}_{A_\mathfrak p}(M_\mathfrak p) +
\dim((A/\mathfrak p)_\mathfrak q) > d + s
$$
\item[(6)] pour $\mathfrak q \in T$, si
$A', \mathfrak m', I', M'$ désignent les complétions $I$-adiques usuelles
de $A_\mathfrak q, \mathfrak qA_\mathfrak q, I_\mathfrak q, M_\mathfrak q$,
on a
$$
\text{depth}(M'_{\mathfrak p'}) > s
$$
pour tout $\mathfrak p' \in \Spec(A') \setminus V(I')$ tel que
$V(\mathfrak p') \cap V(I') = \{\mathfrak m'\}$.
\end{enumerate}
\end{situation}

\noindent
Le lemme suivant explique pourquoi, dans la situation \ref{situation-bootstrap},
il suffit de considérer les triplets
$\mathfrak p \subset \mathfrak r \subset \mathfrak q$ d'idéaux premiers de (4),
bien que l'hypothèse elle-même ne fasse intervenir que $\mathfrak p$ et $\mathfrak q$.

\begin{lemma}
\label{lemma-helper-bootstrap}
Dans la situation \ref{situation-bootstrap}, soient $\mathfrak p \subset \mathfrak q$
des idéaux premiers de $A$ tels que $\mathfrak p \not \in V(I)$ et
$\mathfrak q \in T$. S'il n'existe aucun
$\mathfrak r \in V(I) \setminus T$ tel que
$\mathfrak p \subset \mathfrak r \subset \mathfrak q$
alors $\text{depth}(M_\mathfrak p) > s$.
\end{lemma}

\begin{proof}
Choisissons $\mathfrak q' \in T$ tel que
$\mathfrak p \subset \mathfrak q' \subset \mathfrak q$
et qu'il n'existe aucun idéal premier de $T$ strictement
compris entre $\mathfrak p$ et $\mathfrak q'$. Pour démontrer le lemme,
on peut, et l'on va, remplacer $\mathfrak q$ par $\mathfrak q'$.
Ensuite, soit $\mathfrak p' \subset A_\mathfrak q$ l'idéal premier correspondant à
$\mathfrak p$. On obtient alors
$V(\mathfrak p') \cap V(IA_\mathfrak q) = \{\mathfrak q A_\mathfrak q\}$
du fait qu'il n'existe pas d'idéal premier $\mathfrak r$ comme dans le lemme.
Soient $A', I', \mathfrak m', M'$ les complétions $I$-adiques de
$A_\mathfrak q, I_\mathfrak q, \mathfrak qA_\mathfrak q, M_\mathfrak q$.
Comme $A_\mathfrak q \to A'$ est fidèlement plat
(Algèbre, lemme \ref{algebra-lemma-completion-faithfully-flat}),
on peut choisir $\mathfrak p'' \subset A'$ au-dessus de $\mathfrak p'$
tel que $\dim(A'_{\mathfrak p''}/\mathfrak p' A'_{\mathfrak p''}) = 0$.
On voit alors que
$$
\text{depth}(M'_{\mathfrak p''}) =
\text{depth}((M_\mathfrak q \otimes_{A_\mathfrak q} A')_{\mathfrak p''}) =
\text{depth}(M_\mathfrak p \otimes_{A_\mathfrak p} A'_{\mathfrak p''}) =
\text{depth}(M_\mathfrak p)
$$
par platitude de $A \to A'$ et par notre choix de $\mathfrak p''$, voir
Algèbre, lemme \ref{algebra-lemma-apply-grothendieck-module}.
Comme $\mathfrak p''$ est au-dessus de $\mathfrak p'$, on a
$V(\mathfrak p'') \cap V(I') = \{\mathfrak m'\}$. Ainsi, la
condition (6) de la situation \ref{situation-bootstrap} entraîne que
$\text{depth}(M'_{\mathfrak p''}) > s$, ce qui achève la démonstration.
\end{proof}

\noindent
Le lemme fastidieux suivant explique les relations entre les diverses
familles de conditions que l'on pourrait imposer.

\begin{lemma}
\label{lemma-bootstrap-inherited}
Dans la situation \ref{situation-bootstrap}, on a
\begin{enumerate}
\item[(E)] si $T' \subset T$ est un sous-ensemble plus petit, stable par spécialisation, alors
$A, I, T', M$ satisfont les hypothèses de la situation \ref{situation-bootstrap},
\item[(F)] si $S \subset A$ est une partie multiplicative, alors
$S^{-1}A, S^{-1}I, T', S^{-1}M$
satisfont les hypothèses de la situation \ref{situation-bootstrap},
où $T' \subset V(S^{-1}I)$ est l'image réciproque de $T$,
\item[(G)] le quadruplet $A', I', T', M'$
satisfait les hypothèses de la situation \ref{situation-bootstrap},
où $A', I', M'$ sont les complétions $I$-adiques usuelles de $A, I, M$
et $T' \subset V(I')$ est l'image réciproque de $T$.
\end{enumerate}
Soit $I \subset \mathfrak a \subset A$ un idéal tel que
$V(\mathfrak a) \subset T$. Alors
\begin{enumerate}
\item[(A)] si $I$ est contenu dans le radical de Jacobson de $A$,
alors toutes les hypothèses des
lemmes \ref{lemma-kill-colimit-weak-general} et
\ref{lemma-kill-colimit-support-general} sont satisfaites
pour $A, I, \mathfrak a, M$,
\item[(B)] si $A$ est complet pour la topologie $I$-adique, alors
toutes les hypothèses, sauf éventuellement (5), du
lemme \ref{lemma-algebraize-local-cohomology-general}
sont satisfaites pour $A, I, \mathfrak a, M$,
\item[(C)] si $A$ est local d'idéal maximal $\mathfrak m = \mathfrak a$,
alors toutes les hypothèses des
lemmes \ref{lemma-kill-colimit-weak} et \ref{lemma-kill-colimit-support}
sont satisfaites pour $A, \mathfrak m, I, M$,
\item[(D)] si $A$ est local d'idéal maximal $\mathfrak m = \mathfrak a$
et complet pour la topologie $I$-adique, alors toutes les hypothèses du
lemme \ref{lemma-algebraize-local-cohomology}
sont satisfaites pour $A, \mathfrak m, I, M$,
\end{enumerate}
\end{lemma}

\begin{proof}
Démonstration de (E). Il faut montrer que les hypothèses (1), (3), (4), (6)
de la situation \ref{situation-bootstrap} sont satisfaites pour
$A, I, T, M$. Remplacer $T$ par $T'$
affaiblit l'hypothèse (6) et renforce l'hypothèse (4). Cependant, si l'on a
$\mathfrak p \subset \mathfrak r \subset \mathfrak q$ avec
$\mathfrak p \not \in V(I)$, $\mathfrak r \in V(I) \setminus T'$,
$\mathfrak q \in T'$ comme dans l'hypothèse (4) pour $A, I, T', M$, alors
soit on peut choisir $\mathfrak r \in V(I) \setminus T$ et
la condition (4) pour $A, I, T, M$ s'applique, soit on ne peut
trouver un tel $\mathfrak r$, auquel cas on obtient
$\text{depth}(M_\mathfrak p) > s$ par le lemme \ref{lemma-helper-bootstrap}.
Cela montre, comme voulu, que (4) est satisfaite pour $A, I, T', M$.

\medskip\noindent
Démonstration de (F). C'est immédiat et nous omettons les détails.

\medskip\noindent
Démonstration de (G). Il faut montrer que les hypothèses (1), (3), (4), (6)
de la situation \ref{situation-bootstrap} sont satisfaites pour les complétions $I$-adiques
$A', I', T', M'$. Gardons à l'esprit que
$\Spec(A') \to \Spec(A)$ induit un isomorphisme $V(I') \to V(I)$.

\medskip\noindent
Hypothèse (1) : L'anneau $A'$ possède un complexe dualisant, voir
Complexes dualisants, lemme \ref{dualizing-lemma-ubiquity-dualizing}.

\medskip\noindent
Hypothèse (3) : Comme $I' = IA'$, cela résulte de Cohomologie locale,
lemme \ref{local-cohomology-lemma-cd-change-rings}.

\medskip\noindent
Hypothèse (4) : Si l'on a des idéaux premiers
$\mathfrak p' \subset \mathfrak r' \subset \mathfrak q'$ de $A'$
tels que $\mathfrak p' \not \in V(I')$,
$\mathfrak r' \in V(I') \setminus T'$,
$\mathfrak q' \in T'$, alors leurs images
$\mathfrak p \subset \mathfrak r \subset \mathfrak q$ dans
le spectre de $A$
satisfont
$\mathfrak p \not \in V(I)$, $\mathfrak r \in V(I) \setminus T$,
$\mathfrak q \in T$.
On a alors
$$
\text{depth}_{A_\mathfrak p}(M_\mathfrak p) \geq s
\quad\text{ou}\quad
\text{depth}_{A_\mathfrak p}(M_\mathfrak p) +
\dim((A/\mathfrak p)_\mathfrak q) > d + s
$$
d'après l'hypothèse (4) pour $A, I, T, M$. On a
$\text{depth}(M'_{\mathfrak p'}) \geq \text{depth}(M_\mathfrak p)$ et
$\text{depth}(M'_{\mathfrak p'}) +
\dim((A'/\mathfrak p')_{\mathfrak q'}) =
\text{depth}(M_\mathfrak p) +
\dim((A/\mathfrak p)_\mathfrak q)$
d'après Cohomologie locale, lemme \ref{local-cohomology-lemma-change-completion}.
Ainsi, l'hypothèse (4) est satisfaite pour $A', I', T', M'$.

\medskip\noindent
Hypothèse (6) : Soit $\mathfrak q' \in T'$ au-dessus de l'idéal
premier $\mathfrak q \in T$. Alors $A'_{\mathfrak q'}$
et $A_\mathfrak q$ ont des complétions $I$-adiques isomorphes,
de même que $M_\mathfrak q$ et $M'_{\mathfrak q'}$.
Ainsi, l'hypothèse (6) pour $A', I', T', M'$ équivaut
à l'hypothèse (6) pour $A, I, T, M$.

\medskip\noindent
Démonstration de (A). Il faut vérifier les conditions (1), (2), (3), (4) et (6)
des lemmes \ref{lemma-kill-colimit-weak-general} et
\ref{lemma-kill-colimit-support-general} pour
$(A, I, \mathfrak a, M)$. Attention : l'ensemble $T$ de l'énoncé de
ces lemmes n'est pas le même que l'ensemble $T$ ci-dessus.

\medskip\noindent
Condition (1) : Cela vaut parce que nous avons supposé que $A$ possède un complexe dualisant dans la
situation \ref{situation-bootstrap}.

\medskip\noindent
Condition (2) : Elle est vide.

\medskip\noindent
Condition (3) : Soit $\mathfrak p \subset A$ tel que
$V(\mathfrak p) \cap V(I) \subset V(\mathfrak a)$.
Comme $I$ est contenu dans le radical de Jacobson de $A$, on voit
que $V(\mathfrak p) \cap V(I) \not = \emptyset$.
Soit $\mathfrak q \in V(\mathfrak p) \cap V(I)$ un point générique.
Puisque $\text{cd}(A_\mathfrak q, I_\mathfrak q) \leq d$
(Cohomologie locale, lemme \ref{local-cohomology-lemma-cd-local}) et puisque
$V(\mathfrak p A_\mathfrak q) \cap V(I_\mathfrak q) =
\{\mathfrak q A_\mathfrak q\}$, on obtient
$\dim((A/\mathfrak p)_\mathfrak q) \leq d$ par Cohomologie locale,
lemme \ref{local-cohomology-lemma-cd-bound-dim-local}, ce qui prouve (3).

\medskip\noindent
Condition (4) : Supposons $\mathfrak p \not \in V(I)$ et
$\mathfrak q \in V(\mathfrak p) \cap V(\mathfrak a)$.
Il suffit de montrer
$$
\text{depth}_{A_\mathfrak p}(M_\mathfrak p) \geq s
\quad\text{ou}\quad
\text{depth}_{A_\mathfrak p}(M_\mathfrak p) +
\dim((A/\mathfrak p)_\mathfrak q) > d + s
$$
S'il existe un idéal premier $\mathfrak p \subset \mathfrak r \subset \mathfrak q$
tel que $\mathfrak r \in V(I) \setminus T$, cela résulte alors
immédiatement de l'hypothèse (4) de la situation \ref{situation-bootstrap}.
Sinon, $\text{depth}(M_\mathfrak p) > s$ d'après le
lemme \ref{lemma-helper-bootstrap}.

\medskip\noindent
Condition (6) : Soit $\mathfrak p \not \in V(I)$ tel que
$V(\mathfrak p) \cap V(I) \subset V(\mathfrak a)$.
Comme $I$ est contenu dans le radical de Jacobson de $A$, on voit
que $V(\mathfrak p) \cap V(I) \not = \emptyset$.
Choisissons $\mathfrak q \in V(\mathfrak p) \cap V(I) \subset V(\mathfrak a)$.
Il est clair qu'il n'existe pas d'idéal premier
$\mathfrak p \subset \mathfrak r \subset \mathfrak q$
tel que $\mathfrak r \in V(I) \setminus T$.
D'après le lemme \ref{lemma-helper-bootstrap}, on a
$\text{depth}(M_\mathfrak p) > s$, ce qui prouve (6).

\medskip\noindent
Démonstration de (B). Il faut vérifier les conditions (1), (2), (3), (4) du
lemme \ref{lemma-algebraize-local-cohomology-general}. Attention :
l'ensemble $T$ dans l'énoncé de
ce lemme n'est pas le même que l'ensemble $T$ ci-dessus.

\medskip\noindent
Condition (1) : Cela vaut parce que $A$ est complet et possède un complexe dualisant.

\medskip\noindent
Condition (2) : Elle est vide.

\medskip\noindent
Condition (3) : C'est l'hypothèse (3) de la
situation \ref{situation-bootstrap}.

\medskip\noindent
Condition (4) : C'est l'hypothèse (4) du
lemme \ref{lemma-kill-colimit-weak-general}, que nous avons démontrée dans (A).

\medskip\noindent
Démonstration de (C). Cela vaut parce que les hypothèses des
lemmes \ref{lemma-kill-colimit-weak} et \ref{lemma-kill-colimit-support}
sont les mêmes que celles des
lemmes \ref{lemma-kill-colimit-weak-general} et
\ref{lemma-kill-colimit-support-general} dans le cas local,
et nous avons montré dans (A) qu'elles sont satisfaites.

\medskip\noindent
Démonstration de (D). Cela vaut parce que les hypothèses du
lemme \ref{lemma-algebraize-local-cohomology}
sont les hypothèses (1), (2), (3), (4) du
lemme \ref{lemma-algebraize-local-cohomology-general},
et nous avons montré dans (B) qu'elles sont satisfaites.
\end{proof}

\begin{lemma}
\label{lemma-algebraize-local-cohomology-bis}
Dans la situation \ref{situation-bootstrap}, supposons que $A$ est local d'idéal
maximal $\mathfrak m$ et que $T = \{\mathfrak m\}$. Alors
$H^i_\mathfrak m(M) \to \lim H^i_\mathfrak m(M/I^nM)$
est un isomorphisme pour $i \leq s$, et ces modules sont
annulés par une puissance de $I$.
\end{lemma}

\begin{proof}
Soient $A', I', \mathfrak m', M'$ les complétions $I$-adiques usuelles
de $A, I, \mathfrak m, M$. Rappelons que l'on a
$H^i_\mathfrak m(M) \otimes_A A' = H^i_{\mathfrak m'}(M')$
par platitude de $A \to A'$ et par Complexes dualisants, lemme
\ref{dualizing-lemma-torsion-change-rings}.
Comme $H^i_\mathfrak m(M)$ est de torsion annulée par une puissance de $\mathfrak m$, on a
$H^i_\mathfrak m(M) = H^i_\mathfrak m(M) \otimes_A A'$ ; voir
Compléments d'algèbre, lemme \ref{more-algebra-lemma-neighbourhood-equivalence}.
On en conclut que $H^i_\mathfrak m(M) = H^i_{\mathfrak m'}(M')$.
Les mêmes arguments montrent exactement que
$H^i_\mathfrak m(M/I^nM) =  H^i_{\mathfrak m'}(M'/(I')^nM')$
pour tous $n$ et $i$.

\medskip\noindent
Les lemmes \ref{lemma-algebraize-local-cohomology},
\ref{lemma-kill-colimit-weak} et
\ref{lemma-kill-colimit-support}
s'appliquent à $A', \mathfrak m', I', M'$ d'après les parties (C) et (D) du
lemme \ref{lemma-bootstrap-inherited}.
On obtient donc un isomorphisme
$$
H^i_{\mathfrak m'}(M') \longrightarrow H^i(R\Gamma_{\mathfrak m'}(M')^\wedge)
$$
pour $i \leq s$, où ${}^\wedge$ désigne la complétion $I'$-adique dérivée, et ces
modules sont annulés par une puissance de $I'$.
D'après le lemme \ref{lemma-local-cohomology-derived-completion},
on obtient des isomorphismes
$$
H^i_{\mathfrak m'}(M') \longrightarrow
\lim H^i_{\mathfrak m'}(M'/(I')^nM'))
$$
pour $i \leq s$. En les combinant à la comparaison déjà établie
avec la cohomologie locale sur $A$, on conclut que le lemme est vrai.
\end{proof}

\begin{lemma}
\label{lemma-bootstrap-bis-bis}
Soient $I \subset \mathfrak a$ des idéaux d'un anneau noethérien $A$.
Soit $M$ un $A$-module de type fini. Soient $s$ et $d$ des entiers.
Supposons que
\begin{enumerate}
\item[(a)] $A$ possède un complexe dualisant,
\item[(b)] $\text{cd}(A, I) \leq d$,
\item[(c)] si $\mathfrak p \not \in V(I)$ et
$\mathfrak q \in V(\mathfrak p) \cap V(\mathfrak a)$, alors
$\text{depth}_{A_\mathfrak p}(M_\mathfrak p) > s$ ou
$\text{depth}_{A_\mathfrak p}(M_\mathfrak p) +
\dim((A/\mathfrak p)_\mathfrak q) > d + s$.
\end{enumerate}
Alors $A, I, V(\mathfrak a), M, s, d$ sont comme dans la
situation \ref{situation-bootstrap}.
\end{lemma}

\begin{proof}
Il faut montrer que les hypothèses (1), (3), (4) et (6) de la
situation \ref{situation-bootstrap} sont satisfaites.
Il est clair que (a) $\Rightarrow$ (1),
(b) $\Rightarrow$ (3) et (c) $\Rightarrow$ (4).
Pour achever la démonstration, nous montrons dans le paragraphe suivant que (6) est satisfaite.

\medskip\noindent
Soit $\mathfrak q \in V(\mathfrak a)$.
Notons $A', I', \mathfrak m', M'$
les complétions $I$-adiques de
$A_\mathfrak q, I_\mathfrak q, \mathfrak qA_\mathfrak q, M_\mathfrak q$.
Soit $\mathfrak p' \subset A'$ un idéal premier non maximal tel que
$V(\mathfrak p') \cap V(I') = \{\mathfrak m'\}$.
Remarquons que cela entraîne $\dim(A'/\mathfrak p') \leq d$
d'après Cohomologie locale, lemme \ref{local-cohomology-lemma-cd-bound-dim-local}.
Notons $\mathfrak p \subset A$ l'image de $\mathfrak p'$.
On a
$\text{depth}(M'_{\mathfrak p'}) \geq \text{depth}(M_\mathfrak p)$ et
$\text{depth}(M'_{\mathfrak p'}) +
\dim(A'/\mathfrak p') =
\text{depth}(M_\mathfrak p) +
\dim((A/\mathfrak p)_\mathfrak q)$
d'après Cohomologie locale, lemme \ref{local-cohomology-lemma-change-completion}.
D'après l'hypothèse (c), ou bien on a
$\text{depth}(M'_{\mathfrak p'}) \geq \text{depth}(M_\mathfrak p) > s$
et c'est terminé, ou bien on a
$\text{depth}(M'_{\mathfrak p'}) +
\dim(A'/\mathfrak p') > s + d$, ce qui entraîne
$\text{depth}(M'_{\mathfrak p'}) > s$ par l'inégalité déjà établie
$\dim(A'/\mathfrak p') \leq d$. Dans les deux cas, on
obtient le résultat voulu.
\end{proof}

\begin{lemma}
\label{lemma-bootstrap}
Dans la situation \ref{situation-bootstrap}, les systèmes projectifs
$\{H^i_T(I^nM)\}_{n \geq 0}$ sont essentiellement nuls pour $i \leq s$.
De plus, il existe un entier $m_0$ tel que, pour tout
$m \geq m_0$, il existe un entier $m'(m) \geq m$ tel que, pour
$k \geq m'(m)$, l'image de
$H^{s + 1}_T(I^kM) \to H^{s + 1}_T(I^mM)$
s'envoie injectivement dans $H^{s + 1}_T(I^{m_0}M)$.
\end{lemma}

\begin{proof}
Fixons $m$. Soit $\mathfrak q \in T$.
D'après les lemmes \ref{lemma-bootstrap-inherited} et
\ref{lemma-algebraize-local-cohomology-bis}
on voit que
$$
H^i_\mathfrak q(M_\mathfrak q)
\longrightarrow
\lim H^i_\mathfrak q(M_\mathfrak q/I^nM_\mathfrak q)
$$
est un isomorphisme pour $i \leq s$. Les systèmes projectifs
$\{H^i_\mathfrak q(I^nM_\mathfrak q)\}_{n \geq 0}$ et
$\{H^i_\mathfrak q(M/I^nM)\}_{n \geq 0}$
satisfont la condition de Mittag-Leffler pour tout $i$ ; voir le
lemme \ref{lemma-ML-local}. En considérant donc le système projectif des
suites exactes longues
$$
0 \to H^0_\mathfrak q(I^nM_\mathfrak q) \to
H^0_\mathfrak q(M_\mathfrak q) \to
H^0_\mathfrak q(M_\mathfrak q/I^nM_\mathfrak q) \to
H^1_\mathfrak q(I^nM_\mathfrak q) \to
H^1_\mathfrak q(M_\mathfrak q) \to \ldots
$$
on conclut (certains détails sont omis) qu'il existe un entier
$m'(m, \mathfrak q) \geq m$ tel que, pour tout $k \geq m'(m, \mathfrak q)$,
l'application
$H^i_\mathfrak q(I^kM_\mathfrak q) \to H^i_\mathfrak q(I^mM_\mathfrak q)$
est nulle pour $i \leq s$, et l'image de
$H^{s + 1}_\mathfrak q(I^kM_\mathfrak q) \to
H^{s + 1}_\mathfrak q(I^mM_\mathfrak q)$
ne dépend pas de $k \geq m'(m, \mathfrak q)$ et
s'envoie injectivement dans $H^{s + 1}_\mathfrak q(M_\mathfrak q)$.

\medskip\noindent
Supposons que l'on puisse montrer que $m'(m, \mathfrak q)$ peut être choisi
indépendamment de $\mathfrak q \in T$.
Alors le lemme résulte immédiatement de
Cohomologie locale, lemmes \ref{local-cohomology-lemma-zero} et
\ref{local-cohomology-lemma-essential-image}.

\medskip\noindent
Soit $\omega_A^\bullet$ un complexe dualisant. Soit
$\delta : \Spec(A) \to \mathbf{Z}$ la fonction de dimension
correspondante. Rappelons que $\delta$ ne prend qu'un
nombre fini de valeurs ; voir
Complexes dualisants, lemme \ref{dualizing-lemma-universally-catenary}.
Assertion : pour tout $d \in \mathbf{Z}$, l'entier
$m'(m, \mathfrak q)$ peut être choisi indépendamment
de $\mathfrak q \in T$ tel que $\delta(\mathfrak q) = d$.
Cette assertion entraîne clairement le lemme d'après ce qui précède.

\medskip\noindent
Choisissons $\mathfrak q \in T$ tel que $\delta(\mathfrak q) = d$.
Considérons les modules d'Ext
$$
E(n, j) = \text{Ext}^j_A(I^nM, \omega_A^\bullet)
$$
La propriété essentielle que nous utiliserons est que ce sont des $A$-modules de type fini.
Rappelons que $(\omega_A^\bullet)_\mathfrak q[-d]$ est un complexe
dualisant normalisé pour $A_\mathfrak q$, par définition de la
fonction de dimension associée à un complexe dualisant ; voir
Complexes dualisants, section \ref{dualizing-section-dimension-function}.
Le théorème de dualité locale (Complexes dualisants, lemme
\ref{dualizing-lemma-special-case-local-duality}) nous dit que
le complété $\mathfrak qA_\mathfrak q$-adique de
$E(n, -d - i)_\mathfrak q$ est le dual de Matlis de
$H^i_\mathfrak q(I^nM_\mathfrak q)$. Le choix de
$m'(m, \mathfrak q)$ pour $i \leq s$ dans le premier paragraphe nous dit donc que,
pour $k \geq m'(m, \mathfrak q)$ et $j \geq -d - s$, l'application
$$
E(m, j)_\mathfrak q \to E(k, j)_\mathfrak q
$$
est nulle. Puisque ces modules sont de type fini et ne sont non nuls que
pour un nombre fini de valeurs possibles de $j$ (petit détail omis),
on peut trouver un voisinage ouvert $W \subset \Spec(A)$ de $\mathfrak q$
tel que
$$
E(m, j)_{\mathfrak q'} \to E(m'(m, \mathfrak q), j)_{\mathfrak q'}
$$
soit nulle pour $j \geq -d - s$ et tout $\mathfrak q' \in W$.
Alors, bien entendu, les applications $E(m, j)_{\mathfrak q'} \to E(k, j)_{\mathfrak q'}$
pour $k \geq m'(m, \mathfrak q)$ sont elles aussi nulles.

\medskip\noindent
Pour $i = s + 1$, correspondant à $j = - d - s - 1$, la dualité locale
et les résultats du premier paragraphe montrent que
$$
K_{k, \mathfrak q} =
\Ker(E(m, -d - s - 1)_\mathfrak q \to E(k, -d - s - 1)_\mathfrak q)
$$
est indépendant de $k \geq m'(m, \mathfrak q)$ et que
$$
E(0, -d - s - 1)_\mathfrak q \to
E(m, -d - s - 1)_\mathfrak q/K_{m'(m, \mathfrak q), \mathfrak q}
$$
est surjective. Pour $k \geq m'(m, \mathfrak q)$, posons
$$
K_k = \Ker(E(m, -d - s - 1) \to E(k, -d - s - 1))
$$
Puisque $K_k$ est une suite croissante de sous-modules du module de type fini
$E(m, -d - s - 1)$, on voit que, quitte à augmenter
un peu $m'(m, \mathfrak q)$, on peut supposer
$K_{m'(m, \mathfrak q)} = K_k$ pour $k \geq m'(m, \mathfrak q)$.
Après avoir encore rétréci $W$ si nécessaire, on peut aussi supposer que
$$
E(0, -d - s - 1)_{\mathfrak q'} \to
E(m, -d - s - 1)_{\mathfrak q'}/K_{m'(m, \mathfrak q), \mathfrak q'}
$$
est surjective pour tout $\mathfrak q' \in W$ (comme précédemment, on utilise que
ces modules sont de type fini
et que l'application est surjective après localisation en $\mathfrak q$).

\medskip\noindent
Toute partie, en particulier
$T_d = \{\mathfrak q \in T \text{ tel que }\delta(\mathfrak q) = d\}$,
de l'espace topologique noethérien $\Spec(A)$,
munie de la topologie induite, est noethérienne et donc quasi-compacte.
Nous avons vu ci-dessus que, pour tout $\mathfrak q \in T_d$,
il existe un voisinage ouvert $W$ sur lequel
$m'(m, \mathfrak q)$ convient à tout $\mathfrak q' \in T_d \cap W$.
On en conclut qu'il existe un entier $m'(m, d)$ tel que, pour tout
$\mathfrak q \in T_d$, on ait
$$
E(m, j)_\mathfrak q \to E(m'(m, d), j)_\mathfrak q
$$
qui soit nulle pour $j \geq -d - s$ et, en posant
$K_{m'(m, d)} = \Ker(E(m, -d - s - 1) \to E(m'(m, d), -d - s - 1))$,
on ait
$$
K_{m'(m, d), \mathfrak q} =
\Ker(E(m, -d - s - 1)_{\mathfrak q} \to E(k, -d - s - 1)_{\mathfrak q})
$$
pour tout $k \geq m'(m, d)$, et que l'application
$$
E(0, -d - s - 1)_\mathfrak q \to
E(m, -d - s - 1)_\mathfrak q/K_{m'(m, d), \mathfrak q}
$$
soit surjective. En utilisant de nouveau le théorème de dualité locale (dans le sens
opposé), on en conclut que l'assertion est vraie. Cela achève la démonstration.
\end{proof}

\begin{lemma}
\label{lemma-final-bootstrap}
Dans la situation \ref{situation-bootstrap}, il existe un entier $m_0 \geq 0$
tel que
\begin{enumerate}
\item $H^i_T(M) \to H^i_T(M/I^nM)$ soit injective pour tout $i \leq s$
et tout $n \geq m_0$,
\item $\{H^i_T(M/I^nM)\}_{n \geq 0}$
vérifie la condition de Mittag-Leffler pour $i < s$,
\item $\{H^i_T(I^{m_0}M/I^nM)\}_{n \geq m_0}$
vérifie la condition de Mittag-Leffler pour $i \leq s$,
\item $H^i_T(M) \to \lim H^i_T(M/I^nM)$
soit un isomorphisme pour $i < s$,
\item $H^s_T(I^{m_0}M) \to \lim H^s_T(I^{m_0}M/I^nM)$
soit un isomorphisme pour $i \leq s$,
\item $H^s_T(M) \to \lim H^s_T(M/I^nM)$ soit
injective, de conoyau annulé par $I^{m_0}$, et
\item $R^1\lim H^s_T(M/I^nM)$ soit annulé par $I^{m_0}$.
\end{enumerate}
\end{lemma}

\begin{proof}
Pour démontrer (1), soit $0 \leq i \leq s$ et choisissons des entiers $m > n > 0$ tels que
$H^i_T(I^mM) \to H^i_T(I^nM)$ soit nulle ; c'est possible d'après le
lemme \ref{lemma-bootstrap}. Le diagramme commutatif à lignes exactes
$$
\xymatrix{
H^i_T(I^nM) \ar[r] &
H^i_T(M) \ar[r] &
H^i_T(M/I^nM) \\
H^i_T(I^mM) \ar[r] \ar[u]^0 &
H^i_T(M) \ar[r] \ar[u]^1 &
H^i_T(M/I^mM) \ar[u] \\
}
$$
montre que l'application $H^i_T(M) \to H^i_T(M/I^mM)$ est injective. Ainsi, (1) vaut
pour tout $m_0$ suffisamment grand.

\medskip\noindent
Considérons les suites exactes longues
$$
0 \to H^0_T(I^nM) \to H^0_T(M) \to
H^0_T(M/I^nM) \to H^1_T(I^nM) \to
H^1_T(M) \to \ldots
$$
Les assertions (2) et (4) en résultent avec le lemme \ref{lemma-bootstrap}.

\medskip\noindent
Soient $m_0$ et $m'(-)$ comme dans le lemme \ref{lemma-bootstrap}.
Pour $m \geq m_0$, considérons la suite exacte longue
$$
H^s_T(I^mM) \to H^s_T(I^{m_0}M) \to
H^s_T(I^{m_0}M/I^mM) \to H^{s + 1}_T(I^mM) \to
H^1_T(I^{m_0}M)
$$
Alors, pour $k \geq m'(m)$, l'image de
$H^{s + 1}_T(I^kM) \to H^{s + 1}_T(I^mM)$
se plonge dans $H^{s + 1}_T(I^{m_0}M)$.
Par conséquent, l'image de
$H^s_T(I^{m_0}M/I^kM) \to H^s_T(I^{m_0}M/I^mM)$
a une image nulle dans $H^{s + 1}_T(I^mM)$ pour tout $k \geq m'(m)$.
On en conclut que (3) et (5) sont vraies.

\medskip\noindent
Considérons les suites exactes courtes
$0 \to I^{m_0}M \to M \to M/I^{m_0} M \to 0$ et
$0 \to I^{m_0}M/I^nM \to M/I^nM \to M/I^{m_0} M \to 0$.
On obtient un diagramme
$$
\xymatrix{
H^{s - 1}_T(M/I^{m_0}M) \ar[r] &
\lim H^s_T(I^{m_0}M/I^nM) \ar[r] &
\lim H^s_T(M/I^nM) \ar[r] &
H^s_T(M/I^{m_0}M) \\
H^{s - 1}_T(M/I^{m_0}M) \ar[r] \ar@{=}[u] &
H^s_T(I^{m_0}M) \ar[r] \ar[u]_{\cong} &
H^s_T(M) \ar[r] \ar[u] &
H^s_T(M/I^{m_0}M) \ar@{=}[u]
}
$$
dont la ligne inférieure est exacte. La ligne supérieure est elle aussi exacte
(aux deux places médianes) d'après
Homologie, lemme \ref{homology-lemma-apply-Mittag-Leffler}.
L'assertion (6) en résulte.

\medskip\noindent
Posons $B_n = H^s_T(M/I^nM)$. Soit $A_n \subset B_n$
l'image de $H^s_T(I^{m_0}M/I^nM) \to H^s_T(M/I^nM)$.
Alors $(A_n)$ vérifie la condition de Mittag-Leffler d'après (3) et
Homologie, lemme \ref{homology-lemma-Mittag-Leffler}.
De plus, $C_n = B_n/A_n$ est annulé par $I^{m_0}$. Ainsi,
$R^1\lim B_n \cong R^1\lim C_n$ est annulé par $I^{m_0}$, ce qui donne (7).
\end{proof}

\begin{theorem}
\label{theorem-final-bootstrap}
Dans la situation \ref{situation-bootstrap}, pour $i \leq s$, le système projectif
$\{H^i_T(M/I^nM)\}_{n \geq 0}$ est essentiellement constant de valeur $H^i_T(M)$.
En particulier, $\{H^i_T(M/I^nM)\}_{n \geq 0}$ est de Mittag-Leffler,
$H^i_T(M)$ est annulé par une puissance de $I$, et
$H^i_T(M) = \lim H^i_T(M/I^nM)$.
\end{theorem}

\begin{proof}
Soit $0 \leq i \leq s$.
Si l'on montre que $\{H^i_T(M/I^nM)\}_{n \geq 0}$ est de Mittag-Leffler et que
$H^i_T(M) = \lim H^i_T(M/I^nM)$, alors l'injectivité de
$H^i_T(M) \to H^i_T(M/I^nM)$ pour $n \gg 0$, établie dans le
lemme \ref{lemma-final-bootstrap}, entraînera que
$\{H^i_T(M/I^nM)\}_{n \geq 0}$ est essentiellement constant de valeur $H^i_T(M)$.
Nous omettons un petit détail.

\medskip\noindent
Nous savons déjà que $\{H^i_T(M/I^nM)\}_{n \geq 0}$ est de Mittag-Leffler et que
$H^i_T(M) = \lim H^i_T(M/I^nM)$ pour $i < s$ d'après le
lemme \ref{lemma-final-bootstrap}. Nous savons aussi que ces assertions valent
pour $i = s$ et le module $I^{m_0}M$, pour un certain $m_0 > 0$.
Pour achever la démonstration, nous allons montrer qu'en fait ces
assertions pour $i = s$ valent pour $M$.

\medskip\noindent
Posons $M' = H^0_I(M)$ et $M'' = M/M'$, de sorte que l'on ait une suite
exacte courte
$$
0 \to M' \to M \to M'' \to 0
$$
et $M''$ vérifie $H^0_I(M') = 0$ d'après
Complexes dualisants, lemme \ref{dualizing-lemma-divide-by-torsion}.
D'après Artin-Rees (Algèbre, lemme \ref{algebra-lemma-Artin-Rees}),
on obtient des suites exactes courtes
$$
0 \to M' \to M/I^n M \to M''/I^n M'' \to 0
$$
pour $n$ assez grand. Considérons les suites exactes longues
$$
H^s_T(M') \to
H^s_T(M/I^nM) \to
H^s_T(M''/I^nM'') \to
H^{s + 1}_T(M')
$$
Il est alors aisé de voir que, si le système projectif
$\{H^s_T(M''/I^nM'')\}_{n \geq 0}$ est de Mittag-Leffler,
le système projectif
$\{H^s_T(M/I^nM)\}_{n \geq 0}$ l'est aussi.
(Remarquons que la condition ML pour un système projectif de groupes $G_n$
ne dépend que des valeurs du système projectif pour $n$ assez grand.)
De plus, la suite
$$
H^s_T(M') \to
\lim H^s_T(M/I^nM) \to
\lim H^s_T(M''/I^nM'') \to
H^{s + 1}_T(M')
$$
est exacte parce que la condition ML est satisfaite aux endroits voulus ; voir
Homologie, lemme \ref{homology-lemma-apply-Mittag-Leffler}.
Ainsi, si $H^s_T(M'') \to \lim H^s_T(M''/I^nM'')$
est un isomorphisme, alors
$H^s_T(M) \to \lim H^s_T(M/I^nM)$
est aussi un isomorphisme d'après le lemme des cinq
(Homologie, lemme \ref{homology-lemma-five-lemma}).
On est ainsi ramené au cas traité dans le paragraphe suivant.

\medskip\noindent
Supposons que $H^0_I(M) = 0$. Choisissons des générateurs
$f_1, \ldots, f_r$ de $I^{m_0}$, où $m_0$ est l'entier
obtenu pour $M$ dans le lemme \ref{lemma-final-bootstrap}.
Considérons alors la suite exacte
$$
0 \to M \xrightarrow{f_1, \ldots, f_r}
(I^{m_0}M)^{\oplus r} \to Q \to 0
$$
qui définit $Q$. Faisons quelques remarques : la première application est injective
précisément parce que $H^0_I(M) = 0$. Le conoyau $Q$ de cette injection
est un $A$-module de type fini tel que, pour tout $1 \leq j \leq r$,
on ait $Q_{f_j} \cong (M_{f_j})^{\oplus r - 1}$.
En particulier, pour un idéal premier $\mathfrak p \subset A$
tel que $\mathfrak p \not \in V(I)$, on a
$Q_\mathfrak p \cong (M_\mathfrak p)^{\oplus r - 1}$.
De même, étant donné $\mathfrak q \in T$ et
$\mathfrak p' \subset A' = (A_\mathfrak q)^\wedge$
non contenu dans $V(IA')$, on a
$Q'_{\mathfrak p'} \cong (M'_{\mathfrak p'})^{\oplus r - 1}$
où $Q' = (Q_\mathfrak q)^\wedge$ et $M' = (M_\mathfrak q)^\wedge$.
Ainsi, les conditions de la situation \ref{situation-bootstrap}
sont satisfaites pour $A, I, T, Q$. (Remarquons que $Q$ peut avoir
un $H^0_I(Q)$ non nul, mais cela n'aura pas d'importance.)

\medskip\noindent
Pour tout $n \geq 0$, posons $F^nM = M \cap I^n(I^{m_0}M)^{\oplus r}$,
de sorte que l'on obtienne des suites exactes courtes
$$
0 \to F^nM \to I^n(I^{m_0}M)^{\oplus r} \to I^nQ \to 0
$$
D'après Artin-Rees (Algèbre, lemme \ref{algebra-lemma-Artin-Rees}),
il existe un $c \geq 0$ tel que
$I^n M \subset F^nM \subset I^{n - c}M$ pour tout $n \geq c$.
Soient $m_0$ l'entier et $m'(m)$
la fonction définie pour $m \geq m_0$,
obtenus dans le lemme \ref{lemma-bootstrap}
appliqué à $M$. Remarquons que l'entier $m_0$
est le même que notre entier $m_0$ choisi ci-dessus (il n'est pas nécessaire de
le vérifier : on peut simplement prendre le maximum des deux entiers si
l'on préfère). Enfin, d'après le lemme \ref{lemma-bootstrap}
appliqué à $Q$, pour tout entier $m$, il existe un entier
$m''(m) \geq m$ tel que $H^s_T(I^kQ) \to H^s_T(I^mQ)$
soit nulle pour tout $k \geq m''(m)$.

\medskip\noindent
Fixons $m \geq m_0$. Choisissons $k \geq m'(m''(m + c))$.
Choisissons $\xi \in H^{s + 1}_T(I^kM)$
qui s'envoie sur zéro dans $H^{s + 1}_T(M)$.
Nous voulons montrer que $\xi$ s'envoie sur zéro dans $H^{s + 1}_T(I^mM)$.
En effet, cela montrera que $\{H^s_T(M/I^nM)\}_{n \geq 0}$
est de Mittag-Leffler, exactement comme dans la démonstration du lemme \ref{lemma-final-bootstrap}.
Voici un diagramme qui permet de visualiser l'argument :
$$
\xymatrix{
&
H^{s + 1}_T(I^kM) \ar[r] \ar[d] &
H^{s + 1}_T(I^k(I^{m_0}M)^{\oplus r}) \ar[d] &
\\
H^s_T(I^{m''(m + c)}Q) \ar[r]_-\delta \ar[d] &
H^{s + 1}_T(F^{m''(m + c)}M) \ar[r] \ar[d] &
H^{s + 1}_T(I^{m''(m + c)}(I^{m_0}M)^{\oplus r}) \\
H^s_T(I^{m + c}Q) \ar[r] &
H^{s + 1}_T(F^{m + c}M) \ar[d] &
\\
&
H^{s + 1}_T(I^mM)
}
$$
L'image de $\xi$ dans $H^{s + 1}_T(I^k(I^{m_0}M)^{\oplus r})$
s'envoie sur zéro dans $H^{s + 1}_T((I^{m_0}M)^{\oplus r})$
et s'envoie donc sur zéro dans
$H^{s + 1}_T(I^{m''(m + c)}(I^{m_0}M)^{\oplus r})$
par le choix de $m'(-)$.
Ainsi, l'image $\xi' \in H^{s + 1}_T(F^{m''(m + c)}M)$
s'envoie sur zéro dans $H^{s + 1}_T(I^{m''(m + c)}(I^{m_0}M)^{\oplus r})$
et, par conséquent, $\xi' = \delta(\eta)$ pour un certain
$\eta \in H^s_T(I^{m''(m + c)}Q)$.
D'après notre choix de $m''(-)$, on voit que $\eta$ s'envoie sur
zéro dans $H^s_T(I^{m + c}Q)$.
Cela signifie à son tour que $\xi'$ s'envoie sur zéro dans
$H^{s + 1}_T(F^{m + c}M)$.
Comme $F^{m + c}M \subset I^mM$, on conclut.

\medskip\noindent
Enfin, démontrons l'assertion sur les limites. Considérons les suites
exactes courtes
$$
0 \to M/F^nM \to (I^{m_0}M)^{\oplus r}/I^n (I^{m_0}M)^{\oplus r}
\to Q/I^nQ \to 0
$$
On a $\lim H^s_T(M/I^nM) = \lim H^s_T(M/F^nM)$,
car ces systèmes projectifs sont pro-isomorphes. On obtient un diagramme commutatif
$$
\xymatrix{
H^{s - 1}_T(Q) \ar[r] \ar[d] &
\lim H^{s - 1}_T(Q/I^nQ) \ar[d] \\
H^s_T(M) \ar[r] \ar[d] &
\lim H^s_T(M/I^nM) \ar[d] \\
H^s_T((I^{m_0}M)^{\oplus r}) \ar[r] \ar[d] &
\lim H^s_T((I^{m_0}M)^{\oplus r}/I^n(I^{m_0}M)^{\oplus r}) \ar[d] \\
H^s_T(Q) \ar[r] &
\lim H^s_T(Q/I^nQ)
}
$$
La colonne de droite est exacte parce que la condition ML est satisfaite aux endroits voulus ; voir
Homologie, lemme \ref{homology-lemma-apply-Mittag-Leffler}.
La flèche horizontale la plus basse est injective (!) d'après la
partie (6) du lemme \ref{lemma-final-bootstrap}.
La flèche horizontale située au-dessus est bijective d'après la
partie (5) du lemme \ref{lemma-final-bootstrap}.
Les flèches en degrés cohomologiques $\leq s - 1$ sont des isomorphismes.
On conclut donc que $H^s_T(M) \to \lim H^s_T(M/I^nM)$
est un isomorphisme d'après le lemme des cinq
(Homologie, lemme \ref{homology-lemma-five-lemma}).
Cela achève la démonstration du théorème.
\end{proof}

\begin{lemma}
\label{lemma-combine-two}
Soient $I \subset \mathfrak a \subset A$ des idéaux d'un anneau noethérien $A$,
et soit $M$ un $A$-module de type fini. Soient $s$ et $d$ des entiers.
Supposons que
\begin{enumerate}
\item $A, I, V(\mathfrak a), M$ satisfont les hypothèses de la
situation \ref{situation-bootstrap} pour $s$ et $d$, et
\item $A, I, \mathfrak a, M$ satisfont les conditions du
lemme \ref{lemma-algebraize-local-cohomology-general}
pour $s + 1$ et $d$, avec $J = \mathfrak a$.
\end{enumerate}
Alors il existe un idéal
$J_0 \subset \mathfrak a$ avec $V(J_0) \cap V(I) = V(\mathfrak a)$
tel que, pour tout $J \subset J_0$ avec $V(J) \cap V(I) = V(\mathfrak a)$,
l'application
$$
H^{s + 1}_J(M) \longrightarrow \lim H^{s + 1}_\mathfrak a(M/I^nM)
$$
soit un isomorphisme.
\end{lemma}

\begin{proof}
En effet, l'existence de $J_0$
et l'isomorphisme
$H^{s + 1}_J(M) = H^{s + 1}(R\Gamma_\mathfrak a(M)^\wedge)$
résultent du lemme \ref{lemma-algebraize-local-cohomology-general} ;
on a une suite exacte courte
$$
0 \to R^1\lim H^s_\mathfrak a(M/I^nM) \to
H^{s + 1}(R\Gamma_\mathfrak a(M)^\wedge) \to
\lim H^{s + 1}_\mathfrak a(M/I^nM) \to 0
$$
d'après Complexes dualisants, lemme \ref{dualizing-lemma-completion-local},
et le module $R^1\lim H^s_\mathfrak a(M/I^nM)$ est nul parce que
$\{H^s_\mathfrak a(M/I^nM)\}_{n \geq 0}$ satisfait la condition de Mittag-Leffler
d'après le théorème \ref{theorem-final-bootstrap}.
\end{proof}








\section{Algébrisation des sections formelles, I}
\label{section-algebraization-sections}

\noindent
Dans cette section, nous étudions le problème de l'algébrisation des
sections formelles dans le cas local.
Soit $(A, \mathfrak m)$ un anneau local noethérien.
Soit $I \subset A$ un idéal. Posons
$$
X = \Spec(A) \supset U = \Spec(A) \setminus \{\mathfrak m\}
$$
et notons $Y = V(I)$ le sous-schéma fermé correspondant à $I$.
Soit $\mathcal{F}$ un $\mathcal{O}_U$-module cohérent.
Dans cette section, nous considérons les limites
$$
\lim_n H^i(U, \mathcal{F}/I^n\mathcal{F})
$$
Ceci est étroitement lié à la cohomologie de l'image inverse
de $\mathcal{F}$ sur la complétion formelle de $U$ le long de $Y$ ;
toutefois, comme nous n'avons pas encore introduit les schémas formels,
nous ne pouvons pas employer ici cette terminologie.

\begin{lemma}
\label{lemma-compare-with-derived-completion}
Soit $U$ le spectre épointé d'un anneau local noethérien $A$.
Soit $\mathcal{F}$ un $\mathcal{O}_U$-module cohérent.
Soit $I \subset A$ un idéal. Alors
$$
H^i(R\Gamma(U, \mathcal{F})^\wedge) =
\lim H^i(U, \mathcal{F}/I^n\mathcal{F})
$$
pour tout $i$, où $R\Gamma(U, \mathcal{F})^\wedge$ désigne
la complétion $I$-adique dérivée.
\end{lemma}

\begin{proof}
D'après les lemmes \ref{lemma-formal-functions-general} et
\ref{lemma-derived-completion-pseudo-coherent}, on a
$$
R\Gamma(U, \mathcal{F})^\wedge =
R\Gamma(U, \mathcal{F}^\wedge) =
R\Gamma(U, R\lim \mathcal{F}/I^n\mathcal{F})
$$
On obtient donc des suites exactes courtes
$$
0 \to R^1\lim H^{i - 1}(U, \mathcal{F}/I^n\mathcal{F}) \to
H^i(R\Gamma(U, \mathcal{F})^\wedge) \to
\lim H^i(U, \mathcal{F}/I^n\mathcal{F}) \to 0
$$
d'après Cohomologie, lemme \ref{cohomology-lemma-RGamma-commutes-with-Rlim}.
Les termes $R^1\lim$ s'annulent parce que les systèmes projectifs de groupes
$H^i(U, \mathcal{F}/I^n\mathcal{F})$ satisfont la condition de Mittag-Leffler
d'après le lemme \ref{lemma-ML-local}.
\end{proof}

\begin{theorem}
\label{theorem-algebraization-formal-sections}
\begin{reference}
La méthode de démonstration suit approximativement celle de la
démonstration de \cite[Theorem 1]{Faltings-algebraisation}
et de \cite[Satz 2]{Faltings-uber}.
Le résultat est presque le même que
\cite[Theorem 1.1]{MRaynaud-paper} (cas du complémentaire affine) et
\cite[Theorem 3.9]{MRaynaud-book} (le complémentaire est réunion de peu d'affines).
\end{reference}
Soit $(A, \mathfrak m)$ un anneau local noethérien qui possède un
complexe dualisant et qui est complet pour la topologie définie par un idéal $I$.
Posons $X = \Spec(A)$, $Y = V(I)$ et $U = X \setminus \{\mathfrak m\}$.
Soit $\mathcal{F}$ un faisceau cohérent sur $U$.
Supposons que
\begin{enumerate}
\item $\text{cd}(A, I) \leq d$, c'est-à-dire que
$H^i(X \setminus Y, \mathcal{G}) = 0$ pour $i \geq d$ et tout
$\mathcal{G}$ quasi-cohérent sur $X$,
\item pour tout $x \in X \setminus Y$ dont l'adhérence $\overline{\{x\}}$
dans $X$ rencontre $U \cap Y$, on ait
$$
\text{depth}_{\mathcal{O}_{X, x}}(\mathcal{F}_x) \geq s
\quad\text{ou}\quad
\text{depth}_{\mathcal{O}_{X, x}}(\mathcal{F}_x)
+ \dim(\overline{\{x\}}) > d + s
$$
\end{enumerate}
Alors il existe un ouvert $V_0 \subset U$ contenant $U \cap Y$
tel que, pour tout ouvert $V \subset V_0$ contenant $U \cap Y$,
l'application
$$
H^i(V, \mathcal{F}) \to \lim H^i(U, \mathcal{F}/I^n\mathcal{F})
$$
soit un isomorphisme pour $i < s$. Si, de plus,
$
\text{depth}_{\mathcal{O}_{X, x}}(\mathcal{F}_x) +
\dim(\overline{\{x\}}) > s
$
pour tout $x \in U \cap Y$, alors ces groupes de cohomologie sont des $A$-modules de type fini.
\end{theorem}

\begin{proof}
Choisissons un $A$-module de type fini $M$ tel que $\mathcal{F}$ soit la
restriction à $U$ du
$\mathcal{O}_X$-module cohérent associé à $M$ ; voir Cohomologie locale,
lemme \ref{local-cohomology-lemma-finiteness-pushforwards-and-H1-local}.
Alors les hypothèses du
lemme \ref{lemma-algebraize-local-cohomology}
sont satisfaites.
Choisissons $J_0$ comme dans ce lemme et posons $V_0 = X \setminus V(J_0)$.
Les ouverts $V \subset V_0$ contenant $U \cap Y$
correspondent alors $1$ à $1$ aux idéaux $J \subset J_0$ tels que
$V(J) \cap V(I) = \{\mathfrak m\}$.
De plus, pour un tel choix, on a un triangle distingué
$$
R\Gamma_J(M) \to M \to R\Gamma(V, \mathcal{F}) \to
R\Gamma_J(M)[1]
$$
On a de même un triangle distingué
$$
R\Gamma_\mathfrak m(M)^\wedge \to
M \to
R\Gamma(U, \mathcal{F})^\wedge \to
R\Gamma_\mathfrak m(M)^\wedge[1]
$$
faisant intervenir les complétions $I$-adiques dérivées.
Les groupes de cohomologie de $R\Gamma(U, \mathcal{F})^\wedge$ sont
égaux aux limites figurant dans l'énoncé du théorème d'après le
lemme \ref{lemma-compare-with-derived-completion}.
L'application canonique entre ces triangles
et quelques arguments simples montrent que notre
théorème résulte du lemme principal \ref{lemma-algebraize-local-cohomology}
(remarquons qu'ici on a $i < s$, tandis que dans le lemme on a
$i \leq s$ ; cela provient du décalage).
Le caractère de type fini des groupes de cohomologie
(sous l'hypothèse supplémentaire) résulte du
lemme \ref{lemma-kill-colimit}.
\end{proof}

\begin{lemma}
\label{lemma-application-theorem}
Soit $(A, \mathfrak m)$ un anneau local noethérien qui possède un
complexe dualisant et qui est complet pour la topologie définie par un idéal $I$.
Posons $X = \Spec(A)$, $Y = V(I)$ et $U = X \setminus \{\mathfrak m\}$.
Soit $\mathcal{F}$ un faisceau cohérent sur $U$.
Supposons que, pour tout point associé $x \in U$ de $\mathcal{F}$,
on ait $\dim(\overline{\{x\}}) > \text{cd}(A, I) + 1$,
où $\overline{\{x\}}$ est l'adhérence dans $X$.
Alors l'application
$$
\colim H^0(V, \mathcal{F})
\longrightarrow
\lim H^0(U, \mathcal{F}/I^n\mathcal{F})
$$
est un isomorphisme de $A$-modules de type fini,
où la limite inductive porte sur les ouverts $V \subset U$
contenant $U \cap Y$.
\end{lemma}

\begin{proof}
Appliquons le théorème \ref{theorem-algebraization-formal-sections} avec $s = 1$
(on obtient également le caractère de type fini).
\end{proof}




\section{Algébrisation des sections formelles, II}
\label{section-algebraization-sections-coherent}

\noindent
Il est assez difficile d'énoncer succinctement toutes les conséquences
possibles des résultats des
sections \ref{section-algebraization-sections-general} et
\ref{section-bootstrap}
pour la cohomologie des faisceaux cohérents sur les schémas quasi-affines
et leur complétion relativement à un idéal.
Cette section donne une liste non exhaustive
d'applications à $H^0$. La section suivante contient des
applications à la cohomologie supérieure.

\begin{lemma}
\label{lemma-ses-H0}
Soient $I \subset \mathfrak a$ des idéaux d'un anneau noethérien $A$.
Soit $0 \to \mathcal{F}' \to \mathcal{F} \to \mathcal{F}'' \to 0$
une suite exacte courte de modules cohérents sur
$U = \Spec(A) \setminus V(\mathfrak a)$. Soit $\mathcal{V}$
l'ensemble des sous-schémas ouverts $V \subset U$ contenant
$U \cap V(I)$, ordonné par inclusion inverse.
Considérons le diagramme commutatif
$$
\xymatrix{
\colim_\mathcal{V} H^0(V, \mathcal{F}') \ar[d] \ar[r] &
\colim_\mathcal{V} H^0(V, \mathcal{F}) \ar[d] \ar[r] &
\colim_\mathcal{V} H^0(V, \mathcal{F}'') \ar[d] \\
\lim H^0(U, \mathcal{F}'/I^n\mathcal{F}') \ar[r] &
\lim H^0(U, \mathcal{F}'/I^n\mathcal{F}) \ar[r] &
\lim H^0(U, \mathcal{F}'/I^n\mathcal{F}'')
}
$$
Si les flèches verticales de gauche et de droite sont des isomorphismes, celle du milieu l'est aussi.
Si les flèches verticales du milieu et de gauche sont des isomorphismes, celle de gauche l'est aussi.
\end{lemma}

\begin{proof}
Les suites du diagramme sont exactes au milieu et la première
flèche est injective. La dernière assertion résulte donc d'une simple
chasse au diagramme. Pour le reste de la démonstration, supposons que les flèches verticales
de gauche et de droite soient des isomorphismes. Une chasse au diagramme montre
que la flèche verticale du milieu est injective. Il ne reste
qu'à montrer qu'elle est surjective.

\medskip\noindent
On peut choisir des $A$-modules de type fini $M$ et $M'$ tels que $\mathcal{F}$
et $\mathcal{F}'$ soient les restrictions
de $\widetilde{M}$ et $\widetilde{M}'$ à $U$ ; voir Cohomologie locale, lemme
\ref{local-cohomology-lemma-finiteness-pushforwards-and-H1-local}.
Après avoir remplacé $M'$ par $\mathfrak a^n M'$ pour un certain $n \geq 0$,
on peut supposer que $\mathcal{F}' \to \mathcal{F}$
correspond à un homomorphisme de modules $M' \to M$ ; voir
Cohomologie des schémas, lemme \ref{coherent-lemma-homs-over-open}.
Après avoir remplacé $M'$ par l'image de $M' \to M$
et posé $M'' = M/M'$, on voit que notre suite exacte courte
correspond à la restriction de la suite exacte courte de
modules cohérents associée à la suite exacte courte
$0 \to M' \to M \to M'' \to 0$ de $A$-modules.

\medskip\noindent
Soit $\hat s \in \lim H^0(U, \mathcal{F}/I^n\mathcal{F})$
d'image $\hat s'' \in \lim H^0(U, \mathcal{F}''/I^n\mathcal{F}'')$.
Par hypothèse, on trouve $V \in \mathcal{V}$
et une section $s'' \in \mathcal{F}''(V)$ qui s'envoie sur $\hat s''$.
Soit $J \subset A$ un idéal tel que $V(J) = \Spec(A) \setminus V$.
D'après Cohomologie des schémas, lemme \ref{coherent-lemma-homs-over-open},
après avoir remplacé $J$ par une puissance, on peut supposer
qu'il existe une application $A$-linéaire $\varphi : J \to M''$
correspondant à $s''$. Fixons ce choix de $J$ ; dans la suite
de la démonstration, nous remplacerons $V$ par un $V$ plus petit dans $\mathcal{V}$,
c'est-à-dire que nous aurons $V \cap V(J) = \emptyset$.

\medskip\noindent
Choisissons une présentation $A^{\oplus m} \to A^{\oplus n} \to J \to 0$.
Notons $g_1, \ldots, g_n \in J$ les images des vecteurs de base
de $A^{\oplus n}$, de sorte que $J = (g_1, \ldots, g_n)$. Soit
$A^{\oplus m} \to A^{\oplus n}$ donnée par la matrice $(a_{ji})$,
de sorte que $\sum a_{ji} g_i = 0$, $j = 1, \ldots, m$.
Comme $M \to M''$ est surjective, pour chaque $i$ on peut choisir $m_i \in M$
qui s'envoie sur $\varphi(g_i) \in M''$. Alors l'élément $g_i \hat s - m_i$
de $\lim H^0(U, \mathcal{F}/I^n\mathcal{F})$ appartient au
sous-module $\lim H^0(U, \mathcal{F}'/I^n\mathcal{F}')$.
Par hypothèse, après avoir rétréci $V$, on peut supposer qu'il existe des
$s'_i \in \mathcal{F}'(V)$, $i = 1, \ldots, n$,
tels que les $s'_i$ s'envoient sur $g_i \hat s - m_i$.
Posons $s_i = s'_i + m_i$ dans $\mathcal{F}(V)$.
Remarquons que $\sum a_{ji} s_i$ s'envoie sur $\sum a_{ji}g_i\hat s = 0$
par l'application
$$
\colim_\mathcal{V} \mathcal{F}(V')
\longrightarrow
\lim H^0(U, \mathcal{F}/I^n\mathcal{F})
$$
Comme cette application est injective (voir ci-dessus), après avoir rétréci $V$ on peut supposer
que $\sum a_{ji}s_i = 0$ dans $\mathcal{F}(V)$ pour tout $j = 1, \ldots, m$.
Il s'ensuit que l'on obtient une application de $A$-modules $J \to \mathcal{F}(V)$
qui envoie $g_i$ sur $s_i$. Par la propriété universelle de $\widetilde{J}$,
cette application de $A$-modules correspond à une application de $\mathcal{O}_V$-modules
$\widetilde{J}|_V \to \mathcal{F}$. Cependant, puisque $V(J) \cap V = \emptyset$,
on a $\widetilde{J}|_V = \mathcal{O}_V$. Nous avons donc construit
une section $s \in \mathcal{F}(V)$. Nous omettons le calcul qui
montre que $s$ s'envoie sur $\hat s$ par l'application affichée ci-dessus.
\end{proof}

\noindent
Le lemme suivant sera supplanté par la
proposition \ref{proposition-application-H0}.

\begin{lemma}
\label{lemma-application-H0-pre}
Soient $I \subset \mathfrak a$ des idéaux d'un anneau noethérien $A$.
Soit $\mathcal{F}$ un module cohérent sur
$U = \Spec(A) \setminus V(\mathfrak a)$.
Supposons que
\begin{enumerate}
\item $A$ soit complet pour la topologie $I$-adique et possède un complexe dualisant,
\item si $x \in \text{Ass}(\mathcal{F})$, $x \not \in V(I)$,
$\overline{\{x\}} \cap V(I) \not \subset V(\mathfrak a)$,
et $z \in \overline{\{x\}} \cap V(\mathfrak a)$, alors
$\dim(\mathcal{O}_{\overline{\{x\}}, z}) > \text{cd}(A, I) + 1$,
\item l'une des conditions suivantes soit satisfaite :
\begin{enumerate}
\item la restriction de $\mathcal{F}$ à $U \setminus V(I)$ est $(S_1)$,
\item la dimension de $V(\mathfrak a)$ est au plus $2$\footnote{Au
sens où la différence entre les valeurs maximale et minimale
sur $V(\mathfrak a)$ d'une fonction de dimension sur $\Spec(A)$ est au plus $2$.}.
\end{enumerate}
\end{enumerate}
Alors on obtient un isomorphisme
$$
\colim H^0(V, \mathcal{F})
\longrightarrow
\lim H^0(U, \mathcal{F}/I^n\mathcal{F})
$$
où la limite inductive porte sur les ouverts $V \subset U$ contenant $U \cap V(I)$.
\end{lemma}

\begin{proof}
Choisissons un $A$-module de type fini $M$ tel que $\mathcal{F}$ soit la restriction
à $U$ du module cohérent associé à $M$ ; voir Cohomologie locale,
lemme \ref{local-cohomology-lemma-finiteness-pushforwards-and-H1-local}.
Posons $d = \text{cd}(A, I)$.
Soit $\mathfrak p$ un idéal premier de $A$ n'appartenant pas à $V(I)$,
et soit $\mathfrak q \in V(\mathfrak p) \cap V(\mathfrak a)$.
Alors, ou bien $\mathfrak p$ n'est pas un idéal premier associé de $M$
et donc $\text{depth}(M_\mathfrak p) \geq 1$,
ou bien on a $\dim((A/\mathfrak p)_\mathfrak q) > d + 1$ d'après (2).
Ainsi, les hypothèses du
lemme \ref{lemma-algebraize-local-cohomology-general}
sont satisfaites pour $s = 1$ et $d$ ; on utilise ici la condition (3).
On trouve donc un idéal
$J_0 \subset \mathfrak a$ avec $V(J_0) \cap V(I) = V(\mathfrak a)$
tel que, pour tout $J \subset J_0$ avec $V(J) \cap V(I) = V(\mathfrak a)$,
les applications
$$
H^i_J(M) \longrightarrow H^i(R\Gamma_\mathfrak a(M)^\wedge)
$$
soient des isomorphismes pour $i = 0, 1$. Considérons les morphismes de
triangles exacts
$$
\xymatrix{
R\Gamma_J(M)  \ar[d] \ar[r] &
M \ar[r] \ar[d] &
R\Gamma(V, \mathcal{F}) \ar[d] \ar[r] &
R\Gamma_J(M)[1]  \ar[d] \\
R\Gamma_J(M)^\wedge \ar[r] &
M \ar[r] &
R\Gamma(V, \mathcal{F})^\wedge \ar[r] &
R\Gamma_J(M)^\wedge[1] \\
R\Gamma_\mathfrak a(M)^\wedge \ar[r] \ar[u] &
M \ar[r] \ar[u] &
R\Gamma(U, \mathcal{F})^\wedge \ar[r] \ar[u] &
R\Gamma_\mathfrak a(M)^\wedge[1] \ar[u]
}
$$
où $V = \Spec(A) \setminus V(J)$. Rappelons que
$R\Gamma_\mathfrak a(M)^\wedge \to R\Gamma_J(M)^\wedge$
est un isomorphisme (parce que $\mathfrak a$, $\mathfrak a + I$ et $J + I$
définissent le même sous-schéma fermé ; voir par exemple la
démonstration du lemme \ref{lemma-algebraize-local-cohomology-general}).
Par conséquent,
$R\Gamma(U, \mathcal{F})^\wedge = R\Gamma(V, \mathcal{F})^\wedge$.
On obtient ainsi un diagramme commutatif
$$
\xymatrix{
0 \ar[r] &
H^0_J(M) \ar[r] \ar[d] &
M \ar[r] \ar[d] \ar[r] &
\Gamma(V, \mathcal{F}) \ar[d] \ar[r] &
H^1_J(M) \ar[d] \ar[r] &
0 \\
0 \ar[r] &
H^0(R\Gamma_J(M)^\wedge) \ar[r] &
M \ar[r] &
H^0(R\Gamma(V, \mathcal{F})^\wedge) \ar[r] &
H^1(R\Gamma_J(M)^\wedge) \ar[r] &
0 \\
0 \ar[r] &
H^0(R\Gamma_\mathfrak a(M)^\wedge) \ar[r] \ar[u] &
M \ar[r] \ar[u] &
H^0(R\Gamma(U, \mathcal{F})^\wedge) \ar[r] \ar[u] &
H^1(R\Gamma_\mathfrak a(M)^\wedge) \ar[r] \ar[u] &
0
}
$$
dont les lignes sont exactes et dont les flèches verticales inférieures sont des isomorphismes. Ainsi,
on obtient un isomorphisme
$\Gamma(V, \mathcal{F}) \to H^0(R\Gamma(U, \mathcal{F})^\wedge)$.
D'après les lemmes \ref{lemma-formal-functions-general}
et \ref{lemma-derived-completion-pseudo-coherent}, on a
$$
R\Gamma(U, \mathcal{F})^\wedge =
R\Gamma(U, \mathcal{F}^\wedge) =
R\Gamma(U, R\lim \mathcal{F}/I^n\mathcal{F})
$$
et l'on trouve $H^0(R\Gamma(U, \mathcal{F})^\wedge) =
\lim H^0(U, \mathcal{F}/I^n\mathcal{F})$ d'après
Cohomologie, lemme \ref{cohomology-lemma-RGamma-commutes-with-Rlim}.
\end{proof}

\noindent
Nous appliquons maintenant un argument d'amorçage au lemme précédent afin d'éliminer la condition (3).

\begin{proposition}
\label{proposition-application-H0}
Soient $I \subset \mathfrak a$ des idéaux d'un anneau noethérien $A$.
Soit $\mathcal{F}$ un module cohérent sur
$U = \Spec(A) \setminus V(\mathfrak a)$.
Supposons que
\begin{enumerate}
\item $A$ soit complet pour la topologie $I$-adique et possède un complexe dualisant,
\item si $x \in \text{Ass}(\mathcal{F})$, $x \not \in V(I)$,
$\overline{\{x\}} \cap V(I) \not \subset V(\mathfrak a)$,
et $z \in \overline{\{x\}} \cap V(\mathfrak a)$, alors
$\dim(\mathcal{O}_{\overline{\{x\}}, z}) > \text{cd}(A, I) + 1$.
\end{enumerate}
Alors on obtient un isomorphisme
$$
\colim H^0(V, \mathcal{F})
\longrightarrow
\lim H^0(U, \mathcal{F}/I^n\mathcal{F})
$$
où la limite inductive porte sur les ouverts $V \subset U$ contenant $U \cap V(I)$.
\end{proposition}

\begin{proof}
Soit $T \subset U$ l'ensemble des points $x$ tels que
$\overline{\{x\}} \cap V(I) \subset V(\mathfrak a)$.
Soit $\mathcal{F} \to \mathcal{F}'$ la surjection
de modules cohérents sur $U$ construite dans
Cohomologie locale, lemme \ref{local-cohomology-lemma-get-depth-1-along-Z}.
Comme $\mathcal{F} \to \mathcal{F}'$ est un isomorphisme
au-dessus d'un ouvert $V \subset U$ contenant $U \cap V(I)$,
il suffit de démontrer le lemme après avoir remplacé $\mathcal{F}$
par $\mathcal{F}'$. On peut donc supposer, et l'on suppose, que
pour $x \in U$ tel que $\overline{\{x\}} \cap V(I) \subset V(\mathfrak a)$,
on a $\text{depth}(\mathcal{F}_x) \geq 1$.

\medskip\noindent
Soit $\mathcal{V}$ l'ensemble des sous-schémas ouverts $V \subset U$
contenant $U \cap V(I)$, ordonné par l'inclusion inverse.
C'est un ensemble ordonné filtrant. Nous affirmons d'abord que
$$
\mathcal{F}(V)
\longrightarrow
\lim H^0(U, \mathcal{F}/I^n\mathcal{F})
$$
est injective pour tout $V \in \mathcal{F}$ (et, en particulier, l'application
du lemme est injective). En effet, tout point associé $x$ de $\mathcal{F}$
doit vérifier $\overline{\{x\}} \cap U \cap Y \not = \emptyset$
d'après le paragraphe précédent. Si $y \in \overline{\{x\}} \cap U \cap Y$, alors
$\mathcal{F}_x$ est un localisé de $\mathcal{F}_y$
et $\mathcal{F}_y \subset \lim \mathcal{F}_y/I^n \mathcal{F}_y$
d'après le théorème d'intersection de Krull
(Algèbre, lemme \ref{algebra-lemma-intersect-powers-ideal-module-zero}).
Cela prouve l'assertion, car une section $s \in \mathcal{F}(V)$
appartenant au noyau devrait avoir un support vide et serait donc nulle.

\medskip\noindent
Choisissons un $A$-module de type fini $M$ tel que $\mathcal{F}$ soit la restriction
de $\widetilde{M}$ à $U$ ; voir Cohomologie locale, lemme
\ref{local-cohomology-lemma-finiteness-pushforwards-and-H1-local}.
On peut supposer, et l'on suppose, que $H^0_\mathfrak a(M) = 0$.
Posons $\text{Ass}(M) \setminus V(I) = \{\mathfrak p_1, \ldots, \mathfrak p_n\}$.
Nous allons démontrer le lemme par récurrence sur $n$. Après réindexation, on
peut supposer que $\mathfrak p_n$ est un élément minimal de l'ensemble
$\{\mathfrak p_1, \ldots, \mathfrak p_n\}$ pour l'inclusion, c'est-à-dire que
$\mathfrak p_n$ est un point générique du support de $M$.
Posons
$$
M' = H^0_{\mathfrak p_1 \ldots \mathfrak p_{n - 1} I}(M)
$$
et $M'' = M/M'$. Soient $\mathcal{F}'$ et $\mathcal{F}''$ les
$\mathcal{O}_U$-modules cohérents correspondant à $M'$ et $M''$.
Complexes dualisants, lemme \ref{dualizing-lemma-divide-by-torsion},
entraîne que $M''$ ne possède qu'un seul idéal premier associé, à savoir $\mathfrak p_n$.
Ainsi, $\mathcal{F}''$ ne possède qu'un seul point associé, et
on voit que la condition (3)(a) du lemme \ref{lemma-application-H0-pre} est satisfaite ;
l'application $\colim H^0(V, \mathcal{F}'')
\to \lim H^0(U, \mathcal{F}''/I^n\mathcal{F}'')$ est donc un isomorphisme.
D'autre part, puisque
$\mathfrak p_n \not \in V(\mathfrak p_1 \ldots \mathfrak p_{n - 1} I)$
on voit que $\mathfrak p_n$ n'est pas un idéal premier associé de $M'$.
L'hypothèse de récurrence s'applique donc à $M'$ ; remarquons
que, puisque $\mathcal{F}' \subset \mathcal{F}$,
la condition $\text{depth}(\mathcal{F}'_x) \geq 1$ est satisfaite aux points $x$ tels que
$\overline{\{x\}} \cap V(I) \subset V(\mathfrak a)$ ; voir
Algèbre, lemme \ref{algebra-lemma-depth-in-ses}.
Ainsi, l'application
$\colim H^0(V, \mathcal{F}')
\to \lim H^0(U, \mathcal{F}'/I^n\mathcal{F}')$ est également un isomorphisme.
On conclut par le lemme \ref{lemma-ses-H0}.
\end{proof}

\begin{lemma}
\label{lemma-application-H0}
Soient $I \subset \mathfrak a$ des idéaux d'un anneau noethérien $A$.
Soit $\mathcal{F}$ un module cohérent sur
$U = \Spec(A) \setminus V(\mathfrak a)$.
Supposons que
\begin{enumerate}
\item $A$ soit complet pour la topologie $I$-adique et possède un complexe dualisant,
\item si $x \in \text{Ass}(\mathcal{F})$, $x \not \in V(I)$,
$\overline{\{x\}} \cap V(I) \not \subset V(\mathfrak a)$, et
$z \in V(\mathfrak a) \cap \overline{\{x\}}$, alors
$\dim(\mathcal{O}_{\overline{\{x\}}, z}) > \text{cd}(A, I) + 1$,
\item pour $x \in U$ tel que $\overline{\{x\}} \cap V(I) \subset V(\mathfrak a)$,
on ait $\text{depth}(\mathcal{F}_x) \geq 2$,
\end{enumerate}
alors on obtient un isomorphisme
$$
H^0(U, \mathcal{F})
\longrightarrow
\lim H^0(U, \mathcal{F}/I^n\mathcal{F})
$$
\end{lemma}

\begin{proof}
Soit $\hat s \in \lim H^0(U, \mathcal{F}/I^n\mathcal{F})$.
D'après la proposition \ref{proposition-application-H0},
on trouve que $\hat s$ est l'image d'un élément $s \in \mathcal{F}(V)$
pour un certain ouvert $V \subset U$ contenant $U \cap V(I)$.
Cependant, la condition (3) montre que $\text{depth}(\mathcal{F}_x) \geq 2$
pour tout $x \in U \setminus V$, et l'on trouve donc
$\mathcal{F}(V) = \mathcal{F}(U)$ d'après
Diviseurs, lemme \ref{divisors-lemma-depth-2-hartog},
ce qui achève la démonstration.
\end{proof}

\begin{lemma}
\label{lemma-alternative-colim-H0}
Soit $A$ un anneau noethérien. Soit $f \in \mathfrak a \subset A$
un élément d'un idéal de $A$. Soit $M$ un $A$-module de type fini.
Supposons que
\begin{enumerate}
\item $A$ soit complet pour la topologie $f$-adique,
\item $f$ soit un non-diviseur de zéro dans $M$,
\item $H^1_\mathfrak a(M/fM)$ soit un $A$-module de type fini.
\end{enumerate}
Alors, pour $U = \Spec(A) \setminus V(\mathfrak a)$, l'application
$$
\colim_V \Gamma(V, \widetilde{M})
\longrightarrow
\lim \Gamma(U, \widetilde{M/f^nM})
$$
est un isomorphisme, où la limite inductive porte sur les ouverts $V \subset U$
contenant $U \cap V(f)$.
\end{lemma}

\begin{proof}
Posons $\mathcal{F} = \widetilde{M}|_U$.
Le fait que $H^1_\mathfrak a(M/fM)$ soit de type fini entraîne que
$H^0(U, \mathcal{F}/f\mathcal{F})$ est de type fini ; voir
Cohomologie locale, lemme
\ref{local-cohomology-lemma-finiteness-pushforwards-and-H1-local}.
D'après Cohomologie, lemme \ref{cohomology-lemma-limit-finite}
(qui s'applique puisque $f$ est un non-diviseur de zéro dans $\mathcal{F}$),
on voit que $N = \lim H^0(U, \mathcal{F}/f^n\mathcal{F})$
est un $A$-module de type fini, est sans $f$-torsion, et que
$N/fN \subset H^0(U, \mathcal{F}/f\mathcal{F})$.
D'autre part, on a une application $M \to N$ et une application compatible
$$
M/fM \longrightarrow H^0(U, \mathcal{F}/f\mathcal{F})
$$
Pour $g \in \mathfrak a$, on voit que $(M/fM)_g$ s'envoie isomorphiquement
sur $H^0(U \cap D(f), \mathcal{F}/f\mathcal{F})$, puisque
$\mathcal{F}/f\mathcal{F}$ est la restriction de $\widetilde{M/fM}$ à $U$.
On en conclut que $M_g \to N_g$ induit un isomorphisme
$$
M_g/fM_g = (M/fM)_g \to (N/fN)_g = N_g/fN_g
$$
Puisque $f$ est un non-diviseur de zéro à la fois dans $N$ et dans $M$, on en conclut
que $M_g \to N_g$ induit un isomorphisme sur les complétés $f$-adiques,
ce qui entraîne à son tour que $M_g \to N_g$ est un isomorphisme dans un voisinage
ouvert de $V(f) \cap D(g)$.
Comme $g \in \mathfrak a$ était arbitraire, on conclut que $M$ et $N$
définissent des modules cohérents isomorphes au-dessus d'un ouvert $V$
comme dans l'énoncé du lemme. Cela achève la démonstration.
\end{proof}

\begin{proposition}
\label{proposition-alternative-colim-H0}
Soit $A$ un anneau noethérien. Soit $f \in \mathfrak a \subset A$
un élément d'un idéal de $A$. Soit $\mathcal{F}$ un module cohérent
sur $U = \Spec(A) \setminus V(\mathfrak a)$. Supposons que
\begin{enumerate}
\item $A$ soit complet pour la topologie $f$-adique et possède un complexe dualisant,
\item si $x \in \text{Ass}(\mathcal{F})$, $x \not \in V(f)$,
$\overline{\{x\}} \cap V(f) \not \subset V(\mathfrak a)$,
et $z \in \overline{\{x\}} \cap V(\mathfrak a)$, alors
$\dim(\mathcal{O}_{\overline{\{x\}}, z}) > 2$.
\end{enumerate}
Alors l'application
$$
\colim_V \Gamma(V, \mathcal{F})
\longrightarrow
\lim \Gamma(U, \mathcal{F}/f^n\mathcal{F})
$$
est un isomorphisme, où la limite inductive porte sur les ouverts $V \subset U$
contenant $U \cap V(f)$.
\end{proposition}

\begin{proof}[Première démonstration]
Rappelons que $A$ est universellement caténaire et à fibres formelles
de Gorenstein ; voir Complexes dualisants, lemmes
\ref{dualizing-lemma-dualizing-gorenstein-formal-fibres} et
\ref{dualizing-lemma-universally-catenary}. On peut donc considérer l'application
$\mathcal{F} \to \mathcal{F}'$ construite dans
Cohomologie locale, lemme \ref{local-cohomology-lemma-make-S2-along-Z}
pour la partie fermée $V(f) \cap U$ de $U$.
Remarquons que
\begin{enumerate}
\item Le noyau et le conoyau de $\mathcal{F} \to \mathcal{F}'$
ont leur support dans $V(f) \cap U$.
\item Le module $\mathcal{F}'$ est sans $f$-torsion, car ses germes sont de
profondeur $\geq 1$ en tout point de $V(f) \cap U$, c'est-à-dire que $\mathcal{F}'$
n'a aucun point associé dans $V(f) \cap U$.
\item Si $y \in V(f) \cap U$ est un point associé de
$\mathcal{F}'/f\mathcal{F}'$, alors
$\text{depth}(\mathcal{F}'_y) = 1$
et, par conséquent (d'après la construction de $\mathcal{F}'$),
il existe une spécialisation immédiate $x \leadsto y$,
où $x \not \in V(f)$ est un point associé de $\mathcal{F}$.
Il résulte de notre hypothèse (2) que $y$ ne peut admettre
dans $\Spec(A)$ une spécialisation immédiate
en un point $z \in V(\mathfrak a)$.
\item Il résulte de (3) que $H^0(U, \mathcal{F}'/f\mathcal{F}')$
est un $A$-module de type fini ; voir
Cohomologie locale, lemme \ref{local-cohomology-lemma-finiteness-Rjstar}.
\end{enumerate}
Ces remarques vont nous permettre d'achever la démonstration.

\medskip\noindent
Nous affirmons d'abord que le lemme vaut pour $\mathcal{F}'$.
Choisissons en effet un $A$-module de type fini $M'$ tel que $\mathcal{F}'$
soit la restriction
à $U$ du module cohérent associé à $M'$ ; voir Cohomologie locale,
lemme \ref{local-cohomology-lemma-finiteness-pushforwards-and-H1-local}.
Comme $\mathcal{F}'$ est sans $f$-torsion, on peut également supposer que $M'$
est sans $f$-torsion. La remarque (4) ci-dessus montre que
$H^1_\mathfrak a(M')$ est un $A$-module de type fini ; voir Cohomologie locale,
lemme \ref{local-cohomology-lemma-finiteness-pushforwards-and-H1-local}.
L'assertion résulte donc du lemme \ref{lemma-alternative-colim-H0}.

\medskip\noindent
Deuxièmement, remarquons que le lemme vaut trivialement pour tout
$\mathcal{O}_U$-module cohérent à support dans $V(f) \cap U$.
Soient $\mathcal{K}$, resp.\ $\mathcal{G}$, resp.\ $\mathcal{Q}$
le noyau, resp.\ l'image, resp.\ le conoyau de l'application
$\mathcal{F} \to \mathcal{F}'$.
La suite exacte courte
$0 \to \mathcal{G} \to \mathcal{F}' \to \mathcal{Q} \to 0$
et le lemme \ref{lemma-ses-H0} montrent que le résultat vaut pour
$\mathcal{G}$. On procède ensuite de même avec la suite exacte courte
$0 \to \mathcal{K} \to \mathcal{F} \to \mathcal{G} \to 0$
pour achever la démonstration.
\end{proof}

\begin{proof}[Deuxième démonstration]
La proposition est un cas particulier de la
proposition \ref{proposition-application-H0}.
\end{proof}

\begin{lemma}
\label{lemma-alternative-H0}
Soit $A$ un anneau noethérien. Soit $f \in \mathfrak a \subset A$
un élément d'un idéal de $A$. Soit $M$ un $A$-module de type fini.
Supposons que
\begin{enumerate}
\item $A$ soit complet pour la topologie $f$-adique,
\item $H^1_\mathfrak a(M)$ et $H^2_\mathfrak a(M)$ soient
annulés par une puissance de $f$.
\end{enumerate}
Alors, pour $U = \Spec(A) \setminus V(\mathfrak a)$, l'application
$$
\Gamma(U, \widetilde{M})
\longrightarrow
\lim \Gamma(U, \widetilde{M/f^nM})
$$
est un isomorphisme.
\end{lemma}

\begin{proof}
On peut appliquer le lemme \ref{lemma-formal-functions-principal} à $U$ et à
$\mathcal{F} = \widetilde{M}|_U$, car $\mathcal{F}$ est un objet noethérien de
la catégorie des $\mathcal{O}_U$-modules cohérents.
Puisque $H^1(U, \mathcal{F}) = H^2_\mathfrak a(M)$
(Cohomologie locale, lemme
\ref{local-cohomology-lemma-finiteness-pushforwards-and-H1-local})
est annulé par une puissance de $f$, on voit que
son module de Tate $f$-adique est nul.
Le lemme montre donc que $\lim H^0(U, \mathcal{F}/f^n \mathcal{F})$
est égal au complété $f$-adique usuel de $H^0(U, \mathcal{F})$.
Considérons la suite exacte courte
$$
0 \to M/H^0_\mathfrak a(M) \to H^0(U, \mathcal{F}) \to H^1_\mathfrak a(M) \to 0
$$
de Cohomologie locale, lemme
\ref{local-cohomology-lemma-finiteness-pushforwards-and-H1-local}.
Comme $M/H^0_\mathfrak a(M)$ est un $A$-module de type fini, il est
complet ; voir Algèbre, lemme \ref{algebra-lemma-completion-tensor}.
Comme $H^1_\mathfrak a(M)$ est annulé par une puissance de $f$, on déduit
d'Algèbre, lemme \ref{algebra-lemma-completion-differ-by-torsion},
que $H^0(U, \mathcal{F})$ est également complet. Cela achève la démonstration.
\end{proof}







\section{Algébrisation des sections formelles, III}
\label{section-algebraization-sections-coherent-III}

\noindent
La section suivante contient une liste non exhaustive
d'applications des résultats
sur la complétion de la cohomologie locale à la cohomologie supérieure
des modules cohérents sur les schémas quasi-affines et à leur
complétion relativement à un idéal.

\begin{proposition}
\label{proposition-application-higher}
Soient $I \subset \mathfrak a$ des idéaux d'un anneau noethérien $A$.
Soit $\mathcal{F}$ un module cohérent sur
$U = \Spec(A) \setminus V(\mathfrak a)$.
Soit $s \geq 0$.
Supposons que
\begin{enumerate}
\item $A$ soit complet pour la topologie $I$-adique et possède un complexe dualisant,
\item si $x \in U \setminus V(I)$, alors
$\text{depth}(\mathcal{F}_x) > s$ ou
$$
\text{depth}(\mathcal{F}_x) +
\dim(\mathcal{O}_{\overline{\{x\}}, z}) > \text{cd}(A, I) + s + 1
$$
pour tout $z \in V(\mathfrak a) \cap \overline{\{x\}}$,
\item l'une des conditions suivantes soit satisfaite :
\begin{enumerate}
\item la restriction de $\mathcal{F}$ à $U \setminus V(I)$
soit $(S_{s + 1})$, ou
\item la dimension de $V(\mathfrak a)$ soit au plus $2$\footnote{Au
sens où la différence entre les valeurs maximale et minimale
sur $V(\mathfrak a)$ d'une fonction de dimension sur $\Spec(A)$ est au plus $2$.}.
\end{enumerate}
\end{enumerate}
Alors les applications
$$
H^i(U, \mathcal{F})
\longrightarrow
\lim H^i(U, \mathcal{F}/I^n\mathcal{F})
$$
sont des isomorphismes pour $i < s$. De plus, on a un isomorphisme
$$
\colim H^s(V, \mathcal{F})
\longrightarrow
\lim H^s(U, \mathcal{F}/I^n\mathcal{F})
$$
où la limite inductive porte sur les ouverts $V \subset U$ contenant $U \cap V(I)$.
\end{proposition}

\begin{proof}
On peut supposer $s > 0$, puisque le cas $s = 0$ a été traité dans la
proposition \ref{proposition-application-H0}.

\medskip\noindent
Choisissons un $A$-module de type fini $M$ tel que $\mathcal{F}$ soit la restriction
à $U$ du module cohérent associé à $M$ ; voir Cohomologie locale,
lemme \ref{local-cohomology-lemma-finiteness-pushforwards-and-H1-local}.
Posons $d = \text{cd}(A, I)$.
Soit $\mathfrak p$ un idéal premier de $A$ n'appartenant pas à $V(I)$,
et soit $\mathfrak q \in V(\mathfrak p) \cap V(\mathfrak a)$.
Alors ou bien $\text{depth}(M_\mathfrak p) \geq s + 1 > s$,
ou bien $\dim((A/\mathfrak p)_\mathfrak q) > d + s + 1$ d'après (2).
D'après le lemme \ref{lemma-bootstrap-bis-bis}, on en conclut que les
hypothèses de la situation \ref{situation-bootstrap}
sont satisfaites pour $A, I, V(\mathfrak a), M, s, d$.
D'autre part, les hypothèses du
lemme \ref{lemma-algebraize-local-cohomology-general}
sont satisfaites pour $s + 1$ et $d$ ; c'est ici que l'on utilise la condition (3).

\medskip\noindent
En appliquant le lemme \ref{lemma-algebraize-local-cohomology-general},
on trouve qu'il existe un idéal
$J_0 \subset \mathfrak a$ tel que $V(J_0) \cap V(I) = V(\mathfrak a)$
et que, pour tout $J \subset J_0$ tel que $V(J) \cap V(I) = V(\mathfrak a)$,
les applications
$$
H^i_J(M) \longrightarrow H^i(R\Gamma_\mathfrak a(M)^\wedge)
$$
sont des isomorphismes pour $i \leq s + 1$.

\medskip\noindent
Pour $i \leq s$, l'application $H^i_\mathfrak a(M) \to H^i_J(M)$
est un isomorphisme d'après les lemmes \ref{lemma-bootstrap-inherited} et
\ref{lemma-kill-colimit-support-general}.
En utilisant la comparaison entre cohomologie et cohomologie locale
(Cohomologie locale, lemme \ref{local-cohomology-lemma-local-cohomology}),
on déduit que
$H^i(U, \mathcal{F}) \to H^i(V,\mathcal{F})$
est un isomorphisme pour $V = \Spec(A) \setminus V(J)$ et
$i < s$.

\medskip\noindent
D'après le théorème \ref{theorem-final-bootstrap}, on a
$H^i_\mathfrak a(M) = \lim H^i_\mathfrak a(M/I^nM)$
pour $i \leq s$. D'après le lemme \ref{lemma-combine-two}, on a
$H^{s + 1}_\mathfrak a(M) = \lim H^{s + 1}_\mathfrak a(M/I^nM)$.

\medskip\noindent
L'isomorphisme $H^0(U, \mathcal{F}) = H^0(V, \mathcal{F}) =
\lim H^0(U, \mathcal{F}/I^n\mathcal{F})$ résulte de ce qui précède et de la
proposition \ref{proposition-application-H0}.
Pour $0 < i < s$, on obtient les isomorphismes voulus
$H^i(U, \mathcal{F}) = H^i(V, \mathcal{F}) =
\lim H^i(U, \mathcal{F}/I^n\mathcal{F})$ de la
même manière, en utilisant la relation entre cohomologie locale
et cohomologie ; c'est plus facile que le cas $i = 0$,
car, pour $i > 0$, on a
$$
H^i(U, \mathcal{F}) = H^{i + 1}_\mathfrak a(M),
\quad
H^i(V, \mathcal{F}) = H^{i + 1}_J(M),
\quad
H^i(R\Gamma(U, \mathcal{F})^\wedge) = 
H^{i + 1}(R\Gamma_\mathfrak a(M)^\wedge)
$$
On procède de même pour la dernière assertion.
\end{proof}

\begin{lemma}
\label{lemma-alternative-higher}
Soit $A$ un anneau noethérien. Soit $f \in \mathfrak a \subset A$
un élément d'un idéal de $A$. Soit $M$ un $A$-module de type fini.
Soit $s \geq 0$. Supposons que
\begin{enumerate}
\item $A$ soit complet pour la topologie $f$-adique,
\item $H^i_\mathfrak a(M)$ soit annulé par une puissance de $f$
pour $i \leq s + 1$.
\end{enumerate}
Alors, pour $U = \Spec(A) \setminus V(\mathfrak a)$, l'application
$$
H^i(U, \widetilde{M})
\longrightarrow
\lim H^i(U, \widetilde{M/f^nM})
$$
est un isomorphisme pour $i < s$.
\end{lemma}

\begin{proof}
Procédons par récurrence sur $s$. Si $s = 0$, l'assertion est vide. Si $s = 1$, le
résultat est le lemme \ref{lemma-alternative-H0}. Supposons $s > 1$. Par récurrence,
il suffit de démontrer le résultat pour $i = s - 1 \geq 1$.
On peut appliquer le lemme \ref{lemma-formal-functions-principal}
à $U$ et à $\mathcal{F} = \widetilde{M}|_U$,
car $\mathcal{F}$ est un objet noethérien de
la catégorie des $\mathcal{O}_U$-modules cohérents.
Remarquons que $H^j(U, \mathcal{F}) = H^{j + 1}_\mathfrak a(M)$ pour tout $j$, d'après
Cohomologie locale, lemme
\ref{local-cohomology-lemma-finiteness-pushforwards-and-H1-local}.
Ainsi, pour $j = s = (s - 1) + 1$, ce module est
annulé par une puissance de $f$ d'après l'hypothèse.
Il résulte alors du
lemme \ref{lemma-formal-functions-principal} que
$\lim H^{s - 1}(U, \mathcal{F}/f^n\mathcal{F})$
est le complété $f$-adique usuel de $H^{s - 1}(U, \mathcal{F})$.
En utilisant de nouveau le fait que ce module est annulé par une puissance de
$f$, on voit que le complété est simplement égal à $H^{s - 1}(U, \mathcal{F})$,
comme voulu.
\end{proof}





\section{Application à la connexité}
\label{section-connected}

\noindent
Dans cette section, nous étudions le théorème de connexité de Grothendieck
et ses variantes ; la version originale se trouve dans
\cite[Exposee XIII, théorème 2.1]{SGA2}. Il existe dans la littérature une version
appelée théorème de connexité de Faltings ;
nous supposons qu'il s'agit de \cite[théorème 6]{Faltings-some}.
Énonçons et démontrons la version optimale pour les anneaux locaux
complets donnée dans \cite[théorème 1.6]{Varbaro}.

\begin{lemma}
\label{lemma-punctured-still-connected}
\begin{reference}
\cite[Theorem 1.6]{Varbaro}
\end{reference}
Soit $(A, \mathfrak m)$ un anneau local noethérien complet.
Soit $I$ un idéal propre de $A$.
Posons $X = \Spec(A)$ et $Y = V(I)$.
Notons
\begin{enumerate}
\item $d$ la dimension minimale d'une composante irréductible de $X$, et
\item $c$ la dimension minimale d'une partie fermée $Z \subset X$
telle que $X \setminus Z$ soit non connexe.
\end{enumerate}
Alors, pour toute partie fermée $Z \subset Y$, l'espace $Y \setminus Z$ est connexe si
$\dim(Z) < \min(c, d - 1) - \text{cd}(A, I)$. En particulier, le spectre
épointé de $A/I$ est connexe si $\text{cd}(A, I) < \min(c, d - 1)$.
\end{lemma}

\begin{proof}
Démontrons d'abord la dernière assertion. Supposons premièrement que le spectre
épointé de $A/I$ soit vide. Alors
Cohomologie locale, lemme \ref{local-cohomology-lemma-cd-bound-dim-local},
montre que toute composante irréductible de $X$ est de dimension
$\leq \text{cd}(A, I)$, et l'on obtient $\min(c, d - 1) - \text{cd}(A, I) < 0$,
ce qui entraîne que le lemme vaut dans ce cas. On peut donc supposer que
$U \cap Y$ est non vide, où $U = X \setminus \{\mathfrak m\}$
est le spectre épointé de $A$. On peut remplacer $A$ par son réduit.
Remarquons que $A$ possède un complexe dualisant
(Complexes dualisants, lemme \ref{dualizing-lemma-ubiquity-dualizing})
et que $A$ est complet relativement à $I$
(Algèbre, lemme \ref{algebra-lemma-complete-by-sub}).
Si l'on suppose $d - 1 > \text{cd}(A, I)$, on peut appliquer le
lemme \ref{lemma-application-theorem} pour voir que
$$
\colim H^0(V, \mathcal{O}_V)
\longrightarrow
\lim H^0(U, \mathcal{O}_U/I^n\mathcal{O}_U)
$$
est un isomorphisme, où la limite inductive porte sur les ouverts $V \subset U$
contenant $U \cap Y$. Si $U \cap Y$ n'est pas connexe, alors
son $n$-ième voisinage infinitésimal dans $U$ n'est pas connexe
pour tout $n$, et l'on trouve que le
membre de droite possède un idempotent non trivial (on utilise ici
que $U \cap Y$ est non vide).
On peut donc trouver un $V$ qui n'est pas connexe.
Posons $Z = X \setminus V$. D'après
Cohomologie locale, lemme \ref{local-cohomology-lemma-cd-bound-dim-local},
on voit que toute composante irréductible de $Z$ est de dimension
$\leq \text{cd}(A, I)$. Ainsi $c \leq \text{cd}(A, I)$, ce qui
démontre bien la dernière assertion.

\medskip\noindent
On peut déduire l'énoncé du lemme de ce que l'on vient de démontrer,
comme suit. Supposons que $Z \subset Y$ soit fermé, que $Y \setminus Z$ ne soit pas
connexe et que $\dim(Z) = e$. Rappelons que, par convention, un espace connexe est non vide.
On conclut donc que soit (a) $Y = Z$, soit (b)
$Y \setminus Z = W_1 \amalg W_2$, avec les $W_i$ non vides, ouverts et fermés
dans $Y \setminus Z$. Dans le cas (b), on peut choisir des points $w_i \in W_i$
qui sont fermés dans $U$ ; voir
Morphismes, lemme \ref{morphisms-lemma-ubiquity-Jacobson-schemes}.
On peut alors trouver $f_1, \ldots, f_e \in \mathfrak m$
tels que $V(f_1, \ldots, f_e) \cap Z = \{\mathfrak m\}$
et, dans le cas (b), on peut supposer $w_i \in V(f_1, \ldots, f_e)$.
En effet, par récurrence et en évitant les idéaux premiers, on peut
choisir $f_i$ tel que $\dim V(f_1, \ldots, f_i) \cap Z = e - i$
et tel que, dans le cas (b), on ait $w_1, w_2 \in V(f_i)$.
Il s'ensuit que le spectre épointé de $A/I + (f_1, \ldots, f_e)$
n'est pas connexe (on omet un petit détail). Puisque
$\text{cd}(A, I + (f_1, \ldots, f_e)) \leq \text{cd}(A, I) + e$ d'après
Cohomologie locale, lemmes \ref{local-cohomology-lemma-cd-sum} et
\ref{local-cohomology-lemma-bound-cd}, on conclut que
$$
\text{cd}(A, I) + e \geq \min(c, d - 1)
$$
d'après la première partie de la démonstration. Cela entraîne
$e \geq \min(c, d - 1) - \text{cd}(A, I)$, ce qu'il fallait démontrer.
\end{proof}

\begin{lemma}
\label{lemma-connected}
Soient $I \subset \mathfrak a$ des idéaux d'un anneau noethérien $A$.
Supposons que
\begin{enumerate}
\item $A$ soit complet pour la topologie $I$-adique et possède un complexe dualisant,
\item si $\mathfrak p \subset A$ est un idéal premier minimal n'appartenant pas
à $V(I)$ et si $\mathfrak q \in V(\mathfrak p) \cap V(\mathfrak a)$, alors
$\dim((A/\mathfrak p)_\mathfrak q) > \text{cd}(A, I) + 1$,
\item tout ouvert non vide $V \subset \Spec(A)$ qui contient
$V(I) \setminus V(\mathfrak a)$ soit connexe\footnote{Par exemple,
si $A$ est un domaine.}.
\end{enumerate}
Alors $V(I) \setminus V(\mathfrak a)$ est vide ou connexe.
\end{lemma}

\begin{proof}
On peut remplacer $A$ par son réduit. On dispose alors de l'inégalité
de (2) pour tous les idéaux premiers associés de $A$. D'après la
proposition \ref{proposition-application-H0}, on voit que
$$
\colim H^0(V, \mathcal{O}_V) = \lim H^0(T_n, \mathcal{O}_{T_n})
$$
où la limite inductive porte sur les ouverts $V$ comme dans (3), et où $T_n$ est le
$n$-ième voisinage infinitésimal de $T = V(I) \setminus V(\mathfrak a)$
dans $U = \Spec(A) \setminus V(\mathfrak a)$. Ainsi, $T$ est vide
ou connexe ; sinon, le membre de droite posséderait en effet un
idempotent non trivial, tandis que l'on a supposé que le membre de gauche n'en possède pas.
On omet quelques détails.
\end{proof}

\begin{lemma}
\label{lemma-H0}
Soit $A$ un domaine noethérien qui possède un complexe dualisant
et qui est complet relativement à un élément non nul $f \in A$.
Soit $f \in \mathfrak a \subset A$ un idéal.
Supposons que toute composante irréductible de $Z = V(\mathfrak a)$
soit de codimension $> 2$ dans $X = \Spec(A)$, c'est-à-dire que toute
composante irréductible de $Z$ soit de codimension $> 1$ dans $Y = V(f)$.
Alors $Y \setminus Z$ est connexe.
\end{lemma}

\begin{proof}
C'est un cas particulier du lemme \ref{lemma-connected} (dont la démonstration
repose sur la proposition \ref{proposition-application-H0}). Nous le
démontrons ci-dessous à l'aide de la proposition \ref{proposition-alternative-colim-H0}, plus facile.

\medskip\noindent
Posons $U = X \setminus Z$. D'après la proposition \ref{proposition-alternative-colim-H0},
on a un isomorphisme
$$
\colim \Gamma(V, \mathcal{O}_V) \to
\lim_n \Gamma(U, \mathcal{O}_U/f^n \mathcal{O}_U)
$$
où la limite inductive porte sur les ouverts $V \subset U$ contenant $U \cap Y$.
Ainsi, si $U \cap Y$ n'est pas connexe, il existe pour un certain $V$
un idempotent non trivial dans $\Gamma(V, \mathcal{O}_V)$. C'est impossible,
car $V$ est un schéma intègre puisque $X$ est le spectre d'un domaine.
\end{proof}






\section{Le foncteur de complétion}
\label{section-completion}

\noindent
Soit $X$ un schéma noethérien. Soit $Y \subset X$ un sous-schéma fermé
de faisceau quasi-cohérent d'idéaux $\mathcal{I} \subset \mathcal{O}_X$.
Dans cette section, nous considérons les systèmes projectifs de
$\mathcal{O}_X$-modules cohérents $(\mathcal{F}_n)$ tels que $\mathcal{F}_n$ soit
annulé par $I^n$ et que les applications de transition induisent des
isomorphismes $\mathcal{F}_{n + 1}/I^n\mathcal{F}_{n + 1} \to \mathcal{F}_n$.
La catégorie de ces systèmes projectifs a été notée
$$
\textit{Coh}(X, \mathcal{I})
$$
dans Cohomologie des schémas, section \ref{coherent-section-coherent-formal}.
Cette catégorie est équivalente à la catégorie des modules cohérents
sur le complété formel de $X$ le long de $Y$ ; cependant, comme nous n'avons
pas encore introduit les schémas formels ni les modules cohérents sur ceux-ci,
nous ne pouvons employer ici cette terminologie. Nous nous intéressons tout particulièrement
au foncteur de complétion
$$
\textit{Coh}(\mathcal{O}_X)
\longrightarrow
\textit{Coh}(X, \mathcal{I}),\quad
\mathcal{F} \longmapsto \mathcal{F}^\wedge
$$
Voir
Cohomologie des schémas, équation (\ref{coherent-equation-completion-functor}).

\begin{lemma}
\label{lemma-completion-fully-faithful}
Soit $X$ un schéma noethérien et soit $Y \subset X$ un sous-schéma fermé.
Soit $Y_n \subset X$ le $n$-ième voisinage infinitésimal de $Y$ dans $X$.
Considérons les conditions suivantes :
\begin{enumerate}
\item $X$ est quasi-affine et
$\Gamma(X, \mathcal{O}_X) \to \lim \Gamma(Y_n, \mathcal{O}_{Y_n})$
est un isomorphisme,
\item $X$ possède un module inversible ample $\mathcal{L}$ et
$\Gamma(X, \mathcal{L}^{\otimes m}) \to
\lim \Gamma(Y_n, \mathcal{L}^{\otimes m}|_{Y_n})$
est un isomorphisme pour tout $m \gg 0$,
\item pour tout $\mathcal{O}_X$-module de type fini et localement libre
$\mathcal{E}$, l'application
$\Gamma(X, \mathcal{E}) \to \lim \Gamma(Y_n, \mathcal{E}|_{Y_n})$
est un isomorphisme, et
\item le foncteur de complétion
$\textit{Coh}(\mathcal{O}_X) \to \textit{Coh}(X, \mathcal{I})$
est pleinement fidèle sur la sous-catégorie pleine des
objets de type fini et localement libres.
\end{enumerate}
Alors (1) $\Rightarrow$ (2) $\Rightarrow$ (3) $\Rightarrow$ (4)
et (4) $\Rightarrow$ (3).
\end{lemma}

\begin{proof}
Preuve de (3) $\Rightarrow$ (4). Si $\mathcal{F}$ et $\mathcal{G}$
sont de type fini et localement libres sur $X$, alors, en considérant
$\mathcal{H} = \SheafHom_{\mathcal{O}_X}(\mathcal{G}, \mathcal{F})$
et en utilisant Cohomologie des schémas, lemme
\ref{coherent-lemma-completion-internal-hom}
on voit que (3) entraîne (4).

\medskip\noindent
Preuve de (2) $\rightarrow$ (3). Soit en effet $\mathcal{L}$ ample
sur $X$, et supposons que $\mathcal{E}$ soit un
$\mathcal{O}_X$-module de type fini et localement libre.
Nous affirmons que l'on peut trouver une suite universellement exacte
$$
0 \to \mathcal{E} \to
(\mathcal{L}^{\otimes p})^{\oplus r} \to
(\mathcal{L}^{\otimes q})^{\oplus s}
$$
pour certains $r, s \geq 0$ et $0 \ll p \ll q$. Si tel est le cas, alors,
en utilisant la suite exacte
$$
0 \to \lim \Gamma(\mathcal{E}|_{Y_n}) \to
\lim \Gamma((\mathcal{L}^{\otimes p})^{\oplus r}|_{Y_n}) \to
\lim \Gamma((\mathcal{L}^{\otimes q})^{\oplus s}|_{Y_n})
$$
et les isomorphismes de (2), on obtient l'isomorphisme de (3).
Pour démontrer l'assertion, considérons le module localement libre dual
$\SheafHom_{\mathcal{O}_X}(\mathcal{E}, \mathcal{O}_X)$
et appliquons
Propriétés, proposition \ref{properties-proposition-characterize-ample},
pour trouver une surjection
$$
(\mathcal{L}^{\otimes -p})^{\oplus r}
\longrightarrow
\SheafHom_{\mathcal{O}_X}(\mathcal{E}, \mathcal{O}_X)
$$
En prenant les duaux, on obtient la première application de la suite exacte
(elle est universellement injective, car le fait d'être une surjection est universel).
En répétant l'opération avec le conoyau, on obtient la seconde. On omet quelques détails.

\medskip\noindent
Preuve de (1) $\Rightarrow$ (2). C'est vrai parce que, si $X$ est quasi-affine,
alors $\mathcal{O}_X$ est un module inversible ample ; voir
Propriétés, lemme \ref{properties-lemma-quasi-affine-O-ample}.

\medskip\noindent
Nous omettons la preuve de (4) $\Rightarrow$ (3).
\end{proof}

\noindent
Étant donnés un schéma noethérien et un faisceau quasi-cohérent d'idéaux
$\mathcal{I} \subset \mathcal{O}_X$, nous dirons
qu'un objet $(\mathcal{F}_n)$ de $\textit{Coh}(X, \mathcal{I})$
est {\it de type fini et localement libre} si chaque $\mathcal{F}_n$ est un
$\mathcal{O}_X/\mathcal{I}^n$-module de type fini et localement libre.

\begin{lemma}
\label{lemma-completion-fully-faithful-general}
Soit $X$ un schéma noethérien et soit $Y \subset X$ un sous-schéma fermé
de faisceau d'idéaux $\mathcal{I} \subset \mathcal{O}_X$.
Soit $Y_n \subset X$ le $n$-ième voisinage infinitésimal de $Y$ dans $X$.
Soit $\mathcal{V}$ l'ensemble des sous-schémas ouverts $V \subset X$ contenant $Y$,
ordonné par l'inclusion inverse.
\begin{enumerate}
\item $X$ est quasi-affine et
$$
\colim_\mathcal{V} \Gamma(V, \mathcal{O}_V)
\longrightarrow
\lim \Gamma(Y_n, \mathcal{O}_{Y_n})
$$
est un isomorphisme,
\item $X$ possède un module inversible ample $\mathcal{L}$ et
$$
\colim_\mathcal{V} \Gamma(V, \mathcal{L}^{\otimes m})
\longrightarrow
\lim \Gamma(Y_n, \mathcal{L}^{\otimes m}|_{Y_n})
$$
est un isomorphisme pour tout $m \gg 0$,
\item pour tout $V \in \mathcal{V}$ et tout
$\mathcal{O}_V$-module $\mathcal{E}$ de type fini et localement libre, l'application
$$
\colim_{V' \geq V} \Gamma(V', \mathcal{E}|_{V'})
\longrightarrow
\lim \Gamma(Y_n, \mathcal{E}|_{Y_n})
$$
est un isomorphisme, et
\item le foncteur de complétion
$$
\colim_\mathcal{V} \textit{Coh}(\mathcal{O}_V)
\longrightarrow
\textit{Coh}(X, \mathcal{I}),
\quad
\mathcal{F} \longmapsto \mathcal{F}^\wedge
$$
est pleinement fidèle sur la sous-catégorie pleine des
objets de type fini et localement libres (voir l'explication ci-dessus).
\end{enumerate}
Alors (1) $\Rightarrow$ (2) $\Rightarrow$ (3) $\Rightarrow$ (4)
et (4) $\Rightarrow$ (3).
\end{lemma}

\begin{proof}
Remarquons que $\mathcal{V}$ est un ensemble ordonné filtrant ; les limites inductives sont donc
comme dans Catégories, section \ref{categories-section-directed-colimits}.
La suite de l'argument est presque exactement la même que l'argument
de la démonstration du lemme \ref{lemma-completion-fully-faithful} ; nous invitons
le lecteur à l'omettre.

\medskip\noindent
Preuve de (3) $\Rightarrow$ (4). Si $\mathcal{F}$ et $\mathcal{G}$
sont de type fini et localement libres sur $V \in \mathcal{V}$, alors, en considérant
$\mathcal{H} = \SheafHom_{\mathcal{O}_V}(\mathcal{G}, \mathcal{F})$
et en utilisant Cohomologie des schémas, lemme
\ref{coherent-lemma-completion-internal-hom}
on voit que (3) entraîne (4).

\medskip\noindent
Preuve de (2) $\Rightarrow$ (3). Soit $\mathcal{L}$ ample
sur $X$, et supposons que $\mathcal{E}$ soit un
$\mathcal{O}_V$-module de type fini et localement libre
pour un certain $V \in \mathcal{V}$.
Nous affirmons que l'on peut trouver une suite universellement exacte
$$
0 \to \mathcal{E} \to
(\mathcal{L}^{\otimes p})^{\oplus r}|_{V} \to
(\mathcal{L}^{\otimes q})^{\oplus s}|_{V}
$$
pour certains $r, s \geq 0$ et $0 \ll p \ll q$. Si tel est le cas,
l'isomorphisme de (2) entraînera l'isomorphisme de (3).
Pour démontrer l'assertion, rappelons que $\mathcal{L}|_V$ est ample ; voir
Propriétés, lemme \ref{properties-lemma-ample-on-locally-closed}.
Considérons le module localement libre dual
$\SheafHom_{\mathcal{O}_V}(\mathcal{E}, \mathcal{O}_V)$
et appliquons
Propriétés, proposition \ref{properties-proposition-characterize-ample},
pour trouver une surjection
$$
(\mathcal{L}^{\otimes -p})^{\oplus r}|_V \longrightarrow
\SheafHom_{\mathcal{O}_V}(\mathcal{E}, \mathcal{O}_V)
$$
(elle est universellement injective, car le fait d'être une surjection est universel).
En prenant les duaux, on obtient la première application de la suite exacte.
En répétant l'opération avec le conoyau, on obtient la seconde. On omet quelques détails.

\medskip\noindent
Preuve de (1) $\Rightarrow$ (2). C'est vrai parce que, si $X$ est quasi-affine,
alors $\mathcal{O}_X$ est un module inversible ample ; voir
Propriétés, lemme \ref{properties-lemma-quasi-affine-O-ample}.

\medskip\noindent
Nous omettons la preuve de (4) $\Rightarrow$ (3).
\end{proof}

\begin{lemma}
\label{lemma-recognize-formal-coherent-modules}
Soit $X$ un schéma noethérien. Soit $\mathcal{I} \subset \mathcal{O}_X$
un faisceau quasi-cohérent d'idéaux. Le foncteur
$$
\textit{Coh}(X, \mathcal{I}) \longrightarrow \text{Pro-}\QCoh(\mathcal{O}_X)
$$
est pleinement fidèle ; voir Catégories, remarque \ref{categories-remark-pro-category}.
\end{lemma}

\begin{proof}
Soient $(\mathcal{F}_n)$ et $(\mathcal{G}_n)$ des objets de
$\textit{Coh}(X, \mathcal{I})$. Un morphisme de pro-objets
$\alpha$ de $(\mathcal{F}_n)$ vers $(\mathcal{G}_n)$ est donné
par un système d'applications
$\alpha_n : \mathcal{F}_{m(n)} \to \mathcal{G}_n$ compatibles avec
les applications de transition, où $\mathbf{N} \to \mathbf{N}$, $n \mapsto m(n)$,
est une fonction croissante (en particulier $m(n) \geq n$). Puisque
$\mathcal{F}_n = \mathcal{F}_{m(n)}/\mathcal{I}^n\mathcal{F}_{m(n)}$
et puisque $\mathcal{G}_n$ est annulé par $\mathcal{I}^n$,
on voit que $\alpha_n$ induit une application $\mathcal{F}_n \to \mathcal{G}_n$.
\end{proof}

\noindent
Nous donnons ensuite quelques exemples du type de résultat de pleine fidélité
que nous pourrons démontrer à l'aide du travail effectué plus haut dans ce chapitre.

\begin{lemma}
\label{lemma-fully-faithful}
Soient $I \subset \mathfrak a$ des idéaux d'un anneau noethérien $A$.
Soit $U = \Spec(A) \setminus V(\mathfrak a)$. Supposons que
\begin{enumerate}
\item $A$ soit complet pour la topologie $I$-adique et possède un complexe dualisant,
\item pour tout idéal premier associé $\mathfrak p \subset A$ tel que
$\mathfrak p \not \in V(I)$ et
$V(\mathfrak p) \cap V(I) \not \subset V(\mathfrak a)$ et
$\mathfrak q \in V(\mathfrak p) \cap V(\mathfrak a)$, on ait
$\dim((A/\mathfrak p)_\mathfrak q) > \text{cd}(A, I) + 1$,
\item pour $\mathfrak p \subset A$ tel que $\mathfrak p \not \in V(I)$
et $V(\mathfrak p) \cap V(I) \subset V(\mathfrak a)$,
on ait $\text{depth}(A_\mathfrak p) \geq 2$.
\end{enumerate}
Alors le foncteur de complétion
$$
\textit{Coh}(\mathcal{O}_U)
\longrightarrow
\textit{Coh}(U, I\mathcal{O}_U),
\quad
\mathcal{F} \longmapsto \mathcal{F}^\wedge
$$
est pleinement fidèle sur la sous-catégorie pleine des
objets de type fini et localement libres.
\end{lemma}

\begin{proof}
D'après le lemme \ref{lemma-completion-fully-faithful},
il suffit de montrer que
$$
\Gamma(U, \mathcal{O}_U) =
\lim \Gamma(U, \mathcal{O}_U/I^n\mathcal{O}_U)
$$
Cela résulte immédiatement du
lemme \ref{lemma-application-H0}.
\end{proof}

\begin{lemma}
\label{lemma-fully-faithful-alternative}
Soit $A$ un anneau noethérien. Soit $f \in \mathfrak a \subset A$
un élément d'un idéal de $A$. Soit $U = \Spec(A) \setminus V(\mathfrak a)$.
Supposons que
\begin{enumerate}
\item $A$ soit complet pour la topologie $f$-adique,
\item $H^1_\mathfrak a(A)$ et $H^2_\mathfrak a(A)$ soient
annulés par une puissance de $f$.
\end{enumerate}
Alors le foncteur de complétion
$$
\textit{Coh}(\mathcal{O}_U)
\longrightarrow
\textit{Coh}(U, I\mathcal{O}_U),
\quad
\mathcal{F} \longmapsto \mathcal{F}^\wedge
$$
est pleinement fidèle sur la sous-catégorie pleine des
objets de type fini et localement libres.
\end{lemma}

\begin{proof}
D'après le lemme \ref{lemma-completion-fully-faithful},
il suffit de montrer que
$$
\Gamma(U, \mathcal{O}_U) =
\lim \Gamma(U, \mathcal{O}_U/I^n\mathcal{O}_U)
$$
Cela résulte immédiatement du
lemme \ref{lemma-alternative-H0}.
\end{proof}

\begin{lemma}
\label{lemma-fully-faithful-simple-two}
Soit $A$ un anneau noethérien. Soit $f \in \mathfrak a$ un élément d'un
idéal de $A$. Soit $U = \Spec(A) \setminus V(\mathfrak a)$. Supposons que
\begin{enumerate}
\item $A$ possède un complexe dualisant et soit complet relativement à $f$,
\item pour tout idéal premier $\mathfrak p \subset A$, $f \not \in \mathfrak p$
et $\mathfrak q \in V(\mathfrak p) \cap V(\mathfrak a)$, on ait
$\text{depth}(A_\mathfrak p) + \dim((A/\mathfrak p)_\mathfrak q) > 2$.
\end{enumerate}
Alors le foncteur de complétion
$$
\textit{Coh}(\mathcal{O}_U)
\longrightarrow
\textit{Coh}(U, I\mathcal{O}_U),
\quad
\mathcal{F} \longmapsto \mathcal{F}^\wedge
$$
est pleinement fidèle sur la sous-catégorie pleine des objets de type fini et localement libres.
\end{lemma}

\begin{proof}
Cela résulte du lemme \ref{lemma-fully-faithful-alternative} et de
Cohomologie locale, proposition \ref{local-cohomology-proposition-annihilator}.
\end{proof}

\begin{lemma}
\label{lemma-fully-faithful-general}
Soient $I \subset \mathfrak a \subset A$ des idéaux d'un anneau noethérien $A$.
Soit $U = \Spec(A) \setminus V(\mathfrak a)$. Soit $\mathcal{V}$
l'ensemble des sous-schémas ouverts de $U$ contenant $U \cap V(I)$,
ordonné par l'inclusion inverse. Supposons que
\begin{enumerate}
\item $A$ soit complet pour la topologie $I$-adique et possède un complexe dualisant,
\item pour tout idéal premier associé
$\mathfrak p \subset A$ tel que
$I \not \subset \mathfrak p$ et
$V(\mathfrak p) \cap V(I) \not \subset V(\mathfrak a)$
et $\mathfrak q \in V(\mathfrak p) \cap V(\mathfrak a)$, on ait
$\dim((A/\mathfrak p)_\mathfrak q) > \text{cd}(A, I) + 1$.
\end{enumerate}
Alors le foncteur de complétion
$$
\colim_\mathcal{V} \textit{Coh}(\mathcal{O}_V)
\longrightarrow
\textit{Coh}(U, I\mathcal{O}_U),
\quad
\mathcal{F} \longmapsto \mathcal{F}^\wedge
$$
est pleinement fidèle sur la sous-catégorie pleine des
objets de type fini et localement libres.
\end{lemma}

\begin{proof}
D'après le lemme \ref{lemma-completion-fully-faithful-general},
il suffit de montrer que
$$
\colim_\mathcal{V} \Gamma(V, \mathcal{O}_V) =
\lim \Gamma(U, \mathcal{O}_U/I^n\mathcal{O}_U)
$$
Cela résulte immédiatement de la proposition \ref{proposition-application-H0}.
\end{proof}

\begin{lemma}
\label{lemma-fully-faithful-general-alternative}
Soit $A$ un anneau noethérien. Soit $f \in \mathfrak a \subset A$
un élément d'un idéal de $A$. Soit $U = \Spec(A) \setminus V(\mathfrak a)$.
Soit $\mathcal{V}$ l'ensemble des sous-schémas ouverts de $U$ contenant $U \cap V(f)$,
ordonné par l'inclusion inverse. Supposons que
\begin{enumerate}
\item $A$ soit complet pour la topologie $f$-adique,
\item $f$ soit un non-diviseur de zéro,
\item $H^1_\mathfrak a(A/fA)$ soit un $A$-module de type fini.
\end{enumerate}
Alors le foncteur de complétion
$$
\colim_\mathcal{V} \textit{Coh}(\mathcal{O}_V)
\longrightarrow
\textit{Coh}(U, f\mathcal{O}_U),
\quad
\mathcal{F} \longmapsto \mathcal{F}^\wedge
$$
est pleinement fidèle sur la sous-catégorie pleine des objets de type fini et localement libres.
\end{lemma}

\begin{proof}
D'après le lemme \ref{lemma-completion-fully-faithful-general},
il suffit de montrer que
$$
\colim_\mathcal{V} \Gamma(V, \mathcal{O}_V) =
\lim \Gamma(U, \mathcal{O}_U/I^n\mathcal{O}_U)
$$
Cela résulte immédiatement du lemme \ref{lemma-alternative-colim-H0}.
\end{proof}

\begin{lemma}
\label{lemma-fully-faithful-very-general}
Soient $I \subset \mathfrak a \subset A$ des idéaux d'un anneau noethérien $A$.
Soit $U = \Spec(A) \setminus V(\mathfrak a)$. Soit $\mathcal{V}$ l'ensemble
des sous-schémas ouverts de $U$ contenant $U \cap V(I)$, ordonné par l'inclusion
inverse. Soient $\mathcal{F}$ et
$\mathcal{G}$ des $\mathcal{O}_V$-modules cohérents pour un certain
$V \in \mathcal{V}$. L'application
$$
\colim_{V' \geq V} \Hom_V(\mathcal{G}|_{V'}, \mathcal{F}|_{V'})
\longrightarrow
\Hom_{\textit{Coh}(U, I\mathcal{O}_U)}(\mathcal{G}^\wedge, \mathcal{F}^\wedge)
$$
est bijective si les hypothèses suivantes sont satisfaites :
\begin{enumerate}
\item $A$ soit complet pour la topologie $I$-adique et possède un complexe dualisant,
\item si $x \in \text{Ass}(\mathcal{F})$, $x \not \in V(I)$,
$\overline{\{x\}} \cap V(I) \not \subset V(\mathfrak a)$
et $z \in \overline{\{x\}} \cap V(\mathfrak a)$, alors
$\dim(\mathcal{O}_{\overline{\{x\}}, z}) > \text{cd}(A, I) + 1$.
\end{enumerate}
\end{lemma}

\begin{proof}
On peut choisir des $\mathcal{O}_U$-modules cohérents
$\mathcal{F}'$ et $\mathcal{G}'$ dont les restrictions à $V$
sont $\mathcal{F}$ et $\mathcal{G}$ ; voir
Propriétés, lemme \ref{properties-lemma-lift-finite-presentation}.
On peut modifier le choix de $\mathcal{F}'$ pour assurer que
$\text{Ass}(\mathcal{F}') \subset V$ ; voir par exemple
Cohomologie locale, lemme \ref{local-cohomology-lemma-get-depth-1-along-Z}.
On peut donc remplacer, et l'on remplace, $V$ par $U$, et $\mathcal{F}$ et $\mathcal{G}$
par $\mathcal{F}'$ et $\mathcal{G}'$.
Posons $\mathcal{H} = \SheafHom_{\mathcal{O}_U}(\mathcal{G}, \mathcal{F})$.
C'est un $\mathcal{O}_U$-module cohérent. On a
$$
\Hom_V(\mathcal{G}|_V, \mathcal{F}|_V) =
H^0(V, \mathcal{H})
\quad\text{et}\quad
\lim H^0(U, \mathcal{H}/\mathcal{I}^n\mathcal{H}) =
\Mor_{\textit{Coh}(U, I\mathcal{O}_U)}
(\mathcal{G}^\wedge, \mathcal{F}^\wedge)
$$
Voir Cohomologie des schémas, lemme \ref{coherent-lemma-completion-internal-hom}.
Ainsi, si l'on peut montrer que les hypothèses de la
proposition \ref{proposition-application-H0}
sont satisfaites pour $\mathcal{H}$, la démonstration est achevée.
C'est le cas parce que
$\text{Ass}(\mathcal{H}) \subset \text{Ass}(\mathcal{F})$.
Voir Cohomologie des schémas, lemme
\ref{coherent-lemma-hom-into-depth}.
\end{proof}








\section{Algébrisation des modules formels cohérents, I}
\label{section-algebraization-modules}

\noindent
La surjectivité essentielle du foncteur de complétion (voir ci-dessous)
a été étudiée systématiquement dans
\cite{SGA2}, \cite{MRaynaud-book} et \cite{MRaynaud-paper}.
Nous nous plaçons dans la situation affine suivante.

\begin{situation}
\label{situation-algebraize}
Ici, $A$ est un anneau noethérien et $I \subset \mathfrak a \subset A$ sont des idéaux.
On pose $X = \Spec(A)$, $Y = V(I) = \Spec(A/I)$ et
$Z = V(\mathfrak a) = \Spec(A/\mathfrak a)$. De plus, $U = X \setminus Z$.
\end{situation}

\noindent
Dans cette section, nous cherchons des conditions assurant qu'un objet
de $\textit{Coh}(U, I\mathcal{O}_U)$ appartient à l'image du foncteur de complétion
$\textit{Coh}(\mathcal{O}_U) \to \textit{Coh}(U, I\mathcal{O}_U)$.
Voir Cohomologie des schémas, section \ref{coherent-section-coherent-formal} et
section \ref{section-completion}.

\begin{lemma}
\label{lemma-system-of-modules}
Dans la situation \ref{situation-algebraize}.
Considérons un système projectif $(M_n)$ de $A$-modules tel
que
\begin{enumerate}
\item $M_n$ soit un $A$-module de type fini,
\item $M_n$ soit annulé par $I^n$,
\item le noyau et le conoyau de $M_{n + 1}/I^nM_{n + 1} \to M_n$
soient de torsion annulée par une puissance de $\mathfrak a$.
\end{enumerate}
Alors $(\widetilde{M}_n|_U)$ appartient à $\textit{Coh}(U, I\mathcal{O}_U)$.
Réciproquement, tout objet de $\textit{Coh}(U, I\mathcal{O}_U)$
s'obtient de cette manière.
\end{lemma}

\begin{proof}
Nous omettons de vérifier que $(\widetilde{M}_n|_U)$ appartient à
$\textit{Coh}(U, I\mathcal{O}_U)$. Soit $(\mathcal{F}_n)$
un objet de $\textit{Coh}(U, I\mathcal{O}_U)$.
D'après Cohomologie locale, lemme
\ref{local-cohomology-lemma-finiteness-pushforwards-and-H1-local}
on voit que $\mathcal{F}_n = \widetilde{M_n}$ pour un certain
$A/I^n$-module de type fini $M_n$. Après avoir quotienté $M_n$ par $H^0_\mathfrak a(M_n)$,
on peut supposer $M_n \subset H^0(U, \mathcal{F}_n)$ ; voir
Complexes dualisants, lemme \ref{dualizing-lemma-divide-by-torsion},
et le lemme déjà cité.
En remplaçant par récurrence $M_{n + 1}$ par l'image réciproque
de $M_n$ par l'application $M_{n + 1} \to H^0(U, \mathcal{F}_{n + 1})
\to H^0(U, \mathcal{F}_n)$, on peut supposer que $M_{n + 1}$ s'envoie dans
$M_n$. On obtient ainsi un système projectif $(M_n)$ satisfaisant (1) et (2)
tel que $\mathcal{F}_n = \widetilde{M_n}$. Pour voir que (3)
est satisfaite, on utilise le fait que $M_{n + 1}/I^nM_{n + 1} \to M_n$ est une application
de $A$-modules de type fini qui induit un isomorphisme après
application de $\widetilde{\ }$ et restriction à $U$
(on utilise ici une fois de plus le premier lemme cité).
\end{proof}

\noindent
Dans la situation \ref{situation-algebraize}, on peut étudier le foncteur de complétion de
Cohomologie des schémas, équation (\ref{coherent-equation-completion-functor}),
\begin{equation}
\label{equation-completion}
\textit{Coh}(\mathcal{O}_U)
\longrightarrow
\textit{Coh}(U, I\mathcal{O}_U),\quad
\mathcal{F} \longmapsto \mathcal{F}^\wedge
\end{equation}
Si $A$ est complet pour la topologie $I$-adique, ce foncteur est pleinement fidèle
sur des sous-catégories appropriées d'après nos travaux antérieurs sur l'algébrisation des
sections formelles ; voir la section \ref{section-completion} et le
lemme \ref{lemma-fully-faithful-inequalities} pour quelques exemples de résultats.
Soit ensuite $(\mathcal{F}_n)$ un objet de $\textit{Coh}(U, I\mathcal{O}_U)$.
En supposant toujours $A$ complet pour la topologie $I$-adique, on peut demander :
quand $(\mathcal{F}_n)$ appartient-il à l'image essentielle du foncteur de complétion
affiché ci-dessus ?

\begin{lemma}
\label{lemma-essential-image-completion}
Dans la situation \ref{situation-algebraize}, soit $(\mathcal{F}_n)$
un objet de $\textit{Coh}(U, I\mathcal{O}_U)$. Considérons les
conditions suivantes :
\begin{enumerate}
\item $(\mathcal{F}_n)$ appartient à l'image essentielle
du foncteur (\ref{equation-completion}),
\item $(\mathcal{F}_n)$ est le complété d'un
$\mathcal{O}_U$-module cohérent,
\item $(\mathcal{F}_n)$ est le complété d'un
$\mathcal{O}_V$-module cohérent pour un ouvert $U \cap Y \subset V \subset U$,
\item $(\mathcal{F}_n)$ est le complété de
la restriction à $U$ d'un $\mathcal{O}_X$-module cohérent,
\item $(\mathcal{F}_n)$ est la restriction à $U$ du
complété d'un $\mathcal{O}_X$-module cohérent,
\item il existe un objet $(\mathcal{G}_n)$ de
$\textit{Coh}(X, I\mathcal{O}_X)$ dont la restriction
à $U$ est $(\mathcal{F}_n)$.
\end{enumerate}
Alors les conditions (1), (2), (3), (4) et (5) sont équivalentes et entraînent (6).
Si $A$ est complet pour la topologie $I$-adique, la condition (6) entraîne les autres.
\end{lemma}

\begin{proof}
Les parties (1) et (2) sont équivalentes, car le complété d'un
$\mathcal{O}_U$-module cohérent $\mathcal{F}$ est par définition l'image de
$\mathcal{F}$ par le foncteur (\ref{equation-completion}).
Si $V \subset U$ est un sous-schéma ouvert contenant $U \cap Y$, alors on a
$$
\textit{Coh}(V, I\mathcal{O}_V) =
\textit{Coh}(U, I\mathcal{O}_U)
$$
car la catégorie des $\mathcal{O}_V$-modules cohérents à support dans
$V \cap Y$ est la même que la catégorie des $\mathcal{O}_U$-modules cohérents
à support dans $U \cap Y$. Ainsi, le complété d'un
$\mathcal{O}_V$-module cohérent est un objet de $\textit{Coh}(U, I\mathcal{O}_U)$.
Ceci étant dit, l'équivalence de (2), (3), (4) et (5)
vaut parce que les foncteurs
$\textit{Coh}(\mathcal{O}_X) \to \textit{Coh}(\mathcal{O}_U) \to
\textit{Coh}(\mathcal{O}_V)$ sont essentiellement surjectifs.
Voir Propriétés, lemme \ref{properties-lemma-lift-finite-presentation}.

\medskip\noindent
La condition (5) entraîne toujours (6). Supposons $A$ complet pour la topologie $I$-adique.
Alors tout objet de $\textit{Coh}(X, I\mathcal{O}_X)$ correspond à un
$A$-module de type fini d'après Cohomologie des schémas, lemme
\ref{coherent-lemma-inverse-systems-affine}.
On voit ainsi que (6) entraîne (5) dans ce cas.
\end{proof}

\begin{example}
\label{example-not-algebraizable}
Soit $k$ un corps. Soit $A = k[x, y][[t]]$, avec $I = (t)$ et
$\mathfrak a = (x, y, t)$. Employons les notations de la
situation \ref{situation-algebraize}. Remarquons que
$U \cap Y = (D(x) \cap Y) \cup (D(y) \cap Y)$ est un recouvrement
ouvert affine. Pour $n \geq 1$, considérons le module inversible
$\mathcal{L}_n$ sur $\mathcal{O}_U/t^n\mathcal{O}_U$
obtenu en recollant $A_x/t^nA_x$ et $A_y/t^nA_y$ au moyen de l'élément inversible
de $A_{xy}/t^nA_{xy}$ qui est l'image d'une série formelle quelconque
de la forme
$$
u = 1 + \frac{t}{xy} + \sum_{n \geq 2} a_n \frac{t^n}{(xy)^{\varphi(n)}}
$$
où $a_n \in k[x, y]$ et $\varphi(n) \in \mathbf{N}$.
Alors $(\mathcal{L}_n)$ est un objet inversible de
$\textit{Coh}(U, I\mathcal{O}_U)$ qui n'est pas le
complété d'un $\mathcal{O}_U$-module cohérent $\mathcal{L}$.
Nous ne faisons qu'esquisser l'argument et omettons la plupart des détails.
Soit $y \in U \cap Y$. Le complété du germe
$\mathcal{L}_y$ serait alors un module inversible, donc $\mathcal{L}_y$
est inversible. Il existerait ainsi un ouvert $V \subset U$
contenant $U \cap Y$ tel que $\mathcal{L}|_V$ soit inversible.
D'après Diviseurs, lemme \ref{divisors-lemma-extend-invertible-module},
on trouve un $A$-module inversible $M$ tel que
$\widetilde{M}|_V \cong \mathcal{L}|_V$. Cependant, l'anneau $A$
est factoriel ; on voit donc que $M \cong A$, ce qui entraînerait
$\mathcal{L}_n \cong \mathcal{O}_U/I^n\mathcal{O}_U$.
Comme $\mathcal{L}_2 \not \cong \mathcal{O}_U/I^2\mathcal{O}_U$
par construction, on obtient la contradiction voulue.

\medskip\noindent
Remarquons que, si l'on prend $a_n = 0$ pour $n \geq 2$, on voit
que $\lim H^0(U, \mathcal{L}_n)$ est non nul : dans ce cas,
la fonction $x$ sur $D(x)$ et la fonction $x + t/y$ sur $D(y)$
se recollent. D'autre part, si l'on prend $a_n = 1$ et $\varphi(n) = 2^n$,
ou même $\varphi(n) = n^2$, le lecteur peut montrer que
$\lim H^0(U, \mathcal{L}_n)$ est nul ; cela fournit une autre démonstration
du fait que $(\mathcal{L}_n)$ n'est pas algébrisable dans ce cas.
\end{example}

\noindent
Si, dans la situation \ref{situation-algebraize}, l'anneau $A$ n'est pas
complet pour la topologie $I$-adique, le lemme \ref{lemma-essential-image-completion}
suggère qu'il convient de demander si $(\mathcal{F}_n)$
appartient à l'image essentielle du foncteur de restriction
$$
\textit{Coh}(X, I\mathcal{O}_X)
\longrightarrow
\textit{Coh}(U, I\mathcal{O}_U)
$$
Cependant, on ne peut plus dire que cela signifie que $(\mathcal{F}_n)$
est algébrisable. Nous introduisons donc la terminologie suivante.

\begin{definition}
\label{definition-algebraizable}
Dans la situation \ref{situation-algebraize}, soit $(\mathcal{F}_n)$ un
objet de $\textit{Coh}(U, I\mathcal{O}_U)$. Nous dirons que
{\it $(\mathcal{F}_n)$ s'étend à $X$} s'il existe un objet
$(\mathcal{G}_n)$ de $\textit{Coh}(X, I\mathcal{O}_X)$ dont la restriction
à $U$ est isomorphe à $(\mathcal{F}_n)$.
\end{definition}

\noindent
Cette notion équivaut à être algébrisable sur la complétion.

\begin{lemma}
\label{lemma-algebraizable}
Dans la situation \ref{situation-algebraize}, soit $(\mathcal{F}_n)$ un
objet de $\textit{Coh}(U, I\mathcal{O}_U)$. Soient $A', I', \mathfrak a'$
les complétions $I$-adiques de $A, I, \mathfrak a$. Posons $X' = \Spec(A')$
et $U' = X' \setminus V(\mathfrak a')$. Les conditions suivantes sont équivalentes :
\begin{enumerate}
\item $(\mathcal{F}_n)$ s'étend à $X$, et
\item l'image réciproque de $(\mathcal{F}_n)$ sur $U'$ est le complété
d'un $\mathcal{O}_{U'}$-module cohérent.
\end{enumerate}
\end{lemma}

\begin{proof}
Rappelons que $A \to A'$ est un homomorphisme plat d'anneaux qui induit un isomorphisme
$A/I \to A'/I'$. Voir
Algèbre, lemmes \ref{algebra-lemma-completion-flat} et
\ref{algebra-lemma-completion-complete}.
Ainsi, $X' \to X$ est un
morphisme plat induisant un isomorphisme $Y' \to Y$. Ainsi, $U' \to U$
est un morphisme plat qui induit un isomorphisme $U' \cap Y' \to U \cap Y$.
Cela entraîne que, dans le diagramme commutatif
$$
\xymatrix{
\textit{Coh}(X', I\mathcal{O}_{X'}) \ar[r] &
\textit{Coh}(U', I\mathcal{O}_{U'}) \\
\textit{Coh}(X, I\mathcal{O}_X) \ar[u] \ar[r] &
\textit{Coh}(U, I\mathcal{O}_U) \ar[u]
}
$$
les foncteurs verticaux sont des équivalences. Voir
Cohomologie des schémas, lemme
\ref{coherent-lemma-inverse-systems-pullback-equivalence}.
Le lemme résulte formellement de ceci et des résultats du
lemme \ref{lemma-essential-image-completion}.
\end{proof}

\noindent
Dans la situation \ref{situation-algebraize}, soit $(\mathcal{F}_n)$ un
objet de $\textit{Coh}(U, I\mathcal{O}_U)$. Pour déterminer si
$(\mathcal{F}_n)$ s'étend à $X$, il est naturel de considérer le $A$-module
\begin{equation}
\label{equation-guess}
M = \lim H^0(U, \mathcal{F}_n)
\end{equation}
Remarquons que $M$ possède une topologie de limite qui est (a priori) moins fine que
la topologie $I$-adique, puisque $M \to H^0(U, \mathcal{F}_n)$ annule $I^nM$.
Il existe des applications canoniques
$$
\widetilde{M}|_U \to \widetilde{M/I^nM}|_U \to
\widetilde{H^0(U, \mathcal{F}_n)}|_U \to
\mathcal{F}_n
$$
On pourrait espérer que $\widetilde{M}$ se restreigne en un module cohérent
sur $U$ et que $(\mathcal{F}_n)$ soit le complété de ce module.
C'est naïf, car cela n'a presque aucune chance d'être vrai
si $A$ n'est pas complet. Mais, même si $A$ est complet pour la topologie $I$-adique,
cette notion est très difficile à manier.
Une approche moins naïve consiste à considérer la condition suivante.

\begin{definition}
\label{definition-canonically-algebraizable}
Dans la situation \ref{situation-algebraize}, soit $(\mathcal{F}_n)$ un
objet de $\textit{Coh}(U, I\mathcal{O}_U)$. Nous dirons que
{\it $(\mathcal{F}_n)$ s'étend canoniquement à $X$} si le
système projectif
$$
\{\widetilde{H^0(U, \mathcal{F}_n)}\}_{n \geq 1}
$$
dans $\QCoh(\mathcal{O}_X)$ est pro-isomorphe à un objet
$(\mathcal{G}_n)$ de $\textit{Coh}(X, I\mathcal{O}_X)$.
\end{definition}

\noindent
Nous verrons dans le lemme \ref{lemma-canonically-algebraizable}
que la condition de la définition \ref{definition-canonically-algebraizable}
est plus forte que celle de la définition \ref{definition-algebraizable}.

\begin{lemma}
\label{lemma-canonically-algebraizable}
Dans la situation \ref{situation-algebraize}, soit $(\mathcal{F}_n)$ un
objet de $\textit{Coh}(U, I\mathcal{O}_U)$. Si $(\mathcal{F}_n)$
s'étend canoniquement à $X$, alors
\begin{enumerate}
\item $(\widetilde{H^0(U, \mathcal{F}_n)})$ est pro-isomorphe à un
objet $(\mathcal{G}_n)$ de $\textit{Coh}(X, I \mathcal{O}_X)$,
unique à isomorphisme unique près,
\item la restriction de $(\mathcal{G}_n)$ à $U$ est isomorphe
à $(\mathcal{F}_n)$, c'est-à-dire que $(\mathcal{F}_n)$ s'étend à $X$,
\item le système projectif $\{H^0(U, \mathcal{F}_n)\}$
satisfait la condition de Mittag-Leffler, et
\item le module $M$ de (\ref{equation-guess}) est de type fini sur le
complété $I$-adique de $A$, et la topologie de limite sur
$M$ est la topologie $I$-adique.
\end{enumerate}
\end{lemma}

\begin{proof}
L'existence de $(\mathcal{G}_n)$ dans (1) résulte de la
définition \ref{definition-canonically-algebraizable}.
L'unicité de $(\mathcal{G}_n)$ dans (1) résulte du
lemme \ref{lemma-recognize-formal-coherent-modules}.
Écrivons $\mathcal{G}_n = \widetilde{M_n}$.
Alors $\{M_n\}$ est un système projectif de $A$-modules de type fini
tel que $M_n = M_{n + 1}/I^n M_{n + 1}$.
D'après la définition \ref{definition-canonically-algebraizable},
le système projectif $\{H^0(U, \mathcal{F}_n)\}$
est pro-isomorphe à $\{M_n\}$.
On voit donc que le système projectif $\{H^0(U, \mathcal{F}_n)\}$
satisfait la condition de Mittag-Leffler et que
$M = \lim M_n$ (comme modules topologiques).
Les propriétés de $M$ de (4) résultent ainsi d'
Algèbre, lemmes \ref{algebra-lemma-limit-complete},
\ref{algebra-lemma-finite-over-complete-ring} et
\ref{algebra-lemma-hathat-finitely-generated}.
Comme $U$ est quasi-affine, les applications canoniques
$$
\widetilde{H^0(U, \mathcal{F}_n)}|_U \to \mathcal{F}_n
$$
sont des isomorphismes (Propriétés, lemme
\ref{properties-lemma-quasi-coherent-quasi-affine}).
On en conclut que $(\mathcal{G}_n|_U)$ et $(\mathcal{F}_n)$ sont
pro-isomorphes, donc isomorphes d'après le
lemme \ref{lemma-recognize-formal-coherent-modules}.
\end{proof}

\begin{lemma}
\label{lemma-canonically-extend-base-change}
Dans la situation \ref{situation-algebraize}, soit $(\mathcal{F}_n)$ un
objet de $\textit{Coh}(U, I\mathcal{O}_U)$. Soit $A \to A'$ un homomorphisme plat
d'anneaux. Posons $X' = \Spec(A')$, soit $U' \subset X'$ l'image réciproque de $U$,
et notons $g : U' \to U$ le morphisme induit. Posons
$(\mathcal{F}'_n) = (g^*\mathcal{F}_n)$ ; voir
Cohomologie des schémas, lemme \ref{coherent-lemma-inverse-systems-pullback}.
Si $(\mathcal{F}_n)$ s'étend canoniquement à $X$, alors
$(\mathcal{F}'_n)$ s'étend canoniquement à $X'$.
De plus, l'extension obtenue dans le lemme \ref{lemma-canonically-algebraizable}
pour $(\mathcal{F}_n)$ donne par image réciproque l'extension de
$(\mathcal{F}'_n)$.
\end{lemma}

\begin{proof}
Soit $f : X' \to X$ le morphisme induit.
On a $H^0(U', \mathcal{F}'_n) = H^0(U, \mathcal{F}_n) \otimes_A A'$ par
changement de base plat ; voir Cohomologie des schémas, lemme
\ref{coherent-lemma-flat-base-change-cohomology}.
Ainsi, si $(\mathcal{G}_n)$ dans $\textit{Coh}(X, I\mathcal{O}_X)$
est pro-isomorphe à $(\widetilde{H^0(U, \mathcal{F}_n)})$, alors
$(f^*\mathcal{G}_n)$ est pro-isomorphe à
$$
(f^*\widetilde{H^0(U, \mathcal{F}_n)}) =
(\widetilde{H^0(U, \mathcal{F}_n) \otimes_A A'}) =
(\widetilde{H^0(U', \mathcal{F}'_n)})
$$
Cela achève la démonstration.
\end{proof}

\begin{lemma}
\label{lemma-when-done}
Dans la situation \ref{situation-algebraize}, soit $(\mathcal{F}_n)$ un objet
de $\textit{Coh}(U, I\mathcal{O}_U)$. Soit $M$ comme dans (\ref{equation-guess}).
Supposons que
\begin{enumerate}
\item[(a)] le système projectif $H^0(U, \mathcal{F}_n)$ satisfait la condition de Mittag-Leffler,
\item[(b)] la topologie de limite sur $M$ coïncide avec la topologie $I$-adique, et
\item[(c)] l'image de $M \to H^0(U, \mathcal{F}_n)$ est un $A$-module de type fini
pour tout $n$.
\end{enumerate}
Alors $(\mathcal{F}_n)$ s'étend canoniquement à $X$.
En particulier, si $A$ est complet pour la topologie $I$-adique, alors
$(\mathcal{F}_n)$ est le complété d'un $\mathcal{O}_U$-module cohérent.
\end{lemma}

\begin{proof}
Puisque $H^0(U, \mathcal{F}_n)$ satisfait la condition de Mittag-Leffler
et puisque la topologie de limite sur $M$ est la topologie $I$-adique,
on voit que $\{M/I^nM\}$ et $\{H^0(U, \mathcal{F}_n)\}$
sont des systèmes projectifs pro-isomorphes de $A$-modules.
Ainsi, si l'on pose
$$
\mathcal{G}_n = \widetilde{M/I^n M}
$$
on voit que, pour vérifier la condition de la
définition \ref{definition-canonically-algebraizable},
il suffit de montrer que $M$ est un module de type fini sur le
complété $I$-adique de $A$. Cela résulte
du fait que $M/I^n M$ est de type fini d'après la condition (c),
de ce qui précède et d'
Algèbre, lemme \ref{algebra-lemma-finite-over-complete-ring}.
\end{proof}

\noindent
Le résultat suivant est, en un certain sens, l'application la plus directe possible du
lemme \ref{lemma-when-done} ci-dessus.

\begin{lemma}
\label{lemma-algebraization-principal-variant}
Dans la situation \ref{situation-algebraize}, soit
$(\mathcal{F}_n)$ un objet de $\textit{Coh}(U, I\mathcal{O}_U)$.
Supposons que
\begin{enumerate}
\item $I = (f)$ soit un idéal principal, pour un non-diviseur de zéro $f \in \mathfrak a$,
\item $\mathcal{F}_n$ soit un
$\mathcal{O}_U/f^n\mathcal{O}_U$-module de type fini et localement libre,
\item $H^1_\mathfrak a(A/fA)$ et $H^2_\mathfrak a(A/fA)$
soient des $A$-modules de type fini.
\end{enumerate}
Alors $(\mathcal{F}_n)$ s'étend canoniquement à $X$. En particulier, si $A$
est complet, alors $(\mathcal{F}_n)$ est le complété d'un
$\mathcal{O}_U$-module cohérent.
\end{lemma}

\begin{proof}
Nous allons le démontrer en vérifiant les hypothèses (a), (b) et (c) du
lemme \ref{lemma-when-done}.

\medskip\noindent
Comme $\mathcal{F}_n$ est localement libre sur $\mathcal{O}_U/f^n\mathcal{O}_U$,
on voit que l'on a des suites exactes courtes
$0 \to \mathcal{F}_n \to \mathcal{F}_{n + 1} \to \mathcal{F}_1 \to 0$
pour tout $n$. La condition (b) est donc satisfaite d'après
Cohomologie, lemme \ref{cohomology-lemma-topology-I-adic-f}.

\medskip\noindent
Comme $f$ est un non-diviseur de zéro, on obtient des suites exactes courtes
$$
0 \to A/f^nA \xrightarrow{f} A/f^{n + 1}A \to A/fA \to 0
$$
et l'on a les suites exactes courtes correspondantes
$0 \to \mathcal{F}_n \to \mathcal{F}_{n + 1} \to \mathcal{F}_1 \to 0$.
Nous utiliserons
Cohomologie locale, lemme
\ref{local-cohomology-lemma-finiteness-pushforwards-and-H1-local}
sans autre mention. Nos hypothèses entraînent que
$H^0(U, \mathcal{O}_U/f\mathcal{O}_U)$ et
$H^1(U, \mathcal{O}_U/f\mathcal{O}_U)$
sont des $A$-modules de type fini. Il en va donc de même pour $\mathcal{F}_1$ ; voir
Cohomologie locale, lemme
\ref{local-cohomology-lemma-finiteness-for-finite-locally-free}.
Par récurrence et à l'aide des suites exactes courtes, on trouve que
$H^0(U, \mathcal{F}_n)$ sont des $A$-modules de type fini pour tout $n$.
On voit ainsi que l'hypothèse (c) est satisfaite.

\medskip\noindent
Enfin, comme $H^1(U, \mathcal{F}_1)$ est un $A$-module de type fini,
on peut appliquer Cohomologie, lemme \ref{cohomology-lemma-ML},
pour voir que l'hypothèse (a) est satisfaite.
\end{proof}

\begin{remark}
\label{remark-interesting-case-variant}
Dans le lemme \ref{lemma-algebraization-principal-variant},
si $A$ est universellement caténaire à fibres formelles
de Cohen--Macaulay (par exemple si $A$ possède un complexe dualisant), alors
la condition selon laquelle
$H^1_\mathfrak a(A/fA)$ et $H^2_\mathfrak a(A/fA)$
sont des $A$-modules de type fini équivaut à
$$
\text{depth}((A/f)_\mathfrak p) + \dim((A/\mathfrak p)_\mathfrak q) > 2
$$
pour tout $\mathfrak p \in V(f) \setminus V(\mathfrak a)$
et $\mathfrak q \in V(\mathfrak p) \cap V(\mathfrak a)$,
d'après Cohomologie locale, théorème \ref{local-cohomology-theorem-finiteness}.

\medskip\noindent
Par exemple, si $A/fA$ est $(S_2)$ et si toute composante irréductible
de $Z = V(\mathfrak a)$ est de codimension $\geq 3$
dans $Y = \Spec(A/fA)$, alors on obtient que
$H^1_\mathfrak a(A/fA)$ et $H^2_\mathfrak a(A/fA)$ sont de type fini.
Il faut comparer cela aux conditions légèrement plus faibles
obtenues dans le lemme \ref{lemma-algebraization-principal}
(voir aussi la remarque \ref{remark-interesting-case}).
\end{remark}






\section{Algébrisation des modules formels cohérents, II}
\label{section-algebraization-modules-general}

\noindent
Nous poursuivons la discussion commencée dans la
section \ref{section-algebraization-modules}.
Cette section peut être omise lors d'une première lecture.

\begin{lemma}
\label{lemma-map-kernel-cokernel-on-closed}
Dans la situation \ref{situation-algebraize}. Soit
$(\mathcal{F}_n) \to (\mathcal{F}'_n)$ un morphisme de
$\textit{Coh}(U, I\mathcal{O}_U)$
dont le noyau et le conoyau sont annulés par une puissance de $I$. Alors
\begin{enumerate}
\item $(\mathcal{F}_n)$ s'étend à $X$ si et seulement si
$(\mathcal{F}'_n)$ s'étend à $X$, et
\item $(\mathcal{F}_n)$ est le complété d'un $\mathcal{O}_U$-module cohérent
si et seulement si $(\mathcal{F}'_n)$ l'est.
\end{enumerate}
\end{lemma}

\begin{proof}
La partie (2) résulte immédiatement de
Cohomologie des schémas, lemme \ref{coherent-lemma-existence-easy}.
Pour voir la partie (1), on utilise d'abord le lemme \ref{lemma-algebraizable}
pour se ramener au cas où $A$ est complet pour la topologie $I$-adique.
Cependant, dans ce cas, (1) se ramène à (2) d'après le
lemme \ref{lemma-essential-image-completion}.
\end{proof}

\noindent
Les deux lemmes suivants étaient initialement utilisés dans la démonstration
du lemme \ref{lemma-when-done}. Nous les conservons ici pour le lecteur qui
souhaite connaître les résultats intermédiaires que l'on peut obtenir.

\begin{lemma}
\label{lemma-when-ML}
Dans la situation \ref{situation-algebraize}, soit $(\mathcal{F}_n)$ un objet
de $\textit{Coh}(U, I\mathcal{O}_U)$. Si le système projectif
$H^0(U, \mathcal{F}_n)$ satisfait la condition de Mittag-Leffler, alors les applications canoniques
$$
\widetilde{M/I^nM}|_U \to \mathcal{F}_n
$$
sont surjectives pour tout $n$, où $M$ est comme dans (\ref{equation-guess}).
\end{lemma}

\begin{proof}
La surjectivité peut se vérifier sur le germe en un point $y \in Y \setminus Z$.
Si $y$ correspond à l'idéal premier $\mathfrak q \subset A$, on peut
choisir $f \in \mathfrak a$, $f \not \in \mathfrak q$. Il suffit alors
de montrer que
$$
M_f \longrightarrow H^0(U, \mathcal{F}_n)_f = H^0(D(f), \mathcal{F}_n)
$$
est surjective, puisque $D(f)$ est affine (l'égalité vaut d'après Propriétés,
lemme \ref{properties-lemma-invert-f-sections}). Comme on dispose de la
propriété de Mittag-Leffler, on trouve que
$$
\Im(M \to H^0(U, \mathcal{F}_n)) =
\Im(H^0(U, \mathcal{F}_m) \to H^0(U, \mathcal{F}_n))
$$
pour un certain $m \geq n$. À l'aide de la suite exacte longue de cohomologie, on voit
que
$$
\Coker(H^0(U, \mathcal{F}_m) \to H^0(U, \mathcal{F}_n))
\subset
H^1(U, \Ker(\mathcal{F}_m \to \mathcal{F}_n))
$$
Comme $U = X \setminus V(\mathfrak a)$, ce $H^1$ est de torsion annulée par une puissance de
$\mathfrak a$. Ainsi, après inversion de $f$, le conoyau devient nul.
\end{proof}

\begin{lemma}
\label{lemma-when-topology}
Dans la situation \ref{situation-algebraize}, soit $(\mathcal{F}_n)$ un objet
de $\textit{Coh}(U, I\mathcal{O}_U)$. Soit $M$ comme dans (\ref{equation-guess}).
Posons
$$
\mathcal{G}_n = \widetilde{M/I^nM}.
$$
Si la topologie de limite sur $M$ coïncide avec la topologie $I$-adique, alors
$\mathcal{G}_n|_U$ est un
$\mathcal{O}_U$-module cohérent et l'application de systèmes projectifs
$$
(\mathcal{G}_n|_U) \longrightarrow (\mathcal{F}_n)
$$
est injective dans la catégorie abélienne $\textit{Coh}(U, I\mathcal{O}_U)$.
\end{lemma}

\begin{proof}
Remarquons que $\mathcal{G}_n$ est un $\mathcal{O}_X$-module quasi-cohérent
annulé par $I^n$ et que
$\mathcal{G}_{n + 1}/I^n\mathcal{G}_{n + 1} = \mathcal{G}_n$.
Considérons
$$
M_n = \Im(M \longrightarrow H^0(U, \mathcal{F}_n))
$$
L'hypothèse dit que les systèmes projectifs $(M_n)$ et
$(M/I^nM)$ sont isomorphes comme pro-objets de $\text{Mod}_A$.
Choisissons $f \in \mathfrak a$ de sorte que $D(f) \subset U$ soit un ouvert affine. On a alors
$$
(M_n)_f \subset H^0(U, \mathcal{F}_n)_f = H^0(D(f), \mathcal{F}_n)
$$
L'égalité vaut d'après Propriétés, lemme \ref{properties-lemma-invert-f-sections}.
Ainsi, $\widetilde{M_n}|_U \to \mathcal{F}_n$ est injective.
Il s'ensuit que $\widetilde{M_n}|_U$ est un module cohérent
(Cohomologie des schémas, lemme
\ref{coherent-lemma-coherent-Noetherian-quasi-coherent-sub-quotient}).
Puisque $M \to M/I^nM$ est surjective et se factorise par
$M_{n'} \to M/I^nM$ pour un certain $n' \geq n$, on trouve que $\mathcal{G}_n|_U$
est cohérent comme quotient d'un module cohérent.
Avec les remarques initiales de la démonstration, on en conclut
que $(\mathcal{G}_n|_U)$ forme bien un objet
de $\textit{Coh}(U, I\mathcal{O}_U)$.
Enfin, pour montrer l'injectivité de l'application,
il suffit de montrer que
$$
\lim (M/I^nM)_f = \lim H^0(D(f), \mathcal{G}_n) \to
\lim H^0(D(f), \mathcal{F}_n)
$$
est injective ; voir Cohomologie des schémas, lemmes
\ref{coherent-lemma-inverse-systems-abelian} et
\ref{coherent-lemma-inverse-systems-affine}.
L'injectivité de $\lim (M_n)_f \to \lim H^0(D(f), \mathcal{F}_n)$
est claire (voir ci-dessus), et notre remarque sur les pro-systèmes donne
$\lim (M_n)_f = \lim (M/I^nM)_f$. Cela achève la démonstration.
\end{proof}





\section{Une fonction de distance}
\label{section-distance}

\noindent
Soit $Y$ un schéma noethérien et soit $Z \subset Y$ une partie fermée.
Nous définissons une fonction
\begin{equation}
\label{equation-delta-Z}
\delta^Y_Z = \delta_Z : Y \longrightarrow \mathbf{Z}_{\geq 0} \cup \{\infty\}
\end{equation}
qui mesure la « distance » d'un point de $Y$ à $Z$.
Pour une discussion informelle, voir la remarque \ref{remark-discussion}.
Soit $y \in Y$. On pose $\delta_Z(y) = \infty$ si $y$ appartient
à une composante connexe de $Y$ qui ne rencontre pas $Z$.
Si $y$ appartient à une composante connexe de $Y$ qui rencontre $Z$,
on peut trouver $k \geq 0$ et un système
$$
V_0 \subset W_0 \supset V_1 \subset W_1 \supset \ldots \supset
V_k \subset W_k
$$
de sous-schémas fermés intègres de $Y$ tel que $V_0 \subset Z$
et que $y \in W_k$ soit le point générique. Posons
$c_i = \text{codim}(V_i, W_i)$ pour $i = 0, \ldots, k$
et $b_i = \text{codim}(V_{i + 1}, W_i)$ pour $i = 0, \ldots, k - 1$.
Pour un tel système, on pose
$$
\delta(V_0, W_0, V_1, \ldots, W_k) =
k +
\max_{i = 0, 1, \ldots, k}
(c_i + c_{i + 1} + \ldots + c_k - b_i - b_{i + 1} - \ldots - b_{k - 1})
$$
Ceci est $\geq k$, car on peut prendre $i = k$ et l'on a $c_k \geq 0$.
Enfin, on pose
$$
\delta_Z(y) = \min \delta(V_0, W_0, V_1, \ldots, W_k)
$$
où le minimum porte sur tous les systèmes de sous-schémas fermés intègres de $Y$
comme ci-dessus.

\begin{lemma}
\label{lemma-discussion}
Soit $Y$ un schéma noethérien et soit $Z \subset Y$ une partie fermée.
\begin{enumerate}
\item Pour $y \in Y$, on a $\delta_Z(y) = 0 \Leftrightarrow y \in Z$.
\item Les parties $\{y \in Y \mid \delta_Z(y) \leq k\}$ sont
stables par spécialisation.
\item Pour $y \in Y$ et $z \in \overline{\{y\}} \cap Z$, on a
$\dim(\mathcal{O}_{\overline{\{y\}}, z}) \geq \delta_Z(y)$.
\item Si $\delta$ est une fonction de dimension sur $Y$, alors
$\delta_Z(y) \geq \delta(y) - \delta_{max}$, où $\delta_{max}$
est la valeur maximale de $\delta$ sur $Z$.
\item Si $Y = \Spec(A)$ est le spectre d'un anneau local noethérien caténaire
d'idéal maximal $\mathfrak m$ et si $Z = \{\mathfrak m\}$, alors
$\delta_Z(y) = \dim(\overline{\{y\}})$.
\item Si $Y' \subset Y$ est un sous-schéma ouvert, alors
$\delta^{Y'}_{Y' \cap Z}(y') \geq \delta^Y_Z(y')$ pour $y' \in Y'$.
\end{enumerate}
Supposons $Y$ caténaire. Alors
\begin{enumerate}
\setcounter{enumi}{6}
\item Soit $y' \leadsto y$ une spécialisation immédiate de points de $Y$.
Si $Y$ est caténaire, alors $\delta_Z(y') \leq \delta_Z(y) + 1$.
\item Étant donné un diagramme de spécialisations
$$
\xymatrix{
& y'_0 \ar@{~>}[ld] \ar@{~>}[rd] &
& y'_1 \ar@{~>}[ld] & \ldots
& y'_{k - 1} \ar@{~>}[rd] &
\\
y_0 & &
y_1 & &
\ldots & &
y_k = y
}
$$
entre points de $Y$, avec $y_0 \in Z$ et $y_i' \leadsto y_i$
une spécialisation immédiate, on a $\delta_Z(y_k) \leq k$.
\end{enumerate}
\end{lemma}

\begin{proof}
Preuve de (1). Si $y \in Z$, on peut prendre $k = 0$ et
$V_0 = W_0 = \overline{\{y\}}$, et l'on obtient $\delta(V_0, W_0) = 0$, donc
$\delta_Z(y) = 0$. Si $y \not \in Z$, alors, pour tout système
$V_0 \subset W_0 \supset V_1 \subset W_1 \supset \ldots \subset W_k$
pour $y$, ou bien $k = 0$ et $V_0 \not = W_0$, ou bien $k > 0$.
Dans les deux cas, $\delta(V_0, W_0, \ldots, W_k) > 0$. Ainsi, $\delta_Z(y) > 0$.

\medskip\noindent
Preuve de (2). Soit $y \leadsto y'$ une spécialisation non triviale et soit
$V_0 \subset W_0 \supset V_1 \subset W_1 \supset \ldots \subset W_k$
un système pour $y$. Il y a deux cas.
Cas I : $V_k = W_k$, c'est-à-dire $c_k = 0$. Dans ce cas,
on peut poser $V'_k = W'_k = \overline{\{y'\}}$.
Un calcul facile montre que
$\delta(V_0, W_0, \ldots, V'_k, W'_k) \leq
\delta(V_0, W_0, \ldots, V_k, W_k)$
car seul $b_{k - 1}$ est remplacé par un entier plus grand.
Cas II : $V_k \not = W_k$, c'est-à-dire $c_k > 0$.
Dans ce cas, en posant $V_{k + 1} = W_{k + 1} = \overline{\{y'\}}$,
on voit que $V_0 \subset W_0 \supset \ldots \subset W_k \supset V_{k + 1}
\subset W_{k + 1}$ est un système pour $y'$. Alors $c_{k + 1} = 0$ et
$b_k > 0$, de sorte que l'on obtient
\begin{align*}
& \delta(V_0, \ldots, W_{k + 1}) \\
& =
k + 1 +
\max_{i = 0, 1, \ldots, k + 1}
(c_i + c_{i + 1} + \ldots + c_k + c_{k + 1}
- b_i - b_{i + 1} - \ldots - b_{k - 1} - b_k) \\
& =
k + 1 +
\max_{i = 0, 1, \ldots, k + 1}
(c_i + c_{i + 1} + \ldots + c_k
- b_i - b_{i + 1} - \ldots - b_{k - 1} - b_k) \\
& \leq
k + \max_{i = 0, 1, \ldots, k}
(c_i + c_{i + 1} + \ldots + c_k - b_i - b_{i + 1} - \ldots - b_{k - 1}) \\
& = \delta(V_0, \ldots, W_k)
\end{align*}
L'inégalité vaut parce que $c_k > 0$ et $b_k > 0$, ce qui entraîne en particulier
que $\delta(V_0, \ldots, W_k) \geq k + c_k \geq k + 1$.

\medskip\noindent
Preuve de (3). Étant donnés $y \in Y$ et $z \in \overline{\{y\}} \cap Z$,
on obtient le système
$$
V_0 = \overline{\{z\}} \subset W_0 = \overline{\{y\}}
$$
et $c_0 = \text{codim}(V_0, W_0) = \dim(\mathcal{O}_{\overline{\{y\}}, z})$
d'après Propriétés, lemme \ref{properties-lemma-codimension-local-ring}.
On voit ainsi que $\delta(V_0, W_0) = 0 + c_0 = c_0$, ce qui prouve
ce que l'on veut.

\medskip\noindent
Preuve de (4). Soit $\delta$ une fonction de dimension sur $Y$.
Soit $V_0 \subset W_0 \supset V_1 \subset W_1 \supset \ldots \subset W_k$
un système pour $y$. Soient $y'_i \in W_i$ et $y_i \in V_i$ les
points génériques ; ainsi, $y_0 \in Z$ et $y_k = y$. On voit alors que
$$
\delta(y_i) - \delta(y_{i - 1}) =
\delta(y'_{i - 1}) - \delta(y_{i - 1}) - \delta(y'_{i - 1}) + \delta(y_i) =
c_{i - 1} - b_{i - 1}
$$
Enfin, on a $\delta(y'_k) - \delta(y_{k - 1}) = c_k$.
On voit donc que
$$
\delta(y) - \delta(y_0) =
c_0 + \ldots + c_k - b_0 - \ldots - b_{k - 1}
$$
On en conclut que
$\delta(V_0, W_0, \ldots, W_k) \geq k + \delta(y) - \delta(y_0)$
ce qui prouve ce que l'on veut.

\medskip\noindent
Preuve de (5). La fonction $\delta(y) = \dim(\overline{\{y\}})$
est une fonction de dimension. Ainsi, $\delta(y) \leq \delta_Z(y)$ d'après la
partie (4). D'après la partie (3), on a $\delta_Z(y) \leq \delta(y)$,
et l'on a terminé.

\medskip\noindent
Preuve de (6). C'est clair, puisqu'il y a moins de systèmes à considérer
dans le calcul de $\delta^{Y'}_{Y' \cap Z}$.

\medskip\noindent
Preuve de (7). Soit
$V_0 \subset W_0 \supset V_1 \subset W_1 \supset \ldots \subset W_k$
un système pour $y$. Posons $W'_k = \overline{\{y'\}}$.
Comme $Y$ est caténaire, on voit que
$\text{codim}(V_k, W'_k) = \text{codim}(V_k, W_k) + 1$.
Il s'ensuit facilement que $\delta(V_0, \ldots, V_k, W'_k) \leq
\delta(V_0, \ldots, V_k, W_k) + 1$, ce qui prouve ce que l'on veut.

\medskip\noindent
Preuve de (8) : combinons (7) et (2).
\end{proof}

\begin{lemma}
\label{lemma-change-distance-function}
Soit $Y$ un schéma noethérien universellement caténaire. Soit $Z \subset Y$
un sous-schéma fermé. Soit $f : Y' \to Y$ un morphisme de type fini
dont toutes les fibres sont de dimension $\leq e$. Posons $Z' = f^{-1}(Z)$.
Alors
$$
\delta_Z(y) \leq \delta_{Z'}(y') + e - \text{trdeg}_{\kappa(y)}(\kappa(y'))
$$
pour $y' \in Y'$ d'image $y \in Y$.
\end{lemma}

\begin{proof}
Si $\delta_{Z'}(y') = \infty$, il n'y a rien à démontrer.
Si $\delta_{Z'}(y') < \infty$, choisissons un système
de sous-schémas fermés intègres
$$
V'_0 \subset W'_0 \supset V'_1 \subset W'_1 \supset \ldots \subset W'_k
$$
de $Y'$, avec $V'_0 \subset Z'$ et $y'$ pour point générique de $W'_k$,
tel que $\delta_{Z'}(y') = \delta(V'_0, W'_0, \ldots, W'_k)$.
Notons
$$
V_0 \subset W_0 \supset V_1 \subset W_1 \supset \ldots \subset W_k
$$
les images schématiques des schémas ci-dessus dans $Y$. Remarquons
que $y$ est le point générique de $W_k$ et que $V_0 \subset Z$.
Pour chaque $i$, considérons le diagramme
$$
\xymatrix{
V'_i \ar[r] \ar[d] & W'_i \ar[d] & V'_{i + 1} \ar[l] \ar[d] \\
V_i \ar[r] & W_i & V_{i + 1} \ar[l]
}
$$
Notons $n_i$ la dimension relative de $V'_i/V_i$ et
$m_i$ la dimension relative de $W'_i/W_i$ ; plus précisément,
ce sont les degrés de transcendance des extensions correspondantes
des corps de fonctions. Posons
$c_i = \text{codim}(V_i, W_i)$,
$c'_i = \text{codim}(V'_i, W'_i)$,
$b_i = \text{codim}(V_{i + 1}, W_i)$, et
$b'_i = \text{codim}(V'_{i + 1}, W'_i)$.
D'après la formule de dimension, on a
$$
c_i = c'_i + n_i - m_i
\quad\text{et}\quad
b_i = b'_i + n_{i + 1} - m_i
$$
Voir Morphismes, lemme \ref{morphisms-lemma-dimension-formula}.
Ainsi, $c_i - b_i = c'_i - b'_i + n_i - n_{i + 1}$. On voit donc que
\begin{align*}
& c_i + c_{i + 1} + \ldots + c_k - b_i - b_{i + 1} - \ldots - b_{k - 1} \\
& =
c'_i + c'_{i + 1} + \ldots + c'_k - b'_i - b'_{i + 1} - \ldots - b'_{k - 1}
+ n_i - n_k + c_k - c'_k \\
& =
c'_i + c'_{i + 1} + \ldots + c'_k - b'_i - b'_{i + 1} - \ldots - b'_{k - 1}
+ n_i - m_k
\end{align*}
On voit donc que
\begin{align*}
\max_{i = 0, \ldots, k}
& (c_i + c_{i + 1} + \ldots + c_k - b_i - b_{i + 1} - \ldots - b_{k - 1}) \\
& =
\max_{i = 0, \ldots, k}
(c'_i + c'_{i + 1} + \ldots + c'_k - b'_i - b'_{i + 1} - \ldots - b'_{k - 1}
+ n_i - m_k) \\
& =
\max_{i = 0, \ldots, k}
(c'_i + c'_{i + 1} + \ldots + c'_k - b'_i - b'_{i + 1} - \ldots - b'_{k - 1}
+ n_i) - m_k \\
& \leq
\max_{i = 0, \ldots, k}
(c'_i + c'_{i + 1} + \ldots + c'_k - b'_i - b'_{i + 1} - \ldots - b'_{k - 1})
+ e - m_k
\end{align*}
Comme $m_k = \text{trdeg}_{\kappa(y)}(\kappa(y'))$, on conclut que
$$
\delta(V_0, W_0, \ldots, W_k) \leq
\delta(V'_0, W'_0, \ldots, W'_k) + e - \text{trdeg}_{\kappa(y)}(\kappa(y'))
$$
comme voulu.
\end{proof}

\begin{remark}
\label{remark-discussion}
Soit $Y$ un schéma noethérien caténaire et soit $Z \subset Y$
une partie fermée. D'après le lemme \ref{lemma-discussion}, on a
$$
\delta_Z(y) \leq \min
\left\{ k \middle|
\begin{matrix}
\text{ il existe des spécialisations dans }Y \\
y_0 \leftarrow y'_0 \rightarrow y_1 \leftarrow y'_1 \rightarrow \ldots
\leftarrow y'_{k - 1} \rightarrow y_k = y \\
\text{ avec }y_0 \in Z\text{ et }y_i' \leadsto y_i
\text{ immédiates}
\end{matrix}
\right\}
$$
Nous affirmons que si $Y$ est de type fini sur un corps,
alors l'égalité a lieu. Si nous avons un jour besoin de ce résultat, nous
énoncerons ici un résultat précis et le démontrerons.
Cependant, en général, si nous définissons $\delta_Z$
par le membre de droite de cette inégalité, nous ne
savons pas si le lemme \ref{lemma-change-distance-function} reste valable.
\end{remark}

\begin{example}
\label{example-distance}
Soit $k$ un corps et $Y = \mathbf{A}^n_k$. Notons
$\delta : Y \to \mathbf{Z}_{\geq 0}$ la fonction de dimension usuelle.
\begin{enumerate}
\item Si $Z = \{z\}$ pour un point fermé $z$, alors
\begin{enumerate}
\item $\delta_Z(y) = \delta(y)$ si $y \leadsto z$, et
\item $\delta_Z(y) = \delta(y) + 1$ si $y \not \leadsto z$.
\end{enumerate}
\item Si $Z$ est une sous-variété fermée et $W = \overline{\{y\}}$, alors
\begin{enumerate}
\item $\delta_Z(y) = 0$ si $W \subset Z$,
\item $\delta_Z(y) = \dim(W) - \dim(Z)$ si $Z$ est contenue dans $W$,
\item $\delta_Z(y) = 1$ si $\dim(W) \leq \dim(Z)$ et $W \not \subset Z$,
\item $\delta_Z(y) = \dim(W) - \dim(Z) + 1$ si $\dim(W) > \dim(Z)$
et $Z \not \subset W$.
\end{enumerate}
\end{enumerate}
Une généralisation du cas (1) s'obtient lorsque $Y$ est de type fini sur un corps
et que $Z = \{z\}$ est un point fermé. Alors $\delta_Z(y) = \delta(y) + t$,
où $t$ est la longueur minimale d'une chaîne de courbes reliant
$z$ à un point fermé de $\overline{\{y\}}$.
\end{example}







\section{Algébrisation des modules formels cohérents, III}
\label{section-uniqueness}

\noindent
Nous poursuivons la discussion commencée aux sections
\ref{section-algebraization-modules} et
\ref{section-algebraization-modules-general}.
Nous utiliserons la fonction de distance de la section \ref{section-distance}
pour formuler certaines conditions naturelles sur les
modules formels cohérents dans la situation \ref{situation-algebraize}.

\medskip\noindent
Dans la situation \ref{situation-algebraize}, étant donné un point $y \in U \cap Y$,
nous pouvons considérer la complétion $I$-adique
$$
\mathcal{O}_{X, y}^\wedge = \lim \mathcal{O}_{X, y}/I^n\mathcal{O}_{X, y}
$$
C'est un anneau local noethérien complet pour la topologie
$I\mathcal{O}_{X, y}^\wedge$-adique, d'idéal maximal $\mathfrak m_y^\wedge$ ; voir
Algèbre, section \ref{algebra-section-completion-noetherian}.
Soit $(\mathcal{F}_n)$ un objet de $\textit{Coh}(U, I\mathcal{O}_U)$.
Définissons le « germe » de $(\mathcal{F}_n)$ en $y$ par la formule
$$
\mathcal{F}_y^\wedge = \lim \mathcal{F}_{n, y}
$$
C'est un module de type fini sur $\mathcal{O}_{X, y}^\wedge$. Voir
Algèbre, lemmes \ref{algebra-lemma-limit-complete} et
\ref{algebra-lemma-finite-over-complete-ring}.

\begin{definition}
\label{definition-s-d-inequalities}
Dans la situation \ref{situation-algebraize}, soit $(\mathcal{F}_n)$ un objet
de $\textit{Coh}(U, I\mathcal{O}_U)$. Soient $a, b$ des entiers.
Soit $\delta^Y_Z$ comme dans (\ref{equation-delta-Z}).
Nous disons que
{\it $(\mathcal{F}_n)$ satisfait les inégalités $(a, b)$} si, pour
$y \in U \cap Y$ et un idéal premier $\mathfrak p \subset \mathcal{O}_{X, y}^\wedge$
tel que $\mathfrak p \not \in V(I\mathcal{O}_{X, y}^\wedge)$,
\begin{enumerate}
\item si $V(\mathfrak p) \cap V(I\mathcal{O}_{X, y}^\wedge) \not =
\{\mathfrak m_y^\wedge\}$, alors
$$
\text{depth}((\mathcal{F}^\wedge_y)_\mathfrak p) + \delta^Y_Z(y) \geq a
\quad\text{ou}\quad
\text{depth}((\mathcal{F}^\wedge_y)_\mathfrak p) +
\dim(\mathcal{O}_{X, y}^\wedge/\mathfrak p) + \delta^Y_Z(y) > b
$$
\item si $V(\mathfrak p) \cap V(I\mathcal{O}_{X, y}^\wedge) =
\{\mathfrak m_y^\wedge\}$, alors
$$
\text{depth}((\mathcal{F}^\wedge_y)_\mathfrak p) + \delta^Y_Z(y) > a
$$
\end{enumerate}
Nous disons que {\it $(\mathcal{F}_n)$ satisfait les inégalités strictes $(a, b)$}
si, pour $y \in U \cap Y$ et un idéal premier
$\mathfrak p \subset \mathcal{O}_{X, y}^\wedge$ tel que
$\mathfrak p \not \in V(I\mathcal{O}_{X, y}^\wedge)$,
on a
$$
\text{depth}((\mathcal{F}^\wedge_y)_\mathfrak p) + \delta^Y_Z(y) > a
\quad\text{ou}\quad
\text{depth}((\mathcal{F}^\wedge_y)_\mathfrak p) +
\dim(\mathcal{O}_{X, y}^\wedge/\mathfrak p) + \delta^Y_Z(y) > b
$$
\end{definition}

\noindent
Voici quelques observations élémentaires.

\begin{lemma}
\label{lemma-elementary}
Dans la situation \ref{situation-algebraize}, soit $(\mathcal{F}_n)$ un objet
de $\textit{Coh}(U, I\mathcal{O}_U)$. Soient $a, b$ des entiers.
\begin{enumerate}
\item Si $(\mathcal{F}_n)$ est annulé par une puissance de $I$, alors
$(\mathcal{F}_n)$ satisfait les inégalités $(a, b)$ pour tous $a, b$.
\item Si $(\mathcal{F}_n)$ satisfait les inégalités $(a + 1, b)$, alors
$(\mathcal{F}_n)$ satisfait les inégalités strictes $(a, b)$.
\end{enumerate}
Si $\text{cd}(A, I) \leq d$ et si $A$ possède un complexe dualisant, alors
\begin{enumerate}
\item[(3)] $(\mathcal{F}_n)$ satisfait les inégalités $(s, s + d)$
si et seulement si, pour tout $y \in U \cap Y$, le sextuplet
$\mathcal{O}_{X, y}^\wedge, I\mathcal{O}_{X, y}^\wedge,
\{\mathfrak m_y^\wedge\}, \mathcal{F}_y^\wedge, s - \delta^Y_Z(y), d$
satisfait les hypothèses de la situation \ref{situation-bootstrap}.
\item[(4)]
Si $(\mathcal{F}_n)$ satisfait les inégalités strictes $(s, s + d)$, alors
$(\mathcal{F}_n)$ satisfait les inégalités $(s, s + d)$.
\end{enumerate}
\end{lemma}

\begin{proof}
C'est immédiat, sauf pour la partie (4), qui découle du
lemme \ref{lemma-bootstrap-bis-bis} et de la reformulation de (3).
\end{proof}

\begin{lemma}
\label{lemma-explain-2-3-cd-1}
Dans la situation \ref{situation-algebraize}, soit $(\mathcal{F}_n)$ un objet
de $\textit{Coh}(U, I\mathcal{O}_U)$. Si $\text{cd}(A, I) = 1$, alors
$\mathcal{F}$ satisfait les inégalités $(2, 3)$ si et seulement si
$$
\text{depth}((\mathcal{F}^\wedge_y)_\mathfrak p) +
\dim(\mathcal{O}_{X, y}^\wedge/\mathfrak p) + \delta^Y_Z(y) > 3
$$
pour tout $y \in U \cap Y$ et tout $\mathfrak p \subset \mathcal{O}_{X, y}^\wedge$
tel que $\mathfrak p \not \in V(I\mathcal{O}_{X, y}^\wedge)$.
\end{lemma}

\begin{proof}
Observons que, pour un idéal premier $\mathfrak p \subset \mathcal{O}_{X, y}^\wedge$,
$\mathfrak p \not \in V(I\mathcal{O}_{X, y}^\wedge)$,
nous avons $V(\mathfrak p) \cap V(I\mathcal{O}_{X, y}^\wedge) =
\{\mathfrak m_y^\wedge\}
\Leftrightarrow \dim(\mathcal{O}_{X, y}^\wedge/\mathfrak p) = 1$
car $\text{cd}(A, I) = 1$.
Voir Cohomologie locale, lemmes
\ref{local-cohomology-lemma-cd-change-rings} et
\ref{local-cohomology-lemma-cd-bound-dim-local}.
Considérons les trois nombres
$\alpha = \text{depth}((\mathcal{F}^\wedge_y)_\mathfrak p) \geq 0$,
$\beta = \dim(\mathcal{O}_{X, y}^\wedge/\mathfrak p) \geq 1$, et
$\gamma = \delta^Y_Z(y) \geq 1$.
La définition \ref{definition-s-d-inequalities} exige alors
\begin{enumerate}
\item si $\beta > 1$, alors
$\alpha + \gamma \geq 2$ ou $\alpha + \beta + \gamma > 3$, et
\item si $\beta = 1$, alors $\alpha + \gamma > 2$.
\end{enumerate}
Il est immédiat que cela équivaut à
$\alpha + \beta + \gamma > 3$.
\end{proof}

\noindent
Dans la suite de cette section, que nous suggérons au lecteur de sauter en première
lecture, nous montrerons que, lorsque $A$ est $I$-adiquement complet,
la catégorie des objets $(\mathcal{F}_n)$ de $\textit{Coh}(U, I\mathcal{O}_U)$
qui s'étendent à $X$ et satisfont aux
inégalités $(1, 1 + \text{cd}(A, I))$ strictes
est équivalente à une sous-catégorie pleine de la catégorie des
$\mathcal{O}_U$-modules cohérents.

\begin{lemma}
\label{lemma-sanity}
Dans la situation \ref{situation-algebraize}, soit $\mathcal{F}$ un
$\mathcal{O}_U$-module cohérent et soit $d \geq 1$. Supposons que
\begin{enumerate}
\item $A$ soit $I$-adiquement complet, possède un complexe dualisant, et que
$\text{cd}(A, I) \leq d$,
\item le complété $\mathcal{F}^\wedge$ de $\mathcal{F}$
satisfasse aux inégalités $(1, 1 + d)$ strictes.
\end{enumerate}
Soit $x \in X$ un point. Posons $W = \overline{\{x\}}$.
Si $W \cap Y$ a une composante irréductible contenue dans $Z$
et une qui ne l'est pas, alors $\text{depth}(\mathcal{F}_x) \geq 1$.
\end{lemma}

\begin{proof}
Soit $W \cap Y = W_1 \cup \ldots \cup W_n$ sa décomposition en
composantes irréductibles. Par hypothèse, après renumérotation, nous pouvons trouver
$0 < m < n$ tel que $W_1, \ldots, W_m  \subset Z$ et
$W_{m + 1}, \ldots, W_n \not \subset Z$. Nous en concluons que
$$
W \cap Y \setminus
\left((W_1 \cup \ldots \cup W_m) \cap (W_{m + 1} \cup \ldots \cup W_n)\right)
$$
n'est pas connexe. Par le lemme \ref{lemma-connected}, nous pouvons trouver
$1 \leq i \leq m < j \leq n$ et
$z \in W_i \cap W_j$ tels que $\dim(\mathcal{O}_{W, z}) \leq d + 1$.
Choisissons une spécialisation immédiate $y \leadsto z$ avec
$y \in W_j$, $y \not \in Z$ ; l'existence de $y$ résulte de
Propriétés, lemme \ref{properties-lemma-complement-closed-point-Jacobson}.
Observons que $\delta^Y_Z(y) = 1$ et $\dim(\mathcal{O}_{W, y}) \leq d$.
Soit $\mathfrak p \subset \mathcal{O}_{X, y}$ l'idéal premier correspondant à $x$.
Soit $\mathfrak p' \subset \mathcal{O}_{X, y}^\wedge$ un idéal premier minimal
au-dessus de $\mathfrak p\mathcal{O}_{X, y}^\wedge$. Alors
$$
\text{depth}(\mathcal{F}_x) =
\text{depth}((\mathcal{F}^\wedge_y)_{\mathfrak p'})
\quad\text{et}\quad
\dim(\mathcal{O}_{W, y}) = \dim(\mathcal{O}_{X, y}^\wedge/\mathfrak p')
$$
Voir Algèbre, lemme \ref{algebra-lemma-apply-grothendieck-module}, et
Cohomologie locale, lemme \ref{local-cohomology-lemma-change-completion}.
La conclusion se lit alors directement sur les inégalités données.
\end{proof}

\begin{lemma}
\label{lemma-recover}
Dans la situation \ref{situation-algebraize}, soit $\mathcal{F}$ un
$\mathcal{O}_U$-module cohérent et soit $d \geq 1$. Supposons que
\begin{enumerate}
\item $A$ soit $I$-adiquement complet, possède un complexe dualisant, et que
$\text{cd}(A, I) \leq d$,
\item le complété $\mathcal{F}^\wedge$ de $\mathcal{F}$
satisfasse aux inégalités $(1, 1+ d)$ strictes, et
\item pour $x \in U$ tel que $\overline{\{x\}} \cap Y \subset Z$,
nous ayons $\text{depth}(\mathcal{F}_x) \geq 2$.
\end{enumerate}
Alors $H^0(U, \mathcal{F}) \to \lim H^0(U, \mathcal{F}/I^n\mathcal{F})$
est un isomorphisme.
\end{lemma}

\begin{proof}
Nous allons le démontrer en vérifiant que le lemme \ref{lemma-application-H0} s'applique.
Soit donc $x \in \text{Ass}(\mathcal{F})$ avec $x \not \in Y$.
Posons $W = \overline{\{x\}}$.
La condition (3) montre que $W \cap Y \not \subset Z$.
Le lemme \ref{lemma-sanity} montre qu'aucune composante
irréductible de $W \cap Y$ n'est contenue dans $Z$.
Ainsi, si $z \in W \cap Z$, il existe une
spécialisation immédiate $y \leadsto z$, avec $y \in W \cap Y$ et $y \not \in Z$.
Pour l'existence de $y$, utiliser
Propriétés, lemme \ref{properties-lemma-complement-closed-point-Jacobson}.
Alors $\delta^Y_Z(y) = 1$ et l'hypothèse
implique que $\dim(\mathcal{O}_{W, y}) > d$.
Par conséquent $\dim(\mathcal{O}_{W, z}) > 1 + d$, ce qui conclut.
\end{proof}

\begin{lemma}
\label{lemma-fully-faithful-inequalities}
Dans la situation \ref{situation-algebraize}, soit $\mathcal{F}$ un
$\mathcal{O}_U$-module cohérent et soit $d \geq 1$. Supposons que
\begin{enumerate}
\item $A$ soit $I$-adiquement complet, possède un complexe dualisant, et que
$\text{cd}(A, I) \leq d$,
\item le complété $\mathcal{F}^\wedge$ de $\mathcal{F}$
satisfasse aux inégalités $(1, 1 + d)$ strictes, et
\item pour $x \in U$ tel que $\overline{\{x\}} \cap Y \subset Z$,
nous ayons $\text{depth}(\mathcal{F}_x) \geq 2$.
\end{enumerate}
Alors l'application
$$
\Hom_U(\mathcal{G}, \mathcal{F})
\longrightarrow
\Hom_{\textit{Coh}(U, I\mathcal{O}_U)}(\mathcal{G}^\wedge, \mathcal{F}^\wedge)
$$
est bijective pour tout $\mathcal{O}_U$-module cohérent $\mathcal{G}$. 
\end{lemma}

\begin{proof}
Posons $\mathcal{H} = \SheafHom_{\mathcal{O}_U}(\mathcal{G}, \mathcal{F})$.
À l'aide de Cohomologie des schémas, lemme
\ref{coherent-lemma-hom-into-depth}, ou de
Compléments sur l'algèbre, lemme \ref{more-algebra-lemma-hom-into-depth},
nous voyons que le complété de $\mathcal{H}$
satisfait aux inégalités $(1, 1 + d)$ strictes et que, pour
$x \in U$ tel que $\overline{\{x\}} \cap Y \subset Z$,
nous avons $\text{depth}(\mathcal{H}_x) \geq 2$. Nous omettons les détails.
Par le lemme \ref{lemma-recover}, nous avons donc
$$
\Hom_U(\mathcal{G}, \mathcal{F}) =
H^0(U, \mathcal{H}) =
\lim H^0(U, \mathcal{H}/\mathcal{I}^n\mathcal{H}) =
\Mor_{\textit{Coh}(U, I\mathcal{O}_U)}
(\mathcal{G}^\wedge, \mathcal{F}^\wedge)
$$
Voir Cohomologie des schémas, lemme \ref{coherent-lemma-completion-internal-hom},
pour la dernière égalité.
\end{proof}

\begin{lemma}
\label{lemma-construct-unique}
Dans la situation \ref{situation-algebraize}, soit $(\mathcal{F}_n)$ un
objet de $\textit{Coh}(U, I\mathcal{O}_U)$ et soit $d \geq 1$. Supposons que
\begin{enumerate}
\item $A$ soit $I$-adiquement complet, possède un complexe dualisant, et que
$\text{cd}(A, I) \leq d$,
\item $(\mathcal{F}_n)$ soit le complété d'un $\mathcal{O}_U$-module cohérent,
\item $(\mathcal{F}_n)$ satisfasse aux inégalités $(1, 1 + d)$ strictes.
\end{enumerate}
Alors il existe un unique $\mathcal{O}_U$-module cohérent $\mathcal{F}$
dont le complété est $(\mathcal{F}_n)$ et tel que, pour
$x \in U$ avec $\overline{\{x\}} \cap Y \subset Z$,
nous ayons $\text{depth}(\mathcal{F}_x) \geq 2$.
\end{lemma}

\begin{proof}
Choisissons un $\mathcal{O}_U$-module cohérent $\mathcal{F}$ dont le
complété est $(\mathcal{F}_n)$. Posons
$T = \{x \in U \mid \overline{\{x\}} \cap Y \subset Z\}$.
Nous construirons $\mathcal{F}$ en appliquant Cohomologie locale,
lemme \ref{local-cohomology-lemma-make-S2-along-T-simple},
à $\mathcal{F}$ et $T$.
L'unicité résultera alors de la propriété d'unicité des morphismes
du lemme \ref{lemma-fully-faithful-inequalities}.

\medskip\noindent
Puisque $T$ est stable par spécialisation dans $U$, il suffit de
vérifier ce qui suit. Si $x' \leadsto x$ est une
spécialisation immédiate de points de $U$ avec $x \in T$
et $x' \not \in T$, alors $\text{depth}(\mathcal{F}_{x'}) \geq 1$.
Posons $W = \overline{\{x\}}$ et $W' = \overline{\{x'\}}$.
Comme $x' \not \in T$, l'ensemble $W' \cap Y$ n'est pas contenu dans $Z$.
Si $W' \cap Y$ contient une composante irréductible contenue dans $Z$,
le lemme \ref{lemma-sanity} permet de conclure.
Sinon, choisissons une composante irréductible $W_1$ de $W \cap Y$ et
une composante irréductible $W'_1$ de $W' \cap Y$ telles que $W_1 \subset W'_1$.
Soit $z \in W_1$ le point générique. Soit $y \leadsto z$, avec $y \in W'_1$,
une spécialisation immédiate telle que $y \not \in Z$ ; l'existence de $y$
résulte de $W'_1 \not \subset Z$ (voir ci-dessus) et de
Propriétés, lemme \ref{properties-lemma-complement-closed-point-Jacobson}.
Nous avons alors $z \in Z$, $x \leadsto z$,
$x' \leadsto y \leadsto z$, $y \in Y \setminus Z$, et $\delta^Y_Z(y) = 1$.
Par Cohomologie locale, lemme \ref{local-cohomology-lemma-cd-bound-dim-local},
et puisque $z$
est un point générique de $W \cap Y$, nous avons
$\dim(\mathcal{O}_{W, z}) \leq d$.
Comme $x' \leadsto x$ est une spécialisation immédiate, nous avons
$\dim(\mathcal{O}_{W', z}) \leq d + 1$.
Puisque $y \not = z$, nous en concluons que
$\dim(\mathcal{O}_{W', y}) \leq d$.
Si $\text{depth}(\mathcal{F}_{x'}) = 0$, nous obtiendrions
une contradiction avec l'hypothèse (3) ; nous omettons les détails du passage
de $\mathcal{O}_{X, y}$ à son complété.
Cela achève la démonstration.
\end{proof}








\section{Algébrisation des modules formels cohérents, IV}
\label{section-algebraization-modules-local}

\noindent
Dans cette section, nous démontrons deux versions plus fortes du
lemme \ref{lemma-algebraization-principal-variant}
dans le cas local, à savoir les lemmes \ref{lemma-algebraization-principal}
et \ref{lemma-algebraization-principal-bis}.
Bien que ces lemmes soient rendus caducs par la proposition plus générale
\ref{proposition-cd-1}, leurs démonstrations sont sensiblement
plus simples.

\begin{lemma}
\label{lemma-algebraization-principal}
Dans la situation \ref{situation-algebraize}, soit $(\mathcal{F}_n)$ un objet
de $\textit{Coh}(U, I\mathcal{O}_U)$. Supposons que
\begin{enumerate}
\item $A$ soit local et que $\mathfrak a = \mathfrak m$ soit l'idéal maximal,
\item $A$ possède un complexe dualisant,
\item $I = (f)$ soit un idéal principal engendré par un élément non diviseur de zéro $f \in \mathfrak m$,
\item $\mathcal{F}_n$ soit un
$\mathcal{O}_U/f^n\mathcal{O}_U$-module localement libre de type fini,
\item si $\mathfrak p \in V(f) \setminus \{\mathfrak m\}$, alors
$\text{depth}((A/f)_\mathfrak p) + \dim(A/\mathfrak p) > 1$, et
\item si $\mathfrak p \not \in V(f)$ et
$V(\mathfrak p) \cap V(f) \not = \{\mathfrak m\}$, alors
$\text{depth}(A_\mathfrak p) + \dim(A/\mathfrak p) > 3$.
\end{enumerate}
Alors $(\mathcal{F}_n)$ s'étend canoniquement à $X$. En particulier, si $A$
est complet, alors $(\mathcal{F}_n)$ est le complété d'un
$\mathcal{O}_U$-module cohérent.
\end{lemma}

\begin{proof}
Nous allons le démontrer en vérifiant les hypothèses (a), (b) et (c) du
lemme \ref{lemma-when-done}.

\medskip\noindent
Puisque $\mathcal{F}_n$ est localement libre sur $\mathcal{O}_U/f^n\mathcal{O}_U$,
nous avons des suites exactes courtes
$0 \to \mathcal{F}_n \to \mathcal{F}_{n + 1} \to \mathcal{F}_1 \to 0$
pour tout $n$. La condition (b) résulte donc de
Cohomologie, lemme \ref{cohomology-lemma-topology-I-adic-f}.

\medskip\noindent
Par récurrence sur $n$ et à l'aide des suites exactes courtes
$0 \to A/f^n \to A/f^{n + 1} \to A/f \to 0$, nous voyons que
les idéaux premiers associés de $A/f^nA$ coïncident avec les idéaux premiers
associés de $A/fA$. Comme les points associés de $\mathcal{F}_n$
correspondent aux idéaux premiers associés de $A/f^nA$ distincts de $\mathfrak m$
par l'hypothèse (3), nous en concluons que
$M_n = H^0(U, \mathcal{F}_n)$ est un $A$-module de type fini d'après (5) et
Cohomologie locale, proposition \ref{local-cohomology-proposition-kollar}.
L'hypothèse (c) est donc satisfaite.

\medskip\noindent
Pour achever la démonstration, il suffit de montrer qu'il existe $n > 1$
tel que l'image de
$$
H^1(U, \mathcal{F}_n) \longrightarrow H^1(U, \mathcal{F}_1)
$$
soit de longueur finie comme $A$-module. Cela impliquera en effet l'hypothèse (a)
par Cohomologie, lemme \ref{cohomology-lemma-ML-better}. L'image est indépendante
de $n$ pour $n$ assez grand d'après le lemme \ref{lemma-ML-local}.
Soit $\omega_A^\bullet$ un complexe dualisant normalisé pour $A$.
D'après le théorème de dualité locale et la dualité de Matlis
(Complexes dualisants, lemme \ref{dualizing-lemma-special-case-local-duality}
et proposition \ref{dualizing-proposition-matlis}),
notre assertion équivaut à dire que l'image de
$$
\text{Ext}^{-2}_A(M_1, \omega_A^\bullet) \to
\text{Ext}^{-2}_A(M_n, \omega_A^\bullet)
$$
est de longueur finie pour $n \gg 1$. Les modules considérés sont des
$A$-modules de type fini à support dans $V(f)$. Il suffit donc de montrer que cette
application est nulle après localisation en un idéal premier $\mathfrak q$
contenant $f$ et distinct de $\mathfrak m$.
Soit $\omega_{A_\mathfrak q}^\bullet$ un complexe
dualisant normalisé sur $A_\mathfrak q$ et rappelons que
$\omega_{A_\mathfrak q}^\bullet =
(\omega_A^\bullet)_\mathfrak q[\dim(A/\mathfrak q)]$ d'après
Complexes dualisants, lemme \ref{dualizing-lemma-dimension-function}.
En utilisant la structure locale de $\mathcal{F}_n$ donnée en (4),
nous voyons qu'il suffit de montrer l'annulation de
$$
\text{Ext}^{-2 + \dim(A/\mathfrak q)}_{A_\mathfrak q}(
A_\mathfrak q/f, \omega_{A_\mathfrak q}^\bullet)
\to
\text{Ext}^{-2 + \dim(A/\mathfrak q)}_{A_\mathfrak q}(
A_\mathfrak q/f^n, \omega_{A_\mathfrak q}^\bullet)
$$
pour $n$ assez grand. Si $\dim(A/\mathfrak q) > 3$, cela résulte immédiatement de
Cohomologie locale, lemme \ref{local-cohomology-lemma-sitting-in-degrees}.
Dans les autres cas, nous utiliserons la suite exacte longue
$$
\ldots
\xrightarrow{f^n}
H^{-1}(\omega_{A_\mathfrak q}^\bullet)
\to
\text{Ext}^{-1}_{A_\mathfrak q}(
A_\mathfrak q/f^n, \omega_{A_\mathfrak q}^\bullet) \to
H^0(\omega_{A_\mathfrak q}^\bullet)
\xrightarrow{f^n}
H^0(\omega_{A_\mathfrak q}^\bullet)
\to
\text{Ext}^0_{A_\mathfrak q}(
A_\mathfrak q/f^n, \omega_{A_\mathfrak q}^\bullet) \to 0
$$
Si $\dim(A/\mathfrak q) = 2$, alors
$H^0(\omega_{A_\mathfrak q}^\bullet) = 0$
car $\text{depth}(A_\mathfrak q) \geq 1$, puisque
$f$ est un élément non diviseur de zéro.
La suite exacte longue montre donc que la condition s'écrit
$$
f^{n - 1} :
H^{-1}(\omega_{A_\mathfrak q}^\bullet)/f \to
H^{-1}(\omega_{A_\mathfrak q}^\bullet)/f^n
$$
est nulle. Or $H^{-1}(\omega^\bullet_\mathfrak q)$ est un module de type fini dont le support
est contenu dans les idéaux premiers $\mathfrak p \subset A_\mathfrak q$ tels que
$\text{depth}(A_\mathfrak p) + \dim((A/\mathfrak p)_\mathfrak q) \leq 1$.
Puisque $\dim((A/\mathfrak p)_\mathfrak q) = \dim(A/\mathfrak p) - 2$,
la condition (6) nous dit que ces idéaux premiers sont contenus dans $V(f)$.
D'où l'annulation voulue pour $n$ assez grand.
Enfin, si $\dim(A/\mathfrak q) = 1$, la condition (5), combinée
au fait que $f$ est un élément non diviseur de zéro,
assure que $A_\mathfrak q$ est de profondeur au moins $2$. Par conséquent,
$H^0(\omega_{A_\mathfrak q}^\bullet) =
H^{-1}(\omega_{A_\mathfrak q}^\bullet) = 0$,
et la suite exacte longue montre que l'assertion
équivaut à l'annulation de
$$
f^{n - 1} :
H^{-2}(\omega_{A_\mathfrak q}^\bullet)/f \to
H^{-2}(\omega_{A_\mathfrak q}^\bullet)/f^n
$$
Or $H^{-2}(\omega^\bullet_\mathfrak q)$ est un module de type fini dont le support
est contenu dans les idéaux premiers $\mathfrak p \subset A_\mathfrak q$
tels que $\text{depth}(A_\mathfrak p) + \dim((A/\mathfrak p)_\mathfrak q)
\leq 2$. D'après la condition (6), tous ces idéaux premiers sont contenus dans $V(f)$.
D'où l'annulation voulue pour $n$ assez grand.
\end{proof}

\begin{remark}
\label{remark-interesting-case}
Soit $(A, \mathfrak m)$ un anneau local noethérien, intègre, normal, complet,
de dimension $\geq 4$, et soit $f \in \mathfrak m$ non nul.
Les hypothèses (1), (2), (3), (5) et (6) du
lemme \ref{lemma-algebraization-principal}
sont alors satisfaites. Les fibrés vectoriels
sur la complétion formelle de $U$ le long de $U \cap V(f)$
peuvent donc être algébrisés. Dans le lemme \ref{lemma-algebraization-principal-bis},
nous généraliserons ceci à des modules formels cohérents plus généraux ;
comparer également avec la remarque \ref{remark-interesting-case-bis}.
\end{remark}

\begin{lemma}
\label{lemma-helper-algebraize}
Dans la situation \ref{situation-algebraize}, soit $(M_n)$ un système projectif de
$A$-modules comme dans le lemme \ref{lemma-system-of-modules}, et soit
$(\mathcal{F}_n)$ l'objet correspondant de
$\textit{Coh}(U, I\mathcal{O}_U)$. Soient $d \geq \text{cd}(A, I)$
et $s \geq 0$ des entiers.
Avec les notations ci-dessus, supposons que
\begin{enumerate}
\item $A$ soit local, d'idéal maximal $\mathfrak m = \mathfrak a$,
\item $A$ possède un complexe dualisant, et
\item $(\mathcal{F}_n)$ satisfasse aux inégalités $(s, s + d)$
(définition \ref{definition-s-d-inequalities}).
\end{enumerate}
Soit $E$ une enveloppe injective du corps résiduel de $A$. Alors, pour $i \leq s$,
il existe un $A$-module de type fini $N$, annulé par une puissance
de $I$, et, pour $n \gg 0$, des applications compatibles
$$
H^i_\mathfrak m(M_n) \to \Hom_A(N, E)
$$
dont les conoyaux sont des $A$-modules de longueur finie et dont les noyaux $K_n$
forment un système projectif tel que $\Im(K_{n''} \to K_{n'})$ soit de longueur
finie pour $n'' \gg n' \gg 0$.
\end{lemma}

\begin{proof}
Soit $\omega_A^\bullet$ un complexe dualisant normalisé. Alors
$\delta^Y_Z = \delta$ est la fonction de dimension associée à
ce complexe dualisant.
Observons que $\Ext^{-i}_A(M_n, \omega_A^\bullet)$ est un $A$-module de type fini
annulé par $I^n$. Fixons $0 \leq i \leq s$.
Nous trouverons ci-dessous $n_1 > n_0 > 0$ tels qu'en posant
$$
N = \Im(\Ext^{-i}_A(M_{n_0}, \omega_A^\bullet) \to
\Ext^{-i}_A(M_{n_1}, \omega_A^\bullet))
$$
les noyaux des applications
$$
N \to \Ext^{-i}_A(M_n, \omega_A^\bullet),\quad n \geq n_1
$$
soient des $A$-modules de longueur finie et que les conoyaux $Q_n$ forment un
système tel que $\Im(Q_{n'} \to Q_{n''})$ soit de longueur finie
pour $n'' \gg n' \gg n_1$. Cela revient à dire que
le système $\{\Ext^{-i}_A(M_n, \omega_A^\bullet)\}_{n \geq 1}$
est essentiellement constant dans la catégorie quotient de la catégorie des
$A$-modules de type fini par la sous-catégorie de Serre des $A$-modules de longueur finie.
D'après le théorème de dualité locale
(Complexes dualisants, lemme \ref{dualizing-lemma-special-case-local-duality})
et la dualité de Matlis
(Complexes dualisants, proposition \ref{dualizing-proposition-matlis}),
nous en concluons qu'il existe des applications
$$
H^i_\mathfrak m(M_n) \to \Hom_A(N, E),\quad n \geq n_1
$$
comme dans l'énoncé du lemme.

\medskip\noindent
Choisissons $f \in \mathfrak m$. Soit $B = A_f^\wedge$ le complété $I$-adique
de la localisation $A_f$. Rappelons que
$\omega_{A_f}^\bullet = \omega_A^\bullet \otimes_A A_f$
et $\omega_B^\bullet = \omega_A^\bullet \otimes_A B$ sont des complexes
dualisants (Complexes dualisants, lemme \ref{dualizing-lemma-dualizing-localize}
et \ref{dualizing-lemma-completion-henselization-dualizing}).
Soit $M$ le $B$-module de type fini $\lim M_{n, f}$ (comparer avec la
discussion dans Cohomologie des schémas, lemme
\ref{coherent-lemma-inverse-systems-affine}). Alors
$$
\Ext^{-i}_A(M_n, \omega_A^\bullet)_f =
\Ext^{-i}_{A_f}(M_{n, f}, \omega_{A_f}^\bullet) =
\Ext^{-i}_B(M/I^n M, \omega_B^\bullet)
$$
Comme $\mathfrak m$ peut être engendré par un nombre fini d'éléments $f \in \mathfrak m$,
il suffit de montrer que, pour chacun de ces $f$, le système
$$
\{\Ext^{-i}_B(M/I^n M, \omega_B^\bullet)\}_{n \geq 1}
$$
est essentiellement constant. Nous omettons quelques détails.

\medskip\noindent
Soit $\mathfrak q \subset IB$ un idéal premier. Alors $\mathfrak q$ correspond
à un point $y \in U \cap Y$. Observons que
$\delta(\mathfrak q) = \dim(\overline{\{y\}})$
est aussi la valeur de la fonction de dimension associée à $\omega_B^\bullet$
(nous omettons les détails ; utiliser que $\omega_B^\bullet$ s'obtient de
$\omega_A^\bullet$ par extension des scalaires à $B$). L'hypothèse
(3) garantit, via le lemme \ref{lemma-elementary},
que le lemme \ref{lemma-algebraize-local-cohomology-bis}
s'applique à
$B_\mathfrak q, IB_\mathfrak q, \mathfrak qB_\mathfrak q, M_\mathfrak q$
en remplaçant $s$ par $s - \delta(y)$. Nous obtenons
$$
H^{i - \delta(\mathfrak q)}_{\mathfrak qB_\mathfrak q}(M_\mathfrak q) =
\lim H^{i - \delta(\mathfrak q)}_{\mathfrak qB_\mathfrak q}(
(M/I^nM)_\mathfrak q)
$$
et ce module est annulé par une puissance de $I$.
Le lemme \ref{lemma-terrific} montre que les systèmes projectifs
$H^{i - \delta(\mathfrak q)}_{\mathfrak qB_\mathfrak q}((M/I^nM)_\mathfrak q)$
sont essentiellement constants, de valeur
$H^{i - \delta(\mathfrak q)}_{\mathfrak qB_\mathfrak q}(M_\mathfrak q)$.
Comme $(\omega_B^\bullet)_\mathfrak q[-\delta(\mathfrak q)]$ est un complexe
dualisant normalisé sur $B_\mathfrak q$, le théorème de dualité locale
montre que le système
$$
\Ext^{-i}_B(M/I^n M, \omega_B^\bullet)_\mathfrak q
$$
est essentiellement constant, de valeur
$\Ext^{-i}_B(M, \omega_B^\bullet)_\mathfrak q$.

\medskip\noindent
Pour achever la démonstration, passons au global comme dans la démonstration du
lemme \ref{lemma-bootstrap} ; l'argument est ici plus simple,
car nous connaissons déjà la valeur de notre système. Considérons les applications
$$
\alpha_n :
\Ext^{-i}_B(M/I^n M, \omega_B^\bullet)
\longrightarrow
\Ext^{-i}_B(M, \omega_B^\bullet)
$$
lorsque $n$ varie. D'après ce qui précède, pour tout $\mathfrak q$, il existe
un $n$ tel que $\alpha_n$ soit surjective après localisation en $\mathfrak q$.
Comme $B$ est noethérien et $\Ext^{-i}_B(M, \omega_B^\bullet)$
un module de type fini, il existe un $n$ tel que $\alpha_n$ soit surjective.
Pour tout $n$ tel que $\alpha_n$ soit surjective, étant donné un idéal premier
$\mathfrak q \in V(IB)$, il existe un $n' > n$ tel que
$\Ker(\alpha_n)$ s'envoie sur zéro dans $\Ext^{-i}(M/I^{n'}M, \omega_B^\bullet)$,
au moins après localisation en $\mathfrak q$.
Comme $\Ker(\alpha_n)$ est un $A$-module de type fini et que les supports des
sections sont quasi-compacts, il existe un $n'$ tel que
$\Ker(\alpha_n)$ s'envoie sur zéro dans $\Ext^{-i}(M/I^{n'}M, \omega_B^\bullet)$.
On voit ainsi que $\Ext^{-i}(M/I^n M, \omega_B^\bullet)$
est essentiellement constant, de valeur $\Ext^{-i}(M, \omega_B^\bullet)$.
Ceci achève la démonstration.
\end{proof}

\noindent
Voici une version plus générale du lemme \ref{lemma-algebraization-principal}.

\begin{lemma}
\label{lemma-algebraization-principal-bis}
Dans la situation \ref{situation-algebraize}, soit $(\mathcal{F}_n)$ un objet
de $\textit{Coh}(U, I\mathcal{O}_U)$. Supposons que
\begin{enumerate}
\item $A$ soit local et $\mathfrak a = \mathfrak m$ soit l'idéal maximal,
\item $A$ possède un complexe dualisant,
\item $I = (f)$ soit un idéal principal,
\item $(\mathcal{F}_n)$ vérifie les inégalités $(2, 3)$.
\end{enumerate}
Alors $(\mathcal{F}_n)$ s'étend à $X$. En particulier, si $A$ est
complet pour la topologie $I$-adique, alors $(\mathcal{F}_n)$ est le complété
d'un $\mathcal{O}_U$-module cohérent. 
\end{lemma}

\begin{proof}
Rappelons que $\textit{Coh}(U, I\mathcal{O}_U)$ est une catégorie abélienne, voir
Cohomologie des schémas, lemme \ref{coherent-lemma-inverse-systems-abelian}.
Sur les ouverts affines de $U$, l'objet $(\mathcal{F}_n)$
correspond à un module de type fini sur un anneau noethérien
(Cohomologie des schémas, lemme \ref{coherent-lemma-inverse-systems-affine}).
Ainsi, les noyaux des applications $f^N : (\mathcal{F}_n) \to (\mathcal{F}_n)$
se stabilisent pour $N$ assez grand. D'après les
lemmes \ref{lemma-map-kernel-cokernel-on-closed} et
\ref{lemma-essential-image-completion},
pour démontrer le lemme,
nous pouvons remplacer $(\mathcal{F}_n)$ par l'image d'une telle application.
Nous pouvons donc supposer que $f$ est injectif sur $(\mathcal{F}_n)$.
Après ce remplacement, les conditions équivalentes du
lemme \ref{lemma-equivalent-f-good} sont vérifiées pour le système projectif
$(\mathcal{F}_n)$ sur $U$. Nous l'utiliserons sans autre mention
dans la suite de la démonstration.

\medskip\noindent
Nous allons vérifier les hypothèses (a), (b) et (c) du
lemme \ref{lemma-when-done}. L'hypothèse (b) résulte de
Cohomologie, lemme \ref{cohomology-lemma-topology-I-adic-f}.

\medskip\noindent
Choisissons un système projectif de modules $\{M_n\}$ comme dans le
lemme \ref{lemma-system-of-modules}.
Nous pouvons supposer que $H^0_\mathfrak m(M_n) = 0$ en remplaçant, si nécessaire, $M_n$ par
$M_n/H^0_\mathfrak m(M_n)$. Nous obtenons alors des suites
exactes courtes
$$
0 \to M_n \to H^0(U, \mathcal{F}_n) \to H^1_\mathfrak m(M_n) \to 0
$$
pour tout $n$. Soit $E$ une enveloppe injective du corps résiduel de $A$.
D'après le lemme \ref{lemma-helper-algebraize} et notre hypothèse (4),
nous pouvons choisir un entier $m \geq 0$, des $A$-modules de type fini
$N_1$ et $N_2$ annulés par $f^c$ pour un certain $c \geq 0$, ainsi que des
systèmes compatibles d'applications
$$
H^i_\mathfrak m(M_n) \to \Hom_A(N_i, E), \quad i = 1, 2
$$
pour $n \geq m$,
ayant les propriétés énoncées dans le lemme.

\medskip\noindent
Nous savons que $M = \lim H^0(U, \mathcal{F}_n)$ est un $A$-module dont la
topologie de limite est la topologie $f$-adique. Ainsi, étant donné $n$, le module
$M/f^nM$ est un sous-quotient de $H^0(U, \mathcal{F}_N)$ pour un certain $N \gg n$.
Les informations obtenues ci-dessus montrent que
$f^cM/f^nM$ est un $A$-module de type fini. Comme $f$ est un élément non diviseur de zéro
dans $M$, nous en déduisons que $M/f^{n - c}M$ est un $A$-module de type fini.
On voit ainsi que l'hypothèse (c) du lemme \ref{lemma-when-done} est vérifiée.

\medskip\noindent
Étudions ensuite le module
$$
Ob = \lim H^1(U, \mathcal{F}_n) = \lim H^2_\mathfrak m(M_n)
$$
Pour $n \geq m$, soit $K_n$ le noyau de l'application
$H^2_\mathfrak m(M_n) \to \Hom_A(N_2, E)$.
Posons $K = \lim K_n$. Nous obtenons une suite exacte
$$
0 \to K \to Ob \to \Hom_A(N_2, E)
$$
D'après ce qui précède, la topologie de limite sur $Ob = \lim H^2_\mathfrak m(M_n)$
est la topologie $f$-adique. Comme $N_2$ est annulé par $f^c$,
nous en déduisons qu'il en va de même pour la topologie de limite sur $K = \lim K_n$.
Ainsi, $K/fK$ est un sous-quotient de $K_n$ pour $n \gg 1$.
Cependant, puisque $\{K_n\}$ est pro-isomorphe à un système projectif de
$A$-modules de longueur finie (d'après la conclusion du
lemme \ref{lemma-helper-algebraize}),
nous en déduisons que $K/fK$ est un sous-quotient d'un
$A$-module de longueur finie. Il s'ensuit que $K$ est un $A$-module de type fini, voir
Algèbre, lemme \ref{algebra-lemma-finite-over-complete-ring}.
(En fait, nous voyons même que $\dim(\text{Supp}(K)) = 1$, mais
nous n'en aurons pas besoin.)

\medskip\noindent
Étant donné $n \geq 1$, considérons l'application de bord
$$
\delta_n :
H^0(U, \mathcal{F}_n)
\longrightarrow
\lim_N H^1(U, f^n\mathcal{F}_N) \xrightarrow{f^{-n}} Ob
$$
(la seconde application est un isomorphisme)
provenant des suites exactes courtes
$$
0 \to f^n\mathcal{F}_N \to \mathcal{F}_N \to \mathcal{F}_n \to 0
$$
Pour tout $n$, posons
$$
P_n = \Im(H^0(U, \mathcal{F}_{n + m}) \to H^0(U, \mathcal{F}_n))
$$
où $m$ est comme ci-dessus. Observons que $\{P_n\}$ est un système
projectif et que l'application $f : \mathcal{F}_n \to \mathcal{F}_{n + 1}$
sur les sections globales envoie $P_n$ dans $P_{n + 1}$.
Si $p \in P_n$, alors $\delta_n(p) \in K \subset Ob$,
car $\delta_n(p)$ s'envoie sur zéro dans
$H^1(U, f^n\mathcal{F}_{n + m}) = H^2_\mathfrak m(M_m)$
et la composée de $\delta_n$ et de $Ob \to \Hom_A(N_2, E)$
se factorise par $H^2_\mathfrak m(M_m)$ en vertu de notre choix de $m$.
Par conséquent,
$$
\bigoplus\nolimits_{n \geq 0} \Im(P_n \to Ob)
$$
est un $A[T]$-module gradué de type fini, où $T$ agit par multiplication par $f$.
En effet, c'est un sous-module gradué de $K[T]$ et $K$ est de type fini sur $A$.
En raisonnant comme dans la démonstration de
Cohomologie, lemme \ref{cohomology-lemma-ML-general}\footnote{Choisir
des générateurs homogènes de la forme $\delta_{n_j}(p_j)$ du module
affiché. Si $k = \max(n_j)$, on voit alors que pour $n \geq k$
et tout $p \in P_n$, il existe des $a_j \in A$ tels que
$p - \sum a_j f^{n - n_j} p_j$ appartienne au noyau de $\delta_n$
et donc à l'image de $P_{n'}$ pour tout $n' \geq n$.
Ainsi, $\Im(P_n \to P_{n - k}) = \Im(P_{n'} \to P_{n - k})$
pour tout $n' \geq n$.},
nous voyons que le système projectif $\{P_n\}$ satisfait à la condition ML.
Comme $\{P_n\}$ est pro-isomorphe à $\{H^0(U, \mathcal{F}_n)\}$,
nous en déduisons que $\{H^0(U, \mathcal{F}_n)\}$ satisfait à la condition ML.
Ainsi, l'hypothèse (a) du lemme \ref{lemma-when-done}
est vérifiée et la démonstration est achevée.
\end{proof}

\noindent
Nous pouvons expliciter la condition du
lemme \ref{lemma-algebraization-principal-bis} comme suit.

\begin{lemma}
\label{lemma-unwinding-conditions}
Dans la situation \ref{situation-algebraize}, soit $(\mathcal{F}_n)$ un objet
de $\textit{Coh}(U, I\mathcal{O}_U)$. Supposons que
\begin{enumerate}
\item $A$ soit local d'idéal maximal $\mathfrak a = \mathfrak m$,
\item $\text{cd}(A, I) = 1$.
\end{enumerate}
Alors $(\mathcal{F}_n)$ vérifie les inégalités $(2, 3)$ si et seulement
si, pour tout $y \in U \cap Y$ tel que $\dim(\{y\}) = 1$ et tout idéal premier
$\mathfrak p \subset \mathcal{O}_{X, y}^\wedge$,
$\mathfrak p \not \in V(I\mathcal{O}_{X, y}^\wedge)$, on a
$$
\text{depth}((\mathcal{F}_y^\wedge)_\mathfrak p) +
\dim(\mathcal{O}_{X, y}^\wedge/\mathfrak p) > 2
$$
\end{lemma}

\begin{proof}
Nous utiliserons le lemme \ref{lemma-explain-2-3-cd-1}
sans autre mention. En particulier, nous voyons que la condition est nécessaire.
Réciproquement, supposons que la condition soit vraie.
Notons que $\delta^Y_Z(y) = \dim(\overline{\{y\}})$ d'après le
lemme \ref{lemma-discussion}. Notons $\delta$ cette fonction.
Soit $y \in U \cap Y$. Si $\delta(y) > 2$, alors l'inégalité
du lemme \ref{lemma-explain-2-3-cd-1} est vérifiée.
Supposons enfin que $\delta(y) = 2$. Il faut montrer que
$$
\text{depth}((\mathcal{F}_y^\wedge)_\mathfrak p) +
\dim(\mathcal{O}_{X, y}^\wedge/\mathfrak p) > 1
$$
Choisissons une spécialisation $y \leadsto y'$ telle que $\delta(y') = 1$. Il existe alors
un homomorphisme d'anneaux $\mathcal{O}_{X, y'}^\wedge \to \mathcal{O}_{X, y}^\wedge$
qui identifie le but au complété de la localisation
de $\mathcal{O}_{X, y'}^\wedge$ en un idéal premier $\mathfrak q$
tel que $\dim(\mathcal{O}_{X, y'}^\wedge/\mathfrak q) = 1$.
De plus, nous obtenons alors
$$
\mathcal{F}_y^\wedge =
\mathcal{F}_{y'}^\wedge
\otimes_{\mathcal{O}_{X, y'}^\wedge}
\mathcal{O}_{X, y}^\wedge
$$
Soit $\mathfrak p' \subset \mathcal{O}_{X, y'}^\wedge$ l'image
de $\mathfrak p$.
D'après Cohomologie locale, lemme \ref{local-cohomology-lemma-change-completion},
nous avons
\begin{align*}
\text{depth}((\mathcal{F}_y^\wedge)_\mathfrak p) +
\dim(\mathcal{O}_{X, y}^\wedge/\mathfrak p)
& =
\text{depth}((\mathcal{F}_{y'}^\wedge)_{\mathfrak p'}) +
\dim((\mathcal{O}_{X, y}^\wedge/\mathfrak p)_{\mathfrak p'}) \\
& =
\text{depth}((\mathcal{F}_{y'}^\wedge)_{\mathfrak p'}) +
\dim(\mathcal{O}_{X, y}^\wedge/\mathfrak p') - 1
\end{align*}
la dernière égalité résultant du fait que la spécialisation est immédiate.
Le lemme résulte donc de l'inégalité supposée pour $y', \mathfrak p'$.
\end{proof}

\begin{lemma}
\label{lemma-unwinding-conditions-bis}
Dans la situation \ref{situation-algebraize}, soit $(\mathcal{F}_n)$ un objet
de $\textit{Coh}(U, I\mathcal{O}_U)$. Supposons que
\begin{enumerate}
\item $A$ soit local d'idéal maximal $\mathfrak a = \mathfrak m$,
\item $A$ possède un complexe dualisant,
\item $\text{cd}(A, I) = 1$,
\item pour $y \in U \cap Y$, le module $\mathcal{F}_y^\wedge$
soit localement libre de type fini hors de $V(I\mathcal{O}_{X, y}^\wedge)$,
par exemple si $\mathcal{F}_n$ est un
$\mathcal{O}_U/I^n\mathcal{O}_U$-module localement libre de type fini, et
\item l'une des conditions suivantes soit vérifiée :
\begin{enumerate}
\item $A_f$ soit $(S_2)$ et toute composante irréductible de $X$
non contenue dans $Y$ soit de dimension $\geq 4$, ou
\item si $\mathfrak p \not \in V(f)$ et
$V(\mathfrak p) \cap V(f) \not = \{\mathfrak m\}$, alors
$\text{depth}(A_\mathfrak p) + \dim(A/\mathfrak p) > 3$.
\end{enumerate}
\end{enumerate}
Alors $(\mathcal{F}_n)$ vérifie les inégalités $(2, 3)$.
\end{lemma}

\begin{proof}
Nous allons utiliser le critère du lemme \ref{lemma-unwinding-conditions}.
Soit $y \in U \cap Y$ tel que $\dim(\overline{\{y\}} = 1$ et soit
$\mathfrak p$ un idéal premier $\mathfrak p \subset \mathcal{O}_{X, y}^\wedge$ tel que
$\mathfrak p \not \in V(I\mathcal{O}_{X, y}^\wedge)$.
La condition (4) montre que
$\text{depth}((\mathcal{F}_y^\wedge)_\mathfrak p) =
\text{depth}((\mathcal{O}_{X, y}^\wedge)_\mathfrak p)$.
Il suffit donc de démontrer que
$$
\text{depth}((\mathcal{O}_{X, y}^\wedge)_\mathfrak p) +
\dim(\mathcal{O}_{X, y}^\wedge/\mathfrak p) > 2
$$
Soit $\mathfrak p_0 \subset A$ l'image de $\mathfrak p$.
Soit $\mathfrak q \subset A$ l'idéal premier correspondant à $y$.
D'après Cohomologie locale, lemme
\ref{local-cohomology-lemma-change-completion},
nous avons
\begin{align*}
\text{depth}((\mathcal{O}_{X, y}^\wedge)_\mathfrak p) +
\dim(\mathcal{O}_{X, y}^\wedge/\mathfrak p)
& =
\text{depth}(A_{\mathfrak p_0}) + \dim((A/\mathfrak p_0)_\mathfrak q) \\
& =
\text{depth}(A_{\mathfrak p_0}) + \dim(A/\mathfrak p_0) - 1
\end{align*}
Si (5)(a) est vérifiée, cette quantité est
$$
\geq \min(2, \dim(A_{\mathfrak p_0})) + \dim(A/\mathfrak p_0) - 1
$$
Notons que, dans tous les cas, $\dim(A/\mathfrak p_0) \geq 2$. Par conséquent, si
le minimum vaut $2$, la conclusion est acquise. Sinon, nous obtenons
$$
\dim(A_{\mathfrak p_0}) + \dim(A/\mathfrak p_0) - 1 \geq 4 - 1
$$
car toute composante de $\Spec(A)$ passant par $\mathfrak p_0$
est de dimension $\geq 4$. Si (5)(b) est vérifiée, la conclusion est immédiate.
\end{proof}

\begin{remark}
\label{remark-interesting-case-bis}
Soit $(A, \mathfrak m)$ un anneau local noethérien qui possède un
complexe dualisant et qui est complet pour la topologie définie par $f \in \mathfrak m$.
Soit $(\mathcal{F}_n)$ un objet de $\textit{Coh}(U, f\mathcal{O}_U)$,
où $U$ est le spectre épointé de $A$.
Posons $Y = V(f) \subset X = \Spec(A)$.
Supposons que, pour $y \in U \cap V(f)$ fermé dans $U$, c'est-à-dire tel que
$\dim(\overline{\{y\}}) = 1$, le
$\mathcal{O}_{X, y}^\wedge$-module $\mathcal{F}_y^\wedge$
vérifie les deux conditions suivantes :
\begin{enumerate}
\item $\mathcal{F}_y^\wedge[1/f]$ soit $(S_2)$ comme
$\mathcal{O}_{X, y}^\wedge[1/f]$-module, et
\item pour $\mathfrak p \in \text{Ass}(\mathcal{F}_y^\wedge[1/f])$,
on ait $\dim(\mathcal{O}_{X, y}^\wedge/\mathfrak p) \geq 3$.
\end{enumerate}
Alors $(\mathcal{F}_n)$ est le complété d'un module cohérent sur $U$.
Ceci résulte des lemmes \ref{lemma-algebraization-principal-bis}
et \ref{lemma-unwinding-conditions}.
\end{remark}
































\section{Amélioration des modules formels cohérents}
\label{section-improving-formal-modules}

\noindent
Soit $X$ un schéma noethérien. Soit $Y \subset X$ un sous-schéma fermé,
dont le faisceau quasi-cohérent d'idéaux est $\mathcal{I} \subset \mathcal{O}_X$.
Soit $(\mathcal{F}_n)$ un objet de $\textit{Coh}(X, \mathcal{I})$.
Dans cette section, nous construisons des applications
$(\mathcal{F}_n) \to (\mathcal{F}'_n)$
analogues aux applications construites dans
Cohomologie locale, section \ref{local-cohomology-section-improve},
pour les modules cohérents. Pour un point $y \in Y$, posons
$$
\mathcal{O}_{X, y}^\wedge = \lim \mathcal{O}_{X, y}/\mathcal{I}^n_y, \quad
\mathcal{I}_y^\wedge = \lim \mathcal{I}_y/\mathcal{I}^n_y
\quad\text{et}\quad
\mathfrak m_y^\wedge = \lim \mathfrak m_y/\mathcal{I}_y^n
$$
Alors $\mathcal{O}_{X, y}^\wedge$ est un anneau local noethérien
d'idéal maximal $\mathfrak m_y^\wedge$, complet pour la topologie définie par
$\mathcal{I}_y^\wedge = \mathcal{I}_y\mathcal{O}_{X, y}^\wedge$.
Posons aussi
$$
\mathcal{F}_y^\wedge = \lim \mathcal{F}_{n, y}
$$
Alors $\mathcal{F}_y^\wedge$ est un module de type fini sur
$\mathcal{O}_{X, y}^\wedge$ tel que
$\mathcal{F}_y^\wedge/(\mathcal{I}_y^\wedge)^n\mathcal{F}_y^\wedge =
\mathcal{F}_{n, y}$ pour tout $n$, voir Algèbre, lemmes
\ref{algebra-lemma-limit-complete} et
\ref{algebra-lemma-finite-over-complete-ring}.

\begin{lemma}
\label{lemma-divide-torsion-formal-coherent-module}
Dans la situation ci-dessus, supposons que $X$ possède localement un complexe dualisant.
Soit $T \subset Y$ une partie stable par spécialisation.
Supposons que, pour $y \in T$ et pour un idéal premier non maximal
$\mathfrak p \subset \mathcal{O}_{X, y}^\wedge$ tel que
$V(\mathfrak p) \cap V(\mathcal{I}^\wedge_y) = \{\mathfrak m_y^\wedge\}$,
on ait
$$
\text{depth}_{(\mathcal{O}_{X, y})_\mathfrak p}
((\mathcal{F}^\wedge_y)_\mathfrak p) > 0
$$
Alors il existe une application canonique
$(\mathcal{F}_n) \to (\mathcal{F}_n')$
de systèmes projectifs de $\mathcal{O}_X$-modules cohérents
ayant les propriétés suivantes :
\begin{enumerate}
\item pour $y \in T$, on a $\text{depth}(\mathcal{F}'_{n, y}) \geq 1$,
\item $(\mathcal{F}'_n)$ est isomorphe, en tant que pro-système, à un objet
$(\mathcal{G}_n)$ de $\textit{Coh}(X, \mathcal{I})$,
\item le morphisme induit
$(\mathcal{F}_n) \to (\mathcal{G}_n)$ dans
$\textit{Coh}(X, \mathcal{I})$ est surjectif, de noyau
annulé par une puissance de $\mathcal{I}$.
\end{enumerate}
\end{lemma}

\begin{proof}
Pour tout $n$, soit $\mathcal{F}_n \to \mathcal{F}'_n$ la surjection
construite dans
Cohomologie locale, lemme \ref{local-cohomology-lemma-get-depth-1-along-Z}.
Comme il s'agit du quotient de $\mathcal{F}_n$ par le sous-faisceau
des sections à support dans $T$, nous voyons qu'il existe des applications canoniques
$\mathcal{F}'_{n + 1} \to \mathcal{F}'_n$ telles que l'on obtienne
une application $(\mathcal{F}_n) \to (\mathcal{F}_n')$
de systèmes projectifs de $\mathcal{O}_X$-modules cohérents.
La propriété (1) est vérifiée par construction.

\medskip\noindent
Pour démontrer les propriétés (2) et (3), nous pouvons supposer que $X = \Spec(A_0)$ est
affine et que $A_0$ possède un complexe dualisant. Soit $I_0 \subset A_0$ l'idéal
correspondant à $Y$. Soient $A, I$ les complétés $I_0$-adiques de
$A_0, I_0$. Pour un usage ultérieur, observons que $A$ possède un complexe dualisant
(Complexes dualisants, lemme \ref{dualizing-lemma-ubiquity-dualizing}).
Soit $M$ le $A$-module de type fini correspondant à $(\mathcal{F}_n)$, voir
Cohomologie des schémas, lemme \ref{coherent-lemma-inverse-systems-affine}.
Alors $\mathcal{F}_n$ correspond à $M_n = M/I^nM$. Rappelons que
$\mathcal{F}'_n$ correspond au quotient $M'_n = M_n / H^0_T(M_n)$,
voir Cohomologie locale, lemme \ref{local-cohomology-lemma-get-depth-1-along-Z}
et sa démonstration.

\medskip\noindent
Posons $s = 0$ et $d = \text{cd}(A, I)$.
Nous affirmons que $A, I, T, M, s, d$ vérifient les hypothèses (1), (3), (4), (6)
de la situation \ref{situation-bootstrap}.
En effet, (1) et (3) résultent immédiatement de ce qui précède, (4) est
une condition vide puisque $s = 0$, et (6) est l'hypothèse
faite dans l'énoncé du lemme.

\medskip\noindent
D'après le théorème \ref{theorem-final-bootstrap}, nous voyons que $\{H^0_T(M_n)\}$
est essentiellement constant, de valeur $H^0_T(M)$. Ainsi, la limite des
suites exactes courtes $0 \to H^0_T(M_n) \to M_n \to M'_n \to 0$
est la suite exacte courte
$$
0 \to H^0_T(M) \to M \to \lim M'_n \to 0
$$
Posons $M' = \lim M'_n = M/H^0_T(M)$. C'est un $A$-module de type fini.
Soit $(\mathcal{G}_n)$ l'objet de $\textit{Coh}(X, \mathcal{I})$
correspondant à $M'$. Pour achever la démonstration, il faut montrer que l'application
canonique $\{M'/I^nM'\} \to \{M'_n\}$ est un pro-isomorphisme.
Ceci équivaut à dire que
$\{H^0_T(M) + I^nM\} \to \{\Ker(M \to M'_n)\}$ est un
pro-isomorphisme. En quotientant par $I^nM$, il suffit de montrer que
$\{H^0_T(M)/H^0_T(M) \cap I^nM\} \to \{H^0_T(M_n)\}$
est un pro-isomorphisme. Comme $H^0_T(M)$ est annulé par une puissance
de $I$, le lemme d'Artin-Rees montre que $H^0_T(M) \cap I^nM = 0$
pour tout $n \gg 0$. Nous obtenons donc le pro-isomorphisme voulu,
puisque nous avons vu ci-dessus que $\{H^0_T(M_n)\}$
est essentiellement constant, de valeur $H^0_T(M)$.
\end{proof}

\begin{lemma}
\label{lemma-improvement-formal-coherent-module-better}
Dans la situation ci-dessus, supposons que $X$ possède localement un complexe dualisant.
Soient $T' \subset T \subset Y$ des parties stables par spécialisation.
Soit $d \geq 0$ un entier. Supposons que
\begin{enumerate}
\item[(a)] localement sur des ouverts affines, on ait $X = \Spec(A_0)$ et $Y = V(I_0)$
et $\text{cd}(A_0, I_0) \leq d$,
\item[(b)] pour $y \in T$ et un idéal premier non maximal
$\mathfrak p \subset \mathcal{O}_{X, y}^\wedge$ tel que
$V(\mathfrak p) \cap V(\mathcal{I}_y^\wedge) = \{\mathfrak m_y^\wedge\}$,
on ait
$$
\text{depth}_{(\mathcal{O}_{X, y})_\mathfrak p}
((\mathcal{F}^\wedge_y)_\mathfrak p) > 0
$$
\item[(c)] pour $y \in T'$ et pour un idéal premier
$\mathfrak p \subset \mathcal{O}_{X, y}^\wedge$ tel que
$\mathfrak p \not \in V(\mathcal{I}_y^\wedge)$
et $V(\mathfrak p) \cap V(\mathcal{I}_y^\wedge) \not =
\{\mathfrak m_y^\wedge\}$, on ait
$$
\text{depth}_{(\mathcal{O}_{X, y})_\mathfrak p}
((\mathcal{F}^\wedge_y)_\mathfrak p) \geq 1
\quad\text{ou}\quad
\text{depth}_{(\mathcal{O}_{X, y})_\mathfrak p}
((\mathcal{F}^\wedge_y)_\mathfrak p) +
\dim(\mathcal{O}_{X, y}^\wedge/\mathfrak p) > 1 + d
$$
\item[(d)] pour $y \in T'$ et un idéal premier non maximal
$\mathfrak p \subset \mathcal{O}_{X, y}^\wedge$ tel que
$V(\mathfrak p) \cap V(\mathcal{I}_y^\wedge) = \{\mathfrak m_y^\wedge\}$,
on ait
$$
\text{depth}_{(\mathcal{O}_{X, y})_\mathfrak p}
((\mathcal{F}^\wedge_y)_\mathfrak p) > 1
$$
\item[(e)] si $y \leadsto y'$ est une spécialisation immédiate et
$y' \in T'$, alors $y \in T$.
\end{enumerate}
Alors il existe une application canonique $(\mathcal{F}_n) \to (\mathcal{F}_n'')$
de systèmes projectifs de $\mathcal{O}_X$-modules cohérents
ayant les propriétés suivantes :
\begin{enumerate}
\item pour $y \in T$, on a $\text{depth}(\mathcal{F}''_{n, y}) \geq 1$,
\item pour $y' \in T'$, on a $\text{depth}(\mathcal{F}''_{n, y'}) \geq 2$,
\item $(\mathcal{F}''_n)$ est isomorphe, en tant que pro-système, à un objet
$(\mathcal{H}_n)$ de $\textit{Coh}(X, \mathcal{I})$,
\item le morphisme induit $(\mathcal{F}_n) \to (\mathcal{H}_n)$ dans
$\textit{Coh}(X, \mathcal{I})$ a un noyau et un conoyau
annulés par une puissance de $\mathcal{I}$.
\end{enumerate}
\end{lemma}

\begin{proof}
Comme dans le lemme \ref{lemma-divide-torsion-formal-coherent-module} et sa démonstration,
pour tout $n$, soit $\mathcal{F}_n \to \mathcal{F}'_n$ la surjection
construite dans
Cohomologie locale, lemme \ref{local-cohomology-lemma-get-depth-1-along-Z},
en utilisant $T$.
Ensuite, soit $\mathcal{F}'_n \to \mathcal{F}''_n$ l'injection
construite dans
Cohomologie locale, lemme \ref{local-cohomology-lemma-make-S2-along-T},
ainsi que dans sa démonstration. Les constructions montrent qu'il existe des applications canoniques
$\mathcal{F}''_{n + 1} \to \mathcal{F}''_n$ telles que l'on obtienne des
applications
$$
(\mathcal{F}_n) \longrightarrow (\mathcal{F}_n') \longrightarrow
(\mathcal{F}''_n)
$$
de systèmes projectifs de $\mathcal{O}_X$-modules cohérents.
Les propriétés (1) et (2) sont vérifiées par construction.

\medskip\noindent
Pour démontrer les propriétés (3) et (4), nous pouvons supposer que $X = \Spec(A_0)$ est
affine et que $A_0$ possède un complexe dualisant. Soit $I_0 \subset A_0$ l'idéal
correspondant à $Y$. Soient $A, I$ les complétés $I_0$-adiques de
$A_0, I_0$. Pour un usage ultérieur, observons que $A$ possède un complexe dualisant
(Complexes dualisants, lemme \ref{dualizing-lemma-ubiquity-dualizing}).
Soit $M$ le $A$-module de type fini correspondant à $(\mathcal{F}_n)$, voir
Cohomologie des schémas, lemme \ref{coherent-lemma-inverse-systems-affine}.
Alors $\mathcal{F}_n$ correspond à $M_n = M/I^nM$. Rappelons que
$\mathcal{F}'_n$ correspond au quotient $M'_n = M_n / H^0_T(M_n)$.
Rappelons aussi que $M' = \lim M'_n$ est le quotient de $M$ par
$H^0_T(M)$, que $H^0_T(M)$ est annulé par une puissance de $I$ et que
$\{M'_n\}$ et $\{M'/I^nM'\}$ sont isomorphes en tant que pro-systèmes.
Enfin, nous voyons que $\mathcal{F}''_n$ correspond à une extension
$$
0 \to M'_n \to M''_n \to H^1_{T'}(M'_n) \to 0
$$
voir la démonstration de
Cohomologie locale, lemme \ref{local-cohomology-lemma-make-S2-along-T}.

\medskip\noindent
Posons $s = 1$. Nous affirmons que $A, I, T', M', s, d$ vérifient les hypothèses
(1), (3), (4), (6) de la situation \ref{situation-bootstrap}. En effet, (1) et (3)
sont immédiates, (4) est impliquée par (c) et (6) résulte de (d).
Nous omettons les démonstrations de (c) $\Rightarrow$ (4) et (d) $\Rightarrow$ (6).

\medskip\noindent
D'après le théorème \ref{theorem-final-bootstrap}, nous voyons que $\{H^1_{T'}(M'/I^nM')\}$
est essentiellement constant, de valeur $H^1_{T'}(M')$. Nous en déduisons que
$\{H^1_{T'}(M'_n)\}$ est essentiellement constant, de valeur $H^1_{T'}(M')$.
Ainsi, la limite des suites exactes courtes affichées ci-dessus
est la suite exacte courte
$$
0 \to M' \to \lim M''_n \to H^1_{T'}(M') \to 0
$$
Posons $M'' = \lim M''_n$. Il résulte de
Cohomologie locale, proposition \ref{local-cohomology-proposition-finiteness},
que $H^1_{T'}(M')$, et donc $M''$, sont des $A$-modules de type fini (pour vérifier les
conditions, utiliser que $M'$ est de profondeur au moins $1$ aux idéaux premiers appartenant à $T - T'$).
Soit $(\mathcal{H}_n)$ l'objet de $\textit{Coh}(X, \mathcal{I})$
correspondant au $A$-module de type fini $M''$. Pour achever la démonstration, il faut
montrer que l'application canonique $\{M''/I^nM''\} \to \{M''_n\}$ est un
pro-isomorphisme. Puisque nous savons déjà que $\{M'/I^nM'\}$ est pro-isomorphe
à $\{M'_n\}$, le lecteur vérifiera (détails omis) que ceci équivaut à demander que
$\{H^1_{T'}(M')/I^nH^1_{T'}(M')\} \to \{H^1_{T'}(M'_n)\}$
soit un pro-isomorphisme. Comme $H^1_{T'}(M')/I^nH^1_{T'}(M') = H^1_{T'}(M')$
pour $n$ assez grand, cela résulte du fait que $\{H^1_{T'}(M'_n)\}$
est essentiellement constant, de valeur $H^1_{T'}(M')$, comme vu ci-dessus.
\end{proof}

\begin{lemma}
\label{lemma-improvement-application}
Dans la situation \ref{situation-algebraize}, supposons que $A$ possède
un complexe dualisant. Soit $d \geq \text{cd}(A, I)$. Soit $(\mathcal{F}_n)$
un objet de $\textit{Coh}(U, I\mathcal{O}_U)$. Supposons que
$(\mathcal{F}_n)$ vérifie les inégalités $(2, 2 + d)$, voir la
définition \ref{definition-s-d-inequalities}.
Alors il existe une application canonique $(\mathcal{F}_n) \to (\mathcal{F}_n'')$
de systèmes projectifs de $\mathcal{O}_U$-modules cohérents
ayant les propriétés suivantes :
\begin{enumerate}
\item $\text{depth}(\mathcal{F}''_{n, y}) + \delta^Y_Z(y) \geq 3$
pour tout $y \in U \cap Y$,
\item $(\mathcal{F}''_n)$ est isomorphe, en tant que pro-système, à un objet
$(\mathcal{H}_n)$ de $\textit{Coh}(U, I\mathcal{O}_U)$,
\item le morphisme induit $(\mathcal{F}_n) \to (\mathcal{H}_n)$ dans
$\textit{Coh}(U, I\mathcal{O}_U)$ a un noyau et un conoyau
annulés par une puissance de $I$,
\item les modules $H^0(U, \mathcal{F}''_n)$ et $H^1(U, \mathcal{F}''_n)$
sont des $A$-modules de type fini pour tout $n$.
\end{enumerate}
\end{lemma}

\begin{proof}
L'existence et les propriétés (1), (2), (3) résultent du
lemme \ref{lemma-improvement-formal-coherent-module-better} appliqué
à $U$, $U \cap Y$, $T = \{y \in U \cap Y : \delta^Y_Z(y) \leq 2\}$,
$T' = \{y \in U \cap Y : \delta^Y_Z(y) \leq 1\}$ et $(\mathcal{F}_n)$.
Le fait que les modules $H^0(U, \mathcal{F}''_n)$ et
$H^1(U, \mathcal{F}''_n)$ soient de type fini résulte de
Cohomologie locale, lemme \ref{local-cohomology-lemma-finiteness-Rjstar},
et des propriétés élémentaires de la fonction $\delta^Y_Z(-)$
démontrées dans le lemme \ref{lemma-discussion}.
\end{proof}











\section{Algébrisation des modules formels cohérents, V}
\label{section-algebraization-modules-conclusion}

\noindent
Dans cette section, nous démontrons nos résultats les plus généraux sur l'algébrisation
des modules formels cohérents. Nous les démontrons d'abord dans le cas où
l'idéal est de dimension cohomologique $1$. Nous les appliquons ensuite
à un éclatement pour démontrer un résultat plus général.

\begin{lemma}
\label{lemma-cd-1-canonical}
Dans la situation \ref{situation-algebraize}, soit $(\mathcal{F}_n)$
un objet de $\textit{Coh}(U, I\mathcal{O}_U)$. Supposons que
\begin{enumerate}
\item $A$ possède un complexe dualisant et $\text{cd}(A, I) = 1$,
\item $(\mathcal{F}_n)$ soit pro-isomorphe à un système projectif
$(\mathcal{F}_n'')$ de $\mathcal{O}_U$-modules cohérents tel que
$\text{depth}(\mathcal{F}''_{n, y}) + \delta^Y_Z(y) \geq 3$
pour tout $y \in U \cap Y$.
\end{enumerate}
Alors $(\mathcal{F}_n)$ s'étend canoniquement à $X$, voir la
définition \ref{definition-canonically-algebraizable}.
\end{lemma}

\begin{proof}
Nous allons vérifier les hypothèses (a), (b) et (c) du lemme \ref{lemma-when-done}.
Avant de commencer, signalons que les modules
$H^0(U, \mathcal{F}''_n)$ et $H^1(U, \mathcal{F}''_n)$
sont des $A$-modules de type fini pour tout $n$ d'après
Cohomologie locale, lemme \ref{local-cohomology-lemma-finiteness-Rjstar}.

\medskip\noindent
Observons que, pour tout $p \geq 0$,
la topologie de limite sur $\lim H^p(U, \mathcal{F}_n)$
est la topologie $I$-adique d'après le lemme \ref{lemma-topology-I-adic}.
En particulier, l'hypothèse (b) est vérifiée.

\medskip\noindent
Nous savons que $M = \lim H^0(U, \mathcal{F}_n)$ est un $A$-module dont la
topologie de limite est la topologie $I$-adique. Ainsi, étant donné $n$, le module
$M/I^nM$ est un sous-quotient de $H^0(U, \mathcal{F}_N)$ pour un certain $N \gg n$.
Comme le système projectif $\{H^0(U, \mathcal{F}_N)\}$ est pro-isomorphe à un
système projectif de $A$-modules de type fini, à savoir $\{H^0(U, \mathcal{F}''_N)\}$,
nous en déduisons que $M/I^nM$ est de type fini. Il s'ensuit que $M$ est de type fini, voir
Algèbre, lemme \ref{algebra-lemma-finite-over-complete-ring}.
En particulier, l'hypothèse (c) est vérifiée.

\medskip\noindent
Pour tout $n \geq 0$, notons $Ob_n = \lim_N H^1(U, I^n\mathcal{F}_N)$.
Un cas particulier est $Ob = Ob_0 = \lim_N H^1(U, \mathcal{F}_N)$.
En raisonnant exactement comme au paragraphe précédent, nous trouvons que $Ob$
est un $A$-module de type fini. (En fait, nous savons aussi que $Ob/I Ob$ est annulé
par une puissance de $\mathfrak a$, mais cela semble assez difficile à utiliser.)

\medskip\noindent
Posons $\mathcal{F} = \lim \mathcal{F}_n$, choisissons des générateurs
$f_1, \ldots, f_r$ de $I$, choisissons $c \geq 1$ et choisissons
$\Phi_\mathcal{F}$ comme dans le lemme \ref{lemma-cd-is-one-for-system}.
Nous utiliserons les résultats du lemme \ref{lemma-properties-system}
sans autre mention. En particulier, pour tout $n \geq 1$, il existe des applications
$$
\delta_n :
H^0(U, \mathcal{F}_n)
\longrightarrow
H^1(U, I^n\mathcal{F})
\longrightarrow
Ob_n
$$
La première provient de la suite exacte courte
$0 \to I^n\mathcal{F} \to \mathcal{F} \to \mathcal{F}_n \to 0$
et la seconde de $I^n\mathcal{F} = \lim I^n\mathcal{F}_N$.
Nous utiliserons plus loin que, si $\delta_n(s) = 0$ pour $s \in H^0(U, \mathcal{F}_n)$,
alors, pour tout $n' \geq n$, il existe $s' \in H^0(U, \mathcal{F}_{n'})$
qui s'envoie sur $s$.
Observons qu'il existe des diagrammes commutatifs
$$
\xymatrix{
H^0(U, \mathcal{F}_{nc}) \ar[r] \ar[dd] &
H^1(U, I^{nc}\mathcal{F}) \ar[dd] \ar[rd]^{\Phi_\mathcal{F}} \\
& &
\bigoplus_{e_1 + \ldots + e_r = n}
H^1(U, \mathcal{F}) \cdot T_1^{e_1} \ldots T_r^{e_r} \ar[ld] \\
H^0(U, \mathcal{F}_n) \ar[r] &
H^1(U, I^n\mathcal{F})
}
$$
Nous en déduisons que l'application d'obstruction
$H^0(U, \mathcal{F}_n) \to Ob_n$
envoie l'image de
$H^0(U, \mathcal{F}_{nc}) \to H^0(U, \mathcal{F}_n)$
dans le sous-module
$$
Ob'_n =
\Im\left(
\bigoplus\nolimits_{e_1 + \ldots + e_r = n}
Ob \cdot T_1^{e_1} \ldots T_r^{e_r} \to Ob_n
\right)
$$
où, sur le facteur $Ob \cdot T_1^{e_1} \ldots T_r^{e_r}$,
nous utilisons l'application en cohomologie provenant des réductions modulo des
puissances de $I$ de l'application de multiplication
$f_1^{e_1} \ldots f_r^{e_r} : \mathcal{F} \to I^n\mathcal{F}$.
Par construction,
$$
\bigoplus\nolimits_{n \geq 0} Ob'_n
$$
est un module gradué de type fini sur l'algèbre de Rees $\bigoplus_{n \geq 0} I^n$.
Pour tout $n$, posons
$$
M_n = \{s \in H^0(U, \mathcal{F}_n) \mid \delta_n(s) \in Ob'_n\}
$$
Observons que $\{M_n\}$ est un système projectif et que l'application
$f_j : \mathcal{F}_n \to \mathcal{F}_{n + 1}$ sur les
sections globales envoie $M_n$ dans $M_{n + 1}$.
Par exactement le même argument que dans la démonstration de
Cohomologie, lemme \ref{cohomology-lemma-ML-general},
nous trouvons que $\{M_n\}$ satisfait à la condition ML. En effet, puisque l'algèbre de Rees
est noethérienne, nous pouvons choisir un nombre fini de générateurs homogènes
de la forme $\delta_{n_j}(z_j)$ avec $z_j \in M_{n_j}$ du sous-module gradué
$\bigoplus_{n \geq 0} \Im(M_n \to Ob'_n)$.
Alors, si $k = \max(n_j)$, nous voyons que pour $n \geq k$
et tout $z \in M_n$, il existe des $a_j \in I^{n - n_j}$ tels que
$z - \sum a_j z_j$ appartienne au noyau de $\delta_n$
et donc à l'image de $M_{n'}$ pour tout $n' \geq n$
(car l'annulation de $\delta_n$ signifie que nous pouvons
relever $z - \sum a_j z_j$ en un élément $z' \in H^0(U, \mathcal{F}_{n'c})$
pour tout $n' \ge n$, puis que l'image de $z'$ dans $H^0(U, \mathcal{F}_{n'})$
appartient à $M_{n'}$ d'après ce que nous avons démontré ci-dessus).
Ainsi, $\Im(M_n \to M_{n - k}) = \Im(M_{n'} \to M_{n - k})$
pour tout $n' \geq n$.

\medskip\noindent
Choisissons $n$. Par la propriété de Mittag-Leffler de $\{M_n\}$ établie ci-dessus,
nous pouvons trouver $n' \geq n$ tel que l'image de $M_{n'} \to M_n$
coïncide avec l'image de $M' \to M_n$. D'après ce qui précède, nous voyons que
l'image de $M' \to M_n$ contient l'image de
$H^0(U, \mathcal{F}_{n'c}) \to H^0(U, \mathcal{F}_n)$.
Ainsi, nous voyons que $\{M_n\}$ et $\{H^0(U, \mathcal{F}_n)\}$
sont pro-isomorphes. Par conséquent, $\{H^0(U, \mathcal{F}_n)\}$
satisfait à la condition ML, et nous concluons enfin que l'hypothèse (a) est vérifiée.
Ceci achève la démonstration.
\end{proof}

\begin{proposition}[Algébrisation en dimension cohomologique 1]
\label{proposition-cd-1}
\begin{reference}
Le cas local de ce résultat est \cite[IV Corollaire 2.9]{MRaynaud-book}.
\end{reference}
Dans la situation \ref{situation-algebraize}, soit $(\mathcal{F}_n)$
un objet de $\textit{Coh}(U, I\mathcal{O}_U)$. Supposons que
\begin{enumerate}
\item $A$ possède un complexe dualisant et $\text{cd}(A, I) = 1$,
\item $(\mathcal{F}_n)$ satisfasse aux inégalités $(2, 3)$, voir la
définition \ref{definition-s-d-inequalities}.
\end{enumerate}
Alors $(\mathcal{F}_n)$ s'étend à $X$. En particulier, si $A$ est
complet pour la topologie $I$-adique, alors $(\mathcal{F}_n)$ est le complété
d'un $\mathcal{O}_U$-module cohérent.
\end{proposition}

\begin{proof}
D'après le lemme \ref{lemma-map-kernel-cokernel-on-closed},
nous pouvons remplacer $(\mathcal{F}_n)$ par l'objet $(\mathcal{H}_n)$
de $\textit{Coh}(U, I\mathcal{O}_U)$ obtenu dans le
lemme \ref{lemma-improvement-application}.
Nous pouvons donc supposer que $(\mathcal{F}_n)$ est pro-isomorphe
à un système projectif $(\mathcal{F}_n'')$ possédant les propriétés
mentionnées dans le lemme \ref{lemma-improvement-application}.
Dans le lemme \ref{lemma-cd-1-canonical}, nous avons démontré que
$(\mathcal{F}_n)$ s'étend canoniquement à $X$.
La dernière assertion résulte du lemme \ref{lemma-canonically-algebraizable}.
\end{proof}

\begin{lemma}
\label{lemma-blowup}
Dans la situation \ref{situation-algebraize}, soit $(\mathcal{F}_n)$
un objet de $\textit{Coh}(U, I\mathcal{O}_U)$. Supposons que
\begin{enumerate}
\item $A$ possède un complexe dualisant,
\item toutes les fibres de l'éclatement $b : X' \to X$ de $I$
soient de dimension $\leq d - 1$,
\item l'une des conditions suivantes soit vérifiée :
\begin{enumerate}
\item $(\mathcal{F}_n)$ satisfait aux inégalités $(d + 1, d + 2)$
(définition \ref{definition-s-d-inequalities}), ou
\item pour $y \in U \cap Y$ et un idéal premier
$\mathfrak p \subset \mathcal{O}_{X, y}^\wedge$ tel que
$\mathfrak p \not \in V(I\mathcal{O}_{X, y}^\wedge)$,
nous ayons
$$
\text{depth}((\mathcal{F}^\wedge_y)_\mathfrak p) +
\dim(\mathcal{O}_{X, y}^\wedge/\mathfrak p) + \delta^Y_Z(y) > d + 2
$$
\end{enumerate}
\end{enumerate}
Alors $(\mathcal{F}_n)$ s'étend à $X$.
\end{lemma}

\begin{proof}
Soit $Y' \subset X'$ le diviseur exceptionnel.
Soit $Z' \subset Y'$ l'image réciproque de $Z \subset Y$.
Alors $U' = X' \setminus Z'$ est l'image réciproque de $U$.
Avec $\delta^{Y'}_{Z'}$ comme dans (\ref{equation-delta-Z}), posons
$$
T' = \{y' \in Y' \mid \delta^{Y'}_{Z'}(y') = 1\}
\subset
T = \{y' \in Y' \mid \delta^{Y'}_{Z'}(y') = 1\text{ ou }2\}
$$
D'après les parties (1) et (2) du lemme \ref{lemma-discussion}, ces sous-ensembles sont
stables par spécialisation dans $U' \cap Y' = Y' \setminus Z'$. Considérons
l'objet $(b|_{U'}^*\mathcal{F}_n)$ de $\textit{Coh}(U', I\mathcal{O}_{U'})$,
voir Cohomologie des schémas, lemme \ref{coherent-lemma-inverse-systems-pullback}.
Pour $y' \in U' \cap Y'$, notons
$$
\mathcal{F}_{y'}^\wedge = \lim (b|_{U'}^*\mathcal{F}_n)_{y'}
$$
le « germe » de cette image réciproque en $y'$. Nous affirmons que les conditions
(a), (b), (c), (d) et (e) du
lemme \ref{lemma-improvement-formal-coherent-module-better}
sont vérifiées pour l'objet $(b|_{U'}^*\mathcal{F}_n)$ sur $U'$, avec $d$
remplacé par $1$, et les sous-ensembles $T' \subset T \subset U' \cap Y'$.
La condition (a) est vérifiée parce que $Y'$ est un diviseur de Cartier effectif
et est donc localement défini par $1$ équation. La condition (e) est vérifiée
d'après la partie (7) du lemme \ref{lemma-discussion}.
Pour démontrer (b), (c) et (d), quelques préparatifs sont nécessaires.

\medskip\noindent
Soit $y' \in U' \cap Y'$ et soit
$\mathfrak p' \subset \mathcal{O}_{X', y'}^\wedge$
un idéal premier n'appartenant pas à $V(I\mathcal{O}_{X', y'}^\wedge)$.
Notons $y = b(y') \in U \cap Y$. Choisissons $f \in I$ tel que
$y'$ appartienne au spectre de l'algèbre d'éclatement affine
$A[\frac{I}{f}]$, voir Diviseurs, lemme \ref{divisors-lemma-blowing-up-affine}.
Pour toute $A$-algèbre $B$, notons $B' = B[\frac{IB}{f}]$ l'algèbre d'éclatement
affine correspondante. Désignons la complétion $I$-adique par ${\ }^\wedge$.
Par notre choix de $f$, nous obtenons un homomorphisme d'anneaux
$(\mathcal{O}_{X, y}^\wedge)' \to \mathcal{O}_{X', y'}^\wedge$.
Si $\mathfrak q' \subset (\mathcal{O}_{X, y}^\wedge)'$
désigne l'image réciproque de $\mathfrak m_{y'}^\wedge$, alors
nous voyons que
$((\mathcal{O}_{X, y}^\wedge)'_{\mathfrak q'})^\wedge =
\mathcal{O}_{X', y'}^\wedge$.
Soit $\mathfrak p \subset \mathcal{O}_{X, y}^\wedge$ l'idéal premier
correspondant. Nous avons à ce stade un diagramme commutatif
$$
\xymatrix{
\mathcal{O}_{X, y}^\wedge \ar[d] \ar[r] &
(\mathcal{O}_{X, y}^\wedge)' \ar[d]_\alpha \ar[r] &
(\mathcal{O}_{X, y}^\wedge)'_{\mathfrak q'} \ar[d] \ar[r]_\beta &
\mathcal{O}_{X', y'}^\wedge \ar[d] \\
\mathcal{O}_{X, y}^\wedge/\mathfrak p \ar[r] &
(\mathcal{O}_{X, y}^\wedge/\mathfrak p)' \ar[r] &
(\mathcal{O}_{X, y}^\wedge/\mathfrak p)'_{\mathfrak q'} \ar[r]^\gamma &
((\mathcal{O}_{X, y}^\wedge/\mathfrak p)'_{\mathfrak q'})^\wedge \ar[d] \\
 & & &
\mathcal{O}_{X', y'}^\wedge/\mathfrak p'
}
$$
dont les flèches verticales sont surjectives. D'après
Compléments d'algèbre, lemme \ref{more-algebra-lemma-completion-dimension},
et la formule de dimension
(Algèbre, lemme \ref{algebra-lemma-dimension-formula}),
nous avons
$$
\dim(((\mathcal{O}_{X, y}^\wedge/\mathfrak p)'_{\mathfrak q'})^\wedge) =
\dim((\mathcal{O}_{X, y}^\wedge/\mathfrak p)'_{\mathfrak q'}) =
\dim(\mathcal{O}_{X, y}^\wedge/\mathfrak p)
- \text{trdeg}(\kappa(y')/\kappa(y))
$$
En suivant les définitions des images réciproques, des germes, des localisations
et des complétions, nous trouvons
$$
(\mathcal{F}_y^\wedge)_{\mathfrak p}
\otimes_{(\mathcal{O}_{X, y}^\wedge)_\mathfrak p}
(\mathcal{O}_{X', y'}^\wedge)_{\mathfrak p'}
=
(\mathcal{F}_{y'}^\wedge)_{\mathfrak p'}
$$
Détails omis. Les homomorphismes d'anneaux $\beta$ et $\gamma$ du diagramme
sont plats à fibres de Gorenstein (donc de Cohen--Macaulay), car ce sont des
complétions d'anneaux possédant un complexe dualisant. Voir
Complexes dualisants, lemmes
\ref{dualizing-lemma-formal-fibres-gorenstein} et
\ref{dualizing-lemma-dualizing-gorenstein-formal-fibres},
ainsi que la discussion de Compléments d'algèbre, section
\ref{more-algebra-section-properties-formal-fibres}.
Observons que $(\mathcal{O}_{X, y}^\wedge)_\mathfrak p =
(\mathcal{O}_{X, y}^\wedge)'_{\tilde{\mathfrak p}}$
où $\tilde{\mathfrak p}$ est le noyau de $\alpha$
dans le diagramme. D'autre part,
$(\mathcal{O}_{X, y}^\wedge)'_{\tilde{\mathfrak p}}
\to (\mathcal{O}_{X', y'}^\wedge)_{\mathfrak p'}$
est plat à fibres CM d'après ce qui précède. Par conséquent,
$(\mathcal{O}_{X, y}^\wedge)_\mathfrak p  \to
(\mathcal{O}_{X', y'}^\wedge)_{\mathfrak p'}$ est plat à fibres CM.
En appliquant Algèbre, lemme \ref{algebra-lemma-apply-grothendieck-module},
nous voyons que
$$
\text{depth}((\mathcal{F}_{y'}^\wedge)_{\mathfrak p'}) =
\text{depth}((\mathcal{F}_y^\wedge)_{\mathfrak p}) +
\dim(F_\mathfrak r)
$$
où $F$ est la fibre formelle générique de
$(\mathcal{O}_{X, y}^\wedge/\mathfrak p)'_{\mathfrak q'}$
et $\mathfrak r$ est l'idéal premier correspondant à $\mathfrak p'$.
Puisque $(\mathcal{O}_{X, y}^\wedge/\mathfrak p)'_{\mathfrak q'}$
est un domaine local universellement caténaire, son complété $I$-adique
est équidimensionnel et (universellement) caténaire d'après le théorème de Ratliff
(Compléments d'algèbre, proposition \ref{more-algebra-proposition-ratliff}).
Il s'ensuit alors que
$$
\dim(((\mathcal{O}_{X, y}^\wedge/\mathfrak p)'_{\mathfrak q'})^\wedge) =
\dim(F_\mathfrak r) + \dim(\mathcal{O}_{X', y'}^\wedge/\mathfrak p')
$$
En combinant ceci avec le lemme \ref{lemma-change-distance-function},
nous obtenons
\begin{equation}
\label{equation-one}
\begin{aligned}
&
\text{depth}((\mathcal{F}_{y'}^\wedge)_{\mathfrak p'}) +
\delta^{Y'}_{Z'}(y') \\
& =
\text{depth}((\mathcal{F}_y^\wedge)_{\mathfrak p}) +
\dim(F_\mathfrak r) + \delta^{Y'}_{Z'}(y') \\
& \geq
\text{depth}((\mathcal{F}_y^\wedge)_{\mathfrak p}) + \delta^Y_Z(y) +
\dim(F_\mathfrak r) + \text{trdeg}(\kappa(y')/\kappa(y)) - (d - 1) \\
& =
\text{depth}((\mathcal{F}_y^\wedge)_{\mathfrak p}) + \delta^Y_Z(y) - (d - 1)
+ \dim(\mathcal{O}_{X, y}^\wedge/\mathfrak p) -
\dim(\mathcal{O}_{X', y'}^\wedge/\mathfrak p')
\end{aligned}
\end{equation}
Gardons à l'esprit que
$\dim(\mathcal{O}_{X, y}^\wedge/\mathfrak p) \geq
\dim(\mathcal{O}_{X', y'}^\wedge/\mathfrak p')$. En réécrivant ceci, nous obtenons
\begin{equation}
\label{equation-two}
\begin{aligned}
&
\text{depth}((\mathcal{F}_{y'}^\wedge)_{\mathfrak p'}) +
\dim(\mathcal{O}_{X', y'}^\wedge/\mathfrak p') +
\delta^{Y'}_{Z'}(y') \\
& \geq
\text{depth}((\mathcal{F}_y^\wedge)_{\mathfrak p}) +
\dim(\mathcal{O}_{X, y}^\wedge/\mathfrak p) +
\delta^Y_Z(y) - (d - 1)
\end{aligned}
\end{equation}
Cette inégalité nous permettra de vérifier les conditions restantes.

\medskip\noindent
Conditions (b) et (d) du
lemme \ref{lemma-improvement-formal-coherent-module-better}. Supposons que
$V(\mathfrak p') \cap V(I\mathcal{O}_{X', y'}^\wedge) =
\{\mathfrak m_{y'}^\wedge\}$.
Cela implique que $\dim(\mathcal{O}_{X', y'}^\wedge/\mathfrak p') = 1$
parce que $Z'$ est un diviseur de Cartier effectif.
La conjonction de (b) et (d) équivaut à
$$
\text{depth}((\mathcal{F}_{y'}^\wedge)_{\mathfrak p'}) + \delta^{Y'}_{Z'}(y')
> 2
$$
Si $(\mathcal{F}_n)$ satisfait aux inégalités de (3)(b),
nous concluons immédiatement que ceci est vrai en appliquant (\ref{equation-two}).
Si $(\mathcal{F}_n)$ satisfait à (3)(a), c'est-à-dire aux
inégalités $(d + 1, d + 2)$, alors nous voyons que, dans tous les cas,
$$
\text{depth}((\mathcal{F}_y^\wedge)_{\mathfrak p}) + \delta^Y_Z(y)
\geq d + 1
\quad\text{ou}\quad
\text{depth}((\mathcal{F}_y^\wedge)_{\mathfrak p}) +
\dim(\mathcal{O}_{X, y}^\wedge/\mathfrak p) +
\delta^Y_Z(y) > d + 2
$$
En examinant (\ref{equation-one}) et (\ref{equation-two}) ci-dessus, ceci donne
ce que nous voulons, sauf éventuellement si
$\dim(\mathcal{O}_{X, y}^\wedge/\mathfrak p) = 1$.
Cependant, si $\dim(\mathcal{O}_{X, y}^\wedge/\mathfrak p) = 1$, alors nous avons
$V(\mathfrak p) \cap V(I\mathcal{O}_{X, y}^\wedge) = \{\mathfrak m_y^\wedge\}$
et nous voyons qu'en fait
$$
\text{depth}((\mathcal{F}_y^\wedge)_{\mathfrak p}) + \delta^Y_Z(y) > d + 1
$$
car $(\mathcal{F}_n)$ satisfait aux inégalités $(d + 1, d + 2)$, et nous
concluons de nouveau.

\medskip\noindent
Condition (c) du
lemme \ref{lemma-improvement-formal-coherent-module-better}. Supposons que
$V(\mathfrak p') \cap V(I\mathcal{O}_{X', y'}^\wedge) \not =
\{\mathfrak m_{y'}^\wedge\}$. Alors la condition (c) équivaut à
$$
\text{depth}((\mathcal{F}_{y'}^\wedge)_{\mathfrak p'}) + \delta^{Y'}_{Z'}(y')
\geq 2
\quad\text{ou}\quad
\text{depth}((\mathcal{F}_{y'}^\wedge)_{\mathfrak p'}) +
\dim(\mathcal{O}_{X', y'}^\wedge/\mathfrak p') +
\delta^{Y'}_{Z'}(y') > 3
$$
Si $(\mathcal{F}_n)$ satisfait aux inégalités de (3)(b),
nous voyons que la seconde des deux inégalités affichées est vérifiée
en appliquant (\ref{equation-two}). Si $(\mathcal{F}_n)$ satisfait à (3)(a), c'est-à-dire
aux inégalités $(d + 1, d + 2)$, cela résulte immédiatement de
(\ref{equation-one}) et (\ref{equation-two}).
Ceci achève la démonstration de notre assertion.

\medskip\noindent
Choisissons $(b|_{U'}^*\mathcal{F}_n) \to (\mathcal{F}_n'')$
et $(\mathcal{H}_n)$ dans $\textit{Coh}(U', I\mathcal{O}_{U'})$
comme dans le lemme \ref{lemma-improvement-formal-coherent-module-better}.
Pour tout ouvert affine $W \subset X'$, observons que
$\delta^{W \cap Y'}_{W \cap Z'}(y') \geq \delta^{Y'}_{Z'}(y')$ d'après la
partie (6) du lemme \ref{lemma-discussion}. Nous voyons donc que
$(\mathcal{H}_n|_W)$ satisfait aux hypothèses du
lemme \ref{lemma-cd-1-canonical}.
Ainsi, $(\mathcal{H}_n|_W)$ s'étend canoniquement à $W$.
Soit $(\mathcal{G}_{W, n})$ dans $\textit{Coh}(W, I\mathcal{O}_W)$
l'extension canonique comme dans le
lemme \ref{lemma-canonically-algebraizable}.
D'après le lemme \ref{lemma-canonically-extend-base-change},
nous voyons que, pour $W' \subset W$, il existe un unique isomorphisme
$$
(\mathcal{G}_{W, n}|_{W'}) \longrightarrow
(\mathcal{G}_{W', n})
$$
compatible avec les isomorphismes donnés
$(\mathcal{G}_{W, n}|_{W \cap U}) \cong (\mathcal{H}_n|_{W \cap U})$.
Nous en concluons qu'il existe un objet
$(\mathcal{G}_n)$ de $\textit{Coh}(X', I\mathcal{O}_{X'})$
dont la restriction à $U$ est isomorphe à $(\mathcal{H}_n)$.

\medskip\noindent
Si $A$ est complet pour la topologie $I$-adique, nous pouvons achever la démonstration comme suit.
D'après le théorème d'existence de Grothendieck
(Cohomologie des schémas, lemme \ref{coherent-lemma-existence-projective}),
nous voyons que $(\mathcal{G}_n)$ est le complété d'un
$\mathcal{O}_{X'}$-module cohérent. Puis, d'après
Cohomologie des schémas, lemme \ref{coherent-lemma-existence-easy},
nous voyons que $(b|_{U'}^*\mathcal{F}_n)$
est le complété d'un $\mathcal{O}_{U'}$-module cohérent
$\mathcal{F}'$.
D'après Cohomologie des schémas, lemme \ref{coherent-lemma-inverse-systems-push-pull},
nous voyons qu'il existe une application
$$
(\mathcal{F}_n) \longrightarrow ((b|_{U'})_*\mathcal{F}')^\wedge
$$
dont le noyau et le conoyau sont annulés par une puissance de $I$.
Enfin, nous concluons en appliquant le
lemme \ref{lemma-map-kernel-cokernel-on-closed}.

\medskip\noindent
Si $A$ n'est pas complet, alors, avant de commencer la démonstration, nous pouvons remplacer $A$
par son complété, voir le lemme \ref{lemma-algebraizable}.
Après complétion, les hypothèses restent vérifiées : cela est immédiat
pour la condition (3), résulte du
lemme \ref{dualizing-lemma-ubiquity-dualizing} de Complexes dualisants
pour la condition (1), et du
lemme \ref{divisors-lemma-flat-base-change-blowing-up} de Diviseurs
pour la condition (2).
Ainsi, le cas complet implique le cas général.
\end{proof}

\begin{proposition}[Algébrisation pour les idéaux ayant peu de générateurs]
\label{proposition-d-generators}
Dans la situation \ref{situation-algebraize}, soit $(\mathcal{F}_n)$
un objet de $\textit{Coh}(U, I\mathcal{O}_U)$. Supposons que
\begin{enumerate}
\item $A$ possède un complexe dualisant,
\item $V(I) = V(f_1, \ldots, f_d)$ pour un certain $d \geq 1$ et
$f_1, \ldots, f_d \in A$,
\item l'une des conditions suivantes soit vérifiée :
\begin{enumerate}
\item $(\mathcal{F}_n)$ satisfait aux inégalités $(d + 1, d + 2)$
(définition \ref{definition-s-d-inequalities}), ou
\item pour $y \in U \cap Y$ et un idéal premier
$\mathfrak p \subset \mathcal{O}_{X, y}^\wedge$ tel que
$\mathfrak p \not \in V(I\mathcal{O}_{X, y}^\wedge)$,
nous ayons
$$
\text{depth}((\mathcal{F}^\wedge_y)_\mathfrak p) +
\dim(\mathcal{O}_{X, y}^\wedge/\mathfrak p) + \delta^Y_Z(y) > d + 2
$$
\end{enumerate}
\end{enumerate}
Alors $(\mathcal{F}_n)$ s'étend à $X$. En particulier, si $A$ est
complet pour la topologie $I$-adique, alors $(\mathcal{F}_n)$ est le complété
d'un $\mathcal{O}_U$-module cohérent. 
\end{proposition}

\begin{proof}
Nous pouvons supposer que $I = (f_1, \ldots, f_d)$, voir
Cohomologie des schémas, lemme
\ref{coherent-lemma-inverse-systems-ideals-equivalence}.
Nous voyons alors que
toutes les fibres de l'éclatement de $X$ le long de $I$ sont de dimension au plus $d - 1$.
Nous obtenons donc l'extension par le lemme \ref{lemma-blowup}.
La dernière assertion résulte du lemme \ref{lemma-essential-image-completion}.
\end{proof}

\noindent
Comparons le lemme suivant avec les
remarques \ref{remark-interesting-case-variant},
\ref{remark-interesting-case},
\ref{remark-interesting-case-bis} et
\ref{remark-interesting-case-ter}.

\begin{lemma}
\label{lemma-interesting-case-final}
Dans la situation \ref{situation-algebraize}, soit $(\mathcal{F}_n)$
un objet de $\textit{Coh}(U, I\mathcal{O}_U)$. Supposons que
\begin{enumerate}
\item $A$ soit un anneau local possédant un complexe dualisant,
\item toutes les composantes irréductibles de $X$ aient la même dimension,
\item le schéma $X \setminus Y$ soit de Cohen--Macaulay,
\item $I$ soit engendré par $d$ éléments,
\item $\dim(X) - \dim(Z) > d + 2$, et
\item pour $y \in U \cap Y$, le module $\mathcal{F}_y^\wedge$
soit localement libre de type fini hors de $V(I\mathcal{O}_{X, y}^\wedge)$,
par exemple si $\mathcal{F}_n$ est un
$\mathcal{O}_U/I^n\mathcal{O}_U$-module localement libre de type fini.
\end{enumerate}
Alors $(\mathcal{F}_n)$ s'étend à $X$. En particulier, si $A$ est complet pour la topologie $I$-adique,
alors $(\mathcal{F}_n)$ est le complété d'un
$\mathcal{O}_U$-module cohérent.
\end{lemma}

\begin{proof}
Nous allons montrer que les hypothèses (1), (2) et (3)(b) de la
proposition \ref{proposition-d-generators} sont vérifiées.
C'est clair pour (1) et (2).

\medskip\noindent
Soit $y \in U \cap Y$ et soit $\mathfrak p$ un idéal premier
$\mathfrak p \subset \mathcal{O}_{X, y}^\wedge$ tel que
$\mathfrak p \not \in V(I\mathcal{O}_{X, y}^\wedge)$.
La dernière condition montre que
$\text{depth}((\mathcal{F}_y^\wedge)_\mathfrak p) =
\text{depth}((\mathcal{O}_{X, y}^\wedge)_\mathfrak p)$.
Puisque $X \setminus Y$ est de Cohen--Macaulay, nous voyons que
$(\mathcal{O}_{X, y}^\wedge)_\mathfrak p$ est de Cohen--Macaulay.
Ainsi, nous voyons que
\begin{align*}
& \text{depth}((\mathcal{F}^\wedge_y)_\mathfrak p) +
\dim(\mathcal{O}_{X, y}^\wedge/\mathfrak p) + \delta^Y_Z(y) \\
& =
\dim((\mathcal{O}_{X, y}^\wedge)_\mathfrak p) +
\dim(\mathcal{O}_{X, y}^\wedge/\mathfrak p) + \delta^Y_Z(y) \\
& =
\dim(\mathcal{O}_{X, y}^\wedge) + \delta^Y_Z(y)
\end{align*}
La dernière égalité vaut parce que $\mathcal{O}_{X, y}$ est équidimensionnel
d'après la deuxième condition.
Soit $\delta(y) = \dim(\overline{\{y\}})$. C'est une fonction de dimension
puisque $A$ est un anneau local caténaire.
D'après le lemme \ref{lemma-discussion},
nous avons $\delta^Y_Z(y) \geq \delta(y) - \dim(Z)$. Puisque $X$ est
équidimensionnel, nous obtenons
$$
\dim(\mathcal{O}_{X, y}^\wedge) + \delta^Y_Z(y)
\geq \dim(\mathcal{O}_{X, y}^\wedge) + \delta(y) - \dim(Z)
= \dim(X) - \dim(Z)
$$
Nous obtenons donc l'inégalité voulue, ce qui conclut.
\end{proof}

\begin{remark}
\label{remark-question}
Nous ne savons ni démontrer ni réfuter l'analogue de la
proposition \ref{proposition-d-generators}
où l'hypothèse selon laquelle $I$ possède $d$ générateurs
est remplacée par l'hypothèse $\text{cd}(A, I) \leq d$.
Si vous connaissez une démonstration ou un contre-exemple, veuillez écrire à
\href{mailto:stacks.project@gmail.com}{stacks.project@gmail.com}.
Une autre question évidente est de savoir dans quelle mesure les conditions de la
proposition \ref{proposition-d-generators}
sont nécessaires.
\end{remark}




\section{Algébrisation des modules formels cohérents, VI}
\label{section-algebraization-modules-yet-more}

\noindent
Dans cette section, nous ajoutons quelques cas plus faciles à démontrer.

\begin{proposition}
\label{proposition-algebraization-regular-sequence}
Dans la situation \ref{situation-algebraize}, soit
$(\mathcal{F}_n)$ un objet de $\textit{Coh}(U, I\mathcal{O}_U)$.
Supposons que
\begin{enumerate}
\item il existe $f_1, \ldots, f_d \in I$ tels que,
pour $y \in U \cap Y$, l'idéal $I\mathcal{O}_{X, y}$
soit engendré par $f_1, \ldots, f_d$ et que
$f_1, \ldots, f_d$ forment une suite $\mathcal{F}_y^\wedge$-régulière,
\item $H^0(U, \mathcal{F}_1)$ et $H^1(U, \mathcal{F}_1)$
soient des $A$-modules de type fini.
\end{enumerate}
Alors $(\mathcal{F}_n)$ s'étend canoniquement à $X$. En particulier, si $A$
est complet, alors $(\mathcal{F}_n)$ est le complété d'un
$\mathcal{O}_U$-module cohérent.
\end{proposition}

\begin{proof}
Nous allons le démontrer en vérifiant les hypothèses (a), (b) et (c) du
lemme \ref{lemma-when-done}.
Pour tout $n$, nous avons une suite exacte courte
$$
0 \to I^n\mathcal{F}_{n + 1} \to \mathcal{F}_{n + 1} \to \mathcal{F}_n \to 0
$$
Puisque $f_1, \ldots, f_d$ forment une suite régulière (et donc
quasi-régulière, voir Algèbre, lemme \ref{algebra-lemma-regular-quasi-regular})
sur chacun des « germes » $\mathcal{F}_y^\wedge$, et puisque nous avons
$I\mathcal{F}_n = (f_1, \ldots, f_d)\mathcal{F}_n$ pour tout $n$,
nous trouvons que
$$
I^n\mathcal{F}_{n + 1} =
\bigoplus\nolimits_{e_1 + \ldots + e_d = n} \mathcal{F}_1 \cdot
f_1^{e_1} \ldots f_d^{e_d}
$$
en le vérifiant sur les germes. En utilisant l'hypothèse de finitude de
$H^0(U, \mathcal{F}_1)$ et une récurrence, nous concluons d'abord que
$M_n = H^0(U, \mathcal{F}_n)$ est un $A$-module de type fini pour tout $n$.
Nous voyons ainsi que la condition (c) du lemme \ref{lemma-when-done} est vérifiée.
Nous voyons aussi que
$$
\bigoplus\nolimits_{n \geq 0} H^1(U, I^n\mathcal{F}_{n + 1})
$$
est un $R = \bigoplus I^n/I^{n +1}$-module gradué de type fini.
D'après Cohomologie, lemme \ref{cohomology-lemma-ML-general},
nous concluons que la condition (a) du
lemme \ref{lemma-when-done} est vérifiée. Enfin, la condition (b) du
lemme \ref{lemma-when-done} est vérifiée parce que
$\bigoplus H^0(U, I^n\mathcal{F}_{n + 1})$ est un $R$-module gradué de type fini
et que nous pouvons appliquer
Cohomologie, lemme \ref{cohomology-lemma-topology-I-adic-general}.
\end{proof}

\begin{remark}
\label{remark-interesting-case-ter}
Dans la situation de la
proposition \ref{proposition-algebraization-regular-sequence},
si nous supposons que $A$ possède un complexe dualisant, alors
la condition que $H^0(U, \mathcal{F}_1)$ et
$H^1(U, \mathcal{F}_1)$ soient de type fini équivaut à
$$
\text{depth}(\mathcal{F}_{1, y}) +
\dim(\mathcal{O}_{\overline{\{y\}}, z}) > 2
$$
pour tout $y \in U \cap Y$ et tout $z \in Z \cap \overline{\{y\}}$.
Voir Cohomologie locale, lemme \ref{local-cohomology-lemma-finiteness-Rjstar}.
Cela vaut par exemple si $\mathcal{F}_1$ est un
$\mathcal{O}_{U \cap Y}$-module localement libre de type fini, si $Y$ est $(S_2)$, et si
$\text{codim}(Z', Y') \geq 3$ pour toute paire de composantes irréductibles
$Y'$ de $Y$ et $Z'$ de $Z$ telles que $Z' \subset Y'$.
\end{remark}

\begin{proposition}
\label{proposition-algebraization-flat}
Dans la situation \ref{situation-algebraize}, soit
$(\mathcal{F}_n)$ un objet de $\textit{Coh}(U, I\mathcal{O}_U)$.
Supposons qu'il existe un anneau local noethérien $(R, \mathfrak m)$ et un homomorphisme
d'anneaux $R \to A$ tels que
\begin{enumerate}
\item $I = \mathfrak m A$,
\item pour $y \in U \cap Y$, le germe $\mathcal{F}_y^\wedge$ soit $R$-plat,
\item $H^0(U, \mathcal{F}_1)$ et $H^1(U, \mathcal{F}_1)$ soient des
$A$-modules de type fini.
\end{enumerate}
Alors $(\mathcal{F}_n)$ s'étend canoniquement à $X$. En particulier, si $A$
est complet, alors $(\mathcal{F}_n)$ est le complété d'un
$\mathcal{O}_U$-module cohérent.
\end{proposition}

\begin{proof}
La démonstration est exactement la même que celle de la
proposition \ref{proposition-algebraization-regular-sequence}.
En effet, si $\kappa = R/\mathfrak m$, alors, pour $n \geq 0$,
il existe un isomorphisme
$$
I^n \mathcal{F}_{n + 1} \cong
\mathcal{F}_1 \otimes_\kappa \mathfrak m^n/\mathfrak m^{n + 1}
$$
et le membre de droite est une somme directe finie de copies
de $\mathcal{F}_1$. Cela se vérifie sur les germes.
Tout le reste est exactement identique.
\end{proof}

\begin{remark}
\label{remark-interesting-case-quater}
La proposition \ref{proposition-algebraization-flat} est une version locale
de \cite[théorème 2.10 (i)]{Baranovsky}. Il est immédiat de déduire
les résultats globaux du résultat local ; esquissons l'argument.
En effet, supposons que $(R, \mathfrak m)$
soit un anneau local noethérien complet et que $X \to \Spec(R)$ soit un morphisme propre.
Pour $n \geq 1$, posons $X_n = X \times_{\Spec(R)} \Spec(R/\mathfrak m^n)$.
Soit $Z \subset X_1$ un sous-ensemble fermé de la fibre spéciale.
Posons $U = X \setminus Z$ et notons $j : U \to X$ le morphisme d'inclusion.
Supposons donné un objet
$$
(\mathcal{F}_n) \text{ de } \textit{Coh}(U, \mathfrak m\mathcal{O}_U)
$$
qui soit plat sur $R$ au sens où $\mathcal{F}_n$ est plat sur
$R/\mathfrak m^n$ pour tout $n$.
Supposons que $j_*\mathcal{F}_1$ et $R^1j_*\mathcal{F}_1$ soient des modules
cohérents. Alors, localement sur les ouverts affines de $X$, nous obtenons une extension canonique
de $(\mathcal{F}_n)$ par la
proposition \ref{proposition-algebraization-flat},
et la formation de cette extension commute à la localisation
(d'après le lemme \ref{lemma-algebraization-principal-variant}).
Nous obtenons donc un objet global canonique $(\mathcal{G}_n)$ de
$\textit{Coh}(X, \mathfrak m\mathcal{O}_X)$
dont la restriction à $U$ est $(\mathcal{F}_n)$.
D'après le théorème d'existence de Grothendieck
(Cohomologie des schémas, proposition
\ref{coherent-proposition-existence-proper}),
nous voyons qu'il existe un $\mathcal{O}_X$-module cohérent
$\mathcal{G}$ dont le complété est $(\mathcal{G}_n)$.
Nous voyons ainsi que $(\mathcal{F}_n)$ est algébrisable, c'est-à-dire
qu'il est le complété d'un $\mathcal{O}_U$-module cohérent.

\medskip\noindent
Ajoutons que la cohérence de $j_*\mathcal{F}_1$ et de $R^1j_*\mathcal{F}_1$
est une condition sur la fibre spéciale. En effet, si nous notons
$j_1 : U_1 \to X_1$ la fibre spéciale de $j : U \to X$, alors nous pouvons
considérer $\mathcal{F}_1$ comme un faisceau cohérent sur $U_1$, et nous avons
$j_*\mathcal{F}_1 = j_{1, *}\mathcal{F}_1$ et
$R^1j_*\mathcal{F}_1 = R^1j_{1, *}\mathcal{F}_1$.
Ainsi, par exemple, si $X_1$ est $(S_2)$ et irréductible, si nous avons
$\dim(X_1) - \dim(Z) \geq 3$, et si $\mathcal{F}_1$ est un
$\mathcal{O}_{U_1}$-module localement libre, alors $j_{1, *}\mathcal{F}_1$ et
$R^1j_{1, *}\mathcal{F}_1$ sont des modules cohérents.
\end{remark}








\section{Application au foncteur de complétion}
\label{section-completion-application}

\noindent
Dans cette section, nous combinons seulement quelques résultats déjà obtenus
afin de pouvoir les citer commodément. Nous pourrions énoncer ici de nombreux
résultats (plus forts).

\begin{lemma}
\label{lemma-equivalence-better}
Dans la situation \ref{situation-algebraize}, supposons que
\begin{enumerate}
\item $A$ possède un complexe dualisant et soit complet pour la topologie $I$-adique,
\item $I = (f)$ soit engendré par un seul élément,
\item $A$ soit local d'idéal maximal $\mathfrak a = \mathfrak m$,
\item l'une des conditions suivantes soit vérifiée :
\begin{enumerate}
\item $A_f$ est $(S_2)$ et, pour $\mathfrak p \subset A$
minimal tel que $f \not \in \mathfrak p$, nous avons $\dim(A/\mathfrak p) \geq 4$, ou
\item si $\mathfrak p \not \in V(f)$ et
$V(\mathfrak p) \cap V(f) \not = \{\mathfrak m\}$, alors
$\text{depth}(A_\mathfrak p) + \dim(A/\mathfrak p) > 3$.
\end{enumerate}
\end{enumerate}
Alors, avec $U_0 = U \cap V(f)$, le foncteur de complétion
$$
\colim_{U_0 \subset U' \subset U\text{ ouvert}}
\textit{Coh}(\mathcal{O}_{U'})
\longrightarrow
\textit{Coh}(U, f\mathcal{O}_U)
$$
est une équivalence sur les sous-catégories pleines d'objets localement libres de type fini.
\end{lemma}

\begin{proof}
Il résulte du lemme \ref{lemma-fully-faithful-general}
que le foncteur est pleinement fidèle (détails omis).
Démontrons la surjectivité essentielle. Soit $(\mathcal{F}_n)$ un objet localement
libre de type fini de $\textit{Coh}(U, f\mathcal{O}_U)$. D'après le
lemme \ref{lemma-algebraization-principal-bis} ou la
proposition \ref{proposition-cd-1},
il existe un $\mathcal{O}_U$-module cohérent $\mathcal{F}$
tel que $(\mathcal{F}_n)$ soit le complété de $\mathcal{F}$.
En effet, pour appliquer l'un ou l'autre résultat, la seule chose à
vérifier est que $(\mathcal{F}_n)$ satisfait aux inégalités $(2, 3)$.
Cela est fait dans le lemme \ref{lemma-unwinding-conditions-bis}. Si $y \in U_0$,
alors le complété $f$-adique du germe $\mathcal{F}_y$ est isomorphe à
un module libre de type fini sur le complété $f$-adique de $\mathcal{O}_{U, y}$.
Par conséquent, $\mathcal{F}$ est localement libre de type fini dans un voisinage ouvert
$U'$ de $U_0$. Ceci achève la démonstration.
\end{proof}

\begin{lemma}
\label{lemma-equivalence}
Dans la situation \ref{situation-algebraize}, supposons que
\begin{enumerate}
\item $I = (f)$ soit principal,
\item $A$ soit complet pour la topologie $f$-adique,
\item $f$ soit un non-diviseur de zéro,
\item $H^1_\mathfrak a(A/fA)$ et $H^2_\mathfrak a(A/fA)$
soient des $A$-modules de type fini.
\end{enumerate}
Alors, avec $U_0 = U \cap V(f)$, le foncteur de complétion
$$
\colim_{U_0 \subset U' \subset U\text{ ouvert}}
\textit{Coh}(\mathcal{O}_{U'})
\longrightarrow
\textit{Coh}(U, f\mathcal{O}_U)
$$
est une équivalence sur les sous-catégories pleines d'objets localement libres de type fini.
\end{lemma}

\begin{proof}
Le foncteur est pleinement fidèle d'après le
lemme \ref{lemma-fully-faithful-general-alternative}.
La surjectivité essentielle résulte du
lemme \ref{lemma-algebraization-principal-variant}.
\end{proof}








\section{Triplets cohérents}
\label{section-coherent-triples}

\noindent
Soit $(A, \mathfrak m)$ un anneau local noethérien.
Soit $f \in \mathfrak m$ un non-diviseur de zéro. Posons
$X = \Spec(A)$, $X_0 = \Spec(A/fA)$, $U = X \setminus V(\mathfrak m)$ et
$U_0 = U \cap X_0$.
Nous disons que $(\mathcal{F}, \mathcal{F}_0, \alpha)$ est un {\it triplet cohérent}
si nous avons
\begin{enumerate}
\item $\mathcal{F}$ est un $\mathcal{O}_U$-module cohérent tel que
$f : \mathcal{F} \to \mathcal{F}$ soit injective,
\item $\mathcal{F}_0$ est un $\mathcal{O}_{X_0}$-module cohérent,
\item $\alpha : \mathcal{F}/f\mathcal{F} \to \mathcal{F}_0|_{U_0}$
est un isomorphisme.
\end{enumerate}
Il existe une notion évidente de {\it morphisme de triplets cohérents}
qui fait de la collection de tous les triplets cohérents une catégorie.

\medskip\noindent
La catégorie des triplets cohérents est additive mais non abélienne.
Cependant, la notion de suite exacte courte de triplets
cohérents est claire.

\medskip\noindent
Étant donnés deux triplets cohérents $(\mathcal{F}, \mathcal{F}_0, \alpha)$
et $(\mathcal{G}, \mathcal{G}_0, \beta)$, il se peut que
$(\mathcal{F} \otimes_{\mathcal{O}_U} \mathcal{G},
\mathcal{F}_0  \otimes_{\mathcal{O}_{X_0}} \mathcal{G}_0,
\alpha \otimes \beta)$ ne soit pas un triplet cohérent\footnote{En effet, il
n'est pas nécessairement vrai que $f$
soit injectif sur $\mathcal{F} \otimes_{\mathcal{O}_U} \mathcal{G}$.}.
Cependant, si les germes $\mathcal{G}_x$ sont
libres pour tout $x \in U_0$, alors c'est bien le cas.

\medskip\noindent
Nous dirons que le triplet cohérent $(\mathcal{G}, \mathcal{G}_0, \beta)$
est {\it localement libre}, resp.\ {\it inversible},
si $\mathcal{G}$ et $\mathcal{G}_0$
sont des modules localement libres, resp.\ inversibles. Dans ce cas, tensoriser
par $(\mathcal{G}, \mathcal{G}_0, \beta)$ a un sens (voir ci-dessus)
et transforme les suites exactes courtes de triplets cohérents en suites exactes courtes
de triplets cohérents.

\begin{lemma}
\label{lemma-prepare-chi-triple}
Pour tout triplet cohérent $(\mathcal{F}, \mathcal{F}_0, \alpha)$,
il existe un $\mathcal{O}_X$-module cohérent $\mathcal{F}'$
tel que $f : \mathcal{F}' \to \mathcal{F}'$ soit injective,
un isomorphisme $\alpha' : \mathcal{F}'|_U \to \mathcal{F}$ et une application
$\alpha'_0 : \mathcal{F}'/f\mathcal{F}' \to \mathcal{F}_0$
tel que $\alpha \circ (\alpha' \bmod f) = \alpha'_0|_{U_0}$.
\end{lemma}

\begin{proof}
Choisissons un $A$-module de type fini $M$ tel que $\mathcal{F}$ soit la restriction
à $U$ du $\mathcal{O}_X$-module cohérent associé à $M$, voir
Cohomologie locale, lemme
\ref{local-cohomology-lemma-finiteness-pushforwards-and-H1-local}.
Puisque $\mathcal{F}$ est sans torsion par $f$, nous pouvons remplacer $M$ par
son quotient par la torsion annulée par une puissance de $f$.
D'autre part, posons $M_0 = \Gamma(X_0, \mathcal{F}_0)$,
de sorte que $\mathcal{F}_0$ soit le $\mathcal{O}_{X_0}$-module cohérent
associé au $A/fA$-module de type fini $M_0$.
D'après Cohomologie des schémas, lemme \ref{coherent-lemma-homs-over-open},
il existe un $n$ tel que
l'isomorphisme $\alpha_0$ corresponde à un homomorphisme de $A/fA$-modules
$\mathfrak m^n M/fM \to M_0$ (dont le noyau et le conoyau
sont annulés par une puissance de $\mathfrak m$, mais nous n'en avons pas besoin).
Ainsi, si nous prenons $M' = \mathfrak m^n M$ et si nous posons
$\mathcal{F}'$ égal au $\mathcal{O}_X$-module cohérent
associé à $M'$, alors le lemme est clair.
\end{proof}

\noindent
Soit $(\mathcal{F}, \mathcal{F}_0, \alpha)$ un triplet cohérent.
Choisissons $\mathcal{F}', \alpha', \alpha'_0$ comme dans le
lemme \ref{lemma-prepare-chi-triple}. Posons
\begin{equation}
\label{equation-chi-triple}
\chi(\mathcal{F}, \mathcal{F}_0, \alpha) =
\text{length}_A(\Coker(\alpha'_0)) - 
\text{length}_A(\Ker(\alpha'_0))
\end{equation}
L'expression du membre de droite a un sens, car $\alpha'_0$ est un isomorphisme
sur $U_0$ ; son noyau et son conoyau sont donc des modules cohérents supportés
par $\{\mathfrak m\}$, qui sont par conséquent de longueur finie
(Algèbre, lemme \ref{algebra-lemma-support-point}).

\begin{lemma}
\label{lemma-well-defined-chi-triple}
La quantité $\chi(\mathcal{F}, \mathcal{F}_0, \alpha)$ de
(\ref{equation-chi-triple}) ne dépend pas du choix de
$\mathcal{F}', \alpha', \alpha'_0$ comme dans le lemme \ref{lemma-prepare-chi-triple}.
\end{lemma}

\begin{proof}
Soient $\mathcal{F}', \alpha', \alpha'_0$ et
$\mathcal{F}'', \alpha'', \alpha''_0$ deux tels choix.
Pour $n > 0$, posons $\mathcal{F}'_n = \mathfrak m^n \mathcal{F}'$.
D'après Cohomologie des schémas, lemme \ref{coherent-lemma-homs-over-open},
pour un certain $n$,
il existe un homomorphisme de $\mathcal{O}_X$-modules $\mathcal{F}'_n \to \mathcal{F}''$
qui coïncide avec l'identification
$\mathcal{F}''|_U = \mathcal{F}'|_U$ déterminée par $\alpha'$ et $\alpha''$.
Alors le diagramme
$$
\xymatrix{
\mathcal{F}'_n/f\mathcal{F}'_n \ar[r] \ar[d] &
\mathcal{F}'/f\mathcal{F}' \ar[d]^{\alpha_0'} \\
\mathcal{F}''/f\mathcal{F}'' \ar[r]^{\alpha_0''} &
\mathcal{F}_0
}
$$
est commutatif après restriction à $U_0$. Par conséquent, d'après
Cohomologie des schémas, lemme \ref{coherent-lemma-homs-over-open},
il est commutatif après restriction à
$\mathfrak m^l(\mathcal{F}'_n/f\mathcal{F}'_n)$ pour un certain $l > 0$. Puisque
$\mathcal{F}'_{n + l}/f\mathcal{F}'_{n + l} \to \mathcal{F}'_n/f\mathcal{F}'_n$
se factorise par $\mathfrak m^l(\mathcal{F}'_n/f\mathcal{F}'_n)$,
nous voyons qu'après avoir remplacé $n$ par $n + l$, le diagramme
est commutatif. En d'autres termes, nous avons trouvé un troisième choix
$\mathcal{F}''', \alpha''', \alpha'''_0$
tel qu'il existe des applications $\mathcal{F}''' \to \mathcal{F}''$
et $\mathcal{F}''' \to \mathcal{F}'$ sur $X$,
compatibles avec les applications sur $U$ et $X_0$. Nous sommes ainsi ramenés
au cas examiné dans le paragraphe suivant.

\medskip\noindent
Supposons disposer d'une application $\mathcal{F}'' \to \mathcal{F}'$ sur $X$, compatible avec
$\alpha', \alpha''$ sur $U$ et avec $\alpha'_0, \alpha''_0$ sur $X_0$.
Observons que $\mathcal{F}'' \to \mathcal{F}'$ est injective, car c'est un
isomorphisme sur $U$ et que $f : \mathcal{F}'' \to \mathcal{F}''$
est injective. Il est clair que $\mathcal{F}'/\mathcal{F}''$ est supporté par
$\{\mathfrak m\}$, donc est de longueur finie. Nous avons les applications
de $\mathcal{O}_{X_0}$-modules cohérents
$$
\mathcal{F}''/f\mathcal{F}'' \to
\mathcal{F}'/f\mathcal{F}' \xrightarrow{\alpha'_0}
\mathcal{F}_0
$$
dont la composée est $\alpha''_0$ et qui sont des isomorphismes sur $U_0$.
L'algèbre homologique élémentaire donne une suite exacte à $6$ termes
$$
\begin{matrix}
0 \to
\Ker(\mathcal{F}''/f\mathcal{F}'' \to \mathcal{F}'/f\mathcal{F}') \to
\Ker(\alpha''_0) \to
\Ker(\alpha'_0) \to \\
\Coker(\mathcal{F}''/f\mathcal{F}'' \to \mathcal{F}'/f\mathcal{F}') \to
\Coker(\alpha''_0) \to
\Coker(\alpha'_0) \to 0
\end{matrix}
$$
Par additivité des longueurs (Algèbre, lemme \ref{algebra-lemma-length-additive}),
nous trouvons qu'il suffit de montrer que
$$
\text{length}_A(
\Coker(\mathcal{F}''/f\mathcal{F}'' \to \mathcal{F}'/f\mathcal{F}')) -
\text{length}_A(
\Ker(\mathcal{F}''/f\mathcal{F}'' \to \mathcal{F}'/f\mathcal{F}')) = 0
$$
Cela résulte de l'application du lemme du serpent au
diagramme
$$
\xymatrix{
0 \ar[r] &
\mathcal{F}'' \ar[r]_f \ar[d] &
\mathcal{F}'' \ar[r] \ar[d] &
\mathcal{F}''/f\mathcal{F}'' \ar[r] \ar[d] &
0 \\
0 \ar[r] &
\mathcal{F}' \ar[r]^f &
\mathcal{F}' \ar[r] &
\mathcal{F}'/f\mathcal{F}' \ar[r] &
0
}
$$
et du fait que $\mathcal{F}'/\mathcal{F}''$ est de longueur finie.
\end{proof}

\begin{lemma}
\label{lemma-ses-chi-triple}
Nous avons
$\chi(\mathcal{G}, \mathcal{G}_0, \beta) =
\chi(\mathcal{F}, \mathcal{F}_0, \alpha) +
\chi(\mathcal{H}, \mathcal{H}_0, \gamma)$ si
$$
0 \to
(\mathcal{F}, \mathcal{F}_0, \alpha) \to
(\mathcal{G}, \mathcal{G}_0, \beta) \to
(\mathcal{H}, \mathcal{H}_0, \gamma)
\to 0
$$
est une suite exacte courte de triplets cohérents.
\end{lemma}

\begin{proof}
Choisissons $\mathcal{G}', \beta', \beta'_0$ comme dans le
lemme \ref{lemma-prepare-chi-triple}
pour le triplet $(\mathcal{G}, \mathcal{G}_0, \beta)$.
Notons $j : U \to X$ le morphisme d'inclusion.
Soit $\mathcal{F}' \subset \mathcal{G}'$
le noyau de la composée
$$
\mathcal{G}' \xrightarrow{\beta'} j_*\mathcal{G} \to j_*\mathcal{H}
$$
Observons que $\mathcal{H}' = \mathcal{G}'/\mathcal{F}'$
est un sous-faisceau cohérent de $j_*\mathcal{H}$ et que, par conséquent,
$f : \mathcal{H}' \to \mathcal{H}'$ est injective.
Le lemme du serpent donne donc une suite exacte courte
$$
0 \to \mathcal{F}'/f\mathcal{F}' \to
\mathcal{G}'/f\mathcal{G}' \to
\mathcal{H}'/f\mathcal{H}' \to 0
$$
Nous avons par construction les isomorphismes
$\alpha' : \mathcal{F}'|_U \to \mathcal{F}$,
$\beta' : \mathcal{G}'|_U \to \mathcal{G}$ et
$\gamma' : \mathcal{H}'|_U \to \mathcal{H}$.
Pour achever la démonstration, il nous faudra construire des applications
$\alpha'_0 : \mathcal{F}'/f\mathcal{F}' \to \mathcal{F}_0$ et
$\gamma'_0 : \mathcal{H}'/f\mathcal{H}' \to \mathcal{H}_0$ comme dans le
lemme \ref{lemma-prepare-chi-triple} et s'insérant dans
un diagramme commutatif
$$
\xymatrix{
0 \ar[r] &
\mathcal{F}'/f\mathcal{F}' \ar[r] \ar@{..>}[d]^{\alpha'_0} &
\mathcal{G}'/f\mathcal{G}' \ar[r] \ar[d]^{\beta'_0} &
\mathcal{H}'/f\mathcal{H}' \ar[r] \ar@{..>}[d]^{\gamma'_0} &
0 \\
0 \ar[r] &
\mathcal{F}_0 \ar[r] &
\mathcal{G}_0 \ar[r] &
\mathcal{H}_0 \ar[r] &
0
}
$$
Cependant, cela peut être impossible avec notre choix initial de $\mathcal{G}'$.
Le diagramme affiché montre que l'obstruction est
exactement la composée
$$
\delta :
\mathcal{F}'/f\mathcal{F}' \to
\mathcal{G}'/f\mathcal{G}' \xrightarrow{\beta'_0}
\mathcal{G}_0 \to
\mathcal{H}_0
$$
Notons que la restriction de $\delta$ à $U_0$ est nulle par notre choix de
$\mathcal{F}'$ et $\mathcal{H}'$. Par conséquent, d'après
Cohomologie des schémas, lemme \ref{coherent-lemma-homs-over-open},
il existe un $k > 0$ tel que
$\delta$ s'annule sur $\mathfrak m^k \cdot (\mathcal{F}'/f\mathcal{F}')$.
Pour $n > k$, posons $\mathcal{G}'_n = \mathfrak m^n \mathcal{G}'$,
$\mathcal{F}'_n = \mathcal{G}'_n \cap \mathcal{F}'$ et
$\mathcal{H}'_n = \mathcal{G}'_n/\mathcal{F}'_n$.
Observons que $\beta'_0$ peut être composée avec
$\mathcal{G}'_n/f\mathcal{G}'_n \to \mathcal{G}'/f\mathcal{G}'$
pour donner une application
$\beta'_{n, 0} : \mathcal{G}'_n/f\mathcal{G}'_n \to \mathcal{G}_0$
comme dans le lemme \ref{lemma-prepare-chi-triple}.
D'après Artin--Rees (Algèbre, lemme \ref{algebra-lemma-Artin-Rees}),
nous pouvons choisir $n$ tel que
$\mathcal{F}'_n \subset \mathfrak m^k \mathcal{F}'$.
Comme ci-dessus, les applications
$f : \mathcal{F}'_n \to \mathcal{F}'_n$,
$f : \mathcal{G}'_n \to \mathcal{G}'_n$ et
$f : \mathcal{H}'_n \to \mathcal{H}'_n$ sont injectives
et, comme ci-dessus, le lemme du serpent donne une suite exacte
courte
$$
0 \to \mathcal{F}'_n/f\mathcal{F}'_n \to
\mathcal{G}'_n/f\mathcal{G}'_n \to
\mathcal{H}'_n/f\mathcal{H}'_n \to 0
$$
Comme ci-dessus, nous avons les isomorphismes
$\alpha'_n : \mathcal{F}'_n|_U \to \mathcal{F}$,
$\beta'_n : \mathcal{G}'_n|_U \to \mathcal{G}$ et
$\gamma'_n : \mathcal{H}'_n|_U \to \mathcal{H}$.
Considérons l'obstruction
$$
\delta_n :
\mathcal{F}'_n/f\mathcal{F}'_n \to
\mathcal{G}'_n/f\mathcal{G}'_n
\xrightarrow{\beta'_{n, 0}}
\mathcal{G}_0 \to
\mathcal{H}_0
$$
comme précédemment. Cependant, le diagramme commutatif
$$
\xymatrix{
\mathcal{F}'_n/f\mathcal{F}'_n \ar[r] \ar[d] &
\mathcal{G}'_n/f\mathcal{G}'_n \ar[r]_{\beta'_{n, 0}} \ar[d] &
\mathcal{G}_0 \ar[r] \ar[d] &
\mathcal{H}_0 \ar[d] \\
\mathcal{F}'/f\mathcal{F}' \ar[r] &
\mathcal{G}'/f\mathcal{G}' \ar[r]^{\beta'_0} &
\mathcal{G}_0 \ar[r] &
\mathcal{H}_0
}
$$
notre choix de $n$ et notre observation concernant $\delta$
montrent que $\delta_n = 0$.
Ceci fournit les applications voulues
$\alpha'_{n, 0} : \mathcal{F}'_n/f\mathcal{F}'_n \to \mathcal{F}_0$ et
$\gamma'_{n, 0} : \mathcal{H}'_n/f\mathcal{H}'_n \to \mathcal{H}_0$.
Ainsi, nous pouvons utiliser
$\mathcal{F}'_n, \alpha'_n, \alpha'_{n, 0}$,
$\mathcal{G}'_n, \beta'_n, \beta'_{n, 0}$ et
$\mathcal{H}'_n, \gamma'_n, \gamma'_{n, 0}$
pour calculer
$\chi(\mathcal{F}, \mathcal{F}_0, \alpha)$,
$\chi(\mathcal{G}, \mathcal{G}_0, \beta)$ et
$\chi(\mathcal{H}, \mathcal{H}_0, \gamma)$.
Enfin, le lemme résulte d'une
application du lemme du serpent à
$$
\xymatrix{
0 \ar[r] &
\mathcal{F}'_n/f\mathcal{F}'_n \ar[r] \ar[d] &
\mathcal{G}'_n/f\mathcal{G}'_n \ar[r] \ar[d] &
\mathcal{H}'_n/f\mathcal{H}'_n \ar[r] \ar[d] &
0 \\
0 \ar[r] &
\mathcal{F}_0 \ar[r] &
\mathcal{G}_0 \ar[r] &
\mathcal{H}_0 \ar[r] &
0
}
$$
et de l'additivité des longueurs (Algèbre, lemme \ref{algebra-lemma-length-additive}).
\end{proof}

\begin{proposition}
\label{proposition-hilbert-triple}
Soit $(\mathcal{F}, \mathcal{F}_0, \alpha)$ un triplet cohérent.
Soit $(\mathcal{L}, \mathcal{L}_0, \lambda)$ un triplet cohérent
inversible. Alors la fonction
$$
\mathbf{Z} \longrightarrow \mathbf{Z},\quad
n \longmapsto
\chi((\mathcal{F}, \mathcal{F}_0, \alpha) \otimes
(\mathcal{L}, \mathcal{L}_0, \lambda)^{\otimes n})
$$
est un polynôme de degré $\leq \dim(\text{Supp}(\mathcal{F}))$.
\end{proposition}

\noindent
Plus précisément, si $\mathcal{F} = 0$, alors la fonction est constante.
Si $\mathcal{F}$ est à support fini dans $U$, alors la fonction est constante.
Si le support de $\mathcal{F}$ dans $U$ est de dimension $1$, c'est-à-dire
si l'adhérence du support de $\mathcal{F}$ dans $X$ est de dimension $2$, alors
la fonction est linéaire, et ainsi de suite.

\begin{proof}
Nous allons le démontrer par récurrence sur la dimension du support de
$\mathcal{F}$.

\medskip\noindent
Le cas initial est celui où $\mathcal{F} = 0$. Alors soit
$\mathcal{F}_0$ est nul, soit son support est $\{\mathfrak m\}$.
Dans ce cas, nous avons
$$
(\mathcal{F}, \mathcal{F}_0, \alpha) \otimes
(\mathcal{L}, \mathcal{L}_0, \lambda)^{\otimes n} =
(0, \mathcal{F}_0 \otimes \mathcal{L}_0^{\otimes n}, 0) \cong
(0, \mathcal{F}_0, 0)
$$
Ainsi, la fonction de la proposition est constante, de valeur égale
à la longueur de $\mathcal{F}_0$.

\medskip\noindent
Étape de récurrence. Supposons le support de $\mathcal{F}$ non vide.
Soit $\mathcal{G}_0 \subset \mathcal{F}_0$ le sous-module
des sections supportées par $\{\mathfrak m\}$. Nous obtenons alors une suite
exacte courte
$$
0 \to (0, \mathcal{G}_0, 0) \to
(\mathcal{F}, \mathcal{F}_0, \alpha) \to
(\mathcal{F}, \mathcal{F}_0/\mathcal{G}_0, \alpha) \to 0
$$
Cette suite reste exacte si nous tensorisons par le triplet cohérent
inversible $(\mathcal{L}, \mathcal{L}_0, \lambda)$, voir la
discussion ci-dessus. Ainsi, par additivité de $\chi$
(lemme \ref{lemma-ses-chi-triple})
et par le cas initial expliqué ci-dessus, il suffit de démontrer
l'étape de récurrence pour
$(\mathcal{F}, \mathcal{F}_0/\mathcal{G}_0, \alpha)$.
Nous voyons ainsi que nous pouvons supposer que $\mathfrak m$ n'est pas
un point associé de $\mathcal{F}_0$.

\medskip\noindent
Soit $T = \text{Ass}(\mathcal{F}) \cup \text{Ass}(\mathcal{F}/f\mathcal{F})$.
Puisque $U$ est quasi-affine, nous pouvons trouver $s \in \Gamma(U, \mathcal{L})$
qui ne s'annule en aucun $u \in T$, voir
Propriétés, lemme
\ref{properties-lemma-quasi-affine-invertible-nonvanishing-section}.
Après avoir multiplié $s$ par un élément convenable de $\mathfrak m$,
nous pouvons supposer que $\lambda(s \bmod f) = s_0|_{U_0}$ pour un certain
$s_0 \in \Gamma(X_0, \mathcal{L}_0)$ ; détails omis.
Nous obtenons un morphisme
$$
(s, s_0) :
(\mathcal{O}_U, \mathcal{O}_{X_0}, 1)
\longrightarrow
(\mathcal{L}, \mathcal{L}_0, \lambda)
$$
dans la catégorie des triplets cohérents. Posons
$\mathcal{G} = \Coker(s : \mathcal{F} \to \mathcal{F} \otimes \mathcal{L})$
et
$\mathcal{G}_0 = \Coker(s_0 : \mathcal{F}_0 \to
\mathcal{F}_0 \otimes \mathcal{L}_0)$. Observons que $s_0 : \mathcal{F}_0 \to
\mathcal{F}_0 \otimes \mathcal{L}_0$ est injective, car elle est injective
sur $U_0$ par notre choix de $s$ et que $\mathfrak m$ n'est pas un
point associé de $\mathcal{F}_0$. Par conséquent,
il existe un
isomorphisme $\beta : \mathcal{G}/f\mathcal{G} \to \mathcal{G}_0|_{U_0}$
tel que nous obtenions une suite exacte courte
$$
0 \to
(\mathcal{F}, \mathcal{F}_0, \alpha) \to
(\mathcal{F}, \mathcal{F}_0, \alpha) \otimes
(\mathcal{L}, \mathcal{L}_0, \lambda) \to
(\mathcal{G}, \mathcal{G}_0, \beta) \to 0
$$
Par récurrence sur la dimension du support, nous savons que la proposition
est vraie pour le triplet cohérent $(\mathcal{G}, \mathcal{G}_0, \beta)$.
En utilisant l'additivité du lemme \ref{lemma-ses-chi-triple},
nous voyons que
$$
n \longmapsto 
\chi((\mathcal{F}, \mathcal{F}_0, \alpha) \otimes
(\mathcal{L}, \mathcal{L}_0, \lambda)^{\otimes n + 1})
-
\chi((\mathcal{F}, \mathcal{F}_0, \alpha) \otimes
(\mathcal{L}, \mathcal{L}_0, \lambda)^{\otimes n})
$$
est un polynôme. Nous concluons par une variante du
lemme \ref{algebra-lemma-numerical-polynomial} d'Algèbre
pour les fonctions définies sur tous les entiers (détails omis).
\end{proof}

\begin{lemma}
\label{lemma-nonnegative-chi-triple}
Supposons que $\text{depth}(A) \geq 3$ ou, de manière équivalente, que
$\text{depth}(A/fA) \geq 2$. Soit $(\mathcal{L}, \mathcal{L}_0, \lambda)$
un triplet cohérent inversible. Alors
$$
\chi(\mathcal{L}, \mathcal{L}_0, \lambda) =
\text{length}_A \Coker(\Gamma(U, \mathcal{L}) \to \Gamma(U_0, \mathcal{L}_0))
$$
et, en particulier, ceci est $\geq 0$. De plus, 
$\chi(\mathcal{L}, \mathcal{L}_0, \lambda) = 0$ si et seulement si
$\mathcal{L} \cong \mathcal{O}_U$.
\end{lemma}

\begin{proof}
L'équivalence des conditions de profondeur résulte du
lemme \ref{algebra-lemma-depth-drops-by-one} d'Algèbre.
La condition de profondeur donne
$\Gamma(U, \mathcal{O}_U) = A$ et
$\Gamma(U_0, \mathcal{O}_{U_0}) = A/fA$, voir
Complexes dualisants, lemme \ref{dualizing-lemma-depth}, et
Cohomologie locale, lemme
\ref{local-cohomology-lemma-finiteness-pushforwards-and-H1-local}.
En utilisant Cohomologie locale, lemme
\ref{local-cohomology-lemma-finiteness-for-finite-locally-free},
nous trouvons que $M = \Gamma(U, \mathcal{L})$ est un $A$-module de type fini.
Ceci implique à son tour que $\text{depth}(M) \geq 2$, par exemple d'après la
partie (4) de Cohomologie locale, lemme
\ref{local-cohomology-lemma-finiteness-pushforwards-and-H1-local},
ou d'après Diviseurs, lemme \ref{divisors-lemma-depth-pushforward}.
De plus, nous avons $\mathcal{L}_0 \cong \mathcal{O}_{X_0}$,
car $X_0$ est un schéma local. Nous voyons donc aussi que
$M_0 = \Gamma(X_0, \mathcal{L}_0) = \Gamma(U_0, \mathcal{L}_0|_{U_0})$
et que ce module est isomorphe à $A/fA$.

\medskip\noindent
D'après ce qui précède, $\mathcal{F}' = \widetilde{M}$ est un
$\mathcal{O}_X$-module cohérent dont la restriction à $U$ est isomorphe à $\mathcal{L}$.
L'isomorphisme $\lambda : \mathcal{L}/f\mathcal{L} \to \mathcal{L}_0|_{U_0}$
détermine sur les sections globales une application $M/fM \to M_0$
qui est un isomorphisme sur $U_0$.
Puisque $\text{depth}(M) \geq 2$, nous voyons
que $H^0_\mathfrak m(M/fM) = 0$, et il s'ensuit que
$M/fM \to M_0$ est injective. Ainsi, par définition,
$$
\chi(\mathcal{L}, \mathcal{L}_0, \lambda) =
\text{length}_A \Coker(M/fM \to M_0)
$$
ce qui donne la première assertion du lemme.

\medskip\noindent
Enfin, si cette longueur vaut $0$, alors $M \to M_0$ est surjective.
Nous pouvons donc trouver $s \in M = \Gamma(U, \mathcal{L})$
qui s'envoie sur une section trivialisante de $\mathcal{L}_0$.
Considérons les $A$-modules de type fini $K$, $Q$ définis par la suite
exacte
$$
0 \to K \to A \xrightarrow{s} M \to Q \to 0
$$
Les supports de $K$ et de $Q$ ne rencontrent pas $U_0$, car $s$
est non nul aux points de $U_0$. En utilisant
Algèbre, lemme \ref{algebra-lemma-depth-in-ses},
nous voyons que $\text{depth}(K) \geq 2$ (observons que
$As \subset M$ est de $\text{depth} \geq 1$ comme sous-module de $M$).
Ainsi, s'il est non vide, le support de $K$ est de dimension $\geq 2$ d'après
Algèbre, lemme \ref{algebra-lemma-bound-depth}.
Ceci contredit $\text{Supp}(M) \cap V(f) \subset \{\mathfrak m\}$,
à moins que $K = 0$. Lorsque $K = 0$, nous trouvons que
$\text{depth}(Q) \geq 2$ et concluons
$Q = 0$ comme précédemment. Ainsi $A \cong M$ et
$\mathcal{L}$ est trivial.
\end{proof}







\section{Modules inversibles sur les spectres épointés, I}
\label{section-local-lefschetz-for-pic}

\noindent
Dans cette section, nous démontrons quelques théorèmes de Lefschetz locaux pour le groupe de Picard.
Certaines idées sont tirées de
\cite{Kollar-pic}, \cite{Bhatt-local} et \cite{Kollar-map-pic}.

\begin{lemma}
\label{lemma-injective-torsion-in-pic}
Soit $(A, \mathfrak m)$ un anneau local noethérien. Soit $f \in \mathfrak m$
un non-diviseur de zéro et supposons que $\text{depth}(A/fA) \geq 2$, ou, de manière équivalente,
$\text{depth}(A) \geq 3$. Soient $U$, resp.\ $U_0$, les spectres épointés
de $A$, resp.\ de $A/fA$. L'application
$$
\Pic(U) \to \Pic(U_0)
$$
est injective sur le sous-groupe de torsion.
\end{lemma}

\begin{proof}
Soit $\mathcal{L}$ un $\mathcal{O}_U$-module inversible.
Observons que $\mathcal{L}$ s'envoie sur $0$ dans $\Pic(U_0)$
si et seulement si nous pouvons prolonger $\mathcal{L}$ en un triplet cohérent
inversible $(\mathcal{L}, \mathcal{L}_0, \lambda)$
comme dans la section \ref{section-coherent-triples}.
D'après la proposition \ref{proposition-hilbert-triple},
la fonction
$$
n \longmapsto \chi((\mathcal{L}, \mathcal{L}_0, \lambda)^{\otimes n})
$$
est un polynôme. D'après le lemme \ref{lemma-nonnegative-chi-triple},
la valeur de ce polynôme est nulle si et seulement si
$\mathcal{L}^{\otimes n}$ est trivial.
Ainsi, si $\mathcal{L}$ est de torsion, alors ce
polynôme a une infinité de zéros, donc est
identiquement nul, et par conséquent $\mathcal{L}$ est trivial.
\end{proof}

\begin{proposition}[Koll\'ar]
\label{proposition-injective-pic}
\begin{reference}
\cite[Theorem 1.9]{Kollar-map-pic}
\end{reference}
Soit $(A, \mathfrak m)$ un anneau local noethérien. Soit $f \in \mathfrak m$.
Supposons que
\begin{enumerate}
\item $A$ possède un complexe dualisant,
\item $f$ soit un non-diviseur de zéro,
\item $\text{depth}(A/fA) \geq 2$, ou, de manière équivalente, $\text{depth}(A) \geq 3$,
\item si $f \in \mathfrak p \subset A$ est un idéal premier tel que
$\dim(A/\mathfrak p) = 2$, alors $\text{depth}(A_\mathfrak p) \geq 2$.
\end{enumerate}
Soient $U$, resp.\ $U_0$, les spectres épointés de $A$, resp.\ de $A/fA$. L'application
$$
\Pic(U) \to \Pic(U_0)
$$
est injective. Enfin, si (1), (2) et (3) sont vérifiées, si $A$ est $(S_2)$ et si
$\dim(A) \geq 4$, alors (4) est vérifiée.
\end{proposition}

\begin{proof}
Soit $\mathcal{L}$ un $\mathcal{O}_U$-module inversible.
Observons que $\mathcal{L}$ s'envoie sur $0$ dans $\Pic(U_0)$
si et seulement si nous pouvons prolonger $\mathcal{L}$ en un triplet cohérent
inversible $(\mathcal{L}, \mathcal{L}_0, \lambda)$
comme dans la section \ref{section-coherent-triples}.
D'après la proposition \ref{proposition-hilbert-triple},
la fonction
$$
n \longmapsto \chi((\mathcal{L}, \mathcal{L}_0, \lambda)^{\otimes n})
$$
est un polynôme $P$. D'après le lemme \ref{lemma-nonnegative-chi-triple},
nous avons
$P(n) \geq 0$ pour tout $n \in \mathbf{Z}$, avec égalité si et seulement si
$\mathcal{L}^{\otimes n}$ est trivial. En particulier, $P(0) = 0$
et soit $P$ est identiquement nul, auquel cas nous concluons, soit $P$ est de degré pair $\geq 2$.

\medskip\noindent
Posons $M = \Gamma(U, \mathcal{L})$ et
$M_0 = \Gamma(X_0, \mathcal{L}_0) = \Gamma(U_0, \mathcal{L}_0)$.
Alors $M$ est un $A$-module de type fini de profondeur $\geq 2$
et $M_0 \cong A/fA$, voir la démonstration du lemme \ref{lemma-nonnegative-chi-triple}.
Notons que $H^2_\mathfrak m(M)$ est un $A$-module de type fini d'après
Cohomologie locale, lemme
\ref{local-cohomology-lemma-local-finiteness-for-finite-locally-free},
et le fait que $H^i_\mathfrak m(A) = 0$ pour $i = 0, 1, 2$
puisque $\text{depth}(A) \geq 3$.
Considérons la suite exacte courte
$$
0 \to M/fM \to M_0 \to Q \to 0
$$
Le lemme \ref{lemma-nonnegative-chi-triple} nous dit que $Q$ est de longueur finie
égale à $\chi(\mathcal{L}, \mathcal{L}_0, \lambda)$.
Nous obtenons $Q = H^1_\mathfrak m(M/fM)$ et
$H^i_\mathfrak m(M/fM) = H^i_\mathfrak m(M_0) \cong H^i_\mathfrak m(A/fA)$
pour $i > 1$ par la suite exacte longue de cohomologie locale
associée à la suite exacte courte affichée. Considérons la suite
exacte longue de cohomologie locale associée à la suite
$0 \to M \to M \to M/fM \to 0$. Elle commence par
$$
0 \to Q \to H^2_\mathfrak m(M) \to H^2_\mathfrak m(M) \to
H^2_\mathfrak m(A/fA)
$$
Par additivité des longueurs, nous voyons que
$\chi(\mathcal{L}, \mathcal{L}_0, \lambda)$
est égale à la longueur de l'image de
$H^2_\mathfrak m(M) \to H^2_\mathfrak m(A/fA)$.

\medskip\noindent
Démontrons la proposition dans un cas particulier afin d'éclairer la suite de la démonstration.
Supposons en effet un instant que $H^2_\mathfrak m(A/fA)$ soit
un module de longueur finie. Nous aurions alors
$P(1) \leq \text{length}_A H^2_\mathfrak m(A/fA)$.
Exactement le même argument, appliqué à $\mathcal{L}^{\otimes n}$, montre que
$P(n) \leq \text{length}_A H^2_\mathfrak m(A/fA)$ pour tout $n$.
Ainsi, $P$ ne peut pas être de degré positif, et nous concluons.
Dans la suite de la démonstration, nous allons modifier cet argument afin d'obtenir
une borne supérieure linéaire pour $P(n)$, ce qui suffit.

\medskip\noindent
Étudions l'application
$H^2_\mathfrak m(M) \to H^2_\mathfrak m(M_0) \cong H^2_\mathfrak m(A/fA)$.
Choisissons un complexe dualisant normalisé $\omega_A^\bullet$ pour $A$.
Par dualité locale
(Complexes dualisants, lemme \ref{dualizing-lemma-special-case-local-duality}),
cette application est, au sens de Matlis, duale de l'application
$$
\text{Ext}^{-2}_A(M, \omega_A^\bullet)
\longleftarrow
\text{Ext}^{-2}_A(M_0, \omega_A^\bullet)
$$
dont l'image a donc la même longueur (finie).
Le support (s'il est non vide) du $A$-module de type fini
$\text{Ext}^{-2}_A(M_0, \omega_A^\bullet)$ est constitué de
$\mathfrak m$ et d'un nombre fini d'idéaux premiers
$\mathfrak p_1, \ldots, \mathfrak p_r$ contenant $f$ et tels que
$\dim(A/\mathfrak p_i) = 1$. En effet, d'après
Cohomologie locale, lemme \ref{local-cohomology-lemma-sitting-in-degrees},
le support est contenu dans l'ensemble des idéaux premiers $\mathfrak p \subset A$ tels que
$\text{depth}_{A_\mathfrak p}(M_{0, \mathfrak p}) + \dim(A/\mathfrak p) \leq 2$.
Il suffit donc de montrer qu'il n'existe aucun idéal premier $\mathfrak p$ contenant $f$ tel que
$\dim(A/\mathfrak p) = 2$ et
$\text{depth}_{A_\mathfrak p}(M_{0, \mathfrak p}) = 0$.
Cependant, puisque $M_{0, \mathfrak p} \cong (A/fA)_\mathfrak p$,
cela donnerait $\text{depth}(A_\mathfrak p) = 1$, ce qui contredit
l'hypothèse (4).
Choisissons une section $t \in \Gamma(U, \mathcal{L}^{\otimes -1})$
qui ne s'annule en aucun des points $\mathfrak p_1, \ldots, \mathfrak p_r$, voir
Propriétés, lemme
\ref{properties-lemma-quasi-affine-invertible-nonvanishing-section}.
La multiplication par $t$ sur les sections globales détermine une application $t : M \to A$
qui définit un isomorphisme $M_{\mathfrak p_i} \to A_{\mathfrak p_i}$ pour
$i = 1, \ldots, r$. Notons $t_0 = t|_{U_0}$ la section correspondante
de $\Gamma(U_0, \mathcal{L}_0^{\otimes -1})$, qui détermine de même
une application $t_0 : M_0 \to A/fA$ compatible avec $t$.
Nous en concluons qu'il existe un diagramme commutatif
$$
\xymatrix{
\text{Ext}^{-2}_A(M, \omega_A^\bullet) &
\text{Ext}^{-2}_A(M_0, \omega_A^\bullet) \ar[l] \\
\text{Ext}^{-2}_A(A, \omega_A^\bullet) \ar[u]^t &
\text{Ext}^{-2}_A(A/fA, \omega_A^\bullet) \ar[l] \ar[u]_{t_0}
}
$$
Il s'ensuit que la longueur de l'image de l'application horizontale
supérieure est au plus la longueur de $\text{Ext}^{-2}_A(A/fA, \omega_A^\bullet)$
augmentée de la longueur du conoyau de $t_0$.

\medskip\noindent
Cependant, si nous remplaçons $\mathcal{L}$ par $\mathcal{L}^n$ pour $n > 1$,
alors nous pouvons utiliser
$$
t^n :
M_n = \Gamma(U, \mathcal{L}^{\otimes n})
\longrightarrow
\Gamma(U, \mathcal{O}_U) = A
$$
à la place de $t$. Ceci remplace $t_0$ par sa puissance $n$-ième.
Ainsi, la longueur de l'image de l'application
$\text{Ext}^{-2}_A(M_n, \omega_A^\bullet) \leftarrow
\text{Ext}^{-2}_A(M_{n, 0}, \omega_A^\bullet)$
est au plus la longueur de $\text{Ext}^{-2}_A(A/fA, \omega_A^\bullet)$
augmentée de la longueur du conoyau de
$$
t_0^n : 
\text{Ext}^{-2}_A(A/fA, \omega_A^\bullet)
\longrightarrow
\text{Ext}^{-2}_A(M_{n, 0}, \omega_A^\bullet)
$$
Via l'isomorphisme $M_0 \cong A/fA$, l'application $t_0$ devient
$g : A/fA \to A/fA$ pour un certain $g \in A/fA$ et, via les isomorphismes
correspondants $M_{n, 0} \cong A/fA$, l'application $t_0^n$ devient
$g^n : A/fA \to A/fA$. Ainsi, la longueur du conoyau ci-dessus
est la longueur du quotient de
$\text{Ext}^{-2}_A(A/fA, \omega_A^\bullet)$ par $g^n$.
Puisque $\text{Ext}^{-2}_A(A/fA, \omega_A^\bullet)$ est un $A$-module de type fini
de support $T$ de dimension $1$, et puisque $V(g) \cap T$
est constitué du point fermé par notre choix de $t$,
cette longueur croît linéairement en $n$ d'après
Algèbre, lemme \ref{algebra-lemma-support-dimension-d}.

\medskip\noindent
Pour achever la démonstration, prouvons la dernière assertion. Supposons que
$f \in \mathfrak m \subset A$ satisfasse
à (1), (2) et (3), que $A$ soit $(S_2)$ et que $\dim(A) \geq 4$.
La condition (1) implique que $A$ est caténaire, voir
Complexes dualisants, lemme \ref{dualizing-lemma-universally-catenary}.
Alors $\Spec(A)$ est équidimensionnel d'après Cohomologie locale, lemme
\ref{local-cohomology-lemma-catenary-S2-equidimensional}.
Ainsi, $\dim(A_\mathfrak p) + \dim(A/\mathfrak p) \geq 4$
pour tout idéal premier $\mathfrak p$ de $A$. Alors
$\text{depth}(A_\mathfrak p) \geq \min(2, \dim(A_\mathfrak p))
\geq \min(2, 4 - \dim(A/\mathfrak p))$, et donc (4) est vérifiée.
\end{proof}

\begin{remark}
\label{remark-compare-SGA2}
Dans SGA 2, nous trouvons le résultat suivant. Soit $(A, \mathfrak m)$ un
anneau local noethérien. Soit $f \in \mathfrak m$. Supposons que $A$
soit un quotient d'un anneau régulier, que l'élément
$f$ soit un non-diviseur de zéro et que
\begin{enumerate}
\item[(a)] si $\mathfrak p \subset A$ est un idéal premier tel que
$\dim(A/\mathfrak p) = 1$, alors $\text{depth}(A_\mathfrak p) \geq 2$, et
\item[(b)] $\text{depth}(A/fA) \geq 3$, ou, de manière équivalente,
$\text{depth}(A) \geq 4$.
\end{enumerate}
Soient $U$, resp.\ $U_0$, les spectres épointés de $A$, resp.\ de $A/fA$. Alors
l'application
$$
\Pic(U) \to \Pic(U_0)
$$
est injective. C'est \cite[Exposee XI, lemme 3.16]{SGA2}\footnote{La condition
(a) résulte de la condition (b), voir
Algèbre, lemme \ref{algebra-lemma-depth-localization}.}. Ce résultat
de SGA 2 résulte de la proposition \ref{proposition-injective-pic},
car
\begin{enumerate}
\item un quotient d'un anneau régulier possède un complexe dualisant (voir
Complexes dualisants, lemme \ref{dualizing-lemma-regular-gorenstein} et
proposition \ref{dualizing-proposition-dualizing-essentially-finite-type}), et
\item si $\text{depth}(A) \geq 4$, alors $\text{depth}(A_\mathfrak p) \geq 2$
pour tout idéal premier $\mathfrak p$ tel que $\dim(A/\mathfrak p) = 2$, voir
Algèbre, lemme \ref{algebra-lemma-depth-localization}.
\end{enumerate}
\end{remark}






\section{Modules inversibles sur les spectres épointés, II}
\label{section-local-lefschetz-for-pic-surjective}

\noindent
Passons maintenant à la surjectivité dans le théorème de Lefschetz local pour le groupe de Picard.
Pour commencer, afin de prolonger un module inversible sur $U_0$ à un voisinage ouvert,
nous disposons du critère simple suivant.

\begin{lemma}
\label{lemma-surjective-Pic-first}
Soient $(A, \mathfrak m)$ un anneau local noethérien et $f \in \mathfrak m$.
Supposons que
\begin{enumerate}
\item $A$ soit complet pour la topologie $f$-adique,
\item $f$ soit un non-diviseur de zéro,
\item $H^1_\mathfrak m(A/fA)$ et $H^2_\mathfrak m(A/fA)$
soient des $A$-modules de type fini, et
\item $H^3_\mathfrak m(A/fA) = 0$\footnote{Observons que (3) et (4) sont vérifiées
si $\text{depth}(A/fA) \geq 4$, ou, de manière équivalente, $\text{depth}(A) \geq 5$.}.
\end{enumerate}
Soient $U$, resp.\ $U_0$, les spectres épointés de $A$, resp.\ de $A/fA$.
Alors
$$
\colim_{U_0 \subset U' \subset U\text{ ouvert}} \Pic(U')
\longrightarrow
\Pic(U_0)
$$
est surjective.
\end{lemma}

\begin{proof}
Soit $U_0 \subset U_n \subset U$ le $n$-ième voisinage infinitésimal
de $U_0$. Observons que le faisceau d'idéaux de $U_n$ dans $U_{n + 1}$ est
isomorphe à $\mathcal{O}_{U_0}$, puisque $U_0 \subset U$ est le sous-schéma
fermé principal défini par le non-diviseur de zéro $f$. Nous avons donc
une suite exacte de groupes abéliens
$$
\Pic(U_{n + 1}) \to \Pic(U_n) \to
H^2(U_0, \mathcal{O}_{U_0}) = H^3_\mathfrak m(A/fA) = 0
$$
voir Plus sur les morphismes, lemme
\ref{more-morphisms-lemma-picard-group-first-order-thickening}.
Ainsi, tout $\mathcal{O}_{U_0}$-module inversible est la restriction
d'un module formel cohérent inversible, c'est-à-dire d'un objet inversible de
$\textit{Coh}(U, f\mathcal{O}_U)$. On conclut en appliquant le
lemme \ref{lemma-equivalence}.
\end{proof}

\begin{remark}
\label{remark-surjective-Pic-second}
Soient $(A, \mathfrak m)$ un anneau local noethérien et $f \in \mathfrak m$.
La conclusion du lemme \ref{lemma-surjective-Pic-first} vaut si l'on suppose que
\begin{enumerate}
\item $A$ possède un complexe dualisant,
\item $A$ est complet pour la topologie $f$-adique,
\item $f$ est un non-diviseur de zéro,
\item l'une des conditions suivantes est vérifiée :
\begin{enumerate}
\item $A_f$ est $(S_2)$ et, pour tout $\mathfrak p \subset A$ minimal tel que
$f \not \in \mathfrak p$, on a $\dim(A/\mathfrak p) \geq 4$, ou
\item si $\mathfrak p \not \in V(f)$ et
$V(\mathfrak p) \cap V(f) \not = \{\mathfrak m\}$, alors
$\text{depth}(A_\mathfrak p) + \dim(A/\mathfrak p) > 3$.
\end{enumerate}
\item $H^3_{\mathfrak m}(A/fA) = 0$.
\end{enumerate}
La démonstration est exactement la même que celle du
lemme \ref{lemma-surjective-Pic-first},
en utilisant le lemme \ref{lemma-equivalence-better} à la place du
lemme \ref{lemma-equivalence}.
Il convient de faire ici deux remarques : (a)
il semble difficile de trouver des exemples où l'on sait que
$H^3_{\mathfrak m}(A/fA) = 0$ sans supposer que
$\text{depth}(A/fA) \geq 4$, et
(b) la démonstration du lemme \ref{lemma-equivalence-better} est
nettement plus difficile que celle du lemme \ref{lemma-equivalence}.
\end{remark}

\begin{lemma}
\label{lemma-surjective-Pic-first-better}
Soient $(A, \mathfrak m)$ un anneau local noethérien et $f \in \mathfrak m$.
Supposons que
\begin{enumerate}
\item les conditions du lemme \ref{lemma-surjective-Pic-first} soient vérifiées, et
\item pour tout idéal maximal $\mathfrak p \subset A_f$,
le spectre épointé de $(A_f)_\mathfrak p$ ait un groupe de Picard trivial.
\end{enumerate}
Soient $U$, resp.\ $U_0$, les spectres épointés de $A$, resp.\ de $A/fA$.
Alors
$$
\Pic(U) \longrightarrow \Pic(U_0)
$$
est surjective.
\end{lemma}

\begin{proof}
Soit $\mathcal{L}_0 \in \Pic(U_0)$. D'après le
lemme \ref{lemma-surjective-Pic-first},
il existe un ouvert $U_0 \subset U' \subset U$
et $\mathcal{L}' \in \Pic(U')$ dont la restriction
à $U_0$ est $\mathcal{L}_0$.
Puisque $U' \supset U_0$, on voit que $U \setminus U'$
est formé de points correspondant à des idéaux premiers
$\mathfrak p_1, \ldots, \mathfrak p_n$ comme en (2).
Par hypothèse, on peut trouver des modules inversibles
$\mathcal{L}'_i$ sur $\Spec(A_{\mathfrak p_i})$ qui coïncident avec
$\mathcal{L}'$ sur le spectre épointé
$U' \times_U \Spec(A_{\mathfrak p_i})$, puisque
les modules inversibles triviaux se prolongent toujours.
D'après Limites, lemme \ref{limits-lemma-glueing-near-closed-point-modules},
appliqué $n$ fois, on voit que $\mathcal{L}'$ se prolonge en un
module inversible sur $U$.
\end{proof}

\begin{lemma}
\label{lemma-local-pic-to-completion}
Soit $(A, \mathfrak m)$ un anneau local noethérien de profondeur $\geq 2$.
Soit $A^\wedge$ son complété. Soient $U$, resp.\ $U^\wedge$,
les spectres épointés de $A$, resp.\ de $A^\wedge$. Alors
$\Pic(U) \to \Pic(U^\wedge)$ est injective.
\end{lemma}

\begin{proof}
Soit $\mathcal{L}$ un $\mathcal{O}_U$-module inversible,
et notons $\mathcal{L}^\wedge$ son tiré en arrière sur $U^\wedge$.
On a $H^0(U, \mathcal{O}_U) = A$ par notre hypothèse sur la profondeur et par
Complexes dualisants, lemme \ref{dualizing-lemma-depth}, ainsi que
Cohomologie locale, lemme
\ref{local-cohomology-lemma-finiteness-pushforwards-and-H1-local}.
Ainsi, $\mathcal{L}$ est trivial si et seulement si
$M = H^0(U, \mathcal{L})$ est isomorphe à $A$ comme $A$-module.
(Détails omis.) Puisque $A \to A^\wedge$ est plat,
on a $M \otimes_A A^\wedge = \Gamma(U^\wedge, \mathcal{L}^\wedge)$
par changement de base plat, voir
Cohomologie des schémas, lemme \ref{coherent-lemma-flat-base-change-cohomology}.
Enfin, il est facile de voir que $M \cong A$ si et seulement si
$M \otimes_A A^\wedge \cong A^\wedge$.
\end{proof}

\begin{lemma}
\label{lemma-trivial-local-pic-regular}
Soit $(A, \mathfrak m)$ un anneau local régulier. Alors le groupe de Picard
du spectre épointé de $A$ est trivial.
\end{lemma}

\begin{proof}
Combinons Diviseurs, lemme \ref{divisors-lemma-extend-invertible-module},
avec Plus sur l'algèbre, lemme \ref{more-algebra-lemma-regular-local-UFD}.
\end{proof}

\noindent
Nous pouvons maintenant tirer parti des résultats précédents pour démontrer que
les groupes de Picard des spectres épointés
d'intersections complètes de dimension $\geq 4$ sont triviaux.
Rappelons qu'un anneau local noethérien est appelé une intersection
complète si son complété est le quotient d'un
anneau local régulier par l'idéal engendré par une suite régulière.
Voir la discussion dans Algèbre à puissances divisées, section \ref{dpa-section-lci}.

\begin{proposition}[Grothendieck]
\label{proposition-trivial-local-pic-complete-intersection}
Soit $(A, \mathfrak m)$ un anneau local noethérien. Si $A$ est une
intersection complète de dimension $\geq 4$, alors le groupe de Picard
du spectre épointé de $A$ est trivial.
\end{proposition}

\begin{proof}
D'après le lemme \ref{lemma-local-pic-to-completion}, on peut supposer que $A$ est
un anneau local complet. Par hypothèse, on peut écrire
$A = B/(f_1, \ldots, f_r)$, où $B$ est un anneau local régulier
complet et $f_1, \ldots, f_r$ une suite régulière.
Nous achevons la démonstration par récurrence sur $r$.
Le cas initial $r = 0$ résulte du
lemme \ref{lemma-trivial-local-pic-regular}.

\medskip\noindent
Supposons que $A = B/(f_1, \ldots, f_r)$ et que la proposition
soit vraie pour $r - 1$. Posons $A' = B/(f_1, \ldots, f_{r - 1})$ et appliquons le
lemme \ref{lemma-surjective-Pic-first-better} à $f_r \in A'$.
C'est possible, car :
\begin{enumerate}
\item la condition (1) du lemme \ref{lemma-surjective-Pic-first} est vérifiée
parce que nos anneaux locaux sont complets,
\item la condition (2) du lemme \ref{lemma-surjective-Pic-first} est vérifiée
puisque $f_1, \ldots, f_r$ est une suite régulière,
\item les conditions (3) et (4) du lemme \ref{lemma-surjective-Pic-first} sont vérifiées
puisque $A = A'/f_r A'$ est de Cohen--Macaulay et de dimension $\dim(A) \geq 4$,
\item la condition (2) du lemme \ref{lemma-surjective-Pic-first-better}
est vérifiée par l'hypothèse de récurrence, car
$\dim((A'_{f_r})_\mathfrak p) \geq 4$ pour tout idéal maximal
$\mathfrak p$ de $A'_{f_r}$ et car
$(A'_{f_r})_\mathfrak p = B_\mathfrak q/(f_1, \ldots, f_{r - 1})$
pour un certain idéal premier $\mathfrak q \subset B$, l'anneau $B_\mathfrak q$ étant régulier.
\end{enumerate}
Ceci achève la démonstration.
\end{proof}

\begin{example}
\label{example-grothendieck-sharp}
La borne sur la dimension dans la
proposition \ref{proposition-trivial-local-pic-complete-intersection}
est optimale. Par exemple, le groupe de Picard du spectre épointé de
$A = k[[x, y, z, w]]/(xy - zw)$ n'est pas trivial. En effet,
l'idéal $I = (x, z)$ définit un diviseur de Cartier effectif $D$ sur le spectre
épointé $U$ de $A$, car il est facile de voir que $I_x, I_y, I_z, I_w$
sont des idéaux inversibles de $A_x, A_y, A_z, A_w$. D'autre part,
$A/I$ est de profondeur $\geq 1$ (en fait $2$), donc $I$ est de profondeur $\geq 2$
(en fait $3$), d'où
$I = \Gamma(U, \mathcal{O}_U(-D))$. Ainsi, si $\mathcal{O}_U(-D)$
était trivial, on aurait $I \cong \Gamma(U, \mathcal{O}_U) = A$,
ce qui est faux puisque $I$ n'est pas engendré par $1$ élément.
\end{example}

\begin{example}
\label{example-grothendieck-sharp-bis}
La proposition \ref{proposition-trivial-local-pic-complete-intersection}
ne s'étend pas aux quotients
$$
A = B/(f_1, \ldots, f_r)
$$
où $B$ est régulier et $\dim(B) - r \geq 4$. Autrement dit, la condition selon laquelle
$f_1, \ldots, f_r$ est une suite régulière est, en général, nécessaire
à l'annulation du groupe de Picard du spectre épointé de $A$.
En effet, soit $k$ un corps et posons
$$
A = k[[a, b, x, y, z, u, v, w]]/(a^3, b^3, xa^2 + yab + zb^2, w^2)
$$
Observons que $A = A_0[w]/(w^2)$ avec
$A_0 = k[[a, b, x, y, z, u, v]]/(a^3, b^3, xa^2 + yab + zb^2)$.
Nous montrerons ci-dessous que $A_0$ est de profondeur $2$.
Notons $U$ le spectre épointé de $A$ et $U_0$ le spectre épointé
de $A_0$. Observons qu'il existe une suite exacte courte
$0 \to A_0 \to A \to A_0 \to 0$ dont la première flèche est
donnée par la multiplication par $w$.
D'après Plus sur les morphismes, lemme
\ref{more-morphisms-lemma-picard-group-first-order-thickening},
on obtient une suite exacte
$$
H^0(U, \mathcal{O}_U^*) \to
H^0(U_0, \mathcal{O}_{U_0}^*) \to
H^1(U_0, \mathcal{O}_{U_0}) \to
\Pic(U)
$$
Comme $A_0$, et donc $A$, est de profondeur $2$, on voit que
$H^0(U_0, \mathcal{O}_{U_0}) = A_0$ et $H^0(U, \mathcal{O}_U) = A$,
et que $H^1(U_0, \mathcal{O}_{U_0})$ est non nul, voir
Complexes dualisants, lemme \ref{dualizing-lemma-depth}, et
Cohomologie locale, lemme \ref{local-cohomology-lemma-local-cohomology}.
Ainsi, la dernière flèche affichée ci-dessus est non nulle, et l'on conclut
que $\Pic(U)$ n'est pas nul.

\medskip\noindent
Pour montrer que $A_0$ est de profondeur $2$, il suffit de montrer que
$A_1 = k[[a, b, x, y, z]]/(a^3, b^3, xa^2 + yab + zb^2)$
est de profondeur $0$. En effet, l'image de $a^2b^2$ est
un élément non nul de $A_1$ qui est annulé
par chacune des variables $a, b, x, y, z$. Par exemple,
$ya^2b^2 = (yab)(ab) = - (xa^2 + zb^2)(ab) = -xa^3b - yab^3 = 0$ dans $A_1$.
Les autres cas sont analogues.
\end{example}










\section{Application aux théorèmes de Lefschetz}
\label{section-lefschetz}

\noindent
Dans cette section, nous étudions la relation entre les faisceaux cohérents sur un
schéma projectif $P$ et les modules cohérents sur la complétion formelle le long
d'un diviseur ample $Q$.

\medskip\noindent
Soit $k$ un corps. Soit $P$ un schéma propre sur $k$.
Soit $\mathcal{L}$ un $\mathcal{O}_P$-module inversible ample.
Soit $s \in \Gamma(P, \mathcal{L})$ une section\footnote{Nous n'exigeons pas
que $s$ soit une section régulière. Par conséquent, $Q$ est seulement
un sous-schéma fermé localement principal de $P$, et non nécessairement un diviseur
de Cartier effectif.}, et soit
$Q = Z(s)$ son schéma des zéros, voir
Diviseurs, définition \ref{divisors-definition-zero-scheme-s}.
Pour tout $n \geq 1$, nous notons $Q_n = Z(s^n)$
le $n$-ième voisinage infinitésimal de $Q$.
Si $\mathcal{F}$ est un $\mathcal{O}_P$-module cohérent, nous notons
$\mathcal{F}_n = \mathcal{F}|_{Q_n}$ sa restriction, c'est-à-dire
le tiré en arrière de $\mathcal{F}$ par l'immersion fermée
$Q_n \to P$.

\begin{proposition}
\label{proposition-lefschetz}
Dans la situation ci-dessus, supposons qu'il existe un entier $\sigma$ tel que
pour tout point $p \in P \setminus Q$, on ait
$$
\text{depth}(\mathcal{F}_p) + \dim(\overline{\{p\}}) > \sigma
$$
Alors l'application
$$
H^i(P, \mathcal{F}) \longrightarrow \lim H^i(Q_n, \mathcal{F}_n)
$$
est un isomorphisme pour $0 \leq i < \sigma$.
\end{proposition}

\begin{proof}
Nous utiliserons Plus sur les morphismes, lemme \ref{more-morphisms-lemma-apply-proj-spec},
ainsi que les notations et les résultats de
Plus sur les morphismes, section \ref{more-morphisms-section-proj-spec},
sans plus de précisions ; cette démonstration ne peut être comprise sans
au moins comprendre l'énoncé du lemme.
Observons que, dans notre cas,
$A = \bigoplus_{m \geq 0} \Gamma(P, \mathcal{L}^{\otimes m})$
est une $k$-algèbre de type fini
dont toutes les composantes graduées sont des $k$-espaces vectoriels de dimension finie, voir
Cohomologie des schémas, lemme \ref{coherent-lemma-coherent-proper-ample}.

\medskip\noindent
Nous considérerons désormais $s$ comme un élément $f \in A_1 \subset A$, c'est-à-dire
comme un élément homogène de degré $1$ de $A$. Notons
$Y = V(f) \subset X$ le sous-schéma fermé défini par $f$.
Alors $U \cap Y = (\pi|_U)^{-1}(Q)$ au sens schématique.
Rappelons la notation
$\mathcal{F}_U = \pi^*\mathcal{F}|_U = (\pi|_U)^*\mathcal{F}$.
Il s'agit d'un $\mathcal{O}_U$-module cohérent.
Choisissons un $A$-module de type fini $M$ tel que
$\mathcal{F}_U = \widetilde{M}|_U$
(pour l'existence, voir Cohomologie locale, lemme
\ref{local-cohomology-lemma-finiteness-pushforwards-and-H1-local}).
Nous affirmons que $H^i_Z(M)$ est annulé par
une puissance de $f$ pour $i \leq \sigma + 1$.

\medskip\noindent
Pour démontrer cette assertion, nous appliquerons
Cohomologie locale, proposition \ref{local-cohomology-proposition-annihilator}.
En termes géométriques, on voit qu'il suffit de démontrer,
pour $u \in U$, $u \not \in Y$ et $z \in \overline{\{u\}} \cap Z$,
que
$$
\text{depth}(\mathcal{F}_{U, u}) +
\dim(\mathcal{O}_{\overline{\{u\}}, z}) > \sigma + 1
$$
Cela ne demande qu'un bref raisonnement.

\medskip\noindent
Observons que $Z = \Spec(A_0)$ est un ensemble fini de points fermés de $X$, car
$A_0$ est une $k$-algèbre de dimension finie.
(Le lecteur qui souhaiterait que $Z$ soit un singleton peut remplacer la
$k$-algèbre finie $A_0$ par $k$ ; cela ne change rien d'autre à la démonstration.)

\medskip\noindent
Le morphisme $\pi : L \to P$ et sa restriction $\pi|_U : U \to P$
sont lisses de dimension relative $1$.
Soient $u \in U$, $u \not \in Y$ et $z \in \overline{\{u\}} \cap Z$.
Soit $p = \pi(u) \in P \setminus Q$ son image.
Alors $u$ est soit un point générique
de la fibre de $\pi$ au-dessus de $p$, soit un point fermé de cette fibre.
Si $u$ est un point générique de la fibre, alors
$\text{depth}(\mathcal{F}_{U, u}) = \text{depth}(\mathcal{F}_p)$
et
$\dim(\overline{\{u\}}) = \dim(\overline{\{p\}}) + 1$.
Si $u$ est un point fermé de la fibre, alors
$\text{depth}(\mathcal{F}_{U, u}) = \text{depth}(\mathcal{F}_p) + 1$
et
$\dim(\overline{\{u\}}) = \dim(\overline{\{p\}})$.
Dans les deux cas, on a
$\dim(\overline{\{u\}}) = \dim(\mathcal{O}_{\overline{\{u\}}, z})$
parce que tout point de $Z$ est fermé. L'inégalité voulue
résulte donc de l'hypothèse de l'énoncé de la
proposition.

\medskip\noindent
Soit $A'$ le complété $f$-adique de $A$. Alors $A \to A'$ est plat d'après
Algèbre, lemme \ref{algebra-lemma-completion-flat}.
Notons $U' \subset X' = \Spec(A')$ l'image réciproque de
$U$, et de même pour $Y'$ et $Z'$. Soit $\mathcal{F}'$
sur $U'$ le tiré en arrière de $\mathcal{F}_U$, et posons
$M' = M \otimes_A A'$.
Par changement de base plat pour la cohomologie locale
(Cohomologie locale, lemme \ref{local-cohomology-lemma-torsion-change-rings}),
on a
$$
H^i_{Z'}(M') = H^i_Z(M) \otimes_A A'
$$
et l'on constate que, pour $i \leq \sigma + 1$,
ces groupes sont annulés par une puissance de $f$. Considérons le diagramme
$$
\xymatrix{
& H^i(U, \mathcal{F}_U) \ar[ld] \ar[d] \ar[r] &
\lim H^i(U, \mathcal{F}_U/f^n\mathcal{F}_U) \ar@{=}[d] \\
H^i(U, \mathcal{F}_U) \otimes_A A' \ar@{=}[r] &
H^i(U', \mathcal{F}') \ar[r] &
\lim H^i(U', \mathcal{F}'/f^n\mathcal{F}')
}
$$
La flèche horizontale inférieure est un isomorphisme pour $i < \sigma$ d'après le
lemme \ref{lemma-alternative-higher} et la propriété de torsion
que nous venons de démontrer. Le signe d'égalité horizontal provient du changement de base plat
(Cohomologie des schémas, lemme \ref{coherent-lemma-flat-base-change-cohomology}),
et le signe d'égalité vertical vient de ce que $U \cap Y$ et $U' \cap Y'$
ainsi que leurs $n$-ièmes voisinages infinitésimaux
sont identifiés deux à deux par des isomorphismes (puisque nous
complétons par rapport à $f$).

\medskip\noindent
En appliquant Plus sur les morphismes, équation
(\ref{more-morphisms-equation-cohomology-torsor}),
on obtient des décompositions en somme directe compatibles
$$
\lim H^i(U, \mathcal{F}_U/f^n\mathcal{F}_U) =
\lim
\left(
\bigoplus\nolimits_{m \in \mathbf{Z}}
H^i(Q_n, \mathcal{F}_n \otimes \mathcal{L}^{\otimes m})
\right)
$$
et
$$
H^i(U, \mathcal{F}_U) =
\bigoplus\nolimits_{m \in \mathbf{Z}}
H^i(P, \mathcal{F} \otimes \mathcal{L}^{\otimes m})
$$
On conclut donc par Algèbre, lemme \ref{algebra-lemma-daniel-litt}.
\end{proof}

\begin{lemma}
\label{lemma-lefschetz-addendum}
Soit $k$ un corps. Soit $X$ un schéma propre sur $k$.
Soit $\mathcal{L}$ un $\mathcal{O}_X$-module inversible ample.
Soit $s \in \Gamma(X, \mathcal{L})$. Soit $Y = Z(s)$ le
schéma des zéros de $s$, dont le $n$-ième voisinage infinitésimal est $Y_n = Z(s^n)$.
Soit $\mathcal{F}$ un $\mathcal{O}_X$-module cohérent.
Supposons que, pour tout $x \in X \setminus Y$, on ait
$$
\text{depth}(\mathcal{F}_x) + \dim(\overline{\{x\}}) > 1
$$
Alors $\Gamma(V, \mathcal{F}) \to \lim \Gamma(Y_n, \mathcal{F}|_{Y_n})$
est un isomorphisme pour tout sous-schéma ouvert $V \subset X$ contenant $Y$.
\end{lemma}

\begin{proof}
D'après la proposition \ref{proposition-lefschetz}, c'est vrai pour $V = X$.
Il suffit donc de montrer que l'application
$\Gamma(V, \mathcal{F}) \to \lim \Gamma(Y_n, \mathcal{F}|_{Y_n})$
est injective. Si $\sigma \in \Gamma(V, \mathcal{F})$
s'envoie sur zéro, son support est disjoint de $Y$ (détails omis ; indication :
utiliser le théorème d'intersection de Krull). Alors l'adhérence $T \subset X$
de $\text{Supp}(\sigma)$ est disjointe de $Y$.
Ainsi, $T$ est propre sur $k$ (car fermé dans $X$)
et affine (car fermé dans le schéma affine $X \setminus Y$, voir
Morphismes, lemme \ref{morphisms-lemma-proper-ample-delete-affine}),
donc fini sur $k$
(Morphismes, lemme \ref{morphisms-lemma-finite-proper}).
Ainsi, $T$ est un ensemble fini de points fermés de $X$.
L'hypothèse donne $\text{depth}(\mathcal{F}_x) \geq 2$, donc a fortiori au moins $1$,
pour $x \in T$. On conclut que
$\Gamma(V, \mathcal{F}) \to \Gamma(V \setminus T, \mathcal{F})$
est injective et que $\sigma = 0$, ce qu'il fallait démontrer.
\end{proof}


\begin{example}
\label{example-lefschetz}
Soit $k$ un corps et soit $X$ une variété propre sur $k$.
Soit $Y \subset X$ un diviseur de Cartier effectif tel que
$\mathcal{O}_X(Y)$ soit ample, et
notons $Y_n$ son $n$-ième voisinage infinitésimal.
Soit $\mathcal{E}$ un $\mathcal{O}_X$-module localement libre de type fini.
Voici quelques cas particuliers de la proposition \ref{proposition-lefschetz}.
\begin{enumerate}
\item Si $X$ est une courbe, on n'obtient aucun renseignement.
\item Si $X$ est une surface de Cohen--Macaulay (par exemple normale), alors
$$
H^0(X, \mathcal{E}) \to \lim H^0(Y_n, \mathcal{E}|_{Y_n})
$$
est un isomorphisme.
\item Si $X$ est une variété tridimensionnelle de Cohen--Macaulay, alors
$$
H^0(X, \mathcal{E}) \to \lim H^0(Y_n, \mathcal{E}|_{Y_n})
\quad\text{et}\quad
H^1(X, \mathcal{E}) \to \lim H^1(Y_n, \mathcal{E}|_{Y_n})
$$
sont des isomorphismes.
\end{enumerate}
Le principe général est sans doute clair. Si $X$ est une variété tridimensionnelle normale, alors
on peut conclure pour $H^0$, mais pas pour $H^1$.
\end{example}

\noindent
Avant de démontrer le résultat principal suivant, nous avons besoin d'un lemme.

\begin{lemma}
\label{lemma-Gm-equivariant-extend-canonically}
Dans la situation \ref{situation-algebraize}, soit $(\mathcal{F}_n)$ un
objet de $\textit{Coh}(U, I\mathcal{O}_U)$. Supposons que
\begin{enumerate}
\item $A$ soit un anneau gradué, $\mathfrak a = A_+$, et
$I$ un idéal homogène,
\item $(\mathcal{F}_n) = (\widetilde{M_n}|_U)$, où $(M_n)$
est un système projectif de $A$-modules gradués, et
\item $(\mathcal{F}_n)$ s'étende canoniquement à $X$.
\end{enumerate}
Alors il existe un $A$-module gradué de type fini $N$ tel que
\begin{enumerate}
\item[(a)] les systèmes projectifs $(N/I^nN)$ et $(M_n)$ soient pro-isomorphes
dans la catégorie quotient des $A$-modules gradués par les modules annulés par une puissance de $A_+$,
et
\item[(b)] $(\mathcal{F}_n)$ soit le complété du module cohérent
associé à $N$.
\end{enumerate}
\end{lemma}

\begin{proof}
Soit $(\mathcal{G}_n)$ le prolongement canonique du
lemme \ref{lemma-canonically-algebraizable}.
Les graduations de $A$ et de $M_n$ déterminent une action
$$
a : \mathbf{G}_m \times X \longrightarrow X
$$
du schéma en groupes $\mathbf{G}_m$ sur $X$ telle que
$(\widetilde{M_n})$ devienne un système projectif de
modules quasi-cohérents $\mathbf{G}_m$-équivariants sur $\mathcal{O}_X$, voir
Groupoïdes, exemple \ref{groupoids-example-Gm-on-affine}.
Puisque $\mathfrak a$ et $I$ sont des idéaux homogènes,
les sous-schémas $Z$ et $Y$ sont fermés et le sous-schéma $U$ est ouvert ;
tous trois sont des sous-schémas $\mathbf{G}_m$-invariants.
La restriction $(\mathcal{F}_n)$ de $(\widetilde{M_n})$
est un système projectif de modules cohérents $\mathbf{G}_m$-équivariants
qui sont des $\mathcal{O}_U$-modules. Autrement dit, $(\mathcal{F}_n)$
est un module formel cohérent $\mathbf{G}_m$-équivariant,
au sens où il existe un isomorphisme
$$
\alpha : (a^*\mathcal{F}_n) \longrightarrow (p^*\mathcal{F}_n)
$$
sur $\mathbf{G}_m \times U$ satisfaisant à une condition de cocycle convenable.
Comme $a$ et $p$ sont des morphismes plats de schémas affines,
le lemme \ref{lemma-canonically-extend-base-change}
montre qu'il existe un unique isomorphisme
$$
\beta : (a^*\mathcal{G}_n) \longrightarrow (p^*\mathcal{G}_n)
$$
sur $\mathbf{G}_m \times X$ dont la restriction est $\alpha$ sur
$\mathbf{G}_m \times U$. L'unicité garantit que
$\beta$ satisfait à la condition de cocycle correspondante.
De cette manière, chaque $\mathcal{G}_n$ devient
un module cohérent $\mathbf{G}_m$-équivariant sur $\mathcal{O}_X$,
de façon compatible avec les applications de transition.

\medskip\noindent
D'après Groupoïdes, lemme \ref{groupoids-lemma-Gm-equivariant-module},
on voit que $\mathcal{G}_n$, muni de sa structure $\mathbf{G}_m$-équivariante,
correspond à un $A$-module gradué $N_n$. Les applications de transition
$N_{n + 1} \to N_n$ sont des homomorphismes de modules gradués. Notons que $N_n$ est un
$A$-module de type fini et que $N_n = N_{n + 1}/I^n N_{n + 1}$, puisque
$(\mathcal{G}_n)$ est un objet de $\textit{Coh}(X, I\mathcal{O}_X)$.
Soit $N$ le $A$-module gradué de type fini fourni par
Algèbre, lemme \ref{algebra-lemma-finiteness-graded}.
Alors $N_n = N/I^nN$, de sorte que $(\mathcal{G}_n)$
est le complété du module cohérent
associé à $N$, et l'on voit a fortiori que (b) est vraie.

\medskip\noindent
Pour voir (a), il faut expliciter un peu davantage la situation ci-dessus.
Observons d'abord que le noyau et le conoyau de $M_n \to H^0(U, \mathcal{F}_n)$
sont de torsion annulée par une puissance de $A_+$ (Cohomologie locale, lemme
\ref{local-cohomology-lemma-finiteness-pushforwards-and-H1-local}).
Observons que $H^0(U, \mathcal{F}_n)$ est muni d'une graduation naturelle
pour laquelle ces applications et les applications de transition du système sont
des homomorphismes gradués de $A$-modules ; on peut par exemple utiliser le fait que
$(U \to X)_*\mathcal{F}_n$ est un module $\mathbf{G}_m$-équivariant
sur $X$ et appliquer 
Groupoïdes, lemme \ref{groupoids-lemma-Gm-equivariant-module}.
Rappelons ensuite que $(N_n)$ et $(H^0(U, \mathcal{F}_n))$
sont pro-isomorphes d'après la définition \ref{definition-canonically-algebraizable}
et le lemme \ref{lemma-canonically-algebraizable}.
Nous omettons la vérification que les applications définissant ce
pro-isomorphisme sont des homomorphismes de modules gradués.
Ainsi, $(N_n)$ et $(M_n)$ sont pro-isomorphes dans la catégorie quotient des
$A$-modules gradués par les modules de torsion annulés par une puissance de $A_+$.
\end{proof}

\noindent
Soit $k$ un corps. Soit $P$ un schéma propre sur $k$.
Soit $\mathcal{L}$ un $\mathcal{O}_P$-module inversible ample.
Soit $s \in \Gamma(P, \mathcal{L})$ une section et soit
$Q = Z(s)$ le schéma des zéros, voir
Diviseurs, définition \ref{divisors-definition-zero-scheme-s}.
Soit $\mathcal{I} \subset \mathcal{O}_P$ le faisceau d'idéaux de $Q$.
Nous noterons $\textit{Coh}(P, \mathcal{I})$ la catégorie
des modules formels cohérents introduite dans
Cohomologie des schémas, section \ref{coherent-section-coherent-formal}.

\begin{proposition}
\label{proposition-lefschetz-existence}
Dans la situation ci-dessus, soit $(\mathcal{F}_n)$ un objet de
$\textit{Coh}(P, \mathcal{I})$. Supposons que, pour tout $q \in Q$ et tout
idéal premier $\mathfrak p \in \mathcal{O}_{P, q}^\wedge$ tel que
$\mathfrak p \not \in V(\mathcal{I}_q^\wedge)$, on ait
$$
\text{depth}((\mathcal{F}_q^\wedge)_\mathfrak p) +
\dim(\mathcal{O}_{P, q}^\wedge/\mathfrak p) +
\dim(\overline{\{q\}}) > 2
$$
Alors $(\mathcal{F}_n)$ est le complété d'un
$\mathcal{O}_P$-module cohérent.
\end{proposition}

\begin{proof}
D'après Cohomologie des schémas, lemme \ref{coherent-lemma-existence-easy},
pour démontrer l'énoncé, on peut remplacer $(\mathcal{F}_n)$ par un objet
qui n'en diffère que par de la $\mathcal{I}$-torsion (voir ci-dessous pour plus de précision).
Posons $T' = \{q \in Q \mid \dim(\overline{\{q\}}) = 0\}$
et $T = \{q \in Q \mid \dim(\overline{\{q\}}) \leq 1\}$.
L'hypothèse de la proposition signifie exactement que
$Q \subset P$, $(\mathcal{F}_n)$ et $T' \subset T \subset Q$
satisfont aux conditions du
lemme \ref{lemma-improvement-formal-coherent-module-better}
avec $d = 1$ ; outre des manipulations immédiates d'inégalités, on utilise que
$V(\mathfrak p) \cap V(\mathcal{I}^\wedge_y) = \{\mathfrak m^\wedge_y\}
\Leftrightarrow \dim(\mathcal{O}_{P, q}^\wedge/\mathfrak p) = 1$
puisque $\mathcal{I}_y^\wedge$ est engendré par $1$ élément.
En combinant ces deux remarques, on peut remplacer $(\mathcal{F}_n)$ par
l'objet $(\mathcal{H}_n)$ de $\textit{Coh}(P, \mathcal{I})$ fourni par le
lemme \ref{lemma-improvement-formal-coherent-module-better}.
Ainsi, on peut supposer, et on suppose, que $(\mathcal{F}_n)$ est pro-isomorphe à
un système projectif $(\mathcal{F}_n'')$ de $\mathcal{O}_P$-modules cohérents
tel que $\text{depth}(\mathcal{F}''_{n, q}) + \dim(\overline{\{q\}}) \geq 2$
pour tout $q \in Q$.

\medskip\noindent
Nous utiliserons Plus sur les morphismes, lemme \ref{more-morphisms-lemma-apply-proj-spec},
ainsi que les notations et les résultats de
Plus sur les morphismes, section \ref{more-morphisms-section-proj-spec},
sans plus de précisions ; cette démonstration ne peut être comprise sans
au moins comprendre l'énoncé du lemme.
Observons que, dans notre cas,
$A = \bigoplus_{m \geq 0} \Gamma(P, \mathcal{L}^{\otimes m})$
est une $k$-algèbre de type fini
dont toutes les composantes graduées sont des $k$-espaces vectoriels de dimension finie, voir
Cohomologie des schémas, lemme \ref{coherent-lemma-coherent-proper-ample}.

\medskip\noindent
D'après Cohomologie des schémas, lemme \ref{coherent-lemma-inverse-systems-pullback},
le tiré en arrière par $\pi|_U : U \to P$ est un objet
$(\pi|_U^*\mathcal{F}_n)$ de $\textit{Coh}(U, f\mathcal{O}_U)$
qui est pro-isomorphe au système projectif
$(\pi|_U^*\mathcal{F}_n'')$ de $\mathcal{O}_U$-modules cohérents.
Nous affirmons que
$$
\text{depth}(\pi|_U^*\mathcal{F}''_{n, y}) + \delta_Z^Y(y) \geq 3
$$
pour tout $y \in U \cap Y$. Puisque tous les points de $Z$ sont fermés, on
voit que $\delta_Z^Y(y) \geq \dim(\overline{\{y\}})$ pour tout
$y \in U \cap Y$, voir le lemme \ref{lemma-discussion}.
Soit $q \in Q$ l'image de $y$. Comme le morphisme $\pi : U \to P$ est
lisse de dimension relative $1$, on voit que $y$ est soit un point fermé
d'une fibre de $\pi$, soit un point générique.
On en déduit que
$$
\text{depth}(\pi^*\mathcal{F}''_{n, y}) + \delta_Z^Y(y)
\geq
\text{depth}(\pi^*\mathcal{F}''_{n, y}) + \dim(\overline{\{y\}}) =
\text{depth}(\mathcal{F}''_{n, q}) + \dim(\overline{\{q\}}) + 1
$$
car soit la profondeur, soit la dimension augmente de $1$.
Ceci démontre l'assertion.

\medskip\noindent
D'après le lemme \ref{lemma-cd-1-canonical}, on conclut que
$(\pi|_U^*\mathcal{F}_n)$ se prolonge canoniquement sur $X$.
Observons que
$$
M_n = \Gamma(U, \pi|_U^*\mathcal{F}_n) =
\bigoplus\nolimits_{m \in \mathbf{Z}}
\Gamma(P, \mathcal{F}_n \otimes_{\mathcal{O}_P} \mathcal{L}^{\otimes m})
$$
est canoniquement un $A$-module gradué, voir
Plus sur les morphismes, équation (\ref{more-morphisms-equation-cohomology-torsor}).
D'après Propriétés, lemme \ref{properties-lemma-quasi-coherent-quasi-affine},
on a $\pi|_U^*\mathcal{F}_n = \widetilde{M_n}|_U$.
On peut donc appliquer le lemme \ref{lemma-Gm-equivariant-extend-canonically}
pour trouver un $A$-module gradué de type fini $N$ tel que
$(M_n)$ et $(N/I^nN)$ soient pro-isomorphes dans la catégorie quotient
des $A$-modules gradués par les modules de $A_+$-torsion.
Soit $\mathcal{F}$ le $\mathcal{O}_P$-module cohérent
associé à $N$, voir
Cohomologie des schémas, proposition
\ref{coherent-proposition-coherent-modules-on-proj-general}.
La même proposition montre que $(\mathcal{F}/\mathcal{I}^n\mathcal{F})$
est pro-isomorphe à $(\mathcal{F}_n)$.
Puisque tous deux sont des objets de $\textit{Coh}(P, \mathcal{I})$,
on conclut par le lemme \ref{lemma-recognize-formal-coherent-modules}.
\end{proof}

\begin{example}
\label{example-lefschetz-existence}
Soit $k$ un corps et soit $X$ une variété propre sur $k$.
Soit $Y \subset X$ un diviseur de Cartier effectif tel que
$\mathcal{O}_X(Y)$ soit ample, et notons
$\mathcal{I} \subset \mathcal{O}_X$ le faisceau d'idéaux
correspondant. Soit $(\mathcal{E}_n)$ un objet
de $\textit{Coh}(X, \mathcal{I})$ tel que $\mathcal{E}_n$ soit
localement libre de type fini.
Voici quelques cas particuliers de la
proposition \ref{proposition-lefschetz-existence}.
\begin{enumerate}
\item Si $X$ est une courbe ou une surface, on n'obtient aucun renseignement.
\item Si $X$ est une variété tridimensionnelle de Cohen--Macaulay, alors
$(\mathcal{E}_n)$ est le complété d'un
$\mathcal{O}_X$-module cohérent $\mathcal{E}$.
\item Plus généralement, si $\dim(X) \geq 3$ et si $X$ est $(S_3)$, alors
$(\mathcal{E}_n)$ est le complété d'un
$\mathcal{O}_X$-module cohérent $\mathcal{E}$.
\end{enumerate}
Bien sûr, si $\mathcal{E}$ existe, alors $\mathcal{E}$ est localement libre
de type fini dans un voisinage ouvert de $Y$.
\end{example}

\begin{proposition}
\label{proposition-lefschetz-equivalence}
Soit $k$ un corps. Soit $X$ un schéma propre sur $k$.
Soit $\mathcal{L}$ un $\mathcal{O}_X$-module inversible ample,
et soit $s \in \Gamma(X, \mathcal{L})$. Soit $Y = Z(s)$
le schéma des zéros de $s$, et notons $\mathcal{I} \subset \mathcal{O}_X$
le faisceau d'idéaux correspondant.
Soit $\mathcal{V}$ l'ensemble des sous-schémas ouverts de $X$ contenant $Y$,
ordonné par l'inclusion inverse.
Supposons que, pour tout $x \in X \setminus Y$, on ait
$$
\text{depth}(\mathcal{O}_{X, x}) + \dim(\overline{\{x\}}) > 2
$$
Alors le foncteur de complétion
$$
\colim_\mathcal{V}
\textit{Coh}(\mathcal{O}_V)
\longrightarrow
\textit{Coh}(X, \mathcal{I})
$$
induit une équivalence entre les sous-catégories pleines d'objets localement libres de type fini.
\end{proposition}

\begin{proof}
Pour démontrer la pleine fidélité, il suffit de montrer que
$$
\colim_\mathcal{V} \Gamma(V, \mathcal{L}^{\otimes m})
\longrightarrow
\lim \Gamma(Y_n, \mathcal{L}^{\otimes m}|_{Y_n})
$$
est un isomorphisme pour tout $m$, voir le
lemme \ref{lemma-completion-fully-faithful-general}.
Ceci résulte du lemme \ref{lemma-lefschetz-addendum}.

\medskip\noindent
Surjectivité essentielle. Soit $(\mathcal{F}_n)$ un objet localement
libre de type fini de $\textit{Coh}(X, \mathcal{I})$. Alors, pour $y \in Y$,
$\mathcal{F}_y^\wedge = \lim \mathcal{F}_{n, y}$
est un $\mathcal{O}_{X, y}^\wedge$-module libre de type fini.
Soit $\mathfrak p \subset \mathcal{O}_{X, y}^\wedge$
un idéal premier tel que $\mathfrak p \not \in V(\mathcal{I}_y^\wedge)$.
Alors $\mathfrak p$ est au-dessus d'un idéal premier $\mathfrak p_0 \subset \mathcal{O}_{X, y}$
qui correspond à une spécialisation $x \leadsto y$ avec
$x \not \in Y$. D'après Cohomologie locale, lemme
\ref{local-cohomology-lemma-change-completion}
et des résultats de théorie de la dimension
(voir Variétés, section \ref{varieties-section-algebraic-schemes}),
on a
$$
\text{depth}((\mathcal{O}_{X, y}^\wedge)_\mathfrak p) +
\dim(\mathcal{O}_{X, y}^\wedge/\mathfrak p) =
\text{depth}(\mathcal{O}_{X, x}) +
\dim(\overline{\{x\}}) - \dim(\overline{\{y\}})
$$
Ainsi, nos hypothèses impliquent que les hypothèses de la
proposition \ref{proposition-lefschetz-existence}
sont satisfaites, et l'on obtient que $(\mathcal{F}_n)$
est le complété d'un $\mathcal{O}_X$-module cohérent $\mathcal{F}$.
Il s'ensuit que $\mathcal{F}_y$ est libre de type fini pour tout $y \in Y$,
et donc que $\mathcal{F}$ est localement libre de type fini dans un
voisinage ouvert $V$ de $Y$. Ceci achève la démonstration.
\end{proof}









\begin{multicols}{2}[\section{Autres chapitres}]
\noindent
Préliminaires
\begin{enumerate}
\item \hyperref[introduction-section-phantom]{Introduction}
\item \hyperref[conventions-section-phantom]{Conventions}
\item \hyperref[sets-section-phantom]{Théorie des ensembles}
\item \hyperref[categories-section-phantom]{Catégories}
\item \hyperref[topology-section-phantom]{Topologie}
\item \hyperref[sheaves-section-phantom]{Faisceaux sur les espaces}
\item \hyperref[sites-section-phantom]{Sites et faisceaux}
\item \hyperref[stacks-section-phantom]{Champs}
\item \hyperref[fields-section-phantom]{Corps}
\item \hyperref[algebra-section-phantom]{Algèbre commutative}
\item \hyperref[brauer-section-phantom]{Groupes de Brauer}
\item \hyperref[homology-section-phantom]{Algèbre homologique}
\item \hyperref[derived-section-phantom]{Catégories dérivées}
\item \hyperref[simplicial-section-phantom]{Méthodes simpliciales}
\item \hyperref[more-algebra-section-phantom]{Compléments d'algèbre}
\item \hyperref[smoothing-section-phantom]{Lissage des morphismes d'anneaux}
\item \hyperref[modules-section-phantom]{Faisceaux de modules}
\item \hyperref[sites-modules-section-phantom]{Modules sur les sites}
\item \hyperref[injectives-section-phantom]{Objets injectifs}
\item \hyperref[cohomology-section-phantom]{Cohomologie des faisceaux}
\item \hyperref[sites-cohomology-section-phantom]{Cohomologie sur les sites}
\item \hyperref[dga-section-phantom]{Algèbre différentielle graduée}
\item \hyperref[dpa-section-phantom]{Algèbre à puissances divisées}
\item \hyperref[sdga-section-phantom]{Faisceaux différentiels gradués}
\item \hyperref[hypercovering-section-phantom]{Hyperrecouvrements}
\end{enumerate}
Schémas
\begin{enumerate}
\setcounter{enumi}{25}
\item \hyperref[schemes-section-phantom]{Schémas}
\item \hyperref[constructions-section-phantom]{Constructions de schémas}
\item \hyperref[properties-section-phantom]{Propriétés des schémas}
\item \hyperref[morphisms-section-phantom]{Morphismes de schémas}
\item \hyperref[coherent-section-phantom]{Cohomologie des schémas}
\item \hyperref[divisors-section-phantom]{Diviseurs}
\item \hyperref[limits-section-phantom]{Limites de schémas}
\item \hyperref[varieties-section-phantom]{Variétés}
\item \hyperref[topologies-section-phantom]{Topologies sur les schémas}
\item \hyperref[descent-section-phantom]{Descente}
\item \hyperref[perfect-section-phantom]{Catégories dérivées de schémas}
\item \hyperref[more-morphisms-section-phantom]{Compléments sur les morphismes}
\item \hyperref[flat-section-phantom]{Compléments sur la platitude}
\item \hyperref[groupoids-section-phantom]{Schémas en groupoïdes}
\item \hyperref[more-groupoids-section-phantom]{Compléments sur les schémas en groupoïdes}
\item \hyperref[etale-section-phantom]{Morphismes étales de schémas}
\end{enumerate}
Thèmes de théorie des schémas
\begin{enumerate}
\setcounter{enumi}{41}
\item \hyperref[chow-section-phantom]{Homologie de Chow}
\item \hyperref[intersection-section-phantom]{Théorie de l'intersection}
\item \hyperref[pic-section-phantom]{Schémas de Picard des courbes}
\item \hyperref[weil-section-phantom]{Théories cohomologiques de Weil}
\item \hyperref[adequate-section-phantom]{Modules adéquats}
\item \hyperref[dualizing-section-phantom]{Complexes dualisants}
\item \hyperref[duality-section-phantom]{Dualité pour les schémas}
\item \hyperref[discriminant-section-phantom]{Discriminants et différentes}
\item \hyperref[derham-section-phantom]{Cohomologie de de Rham}
\item \hyperref[local-cohomology-section-phantom]{Cohomologie locale}
\item \hyperref[algebraization-section-phantom]{Géométrie algébrique et formelle}
\item \hyperref[curves-section-phantom]{Courbes algébriques}
\item \hyperref[resolve-section-phantom]{Résolution des surfaces}
\item \hyperref[models-section-phantom]{Réduction semi-stable}
\item \hyperref[functors-section-phantom]{Foncteurs et morphismes}
\item \hyperref[equiv-section-phantom]{Catégories dérivées de variétés}
\item \hyperref[pione-section-phantom]{Groupes fondamentaux de schémas}
\item \hyperref[etale-cohomology-section-phantom]{Cohomologie étale}
\item \hyperref[crystalline-section-phantom]{Cohomologie cristalline}
\item \hyperref[proetale-section-phantom]{Cohomologie pro-étale}
\item \hyperref[relative-cycles-section-phantom]{Cycles relatifs}
\item \hyperref[more-etale-section-phantom]{Compléments de cohomologie étale}
\item \hyperref[trace-section-phantom]{La formule des traces}
\end{enumerate}
Espaces algébriques
\begin{enumerate}
\setcounter{enumi}{64}
\item \hyperref[spaces-section-phantom]{Espaces algébriques}
\item \hyperref[spaces-properties-section-phantom]{Propriétés des espaces algébriques}
\item \hyperref[spaces-morphisms-section-phantom]{Morphismes d'espaces algébriques}
\item \hyperref[decent-spaces-section-phantom]{Espaces algébriques décents}
\item \hyperref[spaces-cohomology-section-phantom]{Cohomologie des espaces algébriques}
\item \hyperref[spaces-limits-section-phantom]{Limites d'espaces algébriques}
\item \hyperref[spaces-divisors-section-phantom]{Diviseurs sur les espaces algébriques}
\item \hyperref[spaces-over-fields-section-phantom]{Espaces algébriques sur des corps}
\item \hyperref[spaces-topologies-section-phantom]{Topologies sur les espaces algébriques}
\item \hyperref[spaces-descent-section-phantom]{Descente et espaces algébriques}
\item \hyperref[spaces-perfect-section-phantom]{Catégories dérivées d'espaces}
\item \hyperref[spaces-more-morphisms-section-phantom]{Compléments sur les morphismes d'espaces}
\item \hyperref[spaces-flat-section-phantom]{Platitude sur les espaces algébriques}
\item \hyperref[spaces-groupoids-section-phantom]{Groupoïdes dans les espaces algébriques}
\item \hyperref[spaces-more-groupoids-section-phantom]{Compléments sur les groupoïdes dans les espaces}
\item \hyperref[bootstrap-section-phantom]{Amorçage}
\item \hyperref[spaces-pushouts-section-phantom]{Sommes amalgamées d'espaces algébriques}
\end{enumerate}
Thèmes de géométrie
\begin{enumerate}
\setcounter{enumi}{81}
\item \hyperref[spaces-chow-section-phantom]{Groupes de Chow des espaces}
\item \hyperref[groupoids-quotients-section-phantom]{Quotients de groupoïdes}
\item \hyperref[spaces-more-cohomology-section-phantom]{Compléments sur la cohomologie des espaces}
\item \hyperref[spaces-simplicial-section-phantom]{Espaces simpliciaux}
\item \hyperref[spaces-duality-section-phantom]{Dualité pour les espaces}
\item \hyperref[formal-spaces-section-phantom]{Espaces algébriques formels}
\item \hyperref[restricted-section-phantom]{Algébrisation des espaces formels}
\item \hyperref[spaces-resolve-section-phantom]{Résolution des surfaces revisitée}
\end{enumerate}
Théorie des déformations
\begin{enumerate}
\setcounter{enumi}{89}
\item \hyperref[formal-defos-section-phantom]{Théorie formelle des déformations}
\item \hyperref[defos-section-phantom]{Théorie des déformations}
\item \hyperref[cotangent-section-phantom]{Le complexe cotangent}
\item \hyperref[examples-defos-section-phantom]{Problèmes de déformation}
\end{enumerate}
Champs algébriques
\begin{enumerate}
\setcounter{enumi}{93}
\item \hyperref[algebraic-section-phantom]{Champs algébriques}
\item \hyperref[examples-stacks-section-phantom]{Exemples de champs}
\item \hyperref[stacks-sheaves-section-phantom]{Faisceaux sur les champs algébriques}
\item \hyperref[criteria-section-phantom]{Critères de représentabilité}
\item \hyperref[artin-section-phantom]{Axiomes d'Artin}
\item \hyperref[quot-section-phantom]{Espaces de Quot et de Hilbert}
\item \hyperref[stacks-properties-section-phantom]{Propriétés des champs algébriques}
\item \hyperref[stacks-morphisms-section-phantom]{Morphismes de champs algébriques}
\item \hyperref[stacks-limits-section-phantom]{Limites de champs algébriques}
\item \hyperref[stacks-cohomology-section-phantom]{Cohomologie des champs algébriques}
\item \hyperref[stacks-perfect-section-phantom]{Catégories dérivées de champs}
\item \hyperref[stacks-introduction-section-phantom]{Introduction aux champs algébriques}
\item \hyperref[stacks-more-morphisms-section-phantom]{Compléments sur les morphismes de champs}
\item \hyperref[stacks-geometry-section-phantom]{La géométrie des champs}
\end{enumerate}
Thèmes de théorie des espaces de modules
\begin{enumerate}
\setcounter{enumi}{107}
\item \hyperref[moduli-section-phantom]{Champs de modules}
\item \hyperref[moduli-curves-section-phantom]{Espaces de modules de courbes}
\end{enumerate}
Divers
\begin{enumerate}
\setcounter{enumi}{109}
\item \hyperref[examples-section-phantom]{Exemples}
\item \hyperref[exercises-section-phantom]{Exercices}
\item \hyperref[guide-section-phantom]{Guide bibliographique}
\item \hyperref[desirables-section-phantom]{Desiderata}
\item \hyperref[coding-section-phantom]{Style de programmation}
\item \hyperref[obsolete-section-phantom]{Obsolète}
\item \hyperref[fdl-section-phantom]{Licence de documentation libre GNU}
\item \hyperref[index-section-phantom]{Index généré automatiquement}
\end{enumerate}
\end{multicols}


\bibliography{my}
\bibliographystyle{amsalpha}

\end{document}
